\chapter{Proper Base Change}\label{chap:proper-base-change}
\section{Target, scope, and proof spine}

Fix a complete discrete valuation ring $\Lambda$ with maximal ideal
$\mathfrak m$, put $\Lambda_n=\Lambda/\mathfrak m^n$ for $n\geq1$, and
assume that $\ell$ is invertible on every scheme under consideration.  All
schemes in the ordinary finite-level statement are noetherian and
finite-dimensional.  The functors are first formed in the ambient derived
categories and are restricted to constructible objects only through the
separately named constructibility inputs.

\begin{openproblem}[Proper base change target]
\label{OBL-PROPER-BASE-CHANGE}
For a Cartesian square
\[
\begin{CD}
X' @>{g'}>> X\\
@V{f'}VV @VV{f}V\\
Y' @>{g}>> Y
\end{CD}
\qquad (f\text{ proper and finitely presented}),
\]
and for $K\in D_c^+(X,\Lambda_n)$,
construct the Beck--Chevalley mate
\[
 \operatorname{Ex}_{\square,K}:Lg^*Rf_*K\longrightarrow Rf'_*Lg'^*K,
\]
prove that it is an isomorphism, and prove naturality, identity-square,
horizontal-pasting, vertical-pasting, coefficient, localization, projection
formula, and Frobenius compatibilities in the coefficient layer where each
expression is typed.
\uses{OBL-BA31-S0-GLOBAL-MATE,OBL-BA31-S1-FINITE,OBL-BA31-S2-STABILITY,OBL-BA31-S3-PROJECTIVE-REDUCTION,OBL-BA31-S4-HENSELIAN-P1,OBL-BA31-S5-STALK-COMPARISON,OBL-BA31-S6-SHEAF-ASSEMBLY,OBL-BA31-S7-BOUNDED-BELOW}
\openobligation{OBL-PROPER-BASE-CHANGE}
\end{openproblem}

\begin{conditionaltheorem}[Reassembly represented by S0--S10]
\label{thm:weil-pbc-reassembly-s0-s10}
Assume every exact terminal input in Appendix~\ref{app:gap-ledger}, together
with every imported interface whose manifest status is not closed.  Then the
ordered S0--S7 argument proves finite-torsion proper base change in
$D^+_c(X_{\mathrm{et}})$; S8 transports the comparison to finite-level
pro-etale objects; and S9--S10 prove the completed comparison under the
displayed perfect-quotient, projection-formula, reconstruction, and
$R\varprojlim$ hypotheses.
\uses{OBL-BA31-FINITE-PBC}
\end{conditionaltheorem}
\begin{proof}
The detailed construction and every intermediate proof are reproduced in
Chapters~\ref{chap:s0-s10} and \ref{chap:compact-support}.  Appendix
\ref{app:gap-ledger} records the exact unproved hypotheses, so this theorem
does not promote the historical owner.
\end{proof}

\section{Ordered dependency spine}
\begin{enumerate}[label=\textbf{S\arabic*.},leftmargin=3em]
\item Construct the single global unit--Cartesian-identification--counit mate.
\item Prove the finite-morphism fibre formula and exactness.
\item Prove composition stability and proper-surjective cancellation.
\item Use Chow reduction and the finite symmetric-product map to reduce to
      $\mathbf P^1$.
\item Prove the strict-henselian $\mathbf P^1$ comparison from degree zero,
      affine Gabber comparison, and finite-etale torsor lifting.
\item Identify the global stalk map using strict henselization and proper-curve
      invariance under separably closed extension.
\item Assemble prime and arbitrary torsion coefficients, then pass to $D^+$ by
      a bounded-below hypercohomology spectral sequence.
\item Transport finite levels to the pro-etale site.
\item Reconstruct completed coefficients by perfect reduction and
      $R\varprojlim$.
\item Verify naturality, identity, pasting, and typed Frobenius coherence.
\end{enumerate}

\statusnote{The proof order above is a reduction order.  Formal status may be
promoted only in the reverse dependency direction: exact leaves first, then
S0--S7, then finite PBC, and only later the completed and compact-support
exports.}

\section{Corrected typed finite-level interface}\label{chap:corrected-pbc}
\statusnote{This chapter is the authoritative typing boundary for the ordinary finite-level proof.  It replaces the ill-typed historical sectionwise S0 formula.  Every named S0--S7 middle owner remains OPEN/INCOMPLETE, and every atomic input is listed in Appendix~\ref{app:gap-ledger}.}
For a Cartesian square of noetherian finite-dimensional schemes
\[
\begin{array}{ccc}
X'&\xrightarrow{g'}&X\\
\mathllap{f'}\downarrow&&\downarrow\mathrlap{f}\\
Y'&\xrightarrow{g}&Y,
\end{array}
\]
and $K\in D_c^+(X,\Lambda_n)$, first construct in the ambient category the
single derived morphism
\[
\begin{split}
 \operatorname{BC}_{f,g,K}:Lg^*Rf_*K
 &\xrightarrow{\eta_{f'}}Rf'_*Lf'^*Lg^*Rf_*K\\
 &\xrightarrow{\mathrm{cart}}Rf'_*Lg'^*Lf^*Rf_*K
 \xrightarrow{Rf'_*(Lg'^*\varepsilon_f)}Rf'_*Lg'^*K.
\end{split}
\tag{BA31-MATE}
\]
Replacement independence and naturality concern this fixed map.  Its identity,
horizontal-pasting, and vertical-pasting equations must be proved by adjunction
and mate calculus.

Two stalk facts have disjoint types and disjoint uses:
\begin{enumerate}
\item \texttt{CAT-BA31-SHEAF-STALK-FAITHFULNESS}: if
      $u,v:\mathcal F\to\mathcal G$ are maps in the abelian category of
      $\Lambda_n$-sheaves and $u_{\bar y}=v_{\bar y}$ at every geometric
      point, then $u=v$.  It is consumed only in S5, after the maps of
      cohomology sheaves have already been constructed.
\item \texttt{CAT-BA31-DERIVED-STALK-ISO-CONSERVATIVITY}: if
      $u:A\to B$ is an already constructed morphism in
      $D^+(Y'_{\acute et},\Lambda_n)$ and every
      $H^i(u)_{\bar y}$ is invertible, then $u$ is invertible.  It is consumed
      only in S6.
\end{enumerate}
No stalk principle in this fragment detects equality of arbitrary derived
morphisms or constructs a derived inverse from pointwise inverses.

\paragraph{S0--S7 middle interfaces.}
The following are literal middle owners and all remain
\texttt{OPEN/INCOMPLETE}.
\begin{itemize}
\item \texttt{OBL-BA31-S0-GLOBAL-MATE} constructs \textup{(BA31-MATE)} and
  proves replacement independence, naturality, identity, and both pasting
  laws from the three named categorical atoms.
\item \texttt{OBL-BA31-S1-FINITE} proves exact finite pushforward and the
  finite geometric-fibre product formula, then identifies the latter with
  the stalk of the already constructed mate.
\item \texttt{OBL-BA31-S2-STABILITY} proves typed composition and
  proper-surjective cancellation diagrams for the actual $Rf_*$
  pseudofunctor; no inverse is supplied as input data.
\item \texttt{OBL-BA31-S3-PROJECTIVE-REDUCTION} states a proper finitely
  presented Chow datum and a finite-surjective retract/cancellation theorem
  reducing projective space to products of projective lines.
\item \texttt{OBL-BA31-S4-HENSELIAN-P1} assembles independently owned
  degree-zero, affine, and torsor/degree-one theorems into the strict-henselian
  projective-line comparison.
\item \texttt{OBL-BA31-S5-STALK-COMPARISON} combines the strict-henselization
  stalk formula and proper-curve field invariance and identifies their
  comparison with $H^i(\operatorname{BC}_{f,g,K})$.  Sheaf-stalk
  faithfulness is used only for this equality of maps of $H^i$-sheaves.
\item \texttt{OBL-BA31-S6-SHEAF-ASSEMBLY} applies only derived-stalk
  isomorphism conservativity to the fixed mate after S1--S5 have proved all
  its cohomology-stalk maps invertible.
\item \texttt{OBL-BA31-S7-BOUNDED-BELOW} uses the selected bounded-below
  Postnikov/Cartan--Eilenberg leaf, including convergence and identification
  of the filtered comparison with \textup{(BA31-MATE)}.
\end{itemize}
The root \texttt{OBL-BA31-FINITE-PBC} remains
\texttt{OPEN/INCOMPLETE} until these middle nodes and every listed atom are
actually proved.  The historical finite stalk-descent certificate is retired
and has no edge into this graph.

\section{Detailed S0--S10 proof text}\label{chap:s0-s10}
\statusnote{The following block preserves the mathematical S0--S10 body and every original \string\uses declaration.  The sole source-level correction is the ill-typed S0 sectionwise formula, replaced in place by the categorical mate and an explicit erratum.  The remaining propositions are conditional reassembly steps whenever their gap-ledger inputs are OPEN.}
% A3/S0--S4: explicit finite-torsion proper base change route.
% Stacks tags below are source locators only; every mathematical assertion is
% expanded or named as an OPEN obligation.

\subsection{Stacks route: comparison through the strictly henselian projective line}

\statusnote{This file is the S0--S4 core of the A3 DAG
\texttt{research/a3\_stacks\_proper\_base\_change\_dag.md}.  The coefficient
sheaf is a torsion sheaf on the small etale site.  A Stacks tag identifies the
source paragraph to be audited; it is never used as an admitted theorem.}

\subsubsection{S0. Comparison data and injective acyclicity}

\begin{definition}[Cartesian square and chosen pullback identification]
\label{def:weil-pbc-cartesian-square}
\uses{}
For schemes and morphisms
\[
\begin{CD}
 X' @>{g'}>> X\\
 @V{f'}VV       @VV{f}V\\
 Y' @>{g}>> Y,
\end{CD}
\]
we write $\square=(f,f',g,g',\alpha)$ when
$\alpha:X'\xrightarrow{\sim}X\times_Y Y'$ is a specified isomorphism and
$f\circ g'=g\circ f'$ as maps to $Y$.  The notation
$g'^{-1}F$ always means inverse image using this specified identification.
\end{definition}

\begin{definition}[Underived exchange map: corrected categorical construction]
\label{def:weil-pbc-underived-exchange}
\uses{def:weil-pbc-cartesian-square}
Let $F$ be a sheaf of abelian groups on $X_{\mathrm{et}}$.  The chosen
Cartesian identification gives a canonical isomorphism of inverse-image
functors
\[
 f'^{-1}g^{-1}\;\xrightarrow{\sim}\;g'^{-1}f^{-1}.
\]
Define $\operatorname{Ex}_{\square,F}$ to be the adjoint of the composite
\[
 f'^{-1}g^{-1}f_*F
 \xrightarrow{\sim}g'^{-1}f^{-1}f_*F
 \xrightarrow{g'^{-1}\varepsilon_f}g'^{-1}F.
\]
Thus
\[
 \operatorname{Ex}_{\square,F}:g^{-1}f_*F\longrightarrow
 f'_*g'^{-1}F
\]
is a globally constructed morphism of sheaves.  Naturality, identity, and
pasting are properties of this fixed mate; no family of stalk maps is used to
construct it.
\end{definition}

\begin{warning}[Erratum to the historical sectionwise presentation]
The frozen draft represented a local section by a section on
$V_i\times_{Y'}Y$.  The square only supplies $g:Y'\to Y$, not a map
$Y\to Y'$, so that fibre product is not defined.  The categorical
unit--Cartesian-identification--counit construction above is the replacement
used throughout this extraction.  The removed historical block has SHA-256
\texttt{E982FFC21CEAF8ABCD945A58DDFADBA4EB898F882D56A00A70BAC9F6365679A0} and is not a proof input.
\end{warning}

\begin{lemma}[Sectionwise exchange is natural]
\label{lemma:weil-pbc-underived-exchange-naturality}
\uses{def:weil-pbc-underived-exchange}
For every morphism $u:F\to G$ the square
\[
\begin{CD}
 g^{-1}f_*F @>{\operatorname{Ex}_{\square,F}}>> f'_*g'^{-1}F\\
 @V{g^{-1}f_*u}VV @VV{f'_*g'^{-1}u}V\\
 g^{-1}f_*G @>{\operatorname{Ex}_{\square,G}}>> f'_*g'^{-1}G
\end{CD}
\]
commutes on every etale object of $Y'$.
\end{lemma}
\begin{proof}
\uses{def:weil-pbc-underived-exchange}
On a representative section $s_i$, both paths first apply $u$ and then
pull back along the same two maps.  Since inverse image is a functor, the two
resulting sections of $g'^{-1}G$ are equal on each $V_i$ and therefore after
descent on $V$.  Equality for all $V$ is equality of sheaf morphisms.
\end{proof}

\begin{lemma}[Finite refinement of a surjective etale morphism]
\label{lemma:weil-pbc-finite-refinement-etale-cover}
\uses{}
Let $U\to T$ be surjective and etale, with $T$ quasi-compact and
quasi-separated.  There is a finite, finitely presented, surjective morphism
$r:T'\to T$ and a finite open covering $T'=\bigcup_aW_a$ such that every
$W_a\to T$ factors through $U\to T$.
\end{lemma}
\begin{proof}
\uses{} 
Choose finitely many affine opens $U_a\subset U$ whose images cover $T$ and
replace $U$ by their disjoint union.  The resulting map is quasi-compact,
separated, etale, and has a finite upper bound $d$ on the cardinalities of
its geometric fibres.  We induct on $d$.  If $d=1$, the map is radicial and
etale, hence an open immersion; surjectivity makes it an isomorphism.

For $d>1$, apply the affine form of Zariski's Main Theorem to the separated
quasi-finite map $U\to T$: adjoining to $\mathcal O_T$ the integral elements
of the quasi-coherent algebra $u_*\mathcal O_U$ gives an integral
$\pi:\widetilde T\to T$ and a quasi-compact open immersion
$j:U\hookrightarrow\widetilde T$.  Replace $\widetilde T$ by the
scheme-theoretic closure of $U$, so $U$ is schematically dense.  The diagonal
section gives an open-and-closed decomposition
\[
 U\times_T\widetilde T=U\amalg W.
\]
The image $V$ of $W\to\widetilde T$ is open and contains
$\widetilde T\setminus U$.  Over a point of the dense open $U$, the fibre of
$W\to V$ is the fibre of $U\to T$ with the diagonal point removed, so it has
fewer than $d$ points.  The number of points in a separated etale fibre
    cannot increase under specialization: after writing the map locally as a
    finite disjoint union of Jacobian etale algebra presentations, distinct idempotents at a
specialization lift to distinct idempotents at a generization.  Thus every
geometric fibre of $W\to V$ has cardinality $<d$.

By induction there is an integral surjection $Z\to V$, locally factoring
through $W$.  Apply the same integral-closure construction to
$Z\to\widetilde T$ and take the closure $Z'$ of $Z$; then
$Z'\times_{\widetilde T}V=Z$.  If $C$ is the reduced closed complement
$\widetilde T\setminus V$, then $C\subset U$, and
$Z'\amalg C\to T$ is integral and surjective.  Near a point of $Z$ it
factors through $W\to U$; near a point of $Z'\setminus Z$ it factors through
$C\subset U$.  Hence it locally factors through the original cover.

Finally write this integral morphism as the inverse limit of its finite
finitely presented $T$-subalgebras.  Choose finitely many affine opens on
which the factorization through $U$ is defined.  Each open, each map to $U$,
the equalities over $T$, and surjectivity descend to a common finite stage:
the first three assertions are finite-presentation data, while failure of
surjectivity has quasi-compact open image and would persist after pullback to
the limit.  That common stage is the required $T'$.  This is the construction
located at Tag \texttt{09Z0}, with the degree induction and final finite stage
made explicit.
\end{proof}

\begin{lemma}[Degree-zero comparison over a strictly henselian base]
\label{lemma:weil-pbc-proper-strict-henselian-h0}
\uses{}

Let $(A,\mathfrak m)$ be strictly henselian local, let
$h:T\to\operatorname{Spec}A$ be proper, and let $T_0$ be the closed
geometric fibre.  For every etale sheaf of sets $E$ on $T$, restriction is a
bijection
\[
 \Gamma(T,E)\xrightarrow{\sim}\Gamma(T_0,E|_{T_0}).
\]
The bijection is functorial in $E$ and in finite morphisms over $T$.
\end{lemma}
\begin{proof}
\uses{lemma:weil-pbc-finite-refinement-etale-cover}

We first prove the corresponding assertion for Zariski sheaves, uniformly
after every finite $T'\to T$.  The space $T'$ is spectral and $T'\to
\operatorname{Spec}A$ is proper.  For an irreducible closed reduced
$C\subset T'$, Stein factorization gives
\[
 C\longrightarrow\operatorname{Spec}B\longrightarrow
 \operatorname{Spec}A,
\]
where $A\to B$ is integral and the first map has geometrically connected
fibres.  The pair $(B,\mathfrak mB)$ is henselian: an idempotent of a finite
$B$-algebra modulo $\mathfrak m$ satisfies $e^2-e=0$, and the simple roots
$0,1$ lift uniquely by the henselian property of $A$ after restriction of
scalars along the integral algebra $B$.  Consequently the closed subset
$\operatorname{Spec}B\cap V(\mathfrak m)$ is connected and nonempty.  Its
inverse image $C\cap T'_0$ is connected because the map is closed,
surjective onto its image, and has connected fibres.

For completeness, let $G$ be a Zariski sheaf on $T'$.  Every point
specializes to $T'_0$, so a global section is determined by its restriction.
For existence, write $G$ as the filtered union of subsheaves generated by
finitely many sections on quasi-compact opens.  Such a finitely generated
subsheaf embeds into a finite product
\(\prod_b(i_b)_*\underline{M_b}\), where $i_b:C_b\hookrightarrow T'$ is
closed and $M_b$ is a finite set: take the finitely many values of the
generators and separate their equalizer loci.  Sections of
$(i_b)_*\underline{M_b}$ are locally constant maps $C_b\to M_b$.
Connectedness of every $C\cap T'_0$ says that the connected components of
$C_b$ and $C_b\cap T'_0$ correspond, so such maps extend uniquely.  Finite
products, intersections with the embedded subsheaf, and filtered colimits
preserve this equality of section sets.  Thus
\(\Gamma(T',G)=\Gamma(T'_0,G|)\) for every Zariski sheaf and every finite
$T'\to T$.

Now let $E$ be an etale sheaf.  The map from the restriction to the Zariski
site followed by closed-fibre inverse image to the Zariski restriction of
$E|_{T_0}$ is stalkwise injective.  Indeed, an elementary etale
neighbourhood $V\to T$ whose closed fibre has a local section has open image
$V_0\subset T$, and $V\to V_0$ is an etale cover; hence
\[
 E(V_0)\longrightarrow E(V)
 \rightrightarrows E(V\times_{V_0}V)
\]
is an equalizer.  Taking the colimit over such neighbourhoods proves the
stalkwise injection.  The Zariski result therefore proves injectivity of
$\Gamma(T,E)\to\Gamma(T_0,E|)$, and the same argument applies after every
finite $T'\to T$.

For surjectivity, take $t_0\in\Gamma(T_0,E|)$.  By the definition of inverse
image and sheafification there are finitely many etale maps $V_a\to T_0$
covering $T_0$, etale $U_a\to T$, maps $V_a\to U_a\times_TT_0$, and
sections $s_a\in E(U_a)$ representing $t_0|_{V_a}$.  Shrink $U_a$ to the
open image of $V_a$ in $U_a\times_TT_0$ and then in $U_a$; the sheaf
equalizer shows that $s_a|_{U_a\times_TT_0}=t_0|_{U_a\times_TT_0}$.  Every
nonempty closed subset of $T$ meets $T_0$ by the connectedness argument
above, so $\coprod_aU_a\to T$ is surjective.

Choose the finite surjective refinement $r:T'\to T$ of the preceding lemma
and opens $W_a\subset T'$ mapping to $U_a$.  Pulling back $s_a$ gives local
sections on $W_a$.  On $W_a\cap W_b$ their closed-fibre restrictions both
equal $r_0^{-1}t_0$; the already proved injectivity for the finite scheme
$T'\to T$ makes them equal.  They glue to $s'\in\Gamma(T',r^{-1}E)$.
The two pullbacks of $s'$ to $T'\times_TT'$ have the same closed-fibre
restriction, and injectivity for the finite morphism
$T'\times_TT'\to T$ makes them equal.

It remains to spell out descent along $r$, which need not be etale.  At a
geometric point $\bar x\to T$, choose a geometric $\bar x'\to T'$ above it.
The germ of $s'$ at $\bar x'$ is an element of the stalk $E_{\bar x}$.
The equalizer equality on $T'\times_TT'$ makes this element independent of
$\bar x'$.  Represent it on an etale neighbourhood of $\bar x$ in $T$.
After shrinking, its pullback equals $s'$ near every point over $\bar x$;
finiteness of $r$ makes that inverse image quasi-compact, so one shrinking
works for all of them.  These local descended sections agree on overlaps
because equality can be checked at geometric stalks, and hence glue to
$s\in\Gamma(T,E)$.  Its restriction is $t_0$ by injectivity on $T_0$.
All pullbacks used are the functorial sheafification units, so the argument
also proves functoriality in $E$ and compatibility with finite pullback.  The
source locators are Tags \texttt{09ZG} and \texttt{09ZF}; neither is used as
an unexpanded conclusion.
\end{proof}

\begin{lemma}[Underived proper exchange is an isomorphism]
\label{lemma:weil-pbc-proper-underived-exchange}
\uses{def:weil-pbc-underived-exchange}
If $f$ is proper, then
$\operatorname{Ex}_{\square,F}:g^{-1}f_*F\to f'_*g'^{-1}F$ is an
isomorphism for every sheaf $F$.
\end{lemma}
\begin{proof}
\uses{lemma:weil-pbc-underived-exchange-naturality,lemma:weil-pbc-proper-strict-henselian-h0}

Fix a geometric point $\bar y'\to Y'$ and let $\bar y\to Y$ be its composite
with $g$.  Let $\mathcal N_{\bar y}^{\mathrm{aff}}$ be the category of
pointed affine etale neighbourhoods $(V,\bar v)$ of $\bar y$.  It is
cofiltered: an affine open containing the lifted point in a fibre product
refines two objects, and an affine neighbourhood of the equalizer of two
maps makes parallel maps equal.  It is cofinal among all pointed etale
neighbourhoods because every lifted point has an affine open neighbourhood.
Thus
\[
 (f_*F)_{\bar y}=
 \varinjlim_{(V,\bar v)}\Gamma(X\times_YV,F|).
\]

Writing $V=\operatorname{Spec}A_V$, the strict henselization is
$Y_{(\bar y)}=\operatorname{Spec}(\varinjlim A_V)$ and
$X_{(\bar y)}=X\times_YY_{(\bar y)}=\varprojlim(X\times_YV)$.  The natural
map from the displayed colimit to $\Gamma(X_{(\bar y)},F|)$ is bijective.
For surjectivity, represent a section on the limit by sections on a finite
cover by quasi-compact etale objects.  Each object, each covering map, and
each section representative descends to a stage because etale morphisms are
locally of finite presentation.  The finitely many equality conditions on
the pairwise fibre products descend after enlarging that stage.  The sheaf
equalizer then glues the descended representatives.  For injectivity,
represent two sections at one stage; equality on the limit is witnessed on
a finite etale cover, and the cover together with its finitely many equality
witnesses descends to a later stage.  Hence the two representatives agree in
the colimit.

The base $Y_{(\bar y)}$ is strictly henselian local and
$X_{(\bar y)}$ is proper over it.  The preceding lemma identifies its
sections with those on the closed geometric fibre:
\[
 (f_*F)_{\bar y}\xrightarrow{\sim}
 \Gamma(X_{(\bar y)},F|)\xrightarrow{\sim}
 \Gamma(X_{\bar y},F|_{X_{\bar y}}).
\]
The identical construction for $f'$ gives the target stalk.  Since $\bar y$
is the composite of the same $\operatorname{Spec}\Omega\to Y'$ with $g$,
the Cartesian isomorphism identifies $X_{\bar y}$ and $X'_{\bar y'}$ with
$X\times_Y\operatorname{Spec}\Omega$.

Finally, a germ represented by $s\in\Gamma(X\times_YV,F|)$ travels around
the comparison by pulling $s$ first to $X'\times_{Y'}(V\times_YY')$ and then
to the geometric fibre.  The two stalk identifications above pull $s$
directly to that same fibre.  Hence the fibre bijection is the stalk of
$\operatorname{Ex}_{\square,F}$, not an unnamed isomorphism.  Geometric
stalks are conservative, so the exchange map is an isomorphism.  This closes
the compatibility omission at Tag \texttt{0A3U}.
\end{proof}

\begin{definition}[Derived exchange map]
\label{def:weil-pbc-derived-exchange}
\uses{def:weil-pbc-underived-exchange}
For $K\in D^+(X_{\mathrm{et}})$ choose an injective resolution
$K\xrightarrow{\simeq}I^\bullet$.  Apply
$\operatorname{Ex}_{\square,I^r}$ in every degree and use the naturality
lemma to obtain a morphism of complexes
\[ 
 g^{-1}f_*I^\bullet\longrightarrow f'_*g'^{-1}I^\bullet.
\]
Choose a K-injective replacement
$g'^{-1}I^\bullet\xrightarrow{\simeq}J^\bullet$ in the Grothendieck
category of etale abelian sheaves on $X'$.  Compose the displayed map with
$f'_*g'^{-1}I^\bullet\to f'_*J^\bullet$ and define
\[
 \operatorname{Ex}_{\square,K}:g^{-1}f_*I^\bullet\longrightarrow f'_*J^\bullet
\]
to be its class in $D^+(Y'_{\mathrm{et}})$.  The target complex is a model
for $Rf'_*g'^{-1}K$, since $J^\bullet$ is K-injective; the source is a model
for $g^{-1}Rf_*K$ because $g^{-1}$ is exact.
Its class in $D^+(Y'_{\mathrm{et}})$ is
\[
 \operatorname{Ex}_{\square,K}:g^{-1}Rf_*K\longrightarrow
 Rf'_*g'^{-1}K.
\]
\end{definition}


% THROWAWAY EARLY-INSERT CANDIDATE.  Insert immediately after
% def:weil-pbc-derived-exchange and before
% lemma:weil-sf-well-defined-def-weil-pbc-derived-exchange.
\section*{Derived Exchange Comparison Chain}

\begin{definition}[Derived-exchange presentation]
\label{def:weil-pbc-derived-exchange-presentation-v5}
\uses{def:weil-pbc-derived-exchange}
For a bounded-below complex $K$, a presentation consists of an injective
resolution $u:K\to I$, a K-injective complex $J$, a quasi-isomorphism
$q:g'^{-1}I\to J$, and the complex map
\[
 E_I:g^{-1}f_*I\longrightarrow f'_*g'^{-1}I
\]
obtained degreewise from the underived exchange.  Write
$E_{I,J}:=f'_*q\,E_I$.
\end{definition}

\begin{definition}[Hom complex and K-injective test]
\label{def:weil-pbc-derived-exchange-hom-test-v5}
\uses{def:weil-pbc-derived-exchange-presentation-v5}
For complexes $A,B$, let $\underline{\operatorname{Hom}}^r(A,B)$ be the
degree-$r$ graded maps, with differential $d(h)=d_Bh-(-1)^rh d_A$.
A complex $B$ is K-injective when this Hom complex is acyclic for every
acyclic $A$.
\end{definition}

\begin{lemma}[The Hom differential squares to zero]
\label{lemma:weil-pbc-derived-exchange-hom-differential-v5}
\uses{def:weil-pbc-derived-exchange-hom-test-v5}
The displayed Hom differential satisfies $d^2=0$.
\end{lemma}
\begin{proof}
\uses{def:weil-pbc-derived-exchange-hom-test-v5}
For homogeneous $h$ of degree $r$, expanding $d(dh)$ gives
$d_B^2h-(-1)^rd_Bhd_A-(-1)^{r+1}d_Bhd_A-hd_A^2=0$.
\end{proof}

\begin{lemma}[Bounded-below injectives are K-injective]
\label{lemma:weil-pbc-derived-exchange-bounded-injective-kinjective-v5}
\uses{def:weil-pbc-derived-exchange-hom-test-v5}
Every bounded-below complex whose terms are injective is K-injective.
\end{lemma}
\begin{proof}
\uses{def:weil-pbc-derived-exchange-hom-test-v5,lemma:weil-pbc-derived-exchange-hom-differential-v5}
After shifting the target, it suffices to treat a degree-zero Hom cocycle
$f:C\to I$.  Choose $m$ with $I^n=0$ for $n<m$, and construct maps
$s^n:C^n\to I^{n-1}$ upward from $n=m$ so that
\[
 f^n=d s^n+s^{n+1}d. \tag{*}
\]
Assume $s^m,\ldots,s^n$ are known (with $s^m=0$ at the initial step).
On $B^{n+1}=\operatorname{im}(d:C^n\to C^{n+1})=ker d$, define
\[
 t(d x):=f^n(x)-d s^n(x)\in I^n.
\]
This is well-defined: if $dx=0$, acyclicity gives $x=dy$, and the cocycle
equation together with (*) in degree $n-1$ gives
$f^n(x)=df^{n-1}(y)=d s^n(dy)=d s^n(x)$.  Injectivity extends the map
$t:B^{n+1}\to I^n$ along the monomorphism $B^{n+1}\hookrightarrow C^{n+1}$
to $s^{n+1}:C^{n+1}\to I^n$.  Its defining equality is (*) in degree $n$.
Thus $f=ds$, so every Hom cocycle is a coboundary.
\end{proof}

\begin{definition}[Target-replacement comparison problem]
\label{def:weil-pbc-derived-exchange-target-comparison-v5}
\uses{def:weil-pbc-derived-exchange-presentation-v5}
For two target replacements $q:A\to J$ and $q':A\to J'$ with $J,J'$
K-injective, a comparison is a chain map $c:J\to J'$ together with a
homotopy $cq\simeq q'$.
\end{definition}

\begin{lemma}[Target-replacement comparisons exist]
\label{lemma:weil-pbc-derived-exchange-target-comparison-exists-v5}
\uses{def:weil-pbc-derived-exchange-target-comparison-v5,def:weil-pbc-derived-exchange-hom-test-v5}
Every target-replacement comparison problem has a comparison.
\end{lemma}
\begin{proof}
\uses{def:weil-pbc-derived-exchange-target-comparison-v5,def:weil-pbc-derived-exchange-hom-test-v5}
The cone $C(q)$ is acyclic because $q$ is a quasi-isomorphism.  K-injectivity
of $J'$ makes $\underline{\operatorname{Hom}}(C(q),J')$ acyclic.  The
cone triangle gives the exact segment
\[
 H^0\underline{\operatorname{Hom}}(J,J')\longrightarrow
 H^0\underline{\operatorname{Hom}}(A,J')\longrightarrow
 H^1\underline{\operatorname{Hom}}(C(q),J')=0.
\]
Thus $[q']$ lifts to $[c]$; equality $[cq]=[q']$ is exactly a homotopy
$cq\simeq q'$.
\end{proof}

\begin{lemma}[Target comparisons are unique up to homotopy]
\label{lemma:weil-pbc-derived-exchange-target-comparison-unique-v5}
\uses{def:weil-pbc-derived-exchange-target-comparison-v5,def:weil-pbc-derived-exchange-hom-test-v5}
Two comparisons $c,d:J\to J'$ are homotopic.  Comparisons in both directions
are inverse to one another in the homotopy category.
\end{lemma}
\begin{proof}
\uses{def:weil-pbc-derived-exchange-target-comparison-v5,def:weil-pbc-derived-exchange-hom-test-v5}
The difference $(c-d)q$ is null-homotopic.  Since
$\underline{\operatorname{Hom}}(C(q),J')$ is acyclic, the cone exact segment
\[
 H^0\underline{\operatorname{Hom}}(C(q),J')=0\longrightarrow
 H^0\underline{\operatorname{Hom}}(J,J')\xrightarrow{q^*}
 H^0\underline{\operatorname{Hom}}(A,J')
\]
makes $q^*$ injective; hence $c-d$ is null-homotopic.  If
$c:J\to J'$ and $c':J'\to J$ are comparisons, both $c'c$ and $1_J$ compare
$q$ with itself, and both $cc'$ and $1_{J'}$ compare $q'$ with itself.
The first assertion therefore gives $c'c\simeq1_J$ and $cc'\simeq1_{J'}$.
\end{proof}

\begin{definition}[Source-resolution comparison problem]
\label{def:weil-pbc-derived-exchange-source-comparison-v5}
\uses{def:weil-pbc-derived-exchange-presentation-v5}
For injective resolutions $u:K\to I$ and $u':K\to I'$, a comparison is a
chain map $a:I\to I'$ with $au\simeq u'$.
\end{definition}

\begin{lemma}[Source-resolution comparisons exist and invert]
\label{lemma:weil-pbc-derived-exchange-source-comparison-v5}
\uses{def:weil-pbc-derived-exchange-source-comparison-v5,def:weil-pbc-derived-exchange-hom-test-v5}
There are comparisons $a:I\to I'$ and $b:I'\to I$, unique up to homotopy,
with $ba\simeq1_I$ and $ab\simeq1_{I'}$.
\end{lemma}
\begin{proof}
\uses{def:weil-pbc-derived-exchange-source-comparison-v5,
  lemma:weil-pbc-derived-exchange-bounded-injective-kinjective-v5,
  lemma:weil-pbc-derived-exchange-target-comparison-exists-v5,
  lemma:weil-pbc-derived-exchange-target-comparison-unique-v5}
The resolutions are bounded below, so the preceding lemma makes $I'$ and $I$
K-injective.  Apply
Lemma~\ref{lemma:weil-pbc-derived-exchange-target-comparison-exists-v5}
with $A=K$, $q=u$, $q'=u'$, $J=I$, and $J'=I'$ to obtain $a$.
Interchanging the two resolutions gives $b$.  Now apply
Lemma~\ref{lemma:weil-pbc-derived-exchange-target-comparison-unique-v5}:
$ba$ and $1_I$ are two comparisons of $u$ with itself, while $ab$ and
$1_{I'}$ compare $u'$ with itself.  It gives uniqueness up to homotopy and
the two displayed inverse equations.
\end{proof}

\begin{lemma}[Exchange naturality is a homotopy square]
\label{lemma:weil-pbc-derived-exchange-naturality-square-v5}
\uses{def:weil-pbc-derived-exchange-presentation-v5,def:weil-pbc-derived-exchange-target-comparison-v5}
Let $(I,J,q)$ and $(I',J',q')$ be presentations and let $a:I\to I'$ be a
source comparison.  A target comparison $c:J\to J'$ for $q'g'^{-1}a$ and
$cq$ gives
\[
 f'_*c\,E_{I,J}\simeq E_{I',J'}\,g^{-1}f_*a.
\]
\end{lemma}
\begin{proof}
\uses{def:weil-pbc-derived-exchange-presentation-v5,lemma:weil-pbc-underived-exchange-naturality,lemma:weil-pbc-derived-exchange-target-comparison-exists-v5}
Underived naturality gives the strict degreewise equality
$f'_*g'^{-1}a\,E_I=E_{I'}g^{-1}f_*a$.  The defining homotopy
$cq\simeq q'g'^{-1}a$ remains a homotopy after $f'_*$; compose it with
$E_I$ and substitute the strict equality.
\end{proof}

\begin{definition}[Identity and pasted exchange presentations]
\label{def:weil-pbc-derived-exchange-pasting-presentations-v5}
\uses{def:weil-pbc-derived-exchange-presentation-v5}
For an identity or a horizontally or vertically pasted Cartesian square, use
one injective resolution and form the degreewise exchange maps by their
unit--pullback-identification--counit formula before selecting targets.
\end{definition}

\begin{lemma}[Identity exchange equation]
\label{lemma:weil-pbc-derived-exchange-identity-v5}
\uses{def:weil-pbc-derived-exchange-pasting-presentations-v5}
The exchange for an identity square is the identity in the derived category.
\end{lemma}
\begin{proof}
\uses{def:weil-pbc-derived-exchange-pasting-presentations-v5,lemma:weil-pbc-derived-exchange-target-comparison-unique-v5}
The unit and counit for the identity pullback are identities on every sheaf
term.  Thus the degreewise exchange is $1$; changing targets leaves its class
unchanged by the target-comparison uniqueness argument above.
\end{proof}

\begin{lemma}[Horizontal and vertical exchange pasting]
\label{lemma:weil-pbc-derived-exchange-pasting-v5}
\uses{def:weil-pbc-derived-exchange-pasting-presentations-v5,def:weil-pbc-underived-exchange,def:weil-pbc-cartesian-square}
The exchange for a horizontal or vertical pasted square equals the composite
of the exchanges for its two constituent squares in the derived category.
\end{lemma}
\begin{proof}
\uses{def:weil-pbc-derived-exchange-pasting-presentations-v5,def:weil-pbc-underived-exchange,def:weil-pbc-cartesian-square,lemma:weil-pbc-derived-exchange-target-comparison-unique-v5}
Before target replacement, expand the three degreewise maps as unit,
pullback identification, and counit.  Associativity of fibre products makes
the two middle identifications equal, and the adjunction triangle identities
cancel the adjacent unit--counit pairs.  The remaining maps are equal as
complex maps.  Any different target replacements represent this same class by
target-comparison uniqueness.
\end{proof}

\begin{lemma}[Well-definedness certificate for def:weil-pbc-derived-exchange]
\label{lemma:weil-sf-well-defined-def-weil-pbc-derived-exchange}
\uses{def:weil-pbc-derived-exchange}
The construction in Definition \ref{def:weil-pbc-derived-exchange} is a typed raw datum.  Its theorem-valued naturality, comparison, descent, or coherence fields are isolated in this property lemma.
\end{lemma}
\begin{proof}
\uses{def:weil-pbc-derived-exchange-presentation-v5,
  lemma:weil-pbc-derived-exchange-target-comparison-unique-v5,
  lemma:weil-pbc-derived-exchange-source-comparison-v5,
  lemma:weil-pbc-derived-exchange-naturality-square-v5,
  lemma:weil-pbc-derived-exchange-identity-v5,
  lemma:weil-pbc-derived-exchange-pasting-v5}
The presentation definition constructs the complex map.  Source comparisons,
the naturality homotopy square, and target-comparison uniqueness identify all
selected presentations in the derived category.  The identity and the two
pasting equations are supplied by the final two lemmas.  Hence the derived
exchange construction is well-defined with all fields stated in the lemma.
\end{proof}
\begin{lemma}[Derived exchange replacement independence]
\label{lemma:weil-pbc-derived-exchange-replacement}
\uses{def:weil-pbc-derived-exchange}
The K-injective replacement used in Definition
\ref{def:weil-pbc-derived-exchange} gives a well-defined derived exchange
map, independent of all replacement choices.
\end{lemma}
\begin{proof}
\uses{} 
Given two K-injective replacements $J^\bullet$ and $J'^\bullet$, the identity
of $g'^{-1}I^\bullet$ lifts to maps $J^\bullet\to J'^\bullet$ and
$J'^\bullet\to J^\bullet$ because K-injective complexes are homotopically
injective.  Their composites are homotopic to the identities by lifting the
homotopies of the replacement quasi-isomorphisms.  Applying $f'_*$ preserves
chain homotopies.  The two resulting maps from
$g^{-1}f_*I^\bullet$ therefore define the same morphism in the derived
category.  A comparison of two injective resolutions of $K$ is handled by
the same lifting argument, and naturality of the underived exchange map
commutes with the comparison maps.  Thus the class is canonical.
\end{proof}

\begin{lemma}[Resolution independence, identities, and pasting]
\label{lemma:weil-pbc-derived-exchange-coherence}
\uses{def:weil-pbc-derived-exchange}
The class $\operatorname{Ex}_{\square,K}$ is independent of the injective
resolution and is compatible with identity squares, horizontal pasting, and
vertical pasting.
\end{lemma}
\begin{proof}
\uses{def:weil-pbc-derived-exchange}
Given two injective resolutions, the comparison theorem supplies chain maps
each way, unique up to chain homotopy, whose composites are homotopic to the
identities.  Naturality of the degreewise exchange map sends these homotopies
to homotopies after both direct images; localization therefore gives the same
derived morphism.  For the identity square, the unit followed by the counit is
the identity sectionwise, hence degreewise.  For pasted squares, expand each
exchange as ``unit, pullback identification, counit''; the two middle
identifications are equal by associativity of fibre products, and the two
adjunction triangle identities cancel the adjacent unit--counit pairs.  The
same calculation with the order reversed proves vertical pasting.
\end{proof}

\begin{definition}[Proper base-change property]
\label{def:weil-pbc-property}
\uses{def:weil-pbc-derived-exchange}
For a proper morphism $f:X\to Y$, define $\mathsf{PBC}(f)$ to mean that for
every morphism $g:Y'\to Y$, every associated Cartesian square $\square$, and
every torsion sheaf $F$ on $X_{\mathrm{et}}$, the map
$\operatorname{Ex}_{\square,F}$ is an isomorphism in the derived category.
The assertion includes the specified map of Definition
\ref{def:weil-pbc-derived-exchange}, not an unnamed isomorphism.
\end{definition}

\begin{definition}[Injective-acyclicity property]
\label{def:weil-pbc-injective-acyclicity-property}
\uses{def:weil-pbc-cartesian-square}
For a proper $f:X\to Y$, define $\mathsf{IA}(f)$ to mean: for every prime
$\ell$, every $g:Y'\to Y$, every square $\square$, and every injective sheaf
$I$ of $\mathbf F_\ell$-modules on $X_{\mathrm{et}}$, the sheaf
$g'^{-1}I$ is $f'_* $-acyclic, i.e.
$R^qf'_*g'^{-1}I=0$ for all $q>0$.
\end{definition}

\begin{lemma}[Filtered torsion union]
\label{lemma:weil-pbc-torsion-filtered-union}
\uses{}
For every torsion sheaf $F$ the inclusions $F[n]\hookrightarrow F[mn]$ form a
filtered system and the canonical map $\varinjlim_nF[n]\to F$ is an
isomorphism.
\end{lemma}
\begin{proof}
\uses{} 
The transition map sends a section killed by $n$ to the same section, now
regarded as killed by $mn$.  The map to $F$ is monic because a section that is
zero in $F$ is already zero in every $F[n]$.  It is epic because every local
section of $F$ is killed by some positive integer and hence is locally in one
of the $F[n]$; sheafification preserves this local epimorphism.
\end{proof}

\begin{lemma}[Filtered colimits and higher direct image]
\label{lemma:weil-pbc-direct-image-filtered-colimit}
\uses{}
For a quasi-compact and quasi-separated morphism of schemes $h:W\to V$ and a
filtered system of torsion sheaves $(F_i)$, the canonical map
\[
 \varinjlim_i R^qh_*F_i\longrightarrow R^qh_*(\varinjlim_iF_i)
\]
is an isomorphism for every $q$.
\end{lemma}
\begin{proof}
\uses{} 
For $q=0$, sections on a quasi-compact etale object are the equalizer of a
finite covering diagram after a finite quasi-compact refinement; filtered
colimits commute with this finite limit.  For $q>0$, represent a class by a
degree-$q$ cocycle on a split etale hypercover with finite $(q+1)$-skeleton.
Each of its finitely many sections is represented in some $F_i$, and the
finite list of face-map equalities holds at one later index because the system
is filtered.  This proves surjectivity of
$\varinjlim_iH^q(W,F_i)\to H^q(W,\varinjlim_iF_i)$.

If a class from $F_i$ dies after the colimit, choose a finite hypercover
refinement and a degree-$(q-1)$ null-homotopy on the limit.  Its finitely many
sections, refinement maps, and equations descend to one index $j\geq i$;
there the class is already a coboundary.  Thus the map is injective.  The
same argument applies after every quasi-compact etale base change
$U\to V$, so the cohomology-sheaf description of $R^qh_*$ identifies the
displayed sheaves.  This is the derived filtered-colimit argument, with no
unbounded injective replacement required.
\end{proof}

\begin{lemma}[Prime-factor induction]
\label{lemma:weil-pbc-prime-factor-induction}
\uses{}
If a statement about torsion sheaves is stable under short exact sequences,
then its verification for all $F[n]$ reduces by induction on the total number
of prime factors of $n$ to its verification for sheaves killed by one prime.
\end{lemma}
\begin{proof}
\uses{} 
Choose a prime $\ell\mid n$ and use the short exact sequence
$0\to F[\ell]\to F[n]\xrightarrow{\ell}\ell F[n]\to0$.
The kernel is killed by $\ell$, while $\ell F[n]$ is killed by $n/\ell$ and
therefore has fewer prime factors in its exponent.  Let
$C(F)=\operatorname{Cone}(\operatorname{Ex}_{\square,F})$ and use the
cohomological cone convention
$d_{C}(y,x)=(d_Yy+u(x),-d_Xx)$.  Applying the two derived functors to the
short exact sequence gives a distinguished triangle of exchange maps and,
by the octahedral axiom, a triangle
\[
 C(F[\ell])\longrightarrow C(F[n])\longrightarrow
 C(\ell F[n])\xrightarrow{+1}C(F[\ell])[1].
\]
The connecting map is the class of $(y,x)\mapsto y$ in the above cone
convention; the sign in the last arrow is forced by
$d_C^2=0$.  Induct on
$\nu(n)=\sum_\ell v_\ell(n)$.  If the two end cones are acyclic, the long
exact cohomology sequence of this triangle makes the middle cone acyclic.
The base $\nu(n)=1$ is the prime-torsion input.  This displays the invariant
and all connecting morphisms.
\end{proof}

\begin{lemma}[Five-lemma transfer for exchange]
\label{lemma:weil-pbc-five-lemma-transfer}
\uses{def:weil-pbc-derived-exchange}
In a short exact sequence of torsion sheaves, if two of the three derived
exchange maps are isomorphisms, then the third is an isomorphism.
\end{lemma}
\begin{proof}
\uses{lemma:weil-pbc-prime-factor-induction}

Apply $g^{-1}Rf_*$ and $Rf'_*g'^{-1}$ to the short exact sequence and write
the natural long exact cohomology sequences.  Naturality of the exchange map
gives a morphism between these sequences.  In each degree the five-lemma
applies to the five adjacent terms, with the boundary maps identified by the
connecting-homomorphism construction from an injective resolution.  Thus the
middle exchange map is an isomorphism in every degree.
\end{proof}

\begin{lemma}[Leray acyclicity criterion]
\label{lemma:weil-pbc-leray-acyclicity}
\uses{def:weil-pbc-injective-acyclicity-property}
If a complex $I^\bullet$ consists of $f'_*$-acyclic sheaves, then the natural
map $\mathcal H^q(f'_*I^\bullet)\to R^qf'_*K$ is an isomorphism whenever
$K\to I^\bullet$ is a quasi-isomorphism.
\end{lemma}
\begin{proof}
\uses{} 
Resolve $K$ by a bounded-below Cartan--Eilenberg double complex
$J^{r,s}$ with vertical differential $d_v:J^{r,s}\to J^{r+1,s}$ resolving
$I^s$ and horizontal differential $d_h:J^{r,s}\to J^{r,s+1}$ lifting the
differential of $I^\bullet$.  Set
$d_{\mathrm{tot}}=d_v+(-1)^r d_h$.  Applying $f'_*$ and filtering the total
complex by $r$ gives
\[
 E_1^{r,s}=R^rf'_*I^s,\qquad
 E_2^{r,s}=R^rf'_*\mathcal H^s(K)\Longrightarrow
 \mathcal H^{r+s}(Rf'_*K).
\]
The assumed acyclicity makes $E_1^{r,s}=0$ for $r>0$, so the edge map from
the $r=0$ column is the cohomology map of the ordinary complex
$f'_*I^\bullet$.  The sign $(-1)^r$ makes the square with the horizontal
differential commute, and the edge morphism is therefore exactly the natural
map $\mathcal H^q(f'_*I^\bullet)\to R^qf'_*K$.  Collapse of the first column
proves it is an isomorphism.
\end{proof}

\begin{lemma}[Prime-torsion devissage]
\label{lemma:weil-pbc-prime-torsion-devissage}
\uses{}

For a torsion sheaf $F$, write $F[n]=\ker(n:F\to F)$.  Then
$F=\varinjlim_nF[n]$; for $m\mid n$ the inclusion $F[m]\subseteq F[n]$
is monic, and if $n=\prod_i\ell_i^{e_i}$, repeated use of
\[
 0\longrightarrow F[\ell_i]\longrightarrow F[n]
 \xrightarrow{\ell_i}\ell_iF[n]\longrightarrow0
\]
reduces assertions about $F[n]$ to assertions about $\ell_i$-torsion.
\end{lemma}
\begin{proof}
\uses{lemma:weil-pbc-torsion-filtered-union,lemma:weil-pbc-direct-image-filtered-colimit,lemma:weil-pbc-prime-factor-induction,lemma:weil-pbc-five-lemma-transfer}

The filtered union is Lemma~\ref{lemma:weil-pbc-torsion-filtered-union}.
The exact sequence is the kernel--image sequence for multiplication by
$\ell_i$, since the kernel in $F[n]$ is $F[\ell_i]$.  Lemma~\ref{lemma:weil-pbc-prime-factor-induction} and the
five-lemma transfer then prove the assertion for every $F[n]$, and
Lemma~\ref{lemma:weil-pbc-direct-image-filtered-colimit} passes to $F$.
\end{proof}

\begin{lemma}[Equivalence of PBC and IA]
\label{lemma:weil-pbc-property-iff-acyclicity}
\uses{def:weil-pbc-property,def:weil-pbc-injective-acyclicity-property}
For proper $f$, $\mathsf{PBC}(f)$ holds if and only if $\mathsf{IA}(f)$ holds.
\end{lemma}
\begin{proof}
\uses{def:weil-pbc-property,def:weil-pbc-injective-acyclicity-property,lemma:weil-pbc-prime-torsion-devissage,lemma:weil-pbc-proper-underived-exchange,lemma:weil-pbc-leray-acyclicity}
Assume $\mathsf{PBC}(f)$.  For an injective $\mathbf F_\ell$-module sheaf
$I$ and a base change $g$, apply
the isomorphism to the complex concentrated in degree zero.  The source has
no higher cohomology because $I$ is injective, so the target has none; this is
the required acyclicity.  Conversely assume $\mathsf{IA}(f)$.  Resolve any
torsion $F$ by injectives $I^\bullet$.  The inverse image complex
$g'^{-1}I^\bullet$ computes $Rf'_*g'^{-1}F$ because each term is acyclic.
For an $\mathbf F_\ell$-module sheaf $E$, choose an injective resolution
$E\to I^\bullet$ in the $\mathbf F_\ell$-module category.  IA makes every
$g'^{-1}I^r$ $f'_*$-acyclic, so Lemma
\ref{lemma:weil-pbc-leray-acyclicity} identifies the target of the derived
exchange map with the cohomology of $f'_*g'^{-1}I^\bullet$.  In each degree,
the underived proper exchange lemma gives an isomorphism
\[
 g^{-1}f_*I^r\xrightarrow{\sim}f'_*g'^{-1}I^r,
\]
and the replacement-independence lemma makes these maps a morphism of the
chosen resolutions.  Thus $\operatorname{Ex}_{\square,E}$ is an isomorphism.
For $F[n]$, the prime-factor induction lemma and the five-lemma transfer
propagate this isomorphism through the exact sequences
$0\to F[\ell]\to F[n]\xrightarrow{\ell}\ell F[n]\to0$; the induction
parameter is $\nu(n)=\sum_\ell v_\ell(n)$, with the prime-torsion case just
proved.  Finally $F=\varinjlim_nF[n]$ by the filtered torsion-union lemma.
The filtered-colimit direct-image lemma applies after every quasi-compact
etale base change, so it identifies both cohomology sheaves and carries the
isomorphisms for the $F[n]$ to the specified map $\operatorname{Ex}_{\square,F}$.
This proves PBC for every torsion sheaf and closes the converse implication.
\end{proof}

\subsubsection{S1. Finite morphisms}

\begin{lemma}[Geometric stalk formula for finite direct image]
\label{lemma:weil-pbc-finite-stalk-formula}
\uses{def:weil-pbc-cartesian-square}
Let $h:W\to V$ be finite and $\bar v\to V$ geometric.  For every sheaf $F$
on $W_{\mathrm{et}}$ there is a canonical bijection
\[
 (h_*F)_{\bar v}\;\cong\;
 \prod_{\bar w\,|\,\bar v}F_{\bar w},
\]
where the product is finite and the index runs over geometric points of
$W$ above $\bar v$ up to the natural equivalence.  For an abelian sheaf the
same object is canonically the finite direct sum.
\end{lemma}
\begin{proof}
\uses{} 
The strict henselian local scheme $V_{(\bar v)}$ has the property that a
finite $V$-scheme becomes a finite disjoint union of finite local schemes
after base change to $V_{(\bar v)}$.  Sections of $F$ over that disjoint union
are products of the sections over the components.  Passing through the
cofiltered system of etale neighbourhoods of $\bar v$ identifies the stalk
with the displayed product.  Finiteness of the fibre makes product and sum
equal in the abelian category.  This is the stalk calculation located at
Stacks Proposition 59.55.2, Tag \texttt{03QP}; no exactness is inferred until
the next lemma.
\end{proof}

\begin{lemma}[Exactness of finite direct image]
\label{lemma:weil-pbc-finite-direct-image-exact}
\uses{}

If $h$ is finite, then $h_*$ on etale abelian sheaves is exact.
\end{lemma}
\begin{proof}
\uses{lemma:weil-pbc-finite-stalk-formula}

For a short exact sequence $0\to F_1\to F_2\to F_3\to0$, apply the stalk
formula at every geometric $\bar v$.  A finite product of short exact
sequences of abelian groups is exact, so the sequence of stalks of the direct
images is exact.  Geometric stalks are conservative for etale sheaves:
injectivity and surjectivity can be checked after an etale neighbourhood of
each geometric point and then sheafified.  Hence the direct-image sequence is
exact.  Consequently $R^qh_*=0$ for $q>0$.
\end{proof}

\begin{lemma}[Finite exchange isomorphism]
\label{lemma:weil-pbc-finite-exchange}
\uses{def:weil-pbc-cartesian-square}
For a finite Cartesian square, the underived exchange map is an isomorphism;
therefore all its derived exchange maps are isomorphisms.
\end{lemma}
\begin{proof}
\uses{lemma:weil-pbc-finite-stalk-formula,lemma:weil-pbc-finite-direct-image-exact}

At a geometric point $\bar y'\to Y'$, the fibre of $f'$ is the base change of
the finite fibre of $f$ along $\kappa(\bar y)\to\kappa(\bar y')$.  The stalk
formula identifies both sides with the same finite product indexed by these
points, and the exchange map is the identity on each factor.  Thus it is a
stalkwise bijection and hence an isomorphism.  Exactness from the preceding
lemma means injective resolutions remain direct-image acyclic, so the same
degreewise argument computes the derived map.
\end{proof}

\begin{proposition}[Finite morphisms satisfy proper base change]
\label{prop:weil-pbc-finite-morphism}
\uses{def:weil-pbc-property}

Every finite morphism satisfies $\mathsf{PBC}$.
\end{proposition}
\begin{proof}
\uses{def:weil-pbc-property,lemma:weil-pbc-finite-exchange,lemma:weil-pbc-property-iff-acyclicity}

The finite exchange lemma supplies the required isomorphism for every base
change and every torsion sheaf, which is exactly Definition
\ref{def:weil-pbc-property}.  Equivalently, the vanishing in the exactness
lemma gives $\mathsf{IA}$ and Lemma \ref{lemma:weil-pbc-property-iff-acyclicity}
gives $\mathsf{PBC}$.
\end{proof}

\subsubsection{S2. Composition and proper-surjective cancellation}

\begin{lemma}[Composition of proper maps]
\label{lemma:weil-pbc-composition}
\uses{}

Suppose $X\xrightarrow{f}Y\xrightarrow{h}Z$ are proper and both satisfy
$\mathsf{IA}$.  Then $hf$ satisfies $\mathsf{IA}$ and hence $\mathsf{PBC}$.
\end{lemma}
\begin{proof}
\uses{lemma:weil-pbc-property-iff-acyclicity,lemma:weil-pbc-proper-underived-exchange}

Let $a:Z'\to Z$ be a base change, $Y'=Y\times_ZZ'$, $X'=X\times_Y Y'$, and
$I$ an injective $\mathbf F_\ell$-module sheaf on $X$.  Since inverse image is exact, $f_*$ (a right adjoint)
sends injectives to injectives: for an epimorphism $A\twoheadrightarrow B$,
the map $\operatorname{Hom}(B,f_*I)\to\operatorname{Hom}(A,f_*I)$ identifies
with $\operatorname{Hom}(f^{-1}B,I)\to\operatorname{Hom}(f^{-1}A,I)$ and is
surjective.  The proper underived exchange lemma identifies
$f'_*a_X^{-1}I\cong b^{-1}f_*I$, where $b:Y'\to Y$ is the base change.
The Grothendieck spectral sequence has terms
\[
 E_2^{p,q}=R^ph'_*\,R^qf'_*\,a_X^{-1}I
 \Longrightarrow R^{p+q}(h'f')_*a_X^{-1}I,
\]
where $a_X:X'\to X$ and $h':Y'\to Z'$ are the base changes.  The
first-quadrant hypotheses are: this sheaf is $h'_*$-acyclic by the assumed
$\mathsf{IA}(h)$, and $a_X^{-1}I$ is
$f'_*$-acyclic by $\mathsf{IA}(f)$.  Thus the $q>0$ rows vanish, then the
$p>0$ column of the surviving row vanishes, so all terms with $p+q>0$
vanish.  The spectral sequence converges because each total degree has only
finitely many pairs $(p,q)$.  This is the calculation located at Tag
\texttt{0A4D}.
\end{proof}

\begin{lemma}[Cancellation along a proper surjection]
\label{lemma:weil-pbc-cancellation}
\uses{}

Let $f:X\to Y$ be proper and surjective and $h:Y\to Z$ proper.  If $f$ and
$hf$ satisfy $\mathsf{IA}$, then $h$ satisfies $\mathsf{IA}$.
\end{lemma}
\begin{proof}
\uses{lemma:weil-pbc-property-iff-acyclicity,lemma:weil-pbc-proper-underived-exchange}

Fix a base change $a:Z'\to Z$ and an injective $\mathbf F_\ell$-module sheaf
$I$ on $Y$.  Choose a
  monomorphism $f^{-1}I\hookrightarrow J$ with $J$ injective on $X$.  Applying
$f_*$ gives $\iota:I\to f_*J$.  If a germ $a\in I_{\bar y}$ maps to zero,
choose a geometric $\bar x\to X$ above $\bar y$ (surjectivity).  The canonical
stalk identification $(f^{-1}I)_{\bar x}\cong I_{\bar y}$ sends $a$ to the
germ of $a$ at $\bar x$; the monomorphism into $J$ makes that germ zero.  The
definition of the direct-image stalk then says that $a$ vanishes on an etale
neighbourhood of $\bar y$, so $a=0$.  Thus $\iota$ is monic.  Injectivity of $I$ supplies a retraction
$r:f_*J\to I$ with $r\circ\iota=\mathrm{id}_I$, making $I$ a direct summand.
After base change, proper underived exchange gives
$b^{-1}f_*J\cong f'_*a_X^{-1}J$.  The spectral sequence
\[
R^ph'_*R^qf'_*a_Y^{-1}I\Longrightarrow R^{p+q}(h'f')_*a_X^{-1}I
\]
is computed using the injective $J$ and the assumed acyclicity of $hf$.
The $q>0$ terms vanish for $J$ by $\mathsf{IA}(f)$, and total-degree
vanishing for $hf$ forces $R^ph'_*a_Y^{-1}(f_*J)=0$ for $p>0$.  Derived
direct image preserves finite direct sums and retracts, so the same vanishing
holds for the direct summand $a_Y^{-1}I$.  This is Tag \texttt{0A4C}; the
stalk injectivity argument is written here because surjectivity is essential.
\end{proof}

\subsubsection{S3. Reduction to the projective line}

\begin{lemma}[Target locality]
\label{lemma:weil-pbc-target-locality}
\uses{def:weil-pbc-property}
For proper $f:X\to Y$, the property $\mathsf{PBC}(f)$ is local for the etale
topology on $Y$.
\end{lemma}
\begin{proof}
\uses{} 
For an etale $V\to Y$ and a torsion sheaf $F$, sections of $R^qf_*F$ on $V$
are the sheafification of $U\mapsto H^q(X\times_YU,F|)$ over etale
$U\to V$.  This description is functorial in $V$ and commutes with further
etale localization.  Therefore an exchange map is an isomorphism if and
only if it is so after an etale cover of $Y$ (and its pullback cover on
$Y'$), because geometric stalks form a conservative family.  The section
description is the content located at Tag \texttt{03Q8}.
\end{proof}

\begin{lemma}[Chow and projective factorization]
\label{lemma:weil-pbc-chow-factorization}
\uses{}

After replacing an affine target by an etale affine cover, every proper
$f:X\to Y$ admits a proper surjection $\pi:X'\to X$ and an integer $n$ with
$X'\hookrightarrow\mathbf P^n_Y$ a closed immersion.  Moreover
$X'\to Y$ factors as that closed immersion followed by
$\mathbf P^n_Y\to Y$.
\end{lemma}
\begin{proof}
\uses{lemma:weil-pbc-target-locality,lemma:weil-pbc-cancellation,
thm:chow-lemma,cor:chow-lemma-proper,
lemma:nagata-projective-space-base-change-proper}

Chow's lemma in the form located at Tag \texttt{0202}, i.e.
\Cref{cor:chow-lemma-proper}, supplies a proper surjective $\pi$ for which
$X'\to Y$ is projective; it applies verbatim once the base is noetherian, so
the only work left below is to reduce to that case.  To make the source
hypotheses explicit, first work over one affine member $Y=\operatorname{Spec}A$
of the selected etale cover.  Write
$A=\varinjlim_{\lambda\in\Lambda}A_\lambda$, where every $A_\lambda$ is a
finitely generated $\mathbf Z$-algebra and the transition maps are directed.
Since a proper morphism is of finite presentation, a finite affine cover
$X=\bigcup_{a=1}^rU_a$, the finitely presented algebras defining the $U_a$,
the pairwise overlap maps, and their cocycle equations descend to one index
$\lambda_0$.  After increasing once to $\lambda_1$, the descended morphism
$f_{\lambda_1}:X_{\lambda_1}\to\operatorname{Spec}A_{\lambda_1}$ is proper:
finite presentation and separatedness descend from the finite affine data,
and the valuative criterion for properness is stable under this inverse-limit
presentation.

The ring $A_{\lambda_1}$ is noetherian.  Noetherian Chow's lemma produces a
proper surjection
$\pi_{\lambda_1}:X'_{\lambda_1}\to X_{\lambda_1}$ and an
$f_{\lambda_1}\pi_{\lambda_1}$-ample invertible sheaf $L_{\lambda_1}$.
Choose a finite affine cover $X'_{\lambda_1}=\bigcup_{b=1}^sW_b$.  On each
$W_b$, ampleness supplies an exponent $m_b>0$ and finitely many global
sections whose nonvanishing opens cover $W_b$.  With
$m=\prod_bm_b$, take the corresponding powers of these sections.  There are
then finitely many sections $s_0,\ldots,s_n$ of $L_{\lambda_1}^{\otimes m}$
whose nonvanishing opens cover $X'_{\lambda_1}$.  Increasing $m$ by one
further common multiple if necessary, the same finite affine-cover test for
very ampleness separates the finitely many overlap generators; hence the
map
\[
 [s_0:\cdots:s_n]:X'_{\lambda_1}\longrightarrow
 \mathbf P^n_{A_{\lambda_1}}
\]
is an immersion.  It is proper, because $X'_{\lambda_1}$ is proper over
$A_{\lambda_1}$ and projective space is separated over that base.  A proper
immersion is closed.

Base change along $A_{\lambda_1}\to A$ defines
$X'=X'_{\lambda_1}\times_{A_{\lambda_1}}A$, $\pi$, and the closed immersion
$X'\hookrightarrow\mathbf P^n_Y$.  Properness and closed immersions are
preserved by base change.  Surjectivity is preserved because every geometric
point of $X$ lifts after base change from the surjective geometric fibre of
$\pi_{\lambda_1}$.  Thus the displayed index $\lambda_1$, exponent $m$, and
integer $n$ give the asserted single factorization.
\end{proof}

\begin{lemma}[Projective embedding over $X$]
\label{lemma:weil-pbc-chow-embedding-over-x}
\uses{}

Let $\iota_Y:X'\hookrightarrow\mathbf P^n_Y$ be the closed immersion and
$\pi:X'\to X$ the Chow modification.  Then
\[
 \iota_X=(\pi,\iota_Y):X'\longrightarrow
 X\times_Y\mathbf P^n_Y=\mathbf P^n_X
\]
is a closed immersion over $X$.
\end{lemma}
\begin{proof}
\uses{lemma:weil-pbc-chow-factorization}

The product map is an immersion because $\iota_Y$ is an immersion and
immersions are preserved by base change and composition with a graph.  The
map $\iota_X$ is proper over $X$: its source $X'$ is proper over $X$, while
$\mathbf P^n_X$ is separated over $X$, so the graph criterion makes every
$X$-morphism from $X'$ proper.  A proper immersion is a closed immersion:
properness makes its image closed after every base change, and an immersion
identifies the source with a locally closed subscheme of that image; universal
closedness forces the locally closed subset to be closed.  The equality with
$\mathbf P^n_X$ is the defining fibre-product identity.
\end{proof}

\begin{lemma}[Symmetric product map]
\label{lemma:weil-pbc-symmetric-product-p1}
\uses{}

For $n\geq1$ the formula
\[
 ((x_1:y_1),\ldots,(x_n:y_n))\longmapsto
 [F_0:\cdots:F_n],\qquad
 \prod_{i=1}^n(x_it+y_i)=\sum_{j=0}^nF_jt^j,
\]
defines a finite surjective $S_n$-invariant morphism
$\sigma_n:(\mathbf P^1_S)^n\to\mathbf P^n_S$.
\end{lemma}
\begin{proof}
\uses{prop:weil-pbc-finite-morphism,lemma:weil-pbc-composition,lemma:weil-pbc-cancellation}

Each $F_j$ is the sum of monomials choosing $j$ of the $x_i$ and the
remaining $y_i$; hence it is homogeneous of degree one in every pair
$(x_i,y_i)$.  On the chart $y_1\cdots y_n\ne0$, the ratios
$z_i=x_i/y_i$ give elementary symmetric polynomials, so the map is induced
by the universal polynomial $\prod(T-z_i)$.  On a general chart, the same
calculation uses the coefficient ring localized at one nonzero $F_j$; the
coefficients generate the unit ideal because if all $F_j$ vanished then the
homogeneous polynomial $\prod_i(x_it+y_i)$ would be zero, forcing one factor
to be $(0,0)$, impossible in projective coordinates.  The coordinate ring is
integral over the symmetric subring: each $z_i$ satisfies the monic equation
$T^n-e_1T^{n-1}+\cdots+(-1)^ne_n$, and the analogous statement holds on all
charts after clearing the chosen $y_i$ or $x_i$.  Thus $\sigma_n$ is finite.
Permutation of the factors leaves every $F_j$ unchanged, proving
$S_n$-invariance.  To make finiteness independent of the abbreviated chart
calculation, let $[a_0:\cdots:a_n]$ be a geometric point of $\mathbf P^n_S$.
Its fibre is the set of ordered $n$-tuples of points of $\mathbf P^1$ whose
linear forms multiply to $\sum a_jt^j$.  Over an algebraically closed field,
the factorization theorem for one-variable polynomials, applied after choosing
the first nonzero coefficient to decide whether $t=\infty$ is a root, gives
only finitely many such tuples (at most $n!$ after the multiplicities are
fixed).  Hence $\sigma_n$ is quasi-finite.

The source $(\mathbf P^1_S)^n$ is projective over $S$, so $\sigma_n$ is
proper.  On an affine target chart, the quasi-finite proper map is covered by
finitely many affine opens on which the coordinate rings are integral over the
target ring: on the chart where the last coefficient is invertible use the
monic equations for the $z_i$, and on the chart where the first coefficient
is invertible use the reversed monic equations for $w_i=y_i/x_i$; on an
intermediate chart, intersect these two covers with the finitely many opens
where a selected coefficient of the binary form is invertible.  The same
monic equations after multiplying by that selected coefficient show that each
coordinate is integral.  Integral algebras glue on overlaps because the
coordinate identifications are induced by the common coefficients $F_j$.
Thus the proper quasi-finite morphism is finite.  Finally a geometric binary
form of degree $n$ factors into $n$ linear forms, including the possible
factor $x$ for $t=\infty$, so every geometric fibre is nonempty; geometric
surjectivity and properness give surjectivity.  This closes the chart and
gluing obligation.
\end{proof}

\begin{proposition}[Reduction of proper base change to $\mathbf P^1$]
\label{prop:weil-pbc-reduction-to-projective-line}
\uses{def:weil-pbc-derived-exchange}

If every projection $\mathbf P^1_S\to S$ satisfies $\mathsf{PBC}$ for every
base scheme $S$, then every proper morphism satisfies $\mathsf{PBC}$.
\end{proposition}
\begin{proof}
\uses{def:weil-pbc-derived-exchange,lemma:weil-pbc-target-locality,lemma:weil-pbc-chow-factorization,lemma:weil-pbc-chow-embedding-over-x,lemma:weil-pbc-symmetric-product-p1,prop:weil-pbc-finite-morphism,lemma:weil-pbc-composition,lemma:weil-pbc-cancellation}

Use target locality and the Chow factorization.  The closed immersion
$X'\hookrightarrow\mathbf P^n_Y$ is proper and finite onto its image, so
finite base change plus composition reduces $X'\to Y$ to
$\mathbf P^n_Y\to Y$.  The finite surjective map $\sigma_n$ and cancellation
reduce this projection to $(\mathbf P^1_Y)^n\to Y$.  Repeated base change and
composition reduce the product projection to the one-dimensional projection.
Finally cancellation along $\pi$ transfers the property from $X'$ and
$X'\to Y$ to $X\to Y$.  Every comparison used is the named exchange map of
Definition \ref{def:weil-pbc-derived-exchange}.
\end{proof}

\subsubsection{S4. The projective line over a strictly henselian local base}

\begin{definition}[Strictly henselian setup]
\label{def:weil-pbc-strict-henselian-setup}
\uses{}
Let $(A,\mathfrak m)$ be strictly henselian local, $S=\operatorname{Spec}A$,
$s=\operatorname{Spec}(A/\mathfrak m)$, $X=\mathbf P^1_S$, and
$X_0=\mathbf P^1_s$.  For a sheaf $F$ on $X_{\mathrm{et}}$ write
$F_0$ for its inverse image on $(X_0)_{\mathrm{et}}$.
\end{definition}

\begin{lemma}[Zariski restriction on a spectral space]
\label{lemma:weil-pbc-spectral-restriction}
\uses{}
If $T$ is spectral, $T_0\subseteq T$ is closed, and every irreducible closed
subset $C\subseteq T$ has $C\cap T_0$ connected and nonempty, then for every
abelian sheaf $F$ on $T$ the restriction $\Gamma(T,F)\to\Gamma(T_0,F|)$ is a
bijective map.
\end{lemma}
\begin{proof}
\uses{} 
Injectivity is pointwise.  If $s\in\Gamma(T,F)$ restricts to zero, its germ
vanishes at every point of $T_0$.  The zero-germ locus is open and stable under
generization.  If its complement contained $x$, then the irreducible closed
$C=\overline{\{x\}}$ would meet $T_0$; the specialization map from the germ
at the generic point of $C$ to the germ at $x$ would force the latter to be
zero, a contradiction.

For surjectivity, write $F$ as the filtered colimit of subsheaves generated
by finitely many sections on quasi-compact opens.  Restriction to $T_0$ and
both section functors commute with this filtered colimit, so it is enough to
handle one finitely generated subsheaf.  Such a subsheaf embeds into a finite
product
\[
 F\hookrightarrow\prod_{b=1}^r(i_b)_*\underline{M_b},
\]
with $i_b:C_b\hookrightarrow T$ closed and $M_b$ finite: the finite germs of
the chosen generators determine finitely many closed equalizer loci, and the
corresponding finite value sets give the displayed embedding.  If $F\subseteq
G$, the injectivity already proved gives
\[
 \Gamma(T,F)=\Gamma(T_0,F|_{T_0})\cap\Gamma(T,G)
 \subseteq\Gamma(T_0,G|_{T_0}).
\]
Thus it suffices to extend sections of one finite constant pushforward.

For $i:C\hookrightarrow T$ closed, the condition on irreducible closures is
inherited by $C$.  A section of $\underline M$ on $C\cap T_0$ is constant on
each connected component.  Every component of $C$ meets its closed subset,
and that intersection is connected by hypothesis, so the same constant gives
a unique section on the component.  These componentwise sections agree on
quasi-compact overlaps and glue by the sheaf equalizer; pushing forward gives
the extension on $T$.  Finite products, subobjects, and the filtered colimit
then give an extension for $F$.  The stalkwise injectivity proves uniqueness,
so restriction is bijective.  This is the expanded spectral argument at Tag
\texttt{09ZG}; no etale lifting is used.
\end{proof}

\begin{lemma}[Stein connectedness for $\mathbf P^1_A$]
\label{lemma:weil-pbc-stein-connectedness}
\uses{def:weil-pbc-strict-henselian-setup}
Every irreducible closed subset of $X$ meets $X_0$ in a connected nonempty
set.  Consequently $\Gamma(X,F)\xrightarrow{\sim}\Gamma(X_0,F_0)$ for every
Zariski sheaf $F$.
\end{lemma}
\begin{proof}
\uses{lemma:weil-pbc-spectral-restriction}

Fix an irreducible closed subset and give it the reduced closed-subscheme
structure $C\subseteq X$.  The map $C\to S$ is proper.  Its Stein
factorization gives $C\xrightarrow{u}\operatorname{Spec}B
\xrightarrow{v}S$, with $A\to B$ integral and $u$ having geometrically
connected fibres.  Since $(A,\mathfrak m)$ is henselian,
$(B,\mathfrak mB)$ is a henselian pair.  The image
$D=v(\operatorname{Spec}B)$ is irreducible and closed because $C\to S$ is
proper.  Henselian connectedness makes $D\cap V(\mathfrak m)$ connected and
nonempty.  The map
$C\cap X_0\to D\cap V(\mathfrak m)$ is closed and surjective with connected
fibres, so the closed-surjection connected-fibre criterion makes
$C\cap X_0$ connected.  This expands Tag
\texttt{0A0B}; the final equality of sections is Lemma
\ref{lemma:weil-pbc-spectral-restriction}.
\end{proof}

\begin{definition}[Pullback of a section along a finite morphism]
\label{def:weil-pbc-finite-section-pullback}
\uses{}
Let $u:T'\to T$ be finite and $F$ an etale sheaf on $T$.  For an etale
$V\to T$ and $s\in F(V)$, define
\[
 u_V^\sharp(s)\in(u^{-1}F)(V\times_TT')
\]
to be the image of $s$ under the unit from the presheaf inverse image to its
sheafification.
\end{definition}

\begin{lemma}[Finite pullback identity]
\label{lemma:weil-pbc-finite-pullback-identity}
\uses{def:weil-pbc-finite-section-pullback}
For $u=\operatorname{id}_T$, $u_V^\sharp(s)=s$ under the canonical identity
$V\times_TT=V$.
\end{lemma}
\begin{proof}
\uses{} 
The presheaf inverse image along the identity is the same presheaf, and its
sheafification unit is the identity on an already sheafified section.
\end{proof}

\begin{lemma}[Finite pullback composition]
\label{lemma:weil-pbc-finite-pullback-composition}
\uses{def:weil-pbc-finite-section-pullback}
For finite $T''\xrightarrow{v}T'\xrightarrow{u}T$,
$(uv)_V^\sharp(s)=v_{V\times_TT'}^\sharp(u_V^\sharp(s))$ under the
associativity isomorphism of fibre products.
\end{lemma}
\begin{proof}
\uses{lemma:weil-pbc-finite-pullback-identity}

Both sections are the image of the same presheaf section $s$ under the
composite inverse-image functor.  Naturality of sheafification and the
associativity isomorphism make the two unit maps equal.
\end{proof}

\begin{lemma}[Finite pullback and closed-fibre restriction]
\label{lemma:weil-pbc-finite-pullback-closed-fibre}
\uses{def:weil-pbc-finite-section-pullback}
For a closed immersion $i:Z\hookrightarrow T$ and
$Z'=Z\times_TT'$, restriction of $u_V^\sharp(s)$ to $Z'$ equals the finite
pullback of the restriction of $s$ to $Z$.
\end{lemma}
\begin{proof}
\uses{lemma:weil-pbc-finite-pullback-composition}

Apply the composition lemma to the two equal composites
$Z'\to T'\to T$ and $Z'\to Z\to T$.  The Cartesian identification makes
their inverse-image sheafification units equal.
\end{proof}

\begin{lemma}[Finite descent compatibility]
\label{lemma:weil-pbc-finite-pullback-descent}
\uses{def:weil-pbc-finite-section-pullback}
Let $u:T'\to T$ be finite surjective.  Suppose $s'\in(u^{-1}F)(T')$ is
equipped with a descended section $s\in F(T)$, meaning
$u_T^\sharp(s)=s'$.  Then the same equality holds after restriction to a
closed fibre.
\end{lemma}
\begin{proof}
\uses{lemma:weil-pbc-finite-pullback-composition,lemma:weil-pbc-finite-pullback-closed-fibre}

The first equality is the descent witness.  Its closed-fibre restriction is
the equality supplied by Lemma
\ref{lemma:weil-pbc-finite-pullback-closed-fibre}.
\end{proof}

\begin{lemma}[Etale henselian restriction in degree zero]
\label{lemma:weil-pbc-etale-degree-zero}
\uses{def:weil-pbc-strict-henselian-setup}
For every etale abelian sheaf $F$ on $X$, restriction induces
$\Gamma(X,F)\xrightarrow{\sim}\Gamma(X_0,F_0)$.
\end{lemma}
\begin{proof}
\uses{lemma:weil-pbc-stein-connectedness,lemma:weil-pbc-finite-pullback-identity,lemma:weil-pbc-finite-pullback-composition,lemma:weil-pbc-finite-pullback-closed-fibre,lemma:weil-pbc-finite-pullback-descent,lemma:weil-pbc-finite-refinement-etale-cover}

The canonical map from the etale to Zariski restriction is injective on
stalks because a geometric point has a strictly local etale neighbourhood.
For surjectivity, lift a fibre section etale-locally; quasi-compactness of
$X_0$ gives finitely many local lifts, and Lemma
\ref{lemma:weil-pbc-finite-refinement-etale-cover} gives a finite surjective
refinement on which all lifts exist.
The two lifts agree on the finite overlap after restriction to $X_0$ by the
Zariski result, and the equalizer sheaf descends the lift to $X$.  Uniqueness
is stalkwise injectivity.  The finite-pullback compatibility required in
Stacks Tag \texttt{09ZF} is recorded as four explicit equalities:
identity pullback preserves the lift, composition of finite pullbacks preserves
the lift, restriction to $X_0$ recovers the given section, and the descended
section pulls back to each local lift; these are the four preceding lemmas.
The finite refinement lemma proves existence and finite-presentation stage,
and its stalkwise equalizer argument proves equality of the local lifts on
the self-intersection.  Thus the four finite-pullback compatibilities
complete the descent.
\end{proof}

\begin{lemma}[Finite-stage killing of a class]
\label{lemma:weil-pbc-finite-stage-killing}
\uses{}
In an affine henselian pair, if a positive-degree class $c$ dies on an
integral strictly-local cover $T_\infty\to T_0$ represented by a filtered
system $\{T_\lambda\}$ of finite finitely presented covers, then there exists
$\lambda_0$ such that the pullback of $c$ to $H^q(T_{\lambda_0},F|)$ is zero.
\end{lemma}
\begin{proof}
\uses{} 
Represent the chosen cochain and its null-homotopy on the limit by finitely
many sections.  Each section and each equality occurs at some index because
the system is filtered and the etale site is quasi-compact and quasi-separated.
Choose one index dominating those finitely many indices.  The descended
cochain and the descended null-homotopy satisfy the coboundary equation at a
further finite index; replacing by that index gives $\lambda_0$.  This exposes the indexing and
vanishing witness requested by Tag \texttt{09ZD}.
\end{proof}

\begin{lemma}[Constructible approximation for effacement]
\label{lemma:weil-pbc-constructible-approximation}
\uses{}
For a torsion sheaf $F$ on an affine henselian pair and a fixed class
$c\in H^q(T_0,F_0)$, there is a quasi-compact etale neighbourhood $U\to T$
and a constructible sheaf $C$ on $U$ with a monomorphism $C\hookrightarrow F|_U$
through which a cocycle representing $c$ factors.
\end{lemma}
\begin{proof}
\uses{} 
Choose a split etale hypercover $U_\bullet\to T_0$ whose $(q+1)$-skeleton is
finite and quasi-compact and a cocycle $c\in\Gamma(U_q,F_0)$.  Write its
finite components as $U_{r,a}$ for $0\leq r\leq q+1$ and denote all face and
degeneracy maps by $\partial^i_{r,a}$ and $s^i_{r,a}$.  Choose quasi-compact
etale lifts $\widetilde U_{r,a}\to U\to T$ and sections of $F$ representing
the finitely many sections in $c$ and in the equations $dc=0$.

Enlarge the finite list by applying every $\partial^i_{r,a}$ and
$s^i_{r,a}$ to every chosen section.  The simplicial identities are now a
finite list of equalities on finite fibre products.  Shrink the lifts once so
that all equalities hold on the closed fibre and all maps factor through one
quasi-compact etale neighbourhood $U\to T$.  On each lifted component let
$C_{r,a}$ be the subsheaf of $F|_U$ generated by this finite list.  The
equalizer loci of pairs of generators give a finite stratification on which
$C_{r,a}$ is locally constant with finite fibres, hence constructible.

Closure under all displayed face and degeneracy maps gives a simplicial
constructible subsheaf $C_\bullet\subseteq F|_{U_\bullet}$.  The cocycle lies
in $\Gamma(U_q,C_q)$ and the equation $dc=0$ remains true in $C_\bullet$.
Take the image of the augmentation $C_0\to F|_U$ and its finite iterated
extensions through degree $q+1$, then sheafify; this gives a constructible
subsheaf $C\hookrightarrow F|_U$ through which $c$ factors.  Monicity is
checked at a geometric stalk: a finite combination of the generators maps to
zero in $F$ precisely when it is already zero in the generated subsheaf.
Thus every face map, cocycle, and equality used by the selected class is
contained in $C$.
\end{proof}

\begin{lemma}[Pushout preserves monomorphisms]
\label{lemma:weil-pbc-pushout-monicity}
\uses{}
Given monomorphisms $A\hookrightarrow B$ and $A\hookrightarrow C$ of sheaves,
the canonical map $C\to B\amalg_A C$ is monic.
\end{lemma}
\begin{proof}
\uses{} 
Let $P=B\amalg_A C$ be the pushout in sheaves.  It is the cokernel sheaf of
$A\to B\oplus C$, $a\mapsto(i_Ba,-i_Ca)$.  If $c\in C(V)$ maps to zero in
$P(V)$, then by the definition of a sheaf cokernel there is an etale cover
$\{V_\alpha\to V\}$ and $a_\alpha\in A(V_\alpha)$ with
$i_C(a_\alpha)=c|_{V_\alpha}$ and $i_B(a_\alpha)=0$.  Since $i_B$ is monic,
$a_\alpha=0$, hence $c|_{V_\alpha}=0$.  The sheaf axiom gives $c=0$.
This proves monicity on arbitrary etale objects, and the same local
equalizer calculation proves naturality of the pushout map.
\end{proof}

\begin{lemma}[Constructible sheaves reduce to finite pushforwards]
\label{lemma:weil-pbc-constructible-finite-pushforward}
\uses{}

If $U$ is quasi-compact and quasi-separated, every constructible torsion
sheaf $C$ admits a monomorphism
\[
 C\longrightarrow\bigoplus_{i=1}^r(h_i)_*\underline{M_i},
\]
where each $h_i:V_i\to U$ is finite and finitely presented and each $M_i$ is
a finite abelian group.
\end{lemma}
\begin{proof}
\uses{lemma:weil-pbc-constructible-approximation}

Write $U=\varprojlim U_\lambda$ with affine transition maps and noetherian,
quasi-compact, quasi-separated $U_\lambda$.  The finite stratification of
$C$, its finite stalk groups, all restriction maps, and their equalizer
relations descend to one stage $U_\lambda$.  We prove the assertion there by
noetherian induction on a closed complement of a dense stratum.  Choose a
dense open $j:V\hookrightarrow U_\lambda$ where $C_\lambda$ is finite
locally constant and a finite etale cover $e:W\to V$ trivializing it.  Take
the integral closure of $U_\lambda$ in $e_*\mathcal O_W$ and discard
components disjoint from $V$; this gives a finite map $h:W'\to U_\lambda$
extending $e$.  The adjunction unit is a monomorphism on $V$,
\[
 C_\lambda|_V\hookrightarrow e_*\underline M
      =(h_*\underline M)|_V,
\]
for a finite group $M$.

Let $Z=U_\lambda\setminus V$ with its reduced closed structure.  By
induction, $C_\lambda|_Z$ embeds into a finite sum
$\bigoplus_b(h_b)_*\underline{M_b}$ with $h_b$ finite over $Z$; composing
with $Z\hookrightarrow U_\lambda$ keeps the maps finite.  Map
$C_\lambda$ to
\[
 h_*\underline M\oplus\bigoplus_b(h_b)_*\underline{M_b}
\]
by the dense-open unit in the first component and restriction to $Z$ in the
remaining components.  At a geometric point of $V$ the first component is
injective, and at a geometric point of $Z$ the second block is injective, so
the combined map is monic.  This is the localization induction, including
the boundary and overlap comparison.  Pull the finite diagram and monomorphism
back to $U$.  Etale inverse image is exact, and finite direct image commutes
with this pullback by the finite geometric-stalk product formula.  This gives
the displayed monomorphism and expands Tag \texttt{09Z7}.
\end{proof}

\begin{lemma}[Finite-cover killing and finite pushforward monicity]
\label{lemma:weil-pbc-finite-cover-killing}
\uses{}

For a fixed $c\in H^q(T_0,C_0)$ with $q>0$ and constructible $C$, there is a
finite finitely presented cover $h:V\to T$ whose closed-fibre pullback
$h_0^*c$ is zero; the induced map $C\hookrightarrow h_*h^*C$ is monic.
\end{lemma}
\begin{proof}
\uses{lemma:weil-pbc-finite-stage-killing,lemma:weil-pbc-finite-stalk-formula}

Use the integral strictly-local cover killing $c$ and represent it by the
filtered system of finite, finitely presented covers
$h_\lambda:V_\lambda\to T$.  Lemma~\ref{lemma:weil-pbc-finite-stage-killing}
selects an index $\lambda_0$ at which $h_{\lambda_0,0}^*c=0$; set
$h=h_{\lambda_0}$.  The map to the finite pushforward is the unit of
$(h^*,h_*)$.  At a geometric point its stalk is the diagonal into the finite
product of stalks from Lemma~\ref{lemma:weil-pbc-finite-stalk-formula}; the
factor corresponding to the chosen point above it is the identity.  Hence the
unit is injective on every geometric stalk and therefore monic.
\end{proof}

\begin{lemma}[Affine Gabber effacement]
\label{lemma:weil-pbc-gabber-effacement}
\uses{}

For an affine henselian pair $(T,T_0)$ and a torsion sheaf $F$ on $T_{\mathrm{et}}$,
every class $c\in H^q(T_0,F_0)$ with $q>0$ admits a monomorphism
$F\hookrightarrow F'$ such that the image of $c$ in $H^q(T_0,F'_0)$ is zero.
\end{lemma}
\begin{proof}
\uses{lemma:weil-pbc-constructible-approximation,lemma:weil-pbc-constructible-finite-pushforward,lemma:weil-pbc-finite-cover-killing,lemma:weil-pbc-pushout-monicity}

Apply constructible approximation to obtain $C\hookrightarrow F$ through
which $c$ factors, and use finite-pushforward reduction to choose
\[
 C\hookrightarrow D_0=\bigoplus_{i=1}^r(h_i)_*\underline{M_i}.
\]
Let $c_i$ be the image of $c$ in
$H^q(T_0,(h_i)_*\underline{M_i}|_{T_0})$.  Finite direct-image exactness
identifies $c_i$ with a class on the closed fibre of $h_i$.  Apply finite-cover
killing to each $c_i$ and take the fibre product of the resulting finite
covers over $T$; after replacing by one finite-presentation stage, call the
common finite cover $h:V\to T$.  Every component of $h_0^*c$ is zero.
The composite unit
\[
 C\hookrightarrow D_0\longrightarrow h_*h^*D_0=:D
\]
is monic by the finite stalk product formula.  Form the pushout square
\[
\begin{CD}
 C @>>> D\\
 @VVV @VVV\\
 F @>>> F':=F\amalg_C D.
\end{CD}
\]
Pushout monicity makes $F\to F'$ injective.  The image of $c$ in
$H^q(T_0,D_0)$ is carried to zero in $H^q(T_0,D)$ because its pullback to
$V_0$ is zero and $h_*$ is exact.  The pushout square therefore carries $c$
to zero in $H^q(T_0,F'_0)$.  This displays both the common dominating cover
and the null-homotopy square and closes Tag \texttt{09ZE}'s effacement step.
\end{proof}

\begin{proposition}[Gabber affine comparison]
\label{prop:weil-pbc-gabber-affine-comparison}
\uses{}
For an affine henselian pair $(T,T_0)$ and torsion $F$,
$H^q(T,F)\xrightarrow{\sim}H^q(T_0,F_0)$ for all $q\geq0$.
\end{proposition}
\begin{proof}
\uses{} 
Degree zero is the etale henselian restriction lemma applied on the affine.
For $q>0$, assume comparison in all degrees below $q$.  To prove injectivity,
take a nonzero $\xi\in H^q(T,F)$ and apply effacement with closed subscheme
$T$ itself, obtaining $F\hookrightarrow F'$ in which $\xi$ dies.  Put
$Q=F'/F$.  Exactness lifts $\xi$ from $H^{q-1}(T,Q)$.  The induction
isomorphism in degree $q-1$ carries any such lift to $T_0$; if $\xi$ restricted
to zero, exactness and the commutative boundary square would force the lift
to come from $H^{q-1}(T,F')$, contradicting the choice of nonzero $\xi$.

For surjectivity, take $\eta\in H^q(T_0,F_0)$ and use effacement with closed
subscheme $T_0$ to choose $F\hookrightarrow F'$ killing $\eta$.  Exactness
lifts $\eta$ from $H^{q-1}(T_0,Q_0)$.  By induction this lift is the
restriction of a class in $H^{q-1}(T,Q)$, whose boundary in $H^q(T,F)$
restricts to $\eta$.  Thus the comparison is separately injective and
surjective.  This expands Tag \texttt{09ZI}; no statement is imported from
that tag.
\end{proof}

\begin{lemma}[Injective restriction acyclicity bound]
\label{lemma:weil-pbc-injective-restriction-acyclicity}
\uses{}

Let a separated scheme with affine diagonal be covered by $n+1$ affine opens.
If $I$ is injective on the scheme over a henselian pair, then the restriction
to the closed fibre has vanishing cohomology in degrees $>n$.
\end{lemma}
\begin{proof}
\uses{prop:weil-pbc-gabber-affine-comparison}

For $n=0$ this is the affine comparison applied to the injective $I$.  For
$n>0$, split the cover as $U_0$ and $V=\bigcup_{i>0}U_i$.  The Mayer--Vietoris
long exact sequence for $U_0\cup V$ has terms on $U_0$, $V$, and
$U_0\cap V$; the last two have covers of size at most $n$ and the affine
comparison kills their groups above the inductive bound.  The long exact
sequence therefore kills the group above degree $n$.  This is the dimension
induction located at Tag \texttt{0A51}.
\end{proof}

\begin{lemma}[Finite locally constant sheaves and finite-etale schemes]
\label{lemma:weil-pbc-finite-locally-constant-etale-equivalence}
\uses{}
For a scheme $T$, finite locally constant sheaves of sets on $T_{\mathrm{et}}$
are equivalent to finite etale $T$-schemes.  A finite abelian group action on
the sheaf corresponds to an action by finite-etale $T$-automorphisms.
\end{lemma}
\begin{proof}
\uses{} 
Given a sheaf $E$, choose a finite etale cover $U\to T$ trivializing it with
finite fibre $S$.  The descent isomorphism on $U\times_TU$ is a permutation
of $S$ satisfying the cocycle equation, hence descends the disjoint union
$\coprod_{s\in S}U$ to a finite etale scheme $P$.  Conversely, the sheaf of
sections of a finite etale $P$ is finite locally constant; the two constructions
are inverse because both are computed after the same trivializing cover.
An $M$-action is a compatible family of permutations, so it descends exactly
when the sheaf action is compatible with the descent cocycle.
\end{proof}

\begin{lemma}[Finite-etale lifting over a henselian pair]
\label{lemma:weil-pbc-finite-etale-henselian-lifting}
\uses{def:weil-pbc-strict-henselian-setup}
If $(A,\mathfrak m)$ is henselian and $Y\to\operatorname{Spec}A$ is proper,
restriction from finite etale $Y$-schemes to finite etale $Y_0$-schemes is an
equivalence, including unique lifting of morphisms and group actions.
\end{lemma}
\begin{proof}
\uses{lemma:weil-pbc-finite-locally-constant-etale-equivalence}

We expand the equivalence used by Tag \texttt{0A5G}; its source locator is
Tag \texttt{0A48}.  First suppose that $A$ is complete noetherian and put
$Y_r=Y\times_AA/\mathfrak m^{r+1}$.  Formal etaleness and nilpotent
invariance of the etale site lift a finite etale $P_0\to Y_0$ uniquely to a
compatible system $P_r\to Y_r$.  Finiteness can be checked modulo the
nilpotent ideal: on a finite affine cover
$Y_r=\bigcup_a\operatorname{Spec}R_{r,a}$, each lifted algebra is generated
by lifts of a finite set of generators of the $r=0$ algebra, and Nakayama
makes the resulting algebra finite.  The transition isomorphisms on the
finitely many overlaps obey the cocycle equation because their restrictions
do and etale morphisms are formally unramified.

Grothendieck existence applied to the compatible finite algebras
$(p_r)_*\mathcal O_{P_r}$ algebraizes them to a coherent
$\mathcal O_Y$-algebra, hence to a finite morphism $P\to Y$ with special
fibre $P_0$.  The etale locus of $P\to Y$ is open and contains $P_0$.
Every point of $P$ specializes to the closed fibre because $P$ is proper
over the local base; therefore the complementary closed non-etale locus is
empty.  This proves essential surjectivity in the complete noetherian case.
For full faithfulness, a map $u_0:P_0\to Q_0$ has graph
\[
 \Gamma_{u_0}:P_0\longrightarrow P_0\times_{Y_0}Q_0,
\]
which is an open-and-closed finite etale subscheme.  Algebraize this graph to
$W\to P\times_YQ$.  The projection $W\to P$ is finite etale and is an
isomorphism on the closed fibre; its isomorphism locus is open, and the
proper closed complement would meet the closed fibre if nonempty.  Hence it
is empty, and $W\cong P$ gives the unique lift $u:P\to Q$.

Now suppose $A$ is a noetherian henselian local G-ring.  Complete it to
$\widehat A$.  The first step gives a finite etale cover over
$Y_{\widehat A}$.  Choose a finite affine cover of $Y$, write the finitely
many cover algebras, overlap maps, and cocycle equations with finitely many
generators and relations, and descend all of them to a finite type
$A$-algebra
$A_i=A[x_1,\ldots,x_N]/(f_1,\ldots,f_M)$ inside $\widehat A$.
After increasing $i$, the descended map is finite etale.  Approximation for
henselian G-rings gives a solution $(b_1,\ldots,b_N)\in A^N$ of all equations
with $b_j\equiv x_j\pmod{\mathfrak m}$.  The induced retraction
$A_i\to A$ base-changes the descended cover to a finite etale cover of $Y$
with special fibre $P_0$.  The graph argument proves full faithfulness.

For a general henselian local $A$, choose noetherian local G-rings $A_i$ from
finite-type $\mathbf Z$-models and write
$A=\varinjlim_iA_i^h$.  Proper finite presentation gives an index $i_0$ and
a proper model $Y_{i_0}\to\operatorname{Spec}A_{i_0}^h$.  A finite etale
cover, a morphism between two such covers, and an equality of two morphisms
are each finite presentation data; using a finite affine cover and the
finitely many overlap equations, they descend to one common $i\geq i_0$.
Thus
\[
 \mathit{FEt}_Y=\varinjlim_i\mathit{FEt}_{Y_i},\qquad
 \mathit{FEt}_{Y_0}=\varinjlim_i\mathit{FEt}_{(Y_i)_0},
\]
including the Hom sets.  The equivalence at every henselian noetherian
G-ring stage gives the desired equivalence after the colimit.  Finally an
action of a finite group consists of finitely many morphisms and finitely
many multiplication, unit, and inverse equations.  Full faithfulness lifts
the morphisms and their equations uniquely, so group actions lift as well.
\end{proof}

\begin{lemma}[Finite-etale torsor comparison]
\label{lemma:weil-pbc-h1-torsor-comparison}
\uses{def:weil-pbc-strict-henselian-setup}
For a finite abelian group $M$ and proper $Y\to\operatorname{Spec}A$,
restriction induces a bijection
\[
 H^1(Y,\underline M)\xrightarrow{\sim}H^1(Y_0,\underline M).
\]
\end{lemma}
\begin{proof}
\uses{lemma:weil-pbc-etale-degree-zero,lemma:weil-pbc-finite-locally-constant-etale-equivalence,lemma:weil-pbc-finite-etale-henselian-lifting}

An $\underline M$-torsor is a finite etale $Y$-scheme $P$ with a simply
transitive $M$-action by the finite-locally-constant equivalence.  Apply the
henselian lifting lemma to lift $P$, every automorphism $m:P\to P$, and the
multiplication diagrams.  Uniqueness of lifted morphisms preserves the group
law and simple transitivity.  Restriction and lifting are inverse on
isomorphism classes, which is the displayed $H^1$ bijection.  This is the
torsor subtree of Tag \texttt{0A5G}.
\end{proof}

\begin{lemma}[Constructible representation of an $H^1$ class]
\label{lemma:weil-pbc-h1-constructible-representation}
\uses{}

For an injective sheaf $I$ of $\mathbf F_\ell$-modules on $\mathbf P^1_A$,
every class
$c\in H^1(X_0,I_0)$ admits finite $\mathbf F_\ell$-modules $M_i$, finite finitely presented
morphisms $h_i:V_i\to X$, and a map
$\bigoplus_i(h_i)_*\underline{M_i}\to I$ such that $c$ is the image of a
tuple of classes in
$\bigoplus_iH^1(X_0,(h_i)_*\underline{M_i}|_{X_0})$.
\end{lemma}
\begin{proof}
\uses{lemma:weil-pbc-constructible-approximation,lemma:weil-pbc-constructible-finite-pushforward}

View $c$ as the isomorphism class of an $I_0$-torsor.  It is trivial on an
etale cover; quasi-compactness of $X_0$ selects finitely many members
$U_a\to X_0$.  Choosing local origins gives a Cech cocycle
$z=(z_{ab})$ with
\[
 z_{ab}\in I_0(U_a\times_{X_0}U_b),\qquad
 z_{bc}-z_{ac}+z_{ab}=0
\]
on every triple overlap.  There are finitely many sections and equations.
Constructible approximation therefore supplies a constructible sheaf $F$, a
map $\alpha:F\to I$, and lifts $\widetilde z_{ab}\in F_0(U_{ab})$
satisfying the same equations and mapping to $z_{ab}$.

The finite-pushforward lemma gives a monomorphism
$\iota:F\hookrightarrow F'=\bigoplus_{i=1}^r(h_i)_*\underline{M_i}$.
Injectivity of $I$ extends $\alpha$ across $\iota$ to a map
$\psi:F'\to I$.  The sheaf square and the three relevant terms of the Cech
complex form the commuting three-row diagram
\[
\begin{CD}
 F @>{\iota}>> F' @>>> \operatorname{coker}(\iota)\\
 @V{\alpha}VV @V{\psi}VV @VVV\\
 I @= I @>>> 0\\[2pt]
 @VVV @VVV @VVV\\
 \prod_a I_0(U_a) @>{d^0}>>
 \prod_{a,b} I_0(U_{ab}) @>{d^1}>>
 \prod_{a,b,c} I_0(U_{abc}).
\end{CD}
\]
Here the bottom vertical maps mean restriction followed by the indicated
Cech cochain inclusion; commutativity in degrees $0,1,2$ follows from
functoriality of restriction.  The cocycle $\iota(\widetilde z)$ defines
$c'\in H^1(X_0,F'_0)$ and $\psi_0(c')=c$.

Finite direct sums commute with $H^1$, while exact finite direct image and
Leray give
\[
 H^1(X_0,F'_0)=
 \bigoplus_{i=1}^rH^1(X_0,(h_i)_*\underline{M_i}|_{X_0})
 =\bigoplus_{i=1}^rH^1((V_i)_0,\underline{M_i}).
\]
Thus $c'$ is the required finite tuple, and the displayed map $F'\to I$
sends that tuple to the original class.
\end{proof}

\begin{lemma}[Acyclicity of injectives on $\mathbf P^1_s$]
\label{lemma:weil-pbc-p1-injective-acyclic}
\uses{}

For an injective sheaf $I$ of $\mathbf F_\ell$-modules on
$X=\mathbf P^1_A$, the restriction
$I_0$ is $\Gamma(X_0,-)$-acyclic: $H^q(X_0,I_0)=0$ for all $q>0$.
\end{lemma}
\begin{proof}
\uses{lemma:weil-pbc-injective-restriction-acyclicity,lemma:weil-pbc-h1-torsor-comparison,lemma:weil-pbc-h1-constructible-representation}

The affine-cover lemma gives vanishing for $q>1$ because $\mathbf P^1$ is
covered by two affines.  For $q=1$, apply the representation lemma.  For each
finite summand, finite-direct-image exactness and Leray identify its $H^1$ on
$X_0$ with $H^1((V_i)_0,\underline{M_i})$.  Torsor comparison lifts every
component to $H^1(V_i,\underline{M_i})$, and Leray returns a class in
$H^1(X,(h_i)_*\underline{M_i})$.  Their finite sum maps to zero in $H^1(X,I)$
because $I$ is injective.  Restriction of the commuting square shows that the
original class is zero.  The source locator is Tag \texttt{0A5H}.
\end{proof}

\begin{proposition}[Strictly henselian $\mathbf P^1$ comparison]
\label{prop:weil-pbc-p1-strict-henselian-comparison}
\uses{def:weil-pbc-strict-henselian-setup}
For every torsion sheaf $F$ on $X$, restriction induces isomorphisms
\[
 H^q(X,F)\xrightarrow{\sim}H^q(X_0,F_0)\qquad(q\geq0).
\]
\end{proposition}
\begin{proof}
\uses{lemma:weil-pbc-etale-degree-zero,lemma:weil-pbc-p1-injective-acyclic}

Choose an injective resolution $F\to I^\bullet$.  Exactness of inverse
image along the closed immersion gives an exact complex $F_0\to I_0^\bullet$.
The preceding lemma makes each $I_0^r$ acyclic, so Leray's acyclicity lemma
identifies $R\Gamma(X_0,F_0)$ with $\Gamma(X_0,I_0^\bullet)$.  Degree-zero
restriction identifies this complex with $\Gamma(X,I^\bullet)$ term by term.
Taking cohomology gives the displayed isomorphisms, and naturality follows
from the naturality of the restriction map and of injective resolutions.  The
source locators are Tags \texttt{0A0C} and \texttt{0A3S}.
\end{proof}

\begin{proposition}[Strict-henselian projective-line proper base change]
\label{prop:weil-pbc-projective-line-strict-henselian}
\uses{}

Let $A\to B$ be a local map of strictly henselian local rings which induces
an isomorphism of their separably closed residue fields.  For the Cartesian
square obtained from $\mathbf P^1_A\to\operatorname{Spec}A$, the derived
exchange map for a torsion sheaf is an isomorphism on the stalk at the closed
geometric point of $\operatorname{Spec}B$.
\end{proposition}
\begin{proof}
\uses{prop:weil-pbc-p1-strict-henselian-comparison}

Write $S=\operatorname{Spec}A$, $T=\operatorname{Spec}B$, and let
$\bar t$ and $\bar s$ be their closed geometric points under the chosen
residue-field identification.  Choose an injective resolution
$F\to I^\bullet$ on $\mathbf P^1_A$.  A strict henselian local scheme is
initial in the opposite category of pointed etale neighbourhoods of its
closed geometric point: the chosen point selects a unique section of the
connected component of any pointed etale neighbourhood.  Consequently the
pointed-neighbourhood stalk colimit has canonical restriction quasi-isomorphisms to global
cohomology, and the following rectangle is defined degree by degree:
\[
\begin{CD}
 (Rp_*F)_{\bar s} @>{(\operatorname{Ex}_{\square,F})_{\bar t}}>>
 (Rp_{T,*}F_T)_{\bar t}\\
 @V{\Phi_A}V{\sim}V @VV{\Phi_B}V\\
 R\Gamma(\mathbf P^1_A,F) @>{\operatorname{res}_{A,B}}>>
 R\Gamma(\mathbf P^1_B,F_T)\\
 @V{\rho_A}V{\sim}V @VV{\rho_B}V{\sim}V\\
 R\Gamma(\mathbf P^1_{k_A},F_{k_A}) @>{\operatorname{id}}>>
 R\Gamma(\mathbf P^1_{k_B},F_{k_B}).
\end{CD}
\]
Here $\Phi_A$ and $\Phi_B$ send a germ represented on a pointed
neighbourhood to its pullback on the corresponding strict local base.  To
verify that they are quasi-isomorphisms without a stalk theorem as a leaf,
compute a cohomology class on a finite split etale hypercover.  Its
$(q+1)$-skeleton, its cocycle, and any null-homotopy contain finitely many
etale objects, face maps, sections, and equalities.  The unique pointed
sections pull this complete finite datum to the strict local base; conversely
the datum descends to one pointed neighbourhood by finite presentation.
Thus the stalk colimit and strict-local cohomology agree in every degree.

The upper square commutes because both routes apply to a representative
$s\in I^r(\mathbf P^1_A\times_SV)$ by pulling it to
$\mathbf P^1_B\times_T(V\times_ST)$ and then to the pointed limit.  This is
exactly the unit--counit formula defining the S0 exchange map.  The lower
square commutes because restriction of that same section through $B$ and
then $k_B$ equals restriction through $k_A\cong k_B$; this also proves the
compatibility with differentials of $I^\bullet$.

The maps $\rho_A$ and $\rho_B$ are quasi-isomorphisms by Proposition
\ref{prop:weil-pbc-p1-strict-henselian-comparison}, and the bottom map is the
identity under $k_A\cong k_B$.  Hence the top exchange stalk is an
isomorphism.  Notice that no assertion is made here for a nontrivial
extension of separably closed residue fields; that is precisely the S5.2
curve-invariance input.
\end{proof}

\statusnote{The unique stable finite-torsion export label is
\nolinkurl{prop:weil-proet-finite-proper-base-change}.  It is deliberately not
declared in this S0--S4 fragment: S5 (strict-henselization stalks and
separably-closed field invariance) and S6 (global assembly and torsion
devissage) must precede that declaration.  The older A3 blueprint already
occupies the stable label and must be replaced, not duplicated, at merge.}
% ===== END A3 proper base change S0--S4 =====

% ===== BEGIN A3 proper base change S5--S6 =====
% A3/S5--S6: arbitrary-base projective line and finite-torsion endpoint.
% This fragment follows tmp/proetale_a3_stacks_pbc_s0_s4.tex.

\subsubsection{S5. The projective line over an arbitrary base}

\paragraph{S5.1. Stalks through strict henselization}

\begin{definition}[Pointed etale neighbourhood category]
\label{def:weil-pbc-pointed-etale-neighbourhoods}
\uses{}
Let $\bar s:\operatorname{Spec}\Omega\to S$ be a geometric point.  Define
$\mathcal N_{\bar s}$ to have objects pairs $(V,\bar v)$ with $V\to S$
etale and $\bar v:\operatorname{Spec}\Omega\to V$ above $\bar s$.
A morphism $(V,\bar v)\to(W,\bar w)$ is an $S$-map $V\to W$ carrying
$\bar v$ to $\bar w$.  We index stalk colimits by
$\mathcal N_{\bar s}^{\mathrm{op}}$ so that restriction maps point forward.
\end{definition}

\begin{lemma}[Affine pointed neighbourhoods are cofinal]
\label{lemma:weil-pbc-affine-neighbourhood-cofinal}
\uses{def:weil-pbc-pointed-etale-neighbourhoods}
The full subcategory $\mathcal N^{\mathrm{aff}}_{\bar s}$ of affine pointed
etale neighbourhoods is initial in $\mathcal N_{\bar s}$ and hence cofinal
for the opposite-indexed stalk colimit.
\end{lemma}
\begin{proof}
\uses{} 
Given $(V,\bar v)$, choose an affine open $V_0\subseteq V$ containing the
image of $\bar v$; then $(V_0,\bar v)\to(V,\bar v)$ is a morphism.  Given two
affine refinements, their fibre product over $V$ contains the lifted geometric
point; an affine open around that point refines both.  These two facts give
nonemptiness and connectedness of every comma category required for
initiality.
\end{proof}

\begin{definition}[Strict henselization from pointed neighbourhoods]
\label{def:weil-pbc-strict-henselization-limit}
\uses{def:weil-pbc-pointed-etale-neighbourhoods}
Let
\[
 \mathcal I_{\bar s}:=
 (\mathcal N^{\mathrm{aff}}_{\bar s})^{\mathrm{op}}.
\]
For $V=\operatorname{Spec}A_V$, define
$D_{\bar s}:\mathcal I_{\bar s}\to\mathbf{CRing}$ by
$D_{\bar s}(V,\bar v)=A_V$.  An arrow
$(W,\bar w)\to(V,\bar v)$ in $\mathcal I_{\bar s}$ represented by a pointed
map $a:(V,\bar v)\to(W,\bar w)$ is sent to
$a^\#:A_W\to A_V$.  Put
\[
 A^{\mathrm{sh}}_{S,\bar s}:=\varinjlim_{\mathcal I_{\bar s}}D_{\bar s},
 \qquad
 S_{(\bar s)}:=\operatorname{Spec}A^{\mathrm{sh}}_{S,\bar s}.
\]
The affine scheme--ring anti-equivalence carries the colimit cocone to the
canonical limit cone
\[
 S_{(\bar s)}\xrightarrow{\sim}
 \varprojlim_{(V,\bar v)\in\mathcal N^{\mathrm{aff}}_{\bar s}}V.
\]
\end{definition}


\begin{lemma}[Well-definedness certificate for def:weil-pbc-strict-henselization-limit]
\label{lemma:weil-sf-well-defined-def-weil-pbc-strict-henselization-limit}
\uses{def:weil-pbc-strict-henselization-limit}
The coordinate-ring rule of Definition
\ref{def:weil-pbc-strict-henselization-limit} is a functor
$D_{\bar s}:\mathcal I_{\bar s}\to\mathbf{CRing}$.  Applying
$\operatorname{Spec}$ to its colimit cocone gives the displayed limit cone,
and affine cofinality gives the same strict-local object when nonaffine
pointed neighbourhoods are admitted.
\end{lemma}
\begin{proof}
\uses{def:weil-pbc-strict-henselization-limit,lemma:weil-pbc-affine-neighbourhood-cofinal}
Coordinate rings carry an identity pointed map to an identity ring map and
reverse composition.  Passage to
$(\mathcal N^{\mathrm{aff}}_{\bar s})^{\mathrm{op}}$ therefore makes
$D_{\bar s}$ covariant.  The fixed universe of the coefficient convention
makes this a small diagram, so its colimit exists in $\mathbf{CRing}$.

The anti-equivalence between affine schemes and commutative rings sends this
colimit, together with its structure maps, to the inverse-limit cone on the
affine pointed neighbourhoods.  The initiality statement in
Lemma~\ref{lemma:weil-pbc-affine-neighbourhood-cofinal} makes restriction to
the affine subcategory limit-cofinal and opposite-colimit-cofinal.  Hence the
same ring and affine scheme compute the indicated strict-local object.
\end{proof}
\begin{lemma}[Base change commutes with the neighbourhood limit]
\label{lemma:weil-pbc-fibre-product-limit}
\uses{def:weil-pbc-strict-henselization-limit}
For a quasi-compact and quasi-separated $f:X\to S$,
\[
 X\times_SS_{(\bar s)}\cong
 \varprojlim_{(V,\bar v)}(X\times_SV).
\]
\end{lemma}
\begin{proof}
\uses{} 
On affine charts $X_0=\operatorname{Spec}B$ over
$S_0=\operatorname{Spec}A$, the left side has coordinate ring
$B\otimes_A\varinjlim A_V$.  Tensor product commutes with filtered colimits,
so this is $\varinjlim(B\otimes_AA_V)$, the coordinate ring of the inverse
limit on the right.  Quasi-compactness gives finitely many charts and
quasi-separatedness gives quasi-compact pairwise intersections; the affine
identifications therefore glue.
\end{proof}

\begin{lemma}[Finite etale hypercover skeleta descend]
\label{lemma:weil-pbc-finite-hypercover-descent}
\uses{}

Let $X_\infty=\varprojlim_iX_i$ be as in the continuity lemma below and fix
$N\geq0$.  Every etale hypercover of $X_\infty$ admits a refinement whose
$N$-skeleton is split, has finitely many quasi-compact and quasi-separated
components in each degree, and is the pullback of an $N$-truncated etale
hypercover of some $X_j$.  A morphism between two such finite skeleta, and
an equality between two such morphisms, also descends after increasing $j$.
\end{lemma}
\begin{proof}
\uses{lemma:weil-pbc-fibre-product-limit}

Construct the refinement inductively.  In degree zero, quasi-compactness of
$X_\infty$ lets us choose a finite subcover by quasi-compact etale objects.
Suppose degrees $<r$ have been chosen.  The $r$th matching object is a finite
limit of lower-degree objects, because the boundary diagram of the simplex
$r$-simplex has finitely many faces.  It is therefore quasi-compact and
quasi-separated.  Cover it by finitely many quasi-compact etale objects and
adjoin the degenerate summands prescribed by the already chosen degeneracy
maps.  This gives a split hypercover through degree $r$.  Iterating to $N$
gives the required finite skeleton.

Every chosen component is of finite presentation over $X_\infty$, so it is
the pullback of an etale component at some stage.  There are only finitely
many components, face maps, degeneracy maps, and augmentations through
degree $N$; finite presentation descends them all to one common stage.
The simplicial identities are finitely many equalities of morphisms between
finitely presented schemes and hence hold after a further common stage.
For each $r\leq N$, the map to the matching object is an etale cover after
pullback to $X_\infty$.  Etaleness descends by a Jacobian etale algebra
presentation
\(B=A[x_1,\ldots,x_m]/(f_1,\ldots,f_m)\) with an invertible Jacobian minor.
Surjectivity also descends: the complement of the image is constructible and
quasi-compact; if it survived at every later stage, its inverse limit would
give a point outside the limiting image.  Thus, after one more stage, all
matching maps are covers.  The same finite-presentation argument descends a
morphism of skeleta.  Equality of two morphisms is the factorization through
their equalizer, a finitely presented closed condition, so it too holds at a
later stage.
\end{proof}

\begin{lemma}[Etale cohomology continuity]
\label{lemma:weil-pbc-etale-cohomology-continuity}
\uses{}

Let $(X_i)$ be a cofiltered inverse system of quasi-compact and
quasi-separated schemes with affine transition maps and limit $X_\infty$.
For a torsion sheaf $F_i$ descending to one stage, the natural map
\[
 \varinjlim_{j\geq i}H^q(X_j,F_j)\longrightarrow
 H^q(X_\infty,F_\infty)
\]
is an isomorphism.
\end{lemma}
\begin{proof}
\uses{lemma:weil-pbc-fibre-product-limit,lemma:weil-pbc-finite-hypercover-descent}

We first recall why finite hypercover data compute the required cohomology.
For an etale hypercover $U_\bullet\to Z$, put
\[
 C^r(U_\bullet,F)=\Gamma(U_r,F),\qquad
 d=\sum_{a=0}^{r+1}(-1)^a\partial_a^*.
\]
Order split hypercovers by common refinement and set
\(H^q_{\mathrm{hyp}}(Z,F)=
\varinjlim_{U_\bullet}H^q(C^\bullet(U_\bullet,F))\).
The sheaf equalizer gives $H^0_{\mathrm{hyp}}(Z,F)=\Gamma(Z,F)$.  A short
exact sequence of sheaves gives connecting maps after refining so that each
of the finitely many local lifts of a cocycle exists.  The resulting functors
form a cohomological delta functor.  They are effaceable in positive degree:
given a $q$-cocycle, cover its finitely many degree-$q$ components so that
they lift through a chosen monomorphism into an injective sheaf, then fill the
matching objects successively through degree $q+1$; on this refinement the
cocycle is a coboundary.  Hence this delta functor is the universal derived
functor of global sections.  Therefore
\[
 H^q_{\mathrm{hyp}}(Z,F)=H^q_{\mathrm{et}}(Z,F).
\]
The construction also shows that a degree-$q$ class and a witness that it is
zero use only the $(q+1)$-skeleton of one split hypercover.

Let $\alpha\in H^q(X_\infty,F_\infty)$.  Represent it by a cocycle
$c\in C^q(U_\bullet,F_\infty)$ on a split hypercover with finite qcqs
$(q+1)$-skeleton.  By the preceding lemma, this skeleton descends to some
$X_j$.  Each component of $c$ is, by the definition of the inverse-image
sheaf, locally represented by a section at a finite stage.  Refine the finite
degree-$q$ cover once and choose one stage dominating all those
representatives.  The equation $dc=0$ consists of finitely many equalities on
the degree-$(q+1)$ components; these equalities descend after one further
increase of the stage.  We obtain a cocycle $c_j$ and hence a class
$\alpha_j\in H^q(X_j,F_j)$ mapping to $\alpha$.  This proves surjectivity.

For injectivity, suppose $\alpha_j\in H^q(X_j,F_j)$ maps to zero on
$X_\infty$.  Choose a finite split hypercover and cocycle $c_j$ representing
it.  Vanishing on the limit supplies a refinement
$V_\bullet\to U_\bullet\times_{X_j}X_\infty$ and a cochain
$b\in C^{q-1}(V_\bullet,F_\infty)$ with $db=c_j|_V$.  Only the
$(q+1)$-skeleton, the refinement maps, the components of $b$, and the finite
list of equations $db=c_j|_V$ are used.  Descend that entire list to one
stage $X_k$ by the preceding lemma and the same section argument.  It gives a
null-homotopy for the image of $c_j$ on $X_k$, so $\alpha_j$ is zero in the
filtered colimit.  The comparison map is therefore injective.

All constructions commute with restriction and with the transition maps,
so the resulting bijection is the canonical map in the statement.  This
expands the continuity input of Tag \texttt{03Q6}; neither the hypercover nor
the common finite stage is suppressed.
\end{proof}

\begin{proposition}[Strict-henselization stalk formula]
\label{prop:weil-pbc-strict-henselization-stalk-formula}
\uses{def:weil-pbc-pointed-etale-neighbourhoods,def:weil-pbc-strict-henselization-limit}
For quasi-compact quasi-separated $f:X\to S$, torsion $F$, and $q\geq0$,
\[
 (R^qf_*F)_{\bar s}\cong
 H^q(X\times_SS_{(\bar s)},F|).
\]
For $K\in D^+(X_{\mathrm{et}})$ this lifts to
\[
 (Rf_*K)_{\bar s}\xrightarrow{\sim}
 R\Gamma(X\times_SS_{(\bar s)},K|)
\]
in $D^+(\mathbf{Ab})$.
\end{proposition}
\begin{proof}
\uses{lemma:weil-pbc-affine-neighbourhood-cofinal,lemma:weil-pbc-fibre-product-limit,lemma:weil-pbc-etale-cohomology-continuity}

The stalk is the colimit over $\mathcal N_{\bar s}^{\mathrm{op}}$ of
$H^q(X\times_SV,F|)$.  Replace the index by affine neighbourhoods, identify
their inverse limit by Lemma~\ref{lemma:weil-pbc-fibre-product-limit}, and
apply cohomology continuity.  For the derived statement, choose a
bounded-below injective complex $J^\bullet$ representing $K$.  The chain map
from the stalk complex $(f_*J^\bullet)_{\bar s}$ to the sections complex on
the strict henselization is obtained by restricting each representative
section.  Its map on every cohomology group is the preceding isomorphism, so
it is a quasi-isomorphism.  This expands Tag \texttt{03Q9}.
\end{proof}

\paragraph{S5.2. Proper curves under separably closed extension}

\begin{lemma}[Canonical etale descent across a universal homeomorphism]
\label{lemma:weil-pbc-universal-homeomorphism-descent-datum}
\uses{}
Let $u:X_0\to X$ be a universal homeomorphism.  For every etale
$V_0\to X_0$, the two pullbacks to $X_0\times_XX_0$ have a unique
isomorphism whose restriction along the diagonal is the identity.  This
isomorphism is a descent datum and is compatible with morphisms, identities,
and composition.
\end{lemma}
\begin{proof}
\uses{} 
A universal homeomorphism is integral, surjective, and universally
injective.  Hence the diagonal
\[
 \Delta:X_0\longrightarrow X_0\times_XX_0
\]
is a closed immersion and is surjective on points.  On affine charts it is
therefore defined by a nil ideal: if
$X=\operatorname{Spec}A$ and $X_0=\operatorname{Spec}B$, the multiplication
$B\otimes_AB\to B$ is surjective and its kernel is contained in every prime
ideal.

We first record the nil-ideal calculation used below.  If $R\to R/J$ has
$J$ nil, reduction gives an equivalence between etale $R$-algebras and
etale $R/J$-algebras.  For a finitely presented affine etale algebra this is
checked on a Jacobian presentation
\[
 ((R/J)[T_1,\ldots,T_m]/(f_1,\ldots,f_m))_g
\]
with the relevant Jacobian minor invertible.  Lift the finitely many
coefficients and $g$ to $R$.  The same minor remains invertible after
localization because an element invertible modulo a nil ideal is invertible.
Existence of a lift of a morphism follows by successively correcting the
images of the $T_i$ modulo the powers of the finitely generated nil ideal
appearing in those coefficients; the inverse Jacobian gives the correction.
The process is finite because that finitely generated subideal is nilpotent.
If two lifts agree modulo $J$, their difference lies in $J$; the same
linearized equations and inverse Jacobian force equality, one nilpotence
power at a time.  General etale schemes are glued from such affine pieces;
the uniqueness just proved makes the overlap maps and their cocycle
equations lift uniquely.

Apply this nil-invariance to $\Delta$.  The two pullbacks
$p_1^*V_0$ and $p_2^*V_0$ restrict along $\Delta$ to the same etale
$X_0$-scheme $V_0$.  Thus there is a unique isomorphism
\[
 \alpha:p_1^*V_0\xrightarrow{\sim}p_2^*V_0,\qquad
 \Delta^*\alpha=\operatorname{id}_{V_0}.
\]
On $X_0\times_XX_0\times_XX_0$, both
$p_{23}^*\alpha\circ p_{12}^*\alpha$ and $p_{13}^*\alpha$ restrict to the
identity along the small diagonal.  That diagonal is again a surjective
closed immersion defined by a nil ideal, so uniqueness gives the cocycle
equation.  If $a:V_0\to W_0$ is a morphism, the two composites
$p_2^*a\circ\alpha_V$ and $\alpha_W\circ p_1^*a$ have the same diagonal
restriction and hence are equal.  Taking $a$ to be an identity or a
composite proves the asserted identity and composition compatibilities.
\end{proof}

\begin{lemma}[Base change equivalence for universal homeomorphisms]
\label{lemma:weil-pbc-universal-homeomorphism-site-equivalence}
\uses{}
\lean{AlgebraicGeometry.Scheme.Hom.isEquivalence_etale}
\statusnote{The statement is formalised as \texttt{AlgebraicGeometry.Scheme.Hom.isEquivalence\_etale}, but its proof is \texttt{sorry} (Tag \texttt{04DZ}).}

Let $u:X_0\to X$ be a universal homeomorphism.  Base change gives an
equivalence between etale $X$-schemes and etale $X_0$-schemes.
\end{lemma}
\begin{proof}
\uses{lemma:weil-pbc-universal-homeomorphism-descent-datum}

We construct a quasi-inverse to base change.  First let
$X=\operatorname{Spec}A$ and let $V_0=\operatorname{Spec}C$ be affine etale
over an affine inverse image $\operatorname{Spec}B\subset X_0$.  The preceding
lemma gives an isomorphism between the two
$B\otimes_AB$-algebras obtained from $C$, with its cocycle equation.  Define
$D$ as the equalizer of the two resulting maps
\[
 C\rightrightarrows C\otimes_AB,
\]
where the second arrow is followed by the descent isomorphism.  On a
Jacobian etale presentation of $C$, this equalizer uses only its finitely
many generators, relations, and descent equations.  The inverse-Jacobian
lifting calculation from the preceding proof gives
\[
 B\otimes_AD\xrightarrow{\sim}C
\]
and shows that $D$ is finitely presented and etale over $A$: surjectivity
lifts each generator of $C$ after imposing its descent equation, injectivity
uses the cocycle equation on $B\otimes_AB\otimes_AB$, and the lifted
Jacobian minor is invertible.  This is the affine effective-descent
calculation, rather than an appeal to an unnamed descended algebra.

For a general etale $V_0\to X_0$, choose affine opens $W_a\subset V_0$,
each mapping into an affine inverse image in $X_0$.  Their descended affine
etale schemes are denoted $V_a\to X$.  The intersections
$W_a\cap W_b$ are open subspaces.  Because $|X_0|\to|X|$ and every one of
its base changes are homeomorphisms, each such open is the inverse image of
a unique open $V_{ab}\subset V_a$.  The two identifications of
$W_a\cap W_b$ descend uniquely to
$V_{ab}\cong V_{ba}$ by the morphism compatibility in the preceding lemma.
On triple intersections the two composites have equal pullback and hence
are equal; thus the $V_a$ glue to an etale $V\to X$ with
$V\times_XX_0\cong V_0$.  The same uniqueness shows that maps descend and
that the unit and counit of the two functors are identities after pullback.
Base change is therefore an equivalence.

\end{proof}

\begin{lemma}[Covering compatibility for universal homeomorphisms]
\label{lemma:weil-pbc-universal-homeomorphism-coverings}
\uses{}
\lean{AlgebraicGeometry.Scheme.Hom.functorPushforward_mem_smallEtaleTopology_iff}
\leanok

Let $u:X_0\to X$ be a universal homeomorphism.  A family $\{U_i\to U\}$ of
etale $X$-schemes is jointly surjective if and only if its base change along
$u$ is jointly surjective.
\end{lemma}
\begin{proof}
\uses{}
\leanok

This is the topology half, the omission in Tag \texttt{03SI}.  For a
family $\{U_i\to U\}$ of etale $X$-schemes, set
$Z=|U|\setminus\bigcup_i\operatorname{im}(|U_i|\to|U|)$.  Base change gives
a homeomorphism $|U\times_XX_0|\to|U|$, and its inverse image of $Z$ is
exactly the complement of the union of the base-changed images.  Hence
$Z=\varnothing$ if and only if that inverse image is empty.  This proves
joint surjectivity in both directions.  The identity and composition
compatibilities already proved show that this correspondence of covering
families respects refinements and composite coverings.
\end{proof}

\begin{proposition}[Topological invariance of the small etale topos]
\label{prop:weil-pbc-universal-homeomorphism-topos}
\uses{}
\lean{AlgebraicGeometry.Etale.sheafEquivalence,
  AlgebraicGeometry.Etale.isEquivalence_upperStar,
  AlgebraicGeometry.Etale.cohomologyEquiv}
\statusnote{The three declarations are formalised, but each is built on \texttt{attribute [local instance] isEquivalence\_etale} and so inherits the \texttt{sorry} in the proof of Lemma~\ref{lemma:weil-pbc-universal-homeomorphism-site-equivalence}.}

Let $u:X_0\to X$ be a universal homeomorphism.  Pullback along $u$ is an
equivalence between the small etale topos of $X$ and that of $X_0$, and for
every abelian sheaf $F$ on $X_{\mathrm{et}}$ and every $n$ it induces an
isomorphism
\[
 H^n(X_{\mathrm{et}},F)\xrightarrow{\sim}H^n(X_{0,\mathrm{et}},u^{-1}F).
\]
This is Tag \texttt{03SI}.
\end{proposition}
\begin{proof}
\uses{lemma:weil-pbc-universal-homeomorphism-site-equivalence,lemma:weil-pbc-universal-homeomorphism-coverings}

Base change is an equivalence of the underlying categories and matches the
covering families, hence is an equivalence of sites; so restriction along it
is an equivalence of the categories of sheaves, and pullback, being its left
adjoint, is a quasi-inverse.  Cohomology is the $\operatorname{Ext}$-groups
out of the constant sheaf $\mathbf Z$, and an exact equivalence carrying the
constant sheaf to the constant sheaf induces an isomorphism on those groups.
\end{proof}

\begin{definition}[The eight curve-cohomology assertions]
\label{def:weil-pbc-eight-curve-assertions}
\uses{}
Let $k$ be algebraically closed, let $X$ be separated of finite type over
$k$ with $\dim X\leq1$, and let $F$ be a torsion sheaf on $X_{\mathrm{et}}$.
Write $\mathcal C_r(X,F)$, $1\leq r\leq8$, for the following assertions:
\begin{enumerate}[label=(\arabic*)]
\item $H^q(X,F)=0$ for $q>2$;
\item $H^q(X,F)=0$ for $q>1$ when $X$ is affine;
\item $H^q(X,F)=0$ for $q>1$ when $\operatorname{char}k=p>0$ and $F$ is
  $p$-power torsion;
\item every $H^q(X,F)$ is finite when $F$ is constructible and its torsion
  is prime to $\operatorname{char}k$;
\item every $H^q(X,F)$ is finite when $X$ is proper and $F$ is
  constructible;
\item for every algebraically closed extension $k'/k$, the restriction
  $H^q(X,F)\to H^q(X_{k'},F|_{X_{k'}})$ is an isomorphism when the torsion
  of $F$ is prime to $\operatorname{char}k$;
\item the map in (6) is an isomorphism when $X$ is proper, without a
  restriction on the torsion;
\item for every open $a:W\hookrightarrow X$, the restriction
  $H^2(X,F)\to H^2(W,a^{-1}F)$ is surjective.
\end{enumerate}
\end{definition}

\begin{lemma}[Filtered cohomology continuity on a curve]
\label{lemma:weil-pbc-curve-filtered-cohomology}
\uses{}
Let $X$ be quasi-compact and quasi-separated and let $(F_\lambda)$ be a
filtered system of torsion etale sheaves.  For every $q\geq0$ the canonical
map
\[
 \varinjlim_\lambda H^q(X,F_\lambda)\longrightarrow
 H^q(X,\varinjlim_\lambda F_\lambda)
\]
is an isomorphism, naturally in $X$, the system, and pullback along a field
extension.
\end{lemma}
\begin{proof}
\uses{} 
For $q=0$, choose a finite etale cover on which a section is represented.
The section and its finitely many equalizer equations on pairwise fibre
products occur at one index, proving surjectivity; if two representatives
become equal in the colimit, the finite cover and equality witnesses occur
at one later index, proving injectivity.

For $q>0$, represent a class on a split etale hypercover whose
$(q+1)$-skeleton has finitely many quasi-compact components.  A cocycle uses
finitely many sections and face-map equations, so all occur at one index
$\lambda$.  This proves surjectivity.  If the class from $F_\lambda$ dies in
the colimit, choose a finite refinement and a degree-$(q-1)$ cochain whose
coboundary is the cocycle.  The refinement maps, cochain sections, and
coboundary equations occur at one $\mu\geq\lambda$, so the class already
dies in $H^q(X,F_\mu)$.  Pullback of the finite hypercover and all its face
maps gives the same construction over an extension field, which proves the
stated naturality.
\end{proof}

\begin{lemma}[Short-exact devissage for the eight assertions]
\label{lemma:weil-pbc-eight-curve-devissage}
\uses{def:weil-pbc-eight-curve-assertions}
Let
\[
 0\longrightarrow F_1\xrightarrow{a}F\xrightarrow{b}F_2\longrightarrow0
\]
be exact.  If all applicable assertions $\mathcal C_r(X,F_i)$ hold for
$i=1,2$, then all applicable assertions hold for $F$.
\end{lemma}
\begin{proof}
\uses{} 
Pullback to an extension field is exact, so the displayed sequence and its
base change give natural long exact cohomology sequences.

For (1), (2), and (3), in every degree above the stated bound the exact
segment
\[
 H^q(X,F_1)\xrightarrow{a_*}H^q(X,F)
 \xrightarrow{b_*}H^q(X,F_2)
\]
has zero end terms.  Hence its middle term is zero.  The hypotheses on
affineness and characteristic are unchanged, and in (3) an extension of two
$p$-power torsion sheaves is again killed by a power of $p$.

For (4) and (5), constructible sheaves form a Serre subcategory: on a common
finite stratification the stalk sequence is an exact sequence of finite
groups.  Prime-to-characteristic torsion is also preserved by extensions.
Exactness identifies $H^q(X,F)$ as an extension of
$\operatorname{im}a_*$ by $\operatorname{im}b_*$; these are respectively a
quotient of the finite group $H^q(X,F_1)$ and a subgroup of the finite group
$H^q(X,F_2)$.  Thus $H^q(X,F)$ is finite.

For (6) and (7), write the long exact sequence for $F_1,F,F_2$ above its
base-changed sequence.  Every square commutes because restriction commutes
with $a$, $b$, and the boundary map obtained from an injective resolution.
For fixed $q$, the four outer vertical arrows in the five-term portion
\[
 H^{q-1}(F_2)\to H^q(F_1)\to H^q(F)\to H^q(F_2)\to H^{q+1}(F_1)
\]
are isomorphisms; the five lemma makes the middle arrow an isomorphism.

For (8), let $c\in H^2(W,F|_W)$ and put $c_2=b_*(c)$.  Choose
$x_2\in H^2(X,F_2)$ restricting to $c_2$.  Its boundary lies in
$H^3(X,F_1)=0$ by (1), so exactness gives
$x\in H^2(X,F)$ with $b_*(x)=x_2$.  The difference
$c-x|_W$ maps to zero in $H^2(W,F_2|_W)$, hence equals $a_*(d_1)$ for some
$d_1\in H^2(W,F_1|_W)$.  Choose $y_1\in H^2(X,F_1)$ restricting to $d_1$.
Then $x+a_*(y_1)$ restricts to $c$.  This proves the eighth assertion and
completes all eight cases compressed in Tag \texttt{0A5C}.
\end{proof}

\begin{lemma}[Constructible and primary reduction on curves]
\label{lemma:weil-pbc-curve-constructible-primary-reduction}
\uses{}

For a separated finite-type curve $X$ over an algebraically closed field and
a torsion sheaf $F$, extension-invariance of $H^q(X,F)$ reduces to the case
where $X$ is reduced and $F$ is a constructible $\mathbf F_\ell$-sheaf.
\end{lemma}
\begin{proof}
\uses{lemma:weil-pbc-universal-homeomorphism-site-equivalence,lemma:weil-pbc-universal-homeomorphism-coverings,lemma:weil-pbc-curve-filtered-cohomology,lemma:weil-pbc-eight-curve-devissage}

The reduction map $r:X_{\mathrm{red}}\to X$ is a universal homeomorphism.
The site equivalence sends $F$ to $r^{-1}F$, preserves covering families, and
therefore identifies the cochain complexes computing all cohomology groups.
It also commutes with the corresponding reduction after extending $k$.

The etale topos of the noetherian scheme $X$ is coherent.  Given finitely
many local sections of $F$ on quasi-compact etale objects and finitely many
relations, their image generates a constructible subsheaf: refine by a
finite stratification on which the finitely generated stalk groups are
constant.  Ordered by inclusion, these constructible subsheaves $F_\lambda$
form a filtered system with $F=\varinjlim_\lambda F_\lambda$.  The preceding
continuity lemma gives the natural square
\[
\begin{CD}
 \varinjlim_\lambda H^q(X,F_\lambda) @>{\sim}>> H^q(X,F)\\
 @VVV @VVV\\
 \varinjlim_\lambda H^q(X_{k'},F_\lambda|) @>{\sim}>>
 H^q(X_{k'},F|).
\end{CD}
\]
Thus comparison for every constructible $F_\lambda$ implies comparison for
$F$.

Fix a constructible $F$.  A finite stratification with finite stalks gives
one exponent $N$ annihilating every stalk, hence $NF=0$.  Write
$N=\prod_{\ell\mid N}\ell^{r_\ell}$.  Bezout supplies integers
$e_\ell$ which are $1$ modulo $\ell^{r_\ell}$ and $0$ modulo the other prime
powers.  The mutually orthogonal endomorphisms $e_\ell:F\to F$ sum to the
identity and give the explicit decomposition
\[
 F\xrightarrow{\sim}\bigoplus_{\ell\mid N}F[\ell^{r_\ell}],
 \qquad s\longmapsto(e_\ell s)_\ell,
\]
whose inverse is addition.  It remains to treat one summand.

If $\ell^rF=0$, use the exact kernel--quotient sequence
\[
 0\longrightarrow F[\ell]\xrightarrow{\iota}F
 \xrightarrow{\pi}F/F[\ell]\longrightarrow0.
\]
The quotient is killed by $\ell^{r-1}$ because
$\ell^{r-1}s\in F[\ell]$.  Pullback to $X_{k'}$ preserves this kernel,
quotient, and both maps.  Induction on $r$, with the base case
$\ell F=0$, and the short-exact devissage lemma prove all applicable eight
assertions.  In particular assertion (7) gives the required proper
field-extension invariance, reduced to constructible
$\mathbf F_\ell$-module sheaves.
\end{proof}

\begin{lemma}[Dense-open localization reduction]
\label{lemma:weil-pbc-curve-dense-open-reduction}
\uses{}

Let $X$ be a reduced curve and $F$ a constructible $\mathbf F_\ell$-sheaf.
There is a dense open immersion $j:U\hookrightarrow X$ with
$U=\coprod_iU_i$ irreducible and $F|_U$ finite locally constant.  If
$i:Z=X\setminus U\hookrightarrow X$, then
\[
 0\longrightarrow j_!j^{-1}F\longrightarrow F
 \longrightarrow i_*i^{-1}F\longrightarrow0
\]
reduces extension-invariance to $j_!G$ and to sheaves on the zero-dimensional
scheme $Z$.
\end{lemma}
\begin{proof}
\uses{lemma:weil-pbc-curve-constructible-primary-reduction,lemma:weil-pbc-eight-curve-devissage}

Constructibility gives a finite stratification.  Remove the finite closed
strata and all intersections of distinct irreducible components.  The
resulting $U=\coprod_aU_a$ meets every component, each $U_a$ is irreducible,
and $G=j^{-1}F$ is finite locally constant.

Let $\epsilon:j_!G\to F$ be extension by zero of the identity of $G$, and
let $\rho:F\to i_*i^{-1}F$ be restriction to $Z$.  At a geometric point of
$U$, the stalks of $(\epsilon,\rho)$ are
$G_{\bar x}\xrightarrow{\mathrm{id}}F_{\bar x}\to0$; at a geometric point
of $Z$ they are $0\to F_{\bar z}\xrightarrow{\mathrm{id}}F_{\bar z}$.
Geometric stalks are exact and conservative, so
\[
 0\longrightarrow j_!G\xrightarrow{\epsilon}F
 \xrightarrow{\rho}i_*i^{-1}F\longrightarrow0
\]
is exact with the displayed maps.

The reduced zero-dimensional scheme $Z$ is a finite disjoint union of
$k$-points.  Exact finite direct image and Leray give
\[
 H^q(X,i_*i^{-1}F)=H^q(Z,i^{-1}F)=0\quad(q>0),\qquad
 H^0(Z,i^{-1}F)=\bigoplus_{\bar z\in Z}F_{\bar z}.
\]
The last group is finite.  The long exact sequence is therefore
\[
\begin{split}
0&\to H^0(X,j_!G)\to H^0(X,F)\to
 \bigoplus_{\bar z\in Z}F_{\bar z}\\
 &\to H^1(X,j_!G)\to H^1(X,F)\to0,
\end{split}
\]
and restriction induces $H^q(X,j_!G)\cong H^q(X,F)$ for $q>1$.

These maps separate the eight assertions as follows.  Assertions (1)--(3)
and (8) use the isomorphisms in degrees at least two, also after replacing
$X$ by an open $W$ and $U$ by $W\cap U$.  Assertions (4) and (5) use the
displayed five-term sequence and finiteness of the boundary-stalk sum.
For (6) and (7), base change preserves $Z$, its stalk groups, $\epsilon$,
$\rho$, and the boundary map; the two five-term sequences form a commuting
diagram, and the five lemma handles degrees $0,1$, while the higher-degree
maps are already identified.  Thus the eight assertions hold for $F$ if
and only if they hold for $j_!G$.

Finally $j_!G=\bigoplus_a j_{a,!}(G|_{U_a})$.  The inclusions and projections
split this finite sum, and repeated split exact sequences
$0\to E_1\to E_1\oplus E_2\to E_2\to0$ together with short-exact devissage
reduce to the individual irreducible components.
\end{proof}

\begin{lemma}[Finite $\ell$-group representation filtration]
\label{lemma:weil-pbc-finite-l-group-filtration}
\uses{}
Let $H$ be a finite $\ell$-group and $V$ a nonzero finite-dimensional
$\mathbf F_\ell[H]$-module.  Then $V^H\ne0$ and $V$ has a finite filtration
by $H$-stable subspaces whose successive quotients are the trivial
one-dimensional representation.
\end{lemma}
\begin{proof}
\uses{} 
The group algebra augmentation ideal $J\subset\mathbf F_\ell[H]$ is nilpotent:
induct on $|H|$, choose a central subgroup $C$ of order $\ell$, and use
$(c-1)^\ell=c^\ell-1=0$ together with nilpotence for $H/C$.  Choose the least
$r$ with $J^rV=0$; then $J^{r-1}V$ is nonzero and fixed by $H$.  Select a
one-dimensional fixed subspace and induct on $\dim V$ after quotienting by it.
This fills the future-reference gap in Tag \texttt{0A3R}.
\end{proof}

\begin{lemma}[Locally constant sheaf reduction]
\label{lemma:weil-pbc-curve-locally-constant-reduction}
\uses{}

For irreducible $U$ and finite locally constant $G$ of
$\mathbf F_\ell$-vector spaces, there is a finite etale
$v:V\to U$ of degree prime to $\ell$ such that $v^{-1}G$ has a filtration by
constant $\mathbf F_\ell$-sheaves, and $G$ is a direct summand of
$v_*v^{-1}G$.
\end{lemma}
\begin{proof}
\uses{lemma:weil-pbc-finite-l-group-filtration,lemma:weil-pbc-finite-exchange}

Choose a geometric point $\bar u$ of $U$.  The locally constant sheaf gives
a representation
\[
 \rho:\pi_1(U,\bar u)\longrightarrow
 \operatorname{GL}_{\mathbf F_\ell}(G_{\bar u})
\]
with finite image $Q$.  If $H\subseteq Q$ is an $\ell$-Sylow subgroup, the
inverse image $\rho^{-1}(H)$ defines a connected finite etale cover
$v:V\to U$ of degree $d=[Q:H]$, which is prime to $\ell$.  On $V$, the
monodromy factors through $H$.

The representation filtration gives
$0=V_0\subset V_1\subset\cdots\subset V_m=G_{\bar u}$ with
$V_a/V_{a-1}\cong\mathbf F_\ell$ and trivial $H$-action.  Passing from
$H$-modules to locally constant sheaves gives, for every $a$, the explicit
short exact sequence
\[
 0\longrightarrow G_{a-1}\longrightarrow G_a\longrightarrow
 \underline{\mathbf F_\ell}_V\longrightarrow0,
 \qquad G_m=v^{-1}G.
\]
All inclusions, quotient maps, and connecting maps commute with extension of
the algebraically closed ground field.

The adjunction unit and the fibrewise sum are
\[
 \eta:G\longrightarrow v_*v^{-1}G,\qquad
 \operatorname{tr}:v_*v^{-1}G\longrightarrow G.
\]
At a geometric stalk, $\eta(s)=(s,\ldots,s)$ in the direct sum indexed by
the $d$ points of the fibre, while
$\operatorname{tr}(s_1,\ldots,s_d)=\sum_as_a$.  Consequently
$\operatorname{tr}\circ\eta=d\,\operatorname{id}_G$.  Multiplication by
$d^{-1}\in\mathbf F_\ell$ defines
$r=d^{-1}\operatorname{tr}$ with $r\eta=\operatorname{id}_G$.
The stalk formulas are preserved by base change, so the direct-summand
decomposition and every filtration exact sequence are compatible with the
cohomology comparison maps.
\end{proof}

\begin{definition}[Finite-open comparison map]
\label{def:weil-pbc-finite-open-comparison}
\uses{}
Consider a Cartesian square
\[
\begin{CD}
 U @>{j_X}>> X\\
 @V{v}VV @VV{\bar v}V\\
 V @>{j_Y}>> Y
\end{CD}
\]
with $j_X,j_Y$ open immersions and $\bar v$ finite.  For an abelian sheaf
$G$ on $U_{\mathrm{et}}$, define
\[
 \beta_{\bar v,j_Y,G}:(j_Y)_!v_*G\longrightarrow
 \bar v_*(j_X)_!G
\]
as the morphism adjoint to
\[
 v_*G\xrightarrow{\sim}
 j_Y^{-1}\bar v_*(j_X)_!G.
\]
The displayed isomorphism is finite base change
$j_Y^{-1}\bar v_*\cong v_*j_X^{-1}$ followed by
$j_X^{-1}(j_X)_!G\cong G$.
\end{definition}

\begin{lemma}[Section and stalk inverse for the finite-open map]
\label{lemma:weil-pbc-finite-open-comparison-inverse}
\uses{def:weil-pbc-finite-open-comparison}
The map $\beta_{\bar v,j_Y,G}$ is an isomorphism.  Its inverse is obtained by
restricting to $U$, collecting the finite-fibre tuple on $V$, and extending
that tuple by zero to $Y$; both inverse equations hold on sections and on
geometric stalks.
\end{lemma}
\begin{proof}
\uses{} 
Let $T\to Y$ be etale and put
$T_V=T\times_YV$, $T_X=T\times_YX$, and
$T_U=T\times_YU$.  A section of $(j_Y)_!v_*G$ over $T$ is etale-locally a
finite tuple
\[
 (s_x)_{x\in T_U/T_V},\qquad s_x\in G
\]
on $T_V$, extended by zero on $T\setminus T_V$.  The section map
$\beta(T)$ sends this tuple to the identical tuple on $T_U\subset T_X$ and
extends it by zero on $T_X\setminus T_U$.  This is a section of
$(j_X)_!G$ because the support remains closed after the finite map
$T_X\to T$.  Restriction along a map $T'\to T$ pulls back every component,
so these section maps define a sheaf morphism and agree with the
adjunction definition.

Conversely, take
$a\in\Gamma(T_X,(j_X)_!G)$.  Restrict $a$ to $T_U$ and group its components
over the finite fibres of $T_U\to T_V$.  Since $a$ vanishes at every
geometric point of $T_X\setminus T_U$, the grouped tuple vanishes at every
geometric point of $T\setminus T_V$ and therefore defines a section
$\gamma(T)(a)\in\Gamma(T,(j_Y)_!v_*G)$.  Restriction preserves the grouping,
so the $\gamma(T)$ define a sheaf morphism
\[
 \gamma_{\bar v,j_Y,G}:\bar v_*(j_X)_!G\longrightarrow(j_Y)_!v_*G.
\]

At a geometric point $\bar y\to Y$, the two stalks are
\[
\begin{cases}
\displaystyle\bigoplus_{\bar x\mapsto\bar y}G_{\bar x},
   &\bar y\in V,\\[6pt]
0,&\bar y\notin V.
\end{cases}
\]
Here a point $\bar x$ above $\bar y\in V$ lies in $U$ by Cartesianity.
Both $\beta_{\bar y}$ and $\gamma_{\bar y}$ are the identity of the displayed
finite sum; outside $V$ both are the zero map of the zero group.  Hence
\[
 \gamma_{\bar v,j_Y,G}\circ\beta_{\bar v,j_Y,G}
 =\operatorname{id}_{(j_Y)_!v_*G},\qquad
 \beta_{\bar v,j_Y,G}\circ\gamma_{\bar v,j_Y,G}
 =\operatorname{id}_{\bar v_*(j_X)_!G}
\]
on all geometric stalks.  Conservativity of geometric stalks proves both
equations as sheaf maps.
\end{proof}

\begin{lemma}[Finite pushforward and extension by zero]
\label{lemma:weil-pbc-finite-pushforward-extension-zero}
\uses{def:weil-pbc-finite-open-comparison}
In a Cartesian square with $j,j'$ open immersions and $\bar v$ finite,
\[
 j_!v_*G\xrightarrow{\sim}\bar v_*j'_!G
\]
is the map of Definition \ref{def:weil-pbc-finite-open-comparison}.  It is
functorial in $G$, compatible with identity and composite finite maps, and
commutes with arbitrary base change.
\end{lemma}
\begin{proof}
\uses{lemma:weil-pbc-curve-locally-constant-reduction,lemma:weil-pbc-finite-open-comparison-inverse}

The preceding lemma proves that the specified map is an isomorphism and
gives both inverse equations.

For a morphism $\phi:G\to H$, naturality is the commutative square
\[
\begin{CD}
 j_!v_*G @>{\beta_G}>> \bar v_*j'_!G\\
 @V{j_!v_*\phi}VV @VV{\bar v_*j'_!\phi}V\\
 j_!v_*H @>{\beta_H}>> \bar v_*j'_!H.
\end{CD}
\]
At $\bar y\in V$, both routes send
$(s_{\bar x})_{\bar x\mapsto\bar y}$ to
$(\phi_{\bar x}(s_{\bar x}))_{\bar x\mapsto\bar y}$; outside $V$ all four
stalks vanish.  Thus the square commutes.  The same calculation with the
inverse maps $\gamma_G,\gamma_H$ proves naturality of the inverse.

If $\bar v=\operatorname{id}_Y$ and $v=\operatorname{id}_V$, every fibre sum
has one summand.  Both the unit--base-change construction of $\beta$ and the
section formula are the identity of $j_!G$.  This proves identity
compatibility.

For composition, let $f_1:X_1\to Y$ and $f_2:X_2\to X_1$ be finite.  Put
$U_i=X_i\times_YV$, let $g_1:U_1\to V$, $g_2:U_2\to U_1$, and denote the
open immersions by $j_0:V\to Y$, $j_i:U_i\to X_i$.  The compatibility to
prove is the outer equality in
\[
\begin{CD}
 (j_0)_!(g_1g_2)_*G
 @>{\beta_{f_1,g_{2,*}G}}>>
 f_{1,*}(j_1)_!g_{2,*}G
 @>{f_{1,*}\beta_{f_2,G}}>>
 f_{1,*}f_{2,*}(j_2)_!G\\
 @V{\beta_{f_1f_2,G}}V{\sim}V
 @. @VV{=}V\\
 (f_1f_2)_*(j_2)_!G @= (f_1f_2)_*(j_2)_!G
 @= (f_1f_2)_*(j_2)_!G .
\end{CD}
\]
At $\bar y\in V$, the upper route first groups a tuple by
$\bar x_1\mapsto\bar y$ and then by
$\bar x_2\mapsto\bar x_1$:
\[
 \bigoplus_{\bar x_1\mapsto\bar y}
 \left(\bigoplus_{\bar x_2\mapsto\bar x_1}G_{\bar x_2}\right)
 \longrightarrow
 \bigoplus_{\bar x_2\mapsto\bar y}G_{\bar x_2}.
\]
It is the canonical flattening of the doubly indexed finite sum.  The lower
route is the same flattening for the composite fibre.  Outside $V$ every
term is zero.  This proves composition compatibility; applying the reverse
grouping proves the corresponding equation for $\gamma$.

Finally let $h:\widetilde Y\to Y$ be arbitrary and put tildes on all four
schemes and maps after base change.  If $h_X,h_V,h_U$ are the projections,
open and finite base change give the vertical isomorphisms in
\[
\begin{CD}
 h^{-1}(j_Y)_!v_*G
 @>{h^{-1}\beta_{\bar v,j_Y,G}}>>
 h^{-1}\bar v_*(j_X)_!G\\
 @V{\sim}VV @VV{\sim}V\\
 (\widetilde j_Y)_!\widetilde v_*h_U^{-1}G
 @>{\beta_{\widetilde{\bar v},\widetilde j_Y,h_U^{-1}G}}>>
 \widetilde{\bar v}_*(\widetilde j_X)_!h_U^{-1}G .
\end{CD}
\]
The left vertical map is
$h^{-1}(j_Y)_!\cong(\widetilde j_Y)_!h_V^{-1}$ followed by
$h_V^{-1}v_*\cong\widetilde v_*h_U^{-1}$; the right vertical map uses the
same two comparisons in the other order.  At a geometric point of
$\widetilde V$, both routes pull each component $s_{\bar x}$ to the
corresponding component over the pulled-back geometric fibre.  Off
$\widetilde V$, both routes are zero.  The rectangle therefore commutes,
and the inverse rectangle follows from the two inverse equations.  This
supplies the functoriality omitted at Tag \texttt{0GJ1}.
\end{proof}

\begin{lemma}[Smooth prime-to-characteristic curve calculation]
\label{lemma:weil-pbc-smooth-curve-kummer-localization}
\uses{def:weil-pbc-eight-curve-assertions}
Let $k$ be algebraically closed, let $C$ be a smooth separated finite-type
scheme of dimension at most one, and let
$\ell\ne\operatorname{char}k$.  All eight applicable assertions
$\mathcal C_r(C,\underline{\mathbf F_\ell})$ hold, and every field-extension
map and open-restriction map is the one induced by pullback of etale
cocycles.
\end{lemma}
\begin{proof}
\uses{} 
Points have no higher etale cohomology, and finite disjoint unions commute
with every map under consideration.  We may therefore take $C$ to be a
smooth connected curve.  Let $\overline C$ be its smooth proper
compactification, of genus $g$, and let
$D=\overline C\setminus C$ consist of $r$ points.

Choose a primitive $\ell$th root in $k$, identifying
$\underline{\mathbf F_\ell}$ with $\mu_\ell$.  The Kummer sequence
\[
 1\longrightarrow\mu_\ell\longrightarrow\mathbf G_m
 \xrightarrow{(\cdot)^\ell}\mathbf G_m\longrightarrow1
\]
gives
\[
 H^1(\overline C,\mu_\ell)=\operatorname{Pic}(\overline C)[\ell],
 \qquad
 0\to\operatorname{Pic}(\overline C)/\ell
 \to H^2(\overline C,\mu_\ell)\to
 \operatorname{Br}(\overline C)[\ell]\to0.
\]
The degree sequence
$0\to\operatorname{Pic}^0(\overline C)(k)\to
\operatorname{Pic}(\overline C)\to\mathbf Z\to0$
and the finite etale isogeny $[\ell]$ on the $g$-dimensional Jacobian give
\[
 \operatorname{Pic}(\overline C)[\ell]\cong(\mathbf Z/\ell)^{2g},
 \qquad
 \operatorname{Pic}(\overline C)/\ell\cong\mathbf Z/\ell.
\]
The Brauer group is zero.  To expose the curve input, restriction injects it
into $\operatorname{Br}(k(\overline C))$ by the discrete-valuation residue
sequence.  The latter group is zero because $k(\overline C)$ is $C_1$:
writing the unknowns as polynomials of degree at most $N$ turns a homogeneous
equation of degree $d$ in more than $d$ variables into fewer coefficient
equations than unknown coefficients for $N$ large; a nonzero solution over
the algebraically closed field gives a zero of the original form.  Applied
to a reduced norm form, this splits every central simple algebra.  The same
residue sequence and $\ell$-cohomological dimension one of this $C_1$ field
give $H^q(\overline C,\mu_\ell)=0$ for $q>2$.  Consequently
\[
 H^0(\overline C,\mathbf F_\ell)=\mathbf F_\ell,\quad
 H^1(\overline C,\mathbf F_\ell)=(\mathbf F_\ell)^{2g},\quad
 H^2(\overline C,\mathbf F_\ell)=\mathbf F_\ell.
\]

Purity for the finite set $D$ is the local DVR calculation: a parameter
$t$ identifies the boundary class with the valuation of an
$\ell$th-root Kummer torsor.  It gives the exact localization segment
\[
\begin{aligned}
0&\longrightarrow H^1(\overline C,\mathbf F_\ell)
 \longrightarrow H^1(C,\mathbf F_\ell)
 \xrightarrow{\partial}H^0(D,\mathbf F_\ell(-1))\\
 &\xrightarrow{\sum\operatorname{tr}_x}
 H^2(\overline C,\mathbf F_\ell)
 \longrightarrow H^2(C,\mathbf F_\ell)\longrightarrow0 .
\end{aligned}
\]
The trace sends the basis vector at each $x\in D$ to the fundamental class,
so when $r>0$ it is the summation map
$(\mathbf F_\ell)^r\to\mathbf F_\ell$ and is surjective.  Thus a nonproper
smooth curve is affine,
\[
 H^1(C,\mathbf F_\ell)\cong(\mathbf F_\ell)^{2g+r-1},
 \qquad H^q(C,\mathbf F_\ell)=0\quad(q>1).
\]

Under an algebraically closed extension $k'/k$, the compactification, genus,
boundary points, Kummer sequence, valuation boundary $\partial$, trace, and
the displayed summation map all pull back to their primed counterparts.
Hence the vertical maps in the localization diagram are isomorphisms.
This proves assertions (1), (2), (4), (5), and (6), and assertion (7) in the
proper case.  For assertion (8), an open in a proper connected curve is
either the whole curve or is affine with zero $H^2$; opens in an affine
curve are affine.  The restriction $H^2(C)\to H^2(W)$ is therefore either
the identity or the unique map onto zero.  All claimed maps are the maps in
the Kummer and localization sequences, so their naturality is fixed.
\end{proof}

\begin{lemma}[Smooth characteristic-\(p\) Artin--Schreier calculation]
\label{lemma:weil-pbc-smooth-curve-artin-schreier}
\uses{def:weil-pbc-eight-curve-assertions}
Let $k$ be algebraically closed of characteristic $p>0$ and let $C$ be a
smooth separated finite-type scheme of dimension at most one.  All
applicable assertions $\mathcal C_r(C,\underline{\mathbf F_p})$ hold.
For proper $C$, invariance under algebraically closed extension is induced
by the kernel of Frobenius minus the identity on $H^1(C,\mathcal O_C)$.
\end{lemma}
\begin{proof}
\uses{} 
The Artin--Schreier sequence on the etale site is
\[
 0\longrightarrow\underline{\mathbf F_p}\longrightarrow\mathcal O_C
 \xrightarrow{F-1}\mathcal O_C\longrightarrow0.
\]
Exactness on the right is checked on an etale neighbourhood by adjoining a
root of $T^p-T-a$; its derivative is $-1$, so that neighbourhood is etale.
Etale cohomology of the quasi-coherent sheaf $\mathcal O_C$ equals its
Zariski cohomology.

If $C$ is affine, $H^q(C,\mathcal O_C)=0$ for $q>0$.  The long exact
sequence gives
\[
 H^1(C,\mathbf F_p)=
 \operatorname{coker}\bigl(F-1:\Gamma(C,\mathcal O_C)\to
 \Gamma(C,\mathcal O_C)\bigr),\qquad
 H^q(C,\mathbf F_p)=0\quad(q>1).
\]
This proves assertions (1)--(3) and (8) in the affine case; no
field-extension assertion for affine $p$-torsion is claimed.

Suppose $C$ is proper and connected.  Then
$H^0(C,\mathcal O_C)=k$, $H^q(C,\mathcal O_C)=0$ for $q>1$, and
$E=H^1(C,\mathcal O_C)$ has dimension $g$.  Since $a\mapsto a^p-a$ is
surjective on the algebraically closed field $k$, the long exact sequence
reduces to
\[
 H^1(C,\mathbf F_p)=\ker(F-1:E\to E),\qquad
 H^2(C,\mathbf F_p)=\operatorname{coker}(F-1:E\to E).
\]
Choose a basis of $E$.  The additive map $F-1$ is a morphism
$\mathbf A_k^g\to\mathbf A_k^g$ whose derivative is $-\operatorname{id}$.
Its kernel is therefore zero-dimensional etale and finite.  The image is an
open subgroup of the connected additive group $\mathbf A_k^g$, hence is the
whole group.  Thus the cokernel is zero and the kernel is a finite
$\mathbf F_p$-vector space.

Proper coherent base change identifies
$E\otimes_kk'$ with $H^1(C_{k'},\mathcal O_{C_{k'}})$ and carries Frobenius
to Frobenius.  The finite etale kernel of $F-1$ is a finite disjoint union
of copies of $\operatorname{Spec}k$; after an algebraically closed extension
it has the same geometric points.  Hence
$H^1(C,\mathbf F_p)\to H^1(C_{k'},\mathbf F_p)$ is an isomorphism, while
degree zero is $\mathbf F_p$ and all degrees above one vanish.  This proves
(1), (3), (5), (7), and (8) in the proper case, with the exact comparison
map specified by the Artin--Schreier sequence.
\end{proof}

\begin{lemma}[Normalization comparison exact sequence]
\label{lemma:weil-pbc-normalization-exact-sequence}
\uses{def:weil-pbc-finite-open-comparison}
Let $X$ be a reduced curve over an algebraically closed field,
$\nu:X^\nu\to X$ its finite normalization, $j:U\hookrightarrow X$ open,
and $j^\nu:U^\nu\hookrightarrow X^\nu$ its pullback.  For every prime
$\ell$, the unit and the finite-open comparison define a monomorphism
\[
 \alpha:j_!\underline{\mathbf F_\ell}\longrightarrow
 \nu_*j^\nu_!\underline{\mathbf F_\ell}
\]
whose cokernel $Q$ is constructible and supported on the finite nonnormal
locus.  The sequence
\[
 0\longrightarrow j_!\underline{\mathbf F_\ell}
 \xrightarrow{\alpha}\nu_*j^\nu_!\underline{\mathbf F_\ell}
 \longrightarrow Q\longrightarrow0
\]
and all its cohomology boundary maps commute with algebraically closed field
extension and restriction to opens.
\end{lemma}
\begin{proof}
\uses{lemma:weil-pbc-finite-pushforward-extension-zero,lemma:weil-pbc-eight-curve-devissage}

Let $v:U^\nu\to U$.  Define $\alpha$ as
\[
 j_!\underline{\mathbf F_\ell}
 \xrightarrow{j_!\eta}
 j_!v_*\underline{\mathbf F_\ell}
 \xrightarrow{\beta_{\nu,j}}
 \nu_*j^\nu_!\underline{\mathbf F_\ell},
\]
where $\eta$ is the unit of $v^{-1}\dashv v_*$ and $\beta$ is the comparison
map already constructed.  At $\bar x\in U$, its stalk is the diagonal
\[
 \mathbf F_\ell\longrightarrow
 \bigoplus_{\bar z\in\nu^{-1}(\bar x)}\mathbf F_\ell,\qquad
 a\longmapsto(a)_{\bar z};
\]
outside $U$, both stalks are zero.  Normalization is surjective, so the
diagonal is injective.  Over the dense open where $\nu$ is an isomorphism it
is an isomorphism.  Therefore the cokernel $Q$ is supported on the finite
nonnormal locus, with stalk the quotient of the displayed finite sum by its
diagonal.  This proves exactness and constructibility.

After an algebraically closed extension, $X^\nu$ remains smooth and normal
and is the normalization of the reduced base-changed curve.  The unit stalk
map remains the diagonal, while the base-change rectangle for $\beta$
identifies the two pulled-back comparison maps.  Thus the short exact
sequence pulls back to the primed sequence.  Naturality of the connecting
homomorphism follows by applying compatible injective resolutions.  For an
open $W\subset X$, finite base change replaces $\nu$ by
$\nu^{-1}W\to W$ and gives the same stalk diagonals, proving compatibility
with open restriction.

Finally $H^q(X,Q)=0$ for $q>0$, and
$H^0(X,Q)$ is the finite direct sum of its stalks; both statements persist
after base change.  Exact finite direct image and Leray identify
$H^q(X,\nu_*j^\nu_!\mathbf F_\ell)$ with
$H^q(X^\nu,j^\nu_!\mathbf F_\ell)$.  Hence the displayed long exact
sequence transfers vanishing, finiteness, field-extension isomorphisms, and
open-restriction surjectivity from the normalized term to
$j_!\mathbf F_\ell$: in degrees at least two the maps are isomorphisms, and
in degrees zero and one the finite sequence
\[
0\to H^0(j_!\mathbf F_\ell)\to H^0(\nu_*j^\nu_!\mathbf F_\ell)
\to H^0(Q)\to H^1(j_!\mathbf F_\ell)\to
H^1(\nu_*j^\nu_!\mathbf F_\ell)\to0
\]
gives the comparison by a five-lemma diagram.
\end{proof}

\begin{lemma}[Passage between separably closed fields]
\label{lemma:weil-pbc-curve-separably-closed-passage}
\uses{}

Let $k\subseteq k'$ be separably closed fields, let $C$ be proper of
dimension at most one over $k$, and let $F$ be a torsion etale sheaf.  If
proper-curve cohomology is invariant for algebraically closed extensions,
then the specified restriction
$H^q(C,F)\to H^q(C_{k'},F|)$ is an isomorphism for every $q$.
\end{lemma}
\begin{proof}
\uses{lemma:weil-pbc-universal-homeomorphism-site-equivalence,lemma:weil-pbc-universal-homeomorphism-coverings}

Choose compatible algebraic closures $\bar k\subseteq\bar k'$.  Since $k$
and $k'$ are separably closed, $\bar k/k$ and $\bar k'/k'$ are purely
inseparable.  Their spectrum maps are universal homeomorphisms; the explicit
site equivalence and covering compatibility give the vertical isomorphisms
in
\[
\begin{CD}
 H^q(C,F) @>>> H^q(C_{k'},F|)\\
 @V{\sim}VV @VV{\sim}V\\
 H^q(C_{\bar k},F|) @>>> H^q(C_{\bar k'},F|).
\end{CD}
\]
The lower arrow is an isomorphism by the algebraically closed case.  The
square commutes because each route pulls a representative cocycle through
the same composite field map; identity and composition compatibility of the
site equivalence identify the two pullbacks.  The upper arrow is therefore
an isomorphism.  In particular spectra of separably closed fields, which are
strictly henselian local schemes, introduce no additional stalk ambiguity.
\end{proof}

\begin{lemma}[Normalization and smooth constant case]
\label{lemma:weil-pbc-curve-normalization-constant}
\uses{}

For a reduced curve, extension-invariance for $j_!\underline{\mathbf F_\ell}$
reduces through the finite normalization to the corresponding calculation on
a smooth curve, both for proper and affine opens and also when
$\ell=\operatorname{char}k$.
\end{lemma}
\begin{proof}
\uses{lemma:weil-pbc-smooth-curve-kummer-localization,lemma:weil-pbc-smooth-curve-artin-schreier,lemma:weil-pbc-normalization-exact-sequence,lemma:weil-pbc-curve-dense-open-reduction}

The normalization $X^\nu$ is a finite disjoint union of smooth curves and
points.  On a component disjoint from $U^\nu$ the extension-by-zero sheaf is
zero and every assertion is immediate; on every other component $U^\nu$ is
dense.  The dense-open localization lemma on $X^\nu$ therefore reduces
$j^\nu_!\underline{\mathbf F_\ell}$ to the constant sheaf on each remaining
smooth component, with the five-term sequence and all comparison maps
already specified there.  If $\ell\ne\operatorname{char}k$, apply the Kummer and
localization calculation; if $\ell=\operatorname{char}k$, apply the
Artin--Schreier calculation.  Thus all applicable eight assertions hold for
$j^\nu_!\underline{\mathbf F_\ell}$.

Exact finite direct image transfers them to
$\nu_*j^\nu_!\underline{\mathbf F_\ell}$.  The normalization exact sequence
then transfers them to $j_!\underline{\mathbf F_\ell}$, including the
low-degree comparison maps and assertion (8).  Properness is preserved by
finite normalization; affineness is preserved because a finite morphism
over an affine is affine.  This completes the normalization, proper/affine
localization, prime-to-characteristic, and characteristic-$\ell$ branches
without an unnamed smooth-curve calculation.
\end{proof}

\begin{proposition}[Proper curve invariance over separably closed fields]
\label{prop:weil-pbc-proper-curve-sepclosed-invariance}
\uses{}

Let $k\subseteq k'$ be separably closed fields, $C$ a proper $k$-scheme of
dimension at most one, and $F$ a torsion etale sheaf.  Then
\[
 H^q(C,F)\xrightarrow{\sim}H^q(C_{k'},F|)
\]
for every $q$.
\end{proposition}
\begin{proof}
\uses{lemma:weil-pbc-universal-homeomorphism-site-equivalence,lemma:weil-pbc-universal-homeomorphism-coverings,lemma:weil-pbc-curve-constructible-primary-reduction,lemma:weil-pbc-curve-dense-open-reduction,lemma:weil-pbc-curve-locally-constant-reduction,lemma:weil-pbc-finite-pushforward-extension-zero,lemma:weil-pbc-curve-normalization-constant,lemma:weil-pbc-curve-separably-closed-passage}

First assume both fields algebraically closed.  Apply constructible/primary
reduction.  The closed immersion $C_{\mathrm{red}}\to C$ is a universal
homeomorphism, so the explicit site equivalence first replaces $C$ by its
reduction.  Then apply dense-open localization, the prime-to-$\ell$ trace
splitting, the
finite-pushforward comparison, and the normalized smooth constant calculation
in that order.  Every step commutes with field extension, giving the result.
For separably closed $k\subseteq k'$, apply Lemma
\ref{lemma:weil-pbc-curve-separably-closed-passage} to the algebraically
closed result just proved.  Its commutative rectangle identifies the
resulting isomorphism with the specified restriction map.  This expands Tags
\texttt{03SC} and \texttt{0A5E}.
\end{proof}

\paragraph{S5.3. Stalkwise projective-line base change}

\begin{lemma}[Exchange on a geometric stalk]
\label{lemma:weil-pbc-p1-exchange-geometric-stalk}
\uses{def:weil-pbc-derived-exchange}
For $p:\mathbf P^1_S\to S$, a base change $g:T\to S$, and
$\bar t\to T$ with image $\bar s\to S$, the stalk of
$\operatorname{Ex}_{\square,F}$ is the restriction map
\[
 H^q(\mathbf P^1_{\kappa(\bar s)},F|)
 \longrightarrow H^q(\mathbf P^1_{\kappa(\bar t)},F|).
\]
\end{lemma}
\begin{proof}
\uses{prop:weil-pbc-strict-henselization-stalk-formula,prop:weil-pbc-p1-strict-henselian-comparison}

Let $\bar t\to T$ be represented by an algebraically closed field $\Omega$,
and let $\bar s\to S$ be its composite.  Put
$S^{\mathrm{sh}}=\operatorname{Spec}\mathcal O^{\mathrm{sh}}_{S,\bar s}$ and
$T^{\mathrm{sh}}=\operatorname{Spec}\mathcal O^{\mathrm{sh}}_{T,\bar t}$.
The map $T\to S$ induces a local map $T^{\mathrm{sh}}\to S^{\mathrm{sh}}$
and a separably closed residue-field extension
$\kappa(\bar s)\hookrightarrow\kappa(\bar t)$.  For an injective
resolution $F\to I^\bullet$, the strict-henselization construction gives
the following diagram in $D^+(\mathbf{Ab})$:
\[
\begin{CD}
 (g^{-1}Rp_*F)_{\bar t}
   @>{(\operatorname{Ex}_{\square,F})_{\bar t}}>>
 (Rp'_*g'^{-1}F)_{\bar t}\\
 @V{\Phi_S}V{\sim}V @VV{\Phi_T}V{\sim}V\\
 R\Gamma(\mathbf P^1_{S^{\mathrm{sh}}},I|)
   @>{\lambda}>>
 R\Gamma(\mathbf P^1_{T^{\mathrm{sh}}},I|)\\
 @V{\rho_S}V{\sim}V @VV{\rho_T}V{\sim}V\\
 R\Gamma(\mathbf P^1_{\kappa(\bar s)},I_{\bar s})
   @>{\kappa(\bar s)\subset\kappa(\bar t)}>>
 R\Gamma(\mathbf P^1_{\kappa(\bar t)},I_{\bar t}).
\end{CD}
\]
The maps $\Phi_S$ and $\Phi_T$ are the stalk-to-strict-local maps from the
strict-henselization formula.  The maps $\rho_S$ and $\rho_T$ are the S4
restriction quasi-isomorphisms, applied after replacing $\mathbf P^1$ by its
base change to the indicated strict local ring.  The bottom arrow is induced
by the residue-field map.

To check the upper square, take a representative section of $I^r$ on an
affine pointed etale neighbourhood $V\to S$.  The left route first applies
the unit $I^r\to g'_*g'^{-1}I^r$ and then restricts along
$\mathbf P^1_T\times_T(V\times_ST)$.  The right route first restricts to
$\mathbf P^1_{S^{\mathrm{sh}}}$ and then along $T^{\mathrm{sh}}\to
S^{\mathrm{sh}}$.  Both are the same section on
$\mathbf P^1_{T^{\mathrm{sh}}}\times_{T^{\mathrm{sh}}}V$, so the square
commutes degree by degree and hence in the derived category.  The lower
square is the identical calculation after passing to the closed geometric
fibres.  Therefore the stalk of the named S0 exchange is exactly the field
extension map displayed in the bottom row, not an abstract replacement.
\end{proof}

\begin{lemma}[Conservativity of geometric stalks]
\label{lemma:weil-pbc-geometric-stalk-conservative}
\uses{}
\lean{AlgebraicGeometry.Scheme.isIso_of_isIso_sheafFiber_map}
\leanok
A morphism of etale abelian sheaves is an isomorphism if and only if it is an
isomorphism at every geometric stalk.
\end{lemma}
\begin{proof}
\uses{} 
The kernel and cokernel have zero geometric stalks.  If a sheaf has zero
stalks, every local section vanishes on an etale neighbourhood of each point;
quasi-compact refinement on an etale object makes it locally zero on a cover,
hence zero by the sheaf axiom.  Thus kernel and cokernel vanish.
\end{proof}

\begin{proposition}[Projective-line proper base change]
\label{prop:weil-pbc-projective-line}
\uses{}

For every scheme $S$, the projection $\mathbf P^1_S\to S$ satisfies
$\mathsf{PBC}$.
\end{proposition}
\begin{proof}
\uses{lemma:weil-pbc-p1-exchange-geometric-stalk,prop:weil-pbc-proper-curve-sepclosed-invariance,lemma:weil-pbc-geometric-stalk-conservative,lemma:weil-pbc-property-iff-acyclicity}

At every geometric $\bar t\to T$, the exchange map is the field-extension
map of the preceding lemma.  The two residue fields are separably closed and
the fibre is a proper curve, so proper-curve invariance makes this map an
isomorphism.  Geometric stalk conservativity gives an isomorphism of every
$R^q$ sheaf.  Specializing to inverse images of injective
$\mathbf F_\ell$-module sheaves gives $\mathsf{IA}$ after every base change;
the S0 equivalence gives $\mathsf{PBC}$.
\end{proof}

\subsubsection{S6. Finite-torsion proper base change}

\begin{lemma}[Prime-coefficient proper base change]
\label{lemma:weil-pbc-prime-proper-base-change}
\uses{}

Every proper morphism satisfies proper base change for sheaves of
$\mathbf F_\ell$-modules.
\end{lemma}
\begin{proof}
\uses{prop:weil-pbc-reduction-to-projective-line,prop:weil-pbc-projective-line}

The projective-line proposition supplies the sole geometric input required by
the S3 reduction.  Finite morphisms, composition, and cancellation in that
reduction preserve the named S0 exchange map, so the resulting isomorphism is
the exchange map for the original proper morphism.
\end{proof}

\begin{lemma}[Torsion coefficient assembly]
\label{lemma:weil-pbc-torsion-assembly}
\uses{}

For every torsion abelian sheaf $F$ and every proper Cartesian square, the
specific map $\operatorname{Ex}_{\square,F}$ is an isomorphism.
\end{lemma}
\begin{proof}
\uses{lemma:weil-pbc-prime-proper-base-change,lemma:weil-pbc-torsion-filtered-union,lemma:weil-pbc-direct-image-filtered-colimit,lemma:weil-pbc-prime-factor-induction,lemma:weil-pbc-five-lemma-transfer}

Write $F=\varinjlim_nF[n]$ and commute both sides of exchange with the
filtered colimit.  Induct on the number of prime factors of $n$ through
$0\to F[\ell]\to F[n]\to F[n]/F[\ell]\to0$; the quotient is killed by
$n/\ell$.  Prime coefficients are handled by the preceding lemma, and the
five-lemma transfer handles the induction step.  Naturality shows at every
stage that the comparison is the S0 exchange map.
\end{proof}

\begin{proposition}[Finite-torsion proper base change]
\label{prop:weil-proet-finite-proper-base-change}
\uses{def:weil-pbc-derived-exchange}
Let
\[
\begin{CD}
X'@>{g'}>>X\\
@V{f'}VV @VV{f}V\\
Y'@>{g}>>Y
\end{CD}
\]
be Cartesian with $f$ proper.  For every torsion sheaf $F$ on
$X_{\mathrm{et}}$, the canonical exchange map
\[
 g^{-1}Rf_*F\xrightarrow{\sim}Rf'_*g'^{-1}F
\]
is an isomorphism.  These isomorphisms are natural in $F$, equal the identity
for identity squares, and commute with horizontal and vertical pasting.
\end{proposition}
\begin{proof}
\uses{lemma:weil-pbc-derived-exchange-coherence,lemma:weil-pbc-torsion-assembly}

The torsion assembly proves the isomorphism in every cohomological degree for
the map constructed in S0.  A morphism in $D^+$ is an isomorphism exactly when
all its cohomology sheaf maps are isomorphisms.  Naturality, identity, and
pasting are the corresponding S0 coherence equations and survive each
devissage step.  This is the expanded endpoint whose source locator is Tag
\texttt{095T}; that tag is not used as a proof leaf.
\end{proof}

\statusnote{Proposition~\ref{prop:weil-proet-finite-proper-base-change} is the
stable A3 finite-torsion export.  It remains conditional until every OPEN in
S0--S5 on which its proof depends is discharged.}
% ===== END A3 proper base change S5--S6 =====

% ===== BEGIN A3 proper base change S7--S10 =====
% A3/S7--S10: bounded complexes, pro-etale transport, completion, and PBC.

\subsubsection{S7. Bounded-below finite-torsion complexes}

\begin{definition}[Hypercohomology filtrations]
\label{def:weil-pbc-hypercohomology-filtrations}
\uses{def:weil-pbc-derived-exchange}
Let $K\in D^+(X_{\mathrm{et}})$ have torsion cohomology sheaves and choose
$a$ with $\mathcal H^q(K)=0$ for $q<a$.  The canonical truncation tower
$\tau_{\leq q}K$ defines spectral sequences
\[
 E_2^{p,q}=R^pf_*\mathcal H^q(K)\Longrightarrow
 \mathcal H^{p+q}(Rf_*K),
\]
and
\[
 {E'}_2^{p,q}=R^pf'_*g'^{-1}\mathcal H^q(K)\Longrightarrow
 \mathcal H^{p+q}(Rf'_*g'^{-1}K).
\]
Their filtration in total degree $n$ is the image filtration induced by
$\mathcal H^n(Rf_*\tau_{\leq q}K)\to\mathcal H^n(Rf_*K)$; the second
filtration is induced by
$\mathcal H^n(Rf'_*g'^{-1}\tau_{\leq q}K)\to
\mathcal H^n(Rf'_*g'^{-1}K)$.
\end{definition}

\begin{lemma}[Half-plane convergence]
\label{lemma:weil-pbc-hypercohomology-convergence}
\uses{def:weil-pbc-hypercohomology-filtrations}
The spectral sequences lie in $p\geq0$, $q\geq a$.  For fixed total degree
$n$, only $a\leq q\leq n$ and $0\leq p\leq n-a$ occur; hence the induced
filtration is finite, exhaustive, separated, and converges to the displayed
abutment.
\end{lemma}
\begin{proof}
\uses{} 
Right-derived direct image has no negative cohomological degree, giving
$p\geq0$.  The lower bound gives $q\geq a$.  On $p+q=n$ these inequalities
give the finite interval.  The truncation triangles show exhaustivity, and
finiteness shows separatedness; no inverse-limit convergence issue occurs.
\end{proof}

\begin{definition}[The two truncation exact couples]
\label{def:weil-pbc-truncation-exact-couples}
\uses{def:weil-pbc-hypercohomology-filtrations,def:weil-pbc-derived-exchange}
For every $q\geq a$, set
\[
\begin{aligned}
A_q&=g^{-1}Rf_*\tau_{\leq q}K,&
G_q&=g^{-1}Rf_*\mathcal H^q(K)[-q],\\
B_q&=Rf'_*g'^{-1}\tau_{\leq q}K,&
H_q&=Rf'_*g'^{-1}\mathcal H^q(K)[-q].
\end{aligned}
\]
The canonical truncation triangle and its exchanged image are
\[
\begin{CD}
A_{q-1}@>{i_q}>>A_q@>{j_q}>>G_q@>{k_q}>>A_{q-1}[1]\\
@V{e_{q-1}}VV@V{e_q}VV@V{\bar e_q}VV@VV{e_{q-1}[1]}V\\
B_{q-1}@>{i'_q}>>B_q@>{j'_q}>>H_q@>{k'_q}>>B_{q-1}[1].
\end{CD}
\]
For total degree $n$, define
\[
\begin{aligned}
D_A^{q,n}&=\mathcal H^n(A_q),&
\mathcal E_A^{p,q}&=\mathcal H^{p+q}(G_q)
 =g^{-1}R^pf_*\mathcal H^q(K),\\
D_B^{q,n}&=\mathcal H^n(B_q),&
\mathcal E_B^{p,q}&=\mathcal H^{p+q}(H_q)
 =R^pf'_*g'^{-1}\mathcal H^q(K).
\end{aligned}
\]
The exact-couple differential is
\[
d_1=(j_{q-1})_*\circ(k_q)_*:
\mathcal E_A^{p,q}\longrightarrow\mathcal E_A^{p+2,q-1},
\]
and likewise $d'_1=(j'_{q-1})_*(k'_q)_*$.  This first
exact-couple page is the page conventionally denoted $E_2$ in the
hypercohomology indexing.
\end{definition}

\begin{lemma}[Hypercover model and \(E_1\)-map]
\label{lemma:weil-pbc-hypercover-exact-couple-map}
\uses{def:weil-pbc-truncation-exact-couples}
On every split etale hypercover, the exchange maps define a morphism between
the two exact couples.  Its first exact-couple page is
\[
\bar e_1^{p,q}:
g^{-1}R^pf_*\mathcal H^q(K)\longrightarrow
R^pf'_*g'^{-1}\mathcal H^q(K),
\]
the S6 finite-torsion map, and
$d'_1\bar e_1^{p,q}=\bar e_1^{p+2,q-1}d_1$ with no residual sign.
\end{lemma}
\begin{proof}
\uses{def:weil-pbc-truncation-exact-couples,prop:weil-proet-finite-proper-base-change}

Evaluate the two rows of truncation triangles on an etale object and choose
a split etale hypercover $U_\bullet$.  A Cartan--Eilenberg representative
gives on each side a double complex $C^{m,r}$, where $m$ is the cosimplicial
degree and $r$ the complex degree.  Write
\[
d_{\mathrm{simp}}=\sum_{t=0}^{m+1}(-1)^t d_t^*,\qquad
d_{\mathrm{tot}}=d_{\mathrm{simp}}+(-1)^m d_{\mathrm{int}}.
\]
The simplicial identities give $d_{\mathrm{simp}}^2=0$; naturality of the
face maps gives
$d_{\mathrm{simp}}d_{\mathrm{int}}=
d_{\mathrm{int}}d_{\mathrm{simp}}$, and the factor $(-1)^m$ makes
$d_{\mathrm{tot}}^2=0$.  The left and right totalizations use this same
sign convention.

The sectionwise exchange map pulls a section on every $U_m$ through the
Cartesian square.  It commutes with every face and degeneracy map and with
$d_{\mathrm{int}}$, hence is a chain map of the two total complexes.
Filtering by canonical truncation gives exactly the objects
$A_q,G_q$ and $B_q,H_q$ of Definition
\ref{def:weil-pbc-truncation-exact-couples}.  Naturality with respect to
$\tau_{\leq q-1}K\to\tau_{\leq q}K$ gives the first three squares of the
displayed triangle.

For the last square, use the cohomological cone convention
\[
\operatorname{Cone}(u)^n=Y^n\oplus X^{n+1},\qquad
d(y,x)=(d_Yy+u(x),-d_Xx).
\]
The connecting map is induced by $(y,x)\mapsto x$.  Applying exchange
before or after this projection gives $e_{q-1}[1](x)$ in both routes; the
minus sign belongs to the shifted differential on both sides and does not
change the map.  Thus
$k'_q\bar e_q=e_{q-1}[1]k_q$, with no extra sign.

Taking cohomology yields a map of exact couples.  On the graded object
$\mathcal H^q(K)[-q]$, degree $p+q$ cohomology is
$g^{-1}R^pf_*\mathcal H^q(K)$ on the left and
$R^pf'_*g'^{-1}\mathcal H^q(K)$ on the right.  The induced map is the S0
exchange for that torsion sheaf and therefore the specified S6 map.  Finally
\[
\begin{aligned}
d'_1\bar e_1
 &=(j'_{q-1})_*(k'_q)_*\bar e_1\\
 &=(j'_{q-1})_*(e_{q-1}[1])_*(k_q)_*
 =\bar e_1(j_{q-1})_*(k_q)_*
 =\bar e_1d_1,
\end{aligned}
\]
which proves differential compatibility.
\end{proof}

\begin{lemma}[Finite rows and the abutment comparison]
\label{lemma:weil-pbc-exact-couple-abutment}
\uses{def:weil-pbc-hypercohomology-filtrations,def:weil-pbc-truncation-exact-couples}
In total degree $n$, the two exact couples have only the finite row range
$a\leq q\leq n$.  Their induced abutment map is the cohomology map
\[
\mathcal H^n(g^{-1}Rf_*K)\longrightarrow
\mathcal H^n(Rf'_*g'^{-1}K)
\]
of the derived exchange morphism, with the image filtrations of Definition
\ref{def:weil-pbc-hypercohomology-filtrations}.
\end{lemma}
\begin{proof}
\uses{lemma:weil-pbc-hypercohomology-convergence}

On $\mathcal E_1^{p,q}$ one has $p\geq0$ and $q\geq a$.  The equality
$p+q=n$ therefore forces $a\leq q\leq n$.  On a hypercover total complex,
only these finitely many filtered quotients contribute to degree $n$.
Consequently direct-sum and product totalization agree in that degree, the
filtration is finite and separated, and passage from a finite skeleton to
the full hypercover changes none of its graded terms.

The maps $e_q:A_q\to B_q$ commute with
$\tau_{\leq q}K\to K$.  For $q\geq n$, the omitted cohomology sheaves of
$K$ have degrees greater than $n$; after a nonnegative right-derived
functor they cannot contribute to total degree $n$.  Hence
\[
\mathcal H^n(A_q)\xrightarrow{\sim}
\mathcal H^n(g^{-1}Rf_*K),\qquad
\mathcal H^n(B_q)\xrightarrow{\sim}
\mathcal H^n(Rf'_*g'^{-1}K).
\]
Under these identifications, $\mathcal H^n(e_q)$ is
$\mathcal H^n(\operatorname{Ex}_{\square,K})$ by naturality of exchange in
the truncation map.  Its filtration is the image of
$\mathcal H^n(A_q)$ and maps to the image of $\mathcal H^n(B_q)$, exactly
the two filtrations defined earlier.  Thus the exact-couple abutment map is
the specified derived exchange map, rather than an unnamed filtered
isomorphism.
\end{proof}

\begin{lemma}[Morphism of hypercohomology spectral sequences]
\label{lemma:weil-pbc-hypercohomology-map}
\uses{def:weil-pbc-hypercohomology-filtrations,def:weil-pbc-truncation-exact-couples}
The S0 exchange map induces a morphism from the inverse image of the first
spectral sequence to the second, whose $E_2$ component is
\[
 g^{-1}R^pf_*\mathcal H^q(K)\longrightarrow
 R^pf'_*g'^{-1}\mathcal H^q(K),
\]
the S6 finite-torsion proper-base-change map.
\end{lemma}
\begin{proof}
\uses{def:weil-pbc-truncation-exact-couples,lemma:weil-pbc-hypercohomology-convergence,lemma:weil-pbc-hypercover-exact-couple-map,lemma:weil-pbc-exact-couple-abutment,prop:weil-proet-finite-proper-base-change}

Definition \ref{def:weil-pbc-truncation-exact-couples} displays both exact
couples and every structure map.  Lemma
\ref{lemma:weil-pbc-hypercover-exact-couple-map} constructs their morphism,
identifies its first page with the S6 map, and checks the differential and
connecting-map sign.  Lemma \ref{lemma:weil-pbc-exact-couple-abutment}
proves finite-row convergence and identifies the abutment map with
$\mathcal H^n(\operatorname{Ex}_{\square,K})$.  Deriving the exact couples
therefore gives the asserted morphism on every page.
\end{proof}

\begin{proposition}[Bounded-below finite derived proper base change]
\label{prop:weil-proet-finite-derived-proper-base-change}
\uses{}

For a proper Cartesian square and $K\in D^+(X_{\mathrm{et}})$ with torsion
cohomology, the exchange map
\[
 g^{-1}Rf_*K\xrightarrow{\sim}Rf'_*g'^{-1}K
\]
is an isomorphism.
\end{proposition}
\begin{proof}
\uses{lemma:weil-pbc-hypercohomology-convergence,lemma:weil-pbc-hypercohomology-map,prop:weil-proet-finite-proper-base-change}

S6 makes every $E_2$ component of the spectral-sequence morphism an
isomorphism.  Induction on the page gives isomorphisms on $E_r$ and
$E_\infty$.  In each total degree the finite filtrations have isomorphic
associated graded pieces, so induction on the filtration length gives an
isomorphism of abutments.  A morphism in the derived category is an
isomorphism when every cohomology sheaf map is an isomorphism.  This expands
Tag \texttt{0DDE}, including its exact-couple comparison.
\end{proof}

\subsubsection{S8. Etale-to-pro-etale finite-level transport}

\begin{definition}[Etale objects inside the pro-etale site]
\label{def:weil-pbc-etale-proetale-inclusion}
\uses{def:weil-small-etale-topology,def:weil-small-etale-pretopology,def:weil-small-proetale-category,def:weil-small-proetale-precoverage,def:weil-small-proetale-topology}
For a scheme $X$, define the assignment
\[
 \iota_X:X_{\mathrm{et}}\longrightarrow X_{\mathrm{proet}}
\]
as follows.  An object $p:U\to X$, etale in the A2 small etale category, is
sent to the same structural morphism $p$.  It is an object of the A2 small
pro-etale category because an etale morphism is weakly etale.  An
$X$-morphism $h:U\to V$ is sent to the same scheme morphism $h$.  Thus both
forgetful functors to $\operatorname{Over}(X)$ agree on this assignment.
\end{definition}

% Atomic leaves for EFI SITE-02 through SITE-04.

\begin{lemma}[Identity law for the etale-to-pro-etale assignment]
\label{lemma:weil-efi-site-inclusion-identity}
\uses{def:weil-pbc-etale-proetale-inclusion}
For every etale object $U\to X$,
\[
 \iota_X(1_U)=1_{\iota_XU}.
\]
\end{lemma}
\begin{proof}
\uses{def:weil-pbc-etale-proetale-inclusion}
Both arrows have underlying scheme morphism $1_U$, and morphisms in the
selected over-category are determined by their underlying scheme maps.
\end{proof}

\begin{lemma}[Composition law for the etale-to-pro-etale assignment]
\label{lemma:weil-efi-site-inclusion-composition}
\uses{def:weil-pbc-etale-proetale-inclusion}
For composable etale-site arrows $U\xrightarrow{h}V\xrightarrow{k}W$,
\[
 \iota_X(kh)=\iota_X(k)\iota_X(h).
\]
\end{lemma}
\begin{proof}
\uses{def:weil-pbc-etale-proetale-inclusion}
Each side has underlying scheme morphism $kh$.  Equality follows from the
faithful forgetful functor to $\operatorname{Over}(X)$.
\end{proof}

\begin{definition}[Functor-law certificate for the site assignment]
\label{def:weil-efi-site-inclusion-functor-certificate}
\uses{def:weil-pbc-etale-proetale-inclusion}
Let $\mathsf{FunctorLaw}(\iota_X)$ be the conjunction of the identity equation
for every etale object and the composition equation for every composable pair
of etale-site arrows.
\end{definition}

\begin{proposition}[The site assignment satisfies its functor laws]
\label{prop:weil-efi-site-inclusion-functor}
\uses{def:weil-efi-site-inclusion-functor-certificate,
  def:weil-pbc-etale-proetale-inclusion}
The predicate $\mathsf{FunctorLaw}(\iota_X)$ holds.  Hence the assignment of
Definition \ref{def:weil-pbc-etale-proetale-inclusion} is a functor.
\end{proposition}
\begin{proof}
\uses{lemma:weil-efi-site-inclusion-identity,
  lemma:weil-efi-site-inclusion-composition}
The first lemma supplies the identity field and the second supplies the
composition field of the displayed predicate.
\end{proof}

\begin{definition}[Image of an etale sieve in the pro-etale site]
\label{def:weil-efi-site-covering-sieve-image}
\uses{def:weil-pbc-etale-proetale-inclusion}
For an etale object $U\to X$ and a sieve $R\subset h_U$, let
$\iota_{X,!}^{\mathrm{sieve}}R$ be the sieve on $\iota_XU$ generated by the
arrows $\iota_XV\to\iota_XU$ with $V\to U$ in $R$.
\end{definition}

\begin{lemma}[Etale covering families remain pro-etale covering]
\label{lemma:weil-efi-site-covering-family-image}
\uses{def:weil-pbc-etale-proetale-inclusion}
The image under $\iota_X$ of an etale covering family is a pro-etale
covering family.
\end{lemma}
\begin{proof}
\uses{def:weil-pbc-etale-proetale-inclusion,
  def:weil-small-etale-pretopology,
  def:weil-small-proetale-precoverage}
Every arrow in the family is etale and hence weakly etale.  The underlying
scheme maps are unchanged, so their joint surjectivity is unchanged.  These
are exactly the membership conditions in the selected pro-etale
precoverage.
\end{proof}

\begin{lemma}[Etale covering sieves remain pro-etale covering]
\label{lemma:weil-efi-site-covering-sieve-image}
\uses{def:weil-efi-site-covering-sieve-image}
If $R$ is covering on $U$ for the selected etale topology, then
$\iota_{X,!}^{\mathrm{sieve}}R$ is covering for the selected pro-etale
topology.
\end{lemma}
\begin{proof}
\uses{lemma:weil-efi-site-covering-family-image,
  def:weil-small-etale-topology,
  def:weil-small-proetale-topology}
The etale topology is generated by its pretopology, so $R$ contains the sieve
generated by an etale covering family.  The preceding lemma sends that family
to a pro-etale covering family.  Its generated sieve is contained in
$\iota_{X,!}^{\mathrm{sieve}}R$, and covering sieves are upward closed.
\end{proof}

\begin{lemma}[Pullback equation for the image-sieve operation]
\label{lemma:weil-efi-site-covering-sieve-pullback}
\uses{def:weil-efi-site-covering-sieve-image}
For an etale arrow $f:V\to U$ and a sieve $R$ on $U$,
\[
 (\iota_Xf)^*(\iota_{X,!}^{\mathrm{sieve}}R)
 =\iota_{X,!}^{\mathrm{sieve}}(f^*R).
\]
\end{lemma}
\begin{proof}
\uses{def:weil-efi-site-covering-sieve-image,
  def:weil-pbc-etale-proetale-inclusion}
Both sieves are generated by the base changes
$V\times_UV'\to V$ as $V'\to U$ ranges through $R$.  The selected etale and
pro-etale over-categories compute these fiber products by the same scheme
$V\times_UV'$, and $\iota_X$ fixes its projections.  Hence their generating
arrows, and therefore the generated sieves, agree.
\end{proof}

\begin{definition}[Continuity predicate for the site inclusion]
\label{def:weil-efi-site-continuity-certificate}
\uses{def:weil-efi-site-covering-sieve-image}
Let $\mathsf{Continuous}(\iota_X)$ mean that the image-sieve operation sends a
covering sieve to a covering sieve and satisfies the displayed pullback
equation.
\end{definition}

\begin{proposition}[Continuity certificate for the site inclusion]
\label{prop:weil-efi-site-continuity-certificate}
\uses{def:weil-efi-site-continuity-certificate}
The predicate $\mathsf{Continuous}(\iota_X)$ holds.
\end{proposition}
\begin{proof}
\uses{lemma:weil-efi-site-covering-sieve-image,
  lemma:weil-efi-site-covering-sieve-pullback}
The first lemma supplies covering preservation and the second supplies the
pullback field.
\end{proof}

\begin{lemma}[The etale/pro-etale assignment is a functor]
\label{lemma:weil-pbc-etale-proetale-functor}
\uses{def:weil-pbc-etale-proetale-inclusion}
The assignment $\iota_X$ preserves identity morphisms and composition.
\end{lemma}
\begin{proof}
\uses{} 
For an etale object $U\to X$, both $\iota_X(1_U)$ and
$1_{\iota_X(U)}$ have underlying scheme morphism $1_U$.  For composable
$X$-morphisms $U\xrightarrow{h}V\xrightarrow{k}W$, both
$\iota_X(kh)$ and $\iota_X(k)\iota_X(h)$ have underlying scheme morphism
$kh$.  Morphisms in either A2 over-category are determined by this underlying
map, which proves both equations.
\end{proof}

\begin{lemma}[Continuity of the etale/pro-etale inclusion]
\label{lemma:weil-pbc-etale-proetale-continuous}
\uses{def:weil-pbc-etale-proetale-inclusion}
The functor $\iota_X:X_{\mathrm{et}}\to X_{\mathrm{proet}}$ is continuous
for the final A2 topologies.
\end{lemma}
\begin{proof}
\uses{lemma:weil-pbc-etale-proetale-functor}

Let $\{U_a\to U\}_{a\in A}$ be a covering family in the A2 small etale
pretopology.  Each $U_a\to U$ is etale, hence weakly etale, and the underlying
family of scheme maps is unchanged by $\iota_X$.  Its images are jointly
surjective, so the same family belongs to the pulled-back pro-etale
precoverage.  Consequently the sieve that it generates is covering for the
A2 pro-etale topology.

More generally, if $R\subset h_U$ is an etale covering sieve, send it to the
sieve on $\iota_XU$ generated by the arrows $\iota_XV\to\iota_XU$ for
$V\to U$ in $R$.  The preceding argument shows that this image sieve is
pro-etale covering.  Formation of the image sieve commutes with pullback,
because the underlying scheme fiber products are unchanged.  Since both A2
topologies are generated by the displayed precoverages, this is the
continuity criterion.
\end{proof}

\begin{lemma}[Finite-limit preservation of the inclusion]
\label{lemma:weil-pbc-etale-proetale-left-exact}
\uses{def:weil-pbc-etale-proetale-inclusion}
The functor $\iota_X$ preserves the terminal object and fiber products, hence
is left exact.
\end{lemma}
\begin{proof}
\uses{lemma:weil-pbc-etale-proetale-functor}

The terminal object in both A2 sites is $1_X:X\to X$, and $\iota_X$ fixes it.
For a diagram $U\to W\leftarrow V$ of etale $X$-objects, its fiber product in
both over-categories has underlying scheme $U\times_WV$ with the same two
projections and structural map to $X$.  Etale morphisms are stable under base
change and composition, so this is again an etale object.  Thus the canonical
map
\[
 \iota_X(U\times_WV)\longrightarrow
 \iota_XU\times_{\iota_XW}\iota_XV
\]
is the identity on $U\times_WV$, and is an isomorphism.  Terminal objects and
fiber products generate finite limits.
\end{proof}

\begin{definition}[Etale/pro-etale comparison morphism]
\label{def:weil-pbc-epsilon-comparison}
\uses{}

The continuous left-exact functor $\iota_X$ induces the geometric morphism
\[
 \epsilon_X:X_{\mathrm{proet}}\longrightarrow X_{\mathrm{et}}.
\]
For every commutative coefficient ring $A$, its inverse- and direct-image
functors have the typed source and target
\[
 \epsilon_X^{-1}:\operatorname{Sh}_A(X_{\mathrm{et}})
   \longrightarrow\operatorname{Sh}_A(X_{\mathrm{proet}}),
 \qquad
 \epsilon_{X,*}:\operatorname{Sh}_A(X_{\mathrm{proet}})
   \longrightarrow\operatorname{Sh}_A(X_{\mathrm{et}}).
\]
For $G\in\operatorname{Sh}_A(X_{\mathrm{proet}})$ one has
$(\epsilon_{X,*}G)(U)=G(\iota_XU)$.  For
$F\in\operatorname{Sh}_A(X_{\mathrm{et}})$, the inverse image is the
pro-etale sheafification of the left Kan extension along $\iota_X^{\mathrm{op}}$;
before sheafification its value at $W\to X$ is
\[
 \underset{(W\to\iota_XU)}{\operatorname{colim}}F(U).
\]
\end{definition}


\begin{definition}[Etale factorization index for a pro-etale object]
\label{def:weil-pbc-epsilon-factorization-index-v1}
\uses{def:weil-pbc-epsilon-comparison}
For a pro-etale object $W\to X$, let $I_W$ be the category whose objects are
factorizations $W\xrightarrow{a}\iota_XU\to X$ with $U\to X$ etale.  A
morphism from $(a,U)$ to $(a',U')$ is an etale-site arrow $h:U'\to U$ such
that $a=\iota_X(h)a'$.  The transition map on a sheaf $F$ is restriction
$F(U)\to F(U')$.
\end{definition}

\begin{lemma}[The etale factorization index is filtered]
\label{lemma:weil-pbc-epsilon-factorization-index-filtered-v1}
\uses{def:weil-pbc-epsilon-factorization-index-v1}
For every $W$, the category $I_W$ is nonempty and filtered.
\end{lemma}
\begin{proof}
\uses{def:weil-pbc-epsilon-factorization-index-v1,lemma:weil-pbc-etale-proetale-left-exact}
The structural map $W\to X=\iota_XX$ gives an object.  Two factorizations
through $U_1$ and $U_2$ are dominated by the factorization through
$U_1\times_XU_2$, since $\iota_X$ preserves this fibre product.  For two
parallel arrows $h_0,h_1:U'\rightrightarrows U$ that agree after precomposition
with $W\to\iota_XU'$, their equalizer is the pullback of the diagonal of the
etale morphism $U\to X$.  This diagonal is an open immersion, so the
equalizer is an etale open subobject of $U'$ through which $W$ factors.  The
resulting object receives a morphism from $(a',U')$ and equalizes the two
arrows in $I_W$.  These are the three filteredness axioms.
\end{proof}

\begin{lemma}[Construction of the comparison geometric morphism]
\label{lemma:weil-pbc-epsilon-geometric-morphism-construction-v1}
\uses{def:weil-pbc-epsilon-comparison,def:weil-pbc-epsilon-factorization-index-v1}
Restriction along $\iota_X^{\mathrm{op}}$, pro-etale sheafification of left
Kan extension, and the displayed evaluation maps define an adjoint pair
$\epsilon_X^{-1}\dashv\epsilon_{X,*}$.  The adjunction bijection is natural
in both variables, and $\epsilon_X^{-1}$ preserves finite limits.
\end{lemma}
\begin{proof}
\uses{lemma:weil-pbc-etale-proetale-continuous,
  lemma:weil-pbc-etale-proetale-left-exact,
  lemma:weil-pbc-epsilon-factorization-index-filtered-v1,
  def:weil-pbc-epsilon-comparison}
Continuity makes restriction of a pro-etale sheaf an etale sheaf.  Presheaf
left Kan extension is left adjoint to presheaf restriction, and sheafification
is left adjoint to inclusion of sheaves.  Composing these two adjunctions gives
\[
 \operatorname{Hom}(\epsilon_X^{-1}F,G)
 \cong\operatorname{Hom}(F,\epsilon_{X,*}G).
\]
The two universal properties make this bijection natural in $F$ and $G$ and
make its identity and composition equations automatic.

Before sheafification, evaluation at $W$ is
$\varinjlim_{I_W}F(U)$.  The preceding lemma makes $I_W$ filtered, so this
colimit of $A$-modules commutes with finite limits.  Sheafification of
$A$-module presheaves is left exact, and the left-exactness of $\iota_X$
identifies the transition diagrams for terminal objects and fibre products.
Thus $\epsilon_X^{-1}$ preserves finite limits.  All maps used are induced by
the same restriction maps, so they are compatible with restriction in $W$ and
with morphisms of coefficient sheaves.
\end{proof}

% Atomic construction and property leaves for EFI SITE-05.

\begin{definition}[Forward Hom map for the comparison]
\label{def:weil-efi-site-adjunction-hom-forward}
\uses{def:weil-pbc-epsilon-comparison}
For $A$-module sheaves $F$ on $X_{\mathrm{et}}$ and $G$ on
$X_{\mathrm{proet}}$, send
$\alpha:\epsilon_X^{-1}F\to G$ to the map $\Phi(\alpha):F\to\epsilon_{X,*}G$
whose component at $U$ is
\[
 s\longmapsto
 \alpha_{\iota_XU}\bigl[(1_{\iota_XU},s)\bigr].
\]
\end{definition}

\begin{definition}[Representative formula for the reverse Hom map]
\label{def:weil-efi-site-adjunction-hom-backward-presheaf}
\uses{def:weil-pbc-epsilon-comparison}
For $\beta:F\to\epsilon_{X,*}G$, assign to a left-Kan-extension
representative $(W\xrightarrow a\iota_XU,s)$ the section
\[
 G(a)(\beta_U(s))\in G(W).
\]
\end{definition}

\begin{lemma}[Descent of the reverse Hom formula]
\label{lemma:weil-efi-site-adjunction-hom-backward-descent}
\uses{def:weil-efi-site-adjunction-hom-backward-presheaf}
The representative formula factors uniquely through pro-etale
sheafification.
\end{lemma}
\begin{proof}
\uses{def:weil-efi-site-adjunction-hom-backward-presheaf,
  def:weil-pbc-epsilon-factorization-index-v1}
For a transition $h:U'\to U$ with $a=\iota_X(h)a'$, naturality of $\beta$
and functoriality of $G$ identify the values of $(a,s)$ and
$(a',F(h)s)$.  Restriction along $b:W'\to W$ changes each formula to
$G(ab)(\beta_U(s))$.  Local equality descends because $G$ is a sheaf.  The
universal property of sheafification gives the unique factorization.
\end{proof}

\begin{definition}[Reverse Hom map for the comparison]
\label{def:weil-efi-site-adjunction-hom-backward}
\uses{def:weil-efi-site-adjunction-hom-backward-presheaf}
Let $\Psi(\beta):\epsilon_X^{-1}F\to G$ be the sheaf map induced from the
representative formula by the universal property of sheafification.
\end{definition}

\begin{lemma}[Left inverse equation for the comparison Hom maps]
\label{lemma:weil-efi-site-adjunction-hom-left-inverse}
\uses{def:weil-efi-site-adjunction-hom-forward,
  def:weil-efi-site-adjunction-hom-backward}
For every $\alpha:\epsilon_X^{-1}F\to G$,
\[
 \Psi(\Phi(\alpha))=\alpha.
\]
\end{lemma}
\begin{proof}
\uses{def:weil-efi-site-adjunction-hom-forward,
  def:weil-efi-site-adjunction-hom-backward}
On a representative $(W\xrightarrow a\iota_XU,s)$, the left side is
$G(a)(\alpha_{\iota_XU}[(1,s)])$.  Naturality of $\alpha$ identifies this
with $\alpha_W[(a,s)]$.  Representatives are locally exhaustive after
sheafification, so the sheaf maps agree.
\end{proof}

\begin{lemma}[Right inverse equation for the comparison Hom maps]
\label{lemma:weil-efi-site-adjunction-hom-right-inverse}
\uses{def:weil-efi-site-adjunction-hom-forward,
  def:weil-efi-site-adjunction-hom-backward}
For every $\beta:F\to\epsilon_{X,*}G$,
\[
 \Phi(\Psi(\beta))=\beta.
\]
\end{lemma}
\begin{proof}
\uses{def:weil-efi-site-adjunction-hom-forward,
  def:weil-efi-site-adjunction-hom-backward}
At an etale object $U$, the left side sends $s$ to the evaluation of
$(1_{\iota_XU},s)$, namely
$G(1_{\iota_XU})(\beta_U(s))=\beta_U(s)$.
\end{proof}

\begin{lemma}[Source naturality of the comparison Hom correspondence]
\label{lemma:weil-efi-site-adjunction-hom-source-naturality}
\uses{def:weil-efi-site-adjunction-hom-forward}
For $u:F'\to F$ and $\alpha:\epsilon_X^{-1}F\to G$,
\[
 \Phi(\alpha\,\epsilon_X^{-1}u)=\Phi(\alpha)u.
\]
\end{lemma}
\begin{proof}
\uses{def:weil-efi-site-adjunction-hom-forward}
At $U$, each side sends $s\in F'(U)$ to
$\alpha_{\iota_XU}[(1_{\iota_XU},u_U(s))]$.
\end{proof}

\begin{lemma}[Target naturality of the comparison Hom correspondence]
\label{lemma:weil-efi-site-adjunction-hom-target-naturality}
\uses{def:weil-efi-site-adjunction-hom-forward}
For $v:G\to G'$ and $\alpha:\epsilon_X^{-1}F\to G$,
\[
 \Phi(v\alpha)=\epsilon_{X,*}(v)\Phi(\alpha).
\]
\end{lemma}
\begin{proof}
\uses{def:weil-efi-site-adjunction-hom-forward}
At $U$, each side sends $s\in F(U)$ to
$v_{\iota_XU}(\alpha_{\iota_XU}[(1_{\iota_XU},s)])$.
\end{proof}

\begin{definition}[Adjunction-certificate predicate for the comparison functors]
\label{def:weil-efi-site-sheaf-adjunction-certificate}
\uses{def:weil-efi-site-adjunction-hom-forward,
  def:weil-efi-site-adjunction-hom-backward}
Let $\mathsf{AdjunctionCert}(\Phi,\Psi)$ be the conjunction of the left
inverse equation, the right inverse equation, source naturality, and target
naturality for the displayed Hom-map pair.
\end{definition}

\begin{proposition}[Adjunction certificate for the comparison functors]
\label{prop:weil-efi-site-sheaf-adjunction-certificate}
\uses{def:weil-efi-site-sheaf-adjunction-certificate}
The predicate $\mathsf{AdjunctionCert}(\Phi,\Psi)$ holds.
\end{proposition}
\begin{proof}
\uses{lemma:weil-efi-site-adjunction-hom-left-inverse,
  lemma:weil-efi-site-adjunction-hom-right-inverse,
  lemma:weil-efi-site-adjunction-hom-source-naturality,
  lemma:weil-efi-site-adjunction-hom-target-naturality}
The four cited lemmas supply the four fields in their displayed order.
\end{proof}

\begin{lemma}[Finite-limit preservation by comparison inverse image]
\label{lemma:weil-efi-site-inverse-image-finite-limits}
\uses{def:weil-pbc-epsilon-comparison}
The functor $\epsilon_X^{-1}$ preserves finite limits of $A$-module sheaves.
\end{lemma}
\begin{proof}
\uses{def:weil-pbc-epsilon-factorization-index-v1,
  lemma:weil-pbc-epsilon-factorization-index-filtered-v1,
  prop:weil-efi-site-continuity-certificate}
Before sheafification, evaluation at $W$ is the filtered colimit
$\varinjlim_{I_W}F(U)$.  Filtered colimits of $A$-modules commute with finite
limits.  Sheafification of $A$-module presheaves is left exact, and the
pullback equation in the continuity certificate identifies the transition
diagrams.  The resulting sheaf functor therefore preserves finite limits.
\end{proof}

\begin{definition}[Geometric-morphism certificate predicate]
\label{def:weil-efi-site-geometric-morphism-certificate}
\uses{def:weil-pbc-epsilon-comparison,
  def:weil-efi-site-sheaf-adjunction-certificate}
Let $\mathsf{GeometricMorphismCert}(\epsilon_X^{-1},\epsilon_{X,*})$ mean
that the Hom correspondence satisfies the adjunction-certificate predicate
and that $\epsilon_X^{-1}$ preserves finite limits.
\end{definition}

\begin{proposition}[Geometric-morphism package for the comparison]
\label{prop:weil-efi-site-geometric-morphism-package}
\uses{def:weil-efi-site-geometric-morphism-certificate}
The predicate
$\mathsf{GeometricMorphismCert}(\epsilon_X^{-1},\epsilon_{X,*})$ holds.
\end{proposition}
\begin{proof}
\uses{prop:weil-efi-site-sheaf-adjunction-certificate,
  lemma:weil-efi-site-inverse-image-finite-limits}
The proposition supplies the adjunction field and the lemma supplies the
finite-limit field.
\end{proof}

\begin{lemma}[Exactness of the packaged comparison inverse image]
\label{lemma:weil-efi-site-inverse-image-exact}
\uses{def:weil-pbc-epsilon-comparison}
The functor $\epsilon_X^{-1}$ is exact on $A$-module sheaves.
\end{lemma}
\begin{proof}
\uses{prop:weil-efi-site-sheaf-adjunction-certificate,
  lemma:weil-efi-site-inverse-image-finite-limits}
Finite-limit preservation gives preservation of kernels.  The adjunction
certificate makes $\epsilon_X^{-1}$ a left adjoint, so it preserves
cokernels.  Kernel and cokernel preservation gives exactness in the abelian
category of $A$-module sheaves.
\end{proof}

\begin{lemma}[Well-definedness certificate for def:weil-pbc-epsilon-comparison]
\label{lemma:weil-sf-well-defined-def-weil-pbc-epsilon-comparison}
\uses{def:weil-pbc-epsilon-comparison}
The construction in Definition \ref{def:weil-pbc-epsilon-comparison} is a typed raw datum.  Its theorem-valued naturality, comparison, descent, or coherence fields are isolated in this property lemma.
\end{lemma}
\begin{proof}
\uses{lemma:weil-pbc-etale-proetale-continuous,
  lemma:weil-pbc-etale-proetale-left-exact,
  lemma:weil-pbc-epsilon-factorization-index-filtered-v1,
  lemma:weil-pbc-epsilon-geometric-morphism-construction-v1}
The inclusion functor is continuous and left exact.  The factorization-index
lemma makes the displayed Kan-extension formula a filtered colimit, and the
construction lemma supplies its sheafification, the adjunction, naturality,
and finite-limit preservation.  The formulas for direct image and inverse
image are therefore independent of representatives and restrictions and
define the stated geometric morphism with all comparison fields.
\end{proof}
\begin{lemma}[The comparison functors form a geometric morphism]
\label{lemma:weil-pbc-epsilon-geometric-morphism}
\uses{def:weil-pbc-epsilon-comparison}
For $A$-module sheaves $F$ on $X_{\mathrm{et}}$ and $G$ on
$X_{\mathrm{proet}}$ there is a natural bijection
\[
 \operatorname{Hom}_{X_{\mathrm{proet}}}(\epsilon_X^{-1}F,G)
 \cong
 \operatorname{Hom}_{X_{\mathrm{et}}}(F,\epsilon_{X,*}G),
\]
and $\epsilon_X^{-1}$ preserves finite limits.
\end{lemma}
\begin{proof}
\uses{lemma:weil-pbc-etale-proetale-continuous,lemma:weil-pbc-etale-proetale-left-exact}

Restriction along $\iota_X^{\mathrm{op}}$ sends pro-etale sheaves to etale
sheaves by Lemma \ref{lemma:weil-pbc-etale-proetale-continuous}.  The left-Kan
extension adjunction, followed by the sheafification adjunction, gives the
displayed bijection.

For a pro-etale object $W$, let $I_W$ be the category of factorizations
$W\to\iota_XU$ with $U$ etale over $X$.  It is filtered: the map to $X$
gives an object, two factorizations are dominated by their fiber product over
$X$, and two parallel arrows are equalized by their scheme-theoretic
equalizer, an open etale subobject through which $W$ factors.  Filtered
colimits of $A$-modules commute with finite limits.  The pointwise Kan
extension therefore preserves finite limits, and $A$-module sheafification
does as well.  This proves left exactness of $\epsilon_X^{-1}$ and hence the
geometric-morphism assertion.
\end{proof}

\begin{lemma}[Exactness of comparison inverse image]
\label{lemma:weil-pbc-epsilon-inverse-image-exact}
\uses{}

For every commutative ring $A$, the functor
$\epsilon_X^{-1}:\operatorname{Sh}_A(X_{\mathrm{et}})\to
\operatorname{Sh}_A(X_{\mathrm{proet}})$ is exact.
\end{lemma}
\begin{proof}
\uses{lemma:weil-pbc-epsilon-geometric-morphism}

It preserves kernels by the finite-limit assertion in
Lemma \ref{lemma:weil-pbc-epsilon-geometric-morphism}.  It preserves
cokernels because it is a left adjoint.  Applying these two facts to a short
exact sequence
\[
 0\longrightarrow F'\longrightarrow F\longrightarrow F''\longrightarrow0
\]
shows that its inverse-image sequence is exact at all three terms.
\end{proof}

\begin{definition}[Unit and counit of the comparison]
\label{def:weil-pbc-epsilon-unit-counit}
\uses{def:weil-pbc-epsilon-comparison}
Write the unit and counit of the preceding adjunction as
\[
 \eta^X:\operatorname{id}_{\operatorname{Sh}_A(X_{\mathrm{et}})}
  \longrightarrow\epsilon_{X,*}\epsilon_X^{-1},
 \qquad
 \varepsilon^X:\epsilon_X^{-1}\epsilon_{X,*}
  \longrightarrow\operatorname{id}_{\operatorname{Sh}_A(X_{\mathrm{proet}})}.
\]
For $F\in\operatorname{Sh}_A(X_{\mathrm{et}})$, the component
$\eta^X_{F,U}$ sends $s\in F(U)$ to the class represented by
$(1_{\iota_XU}:\iota_XU\to\iota_XU,s)$.  For
$G\in\operatorname{Sh}_A(X_{\mathrm{proet}})$, the component
$\varepsilon^X_{G,W}$ sends a representative
$(W\xrightarrow{a}\iota_XU,s)$ to the restriction $G(a)(s)\in G(W)$.
This map on the Kan-extension presheaf factors uniquely through pro-etale
sheafification.
\end{definition}


\begin{lemma}[Representative descent for the comparison counit]
\label{lemma:weil-pbc-epsilon-counit-representative-descent-v1}
\uses{def:weil-pbc-epsilon-unit-counit}
The rule $(W\xrightarrow a\iota_XU,s)\mapsto G(a)(s)$ respects the
left-Kan-extension relations, restriction in $W$, and local equality, hence
factors uniquely through pro-etale sheafification.
\end{lemma}
\begin{proof}
\uses{def:weil-pbc-epsilon-unit-counit,
  def:weil-pbc-epsilon-factorization-index-v1,
  lemma:weil-pbc-epsilon-geometric-morphism-construction-v1}
If $h:U'\to U$ is a transition and $a=\iota_X(h)a'$, then
functoriality of $G$ gives
\[
 G(a')G(\iota_Xh)(s)=G(a)(s),
\]
so the rule is constant on the colimit relation.  For $b:W'\to W$, both
restriction before evaluation and evaluation before restriction send
$(a,s)$ to $G(ab)(s)$.  If two representatives agree after a pro-etale
cover, their images agree on that cover; the sheaf equalizer for $G$ makes
the images equal.  Thus the presheaf map is local and has a unique factor
through sheafification.  This factor is precisely the counit attached to the
adjunction constructed above.
\end{proof}

% Atomic leaves for EFI SITE-06 and SITE-07.

\begin{lemma}[Naturality equation for the comparison unit]
\label{lemma:weil-efi-site-unit-naturality}
\uses{def:weil-pbc-epsilon-unit-counit}
For $\alpha:F\to F'$ on $X_{\mathrm{et}}$,
\[
 (\epsilon_{X,*}\epsilon_X^{-1}\alpha)\eta^X_F
 =\eta^X_{F'}\alpha.
\]
\end{lemma}
\begin{proof}
\uses{def:weil-pbc-epsilon-unit-counit}
At an etale object $U$, each route sends $s\in F(U)$ to
$(1_{\iota_XU},\alpha_U(s))$.
\end{proof}

\begin{lemma}[Naturality equation for the comparison counit]
\label{lemma:weil-efi-site-counit-naturality}
\uses{def:weil-pbc-epsilon-unit-counit}
For $\beta:G\to G'$ on $X_{\mathrm{proet}}$,
\[
 \beta\varepsilon^X_G
 =\varepsilon^X_{G'}(\epsilon_X^{-1}\epsilon_{X,*}\beta).
\]
\end{lemma}
\begin{proof}
\uses{def:weil-pbc-epsilon-unit-counit,
  lemma:weil-pbc-epsilon-counit-representative-descent-v1}
On a representative $(W\xrightarrow a\iota_XU,s)$, each route gives
$\beta_W(G(a)(s))=G'(a)(\beta_{\iota_XU}(s))$.  The representative-descent
lemma carries this equality through sheafification.
\end{proof}

\begin{lemma}[Restriction equation for the comparison unit]
\label{lemma:weil-efi-site-unit-restriction}
\uses{def:weil-pbc-epsilon-unit-counit}
For an etale arrow $b:U'\to U$, restriction carries
$\eta^X_{F,U}(s)$ to $\eta^X_{F,U'}(F(b)s)$.
\end{lemma}
\begin{proof}
\uses{def:weil-pbc-epsilon-unit-counit}
Restriction carries $(1_{\iota_XU},s)$ to $(\iota_Xb,s)$.  The
left-Kan-extension relation identifies this representative with
$(1_{\iota_XU'},F(b)s)$.
\end{proof}

\begin{lemma}[Restriction equation for the comparison counit]
\label{lemma:weil-efi-site-counit-restriction}
\uses{def:weil-pbc-epsilon-unit-counit}
For a pro-etale arrow $b:W'\to W$, restriction carries the counit value of
$(W\xrightarrow a\iota_XU,s)$ to the counit value of
$(W'\xrightarrow{ab}\iota_XU,s)$.
\end{lemma}
\begin{proof}
\uses{def:weil-pbc-epsilon-unit-counit,
  lemma:weil-pbc-epsilon-counit-representative-descent-v1}
The first route gives $G(b)G(a)(s)$ and the second gives $G(ab)(s)$; these
are equal by functoriality of $G$ and descend through sheafification.
\end{proof}

\begin{definition}[Naturality predicate for the comparison unit and counit]
\label{def:weil-efi-site-unit-counit-certificate}
\uses{def:weil-pbc-epsilon-unit-counit}
Let $\mathsf{NaturalUnitCounit}(\eta^X,\varepsilon^X)$ be the conjunction of
unit naturality, counit naturality, the unit restriction equation, and the
counit restriction equation.
\end{definition}

\begin{proposition}[Natural unit-counit package for the comparison]
\label{prop:weil-efi-site-unit-counit-package}
\uses{def:weil-efi-site-unit-counit-certificate}
The predicate $\mathsf{NaturalUnitCounit}(\eta^X,\varepsilon^X)$ holds.
\end{proposition}
\begin{proof}
\uses{lemma:weil-pbc-epsilon-counit-representative-descent-v1,
  lemma:weil-efi-site-unit-naturality,
  lemma:weil-efi-site-counit-naturality,
  lemma:weil-efi-site-unit-restriction,
  lemma:weil-efi-site-counit-restriction}
Representative descent types the counit, and the four atomic equation lemmas
supply the fields of the displayed predicate.
\end{proof}

\begin{lemma}[Left triangle equation for the comparison adjunction]
\label{lemma:weil-efi-site-left-triangle}
\uses{def:weil-pbc-epsilon-unit-counit}
For every etale $A$-module sheaf $F$,
\[
 \varepsilon^X_{\epsilon_X^{-1}F}\,
 \epsilon_X^{-1}(\eta^X_F)=1_{\epsilon_X^{-1}F}.
\]
\end{lemma}
\begin{proof}
\uses{def:weil-pbc-epsilon-unit-counit}
A local section is represented by $(W\xrightarrow a\iota_XU,s)$.  The
composite replaces $s$ by its identity representative and evaluates along
$a$, returning the class $(a,s)$.  Local representatives detect sheaf maps.
\end{proof}

\begin{lemma}[Right triangle equation for the comparison adjunction]
\label{lemma:weil-efi-site-right-triangle}
\uses{def:weil-pbc-epsilon-unit-counit}
For every pro-etale $A$-module sheaf $G$,
\[
 \epsilon_{X,*}(\varepsilon^X_G)\,
 \eta^X_{\epsilon_{X,*}G}=1_{\epsilon_{X,*}G}.
\]
\end{lemma}
\begin{proof}
\uses{def:weil-pbc-epsilon-unit-counit}
At an etale object $U$, a section $s\in G(\iota_XU)$ first becomes
$(1_{\iota_XU},s)$ and is then evaluated along the identity, returning $s$.
\end{proof}

\begin{definition}[Triangle-certificate predicate for the comparison adjunction]
\label{def:weil-efi-site-triangle-certificate}
\uses{def:weil-pbc-epsilon-unit-counit}
Let $\mathsf{TriangleCert}(\eta^X,\varepsilon^X)$ be the conjunction of the
left-triangle equation and the right-triangle equation.
\end{definition}

\begin{proposition}[Triangle certificate for the comparison adjunction]
\label{prop:weil-efi-site-triangle-certificate}
\uses{def:weil-efi-site-triangle-certificate}
The predicate $\mathsf{TriangleCert}(\eta^X,\varepsilon^X)$ holds.
\end{proposition}
\begin{proof}
\uses{lemma:weil-efi-site-left-triangle,
  lemma:weil-efi-site-right-triangle}
The cited lemmas supply the left and right fields, respectively.
\end{proof}

\begin{lemma}[Naturality and triangle equations for the comparison adjunction]
\label{lemma:weil-pbc-epsilon-adjunction-equations-v1}
\uses{def:weil-pbc-epsilon-unit-counit}
The displayed unit and counit are natural in coefficient sheaves and satisfy
the two triangle identities.  Their components commute with restriction on
both sites.
\end{lemma}
\begin{proof}
\uses{def:weil-pbc-epsilon-unit-counit,
  lemma:weil-pbc-epsilon-counit-representative-descent-v1,
  lemma:weil-pbc-epsilon-geometric-morphism-construction-v1}
For $\alpha:F\to F'$, both routes around the unit naturality square send
$s\in F(U)$ to $(1_{\iota_XU},\alpha_U(s))$.  For $\beta:G\to G'$, both
routes around the counit square send $(a,s)$ to
$\beta_W(G(a)(s))=G'(a)(\beta_{\iota_XU}(s))$.  These equalities commute
with restriction and descend to sheaves by the preceding lemma.

A local section of $\epsilon_X^{-1}F$ is represented by $(a,s)$.  The first
triangle composite replaces $s$ by its identity representative and evaluates
along $a$, returning $(a,s)$.  A section
$s\in(\epsilon_{X,*}G)(U)=G(\iota_XU)$ is sent by the second composite to
$(1_{\iota_XU},s)$ and then evaluated along the identity, returning $s$.
Local representatives detect sheaf maps, so both composites are identities.
\end{proof}

\begin{lemma}[Well-definedness certificate for def:weil-pbc-epsilon-unit-counit]
\label{lemma:weil-sf-well-defined-def-weil-pbc-epsilon-unit-counit}
\uses{def:weil-pbc-epsilon-unit-counit}
The construction in Definition \ref{def:weil-pbc-epsilon-unit-counit} is a typed raw datum.  Its theorem-valued naturality, comparison, descent, or coherence fields are isolated in this property lemma.
\end{lemma}
\begin{proof}
\uses{lemma:weil-pbc-epsilon-geometric-morphism-construction-v1,
  lemma:weil-pbc-epsilon-counit-representative-descent-v1,
  lemma:weil-pbc-epsilon-adjunction-equations-v1}
The geometric-morphism construction supplies the adjunction.  Representative
descent makes the displayed counit formula a sheaf morphism, while the unit
formula is the identity-factorization component of the same adjunction.  The
final lemma proves naturality, restriction compatibility, and both triangle
identities, so every theorem-valued field of the raw unit--counit datum is
well-defined.
\end{proof}
\begin{lemma}[Naturality of the comparison unit and counit]
\label{lemma:weil-pbc-epsilon-unit-counit-naturality}
\uses{def:weil-pbc-epsilon-unit-counit}
For $\alpha:F\to F'$ on $X_{\mathrm{et}}$ and $\beta:G\to G'$ on
$X_{\mathrm{proet}}$, one has
\begin{align*}
 (\epsilon_{X,*}\epsilon_X^{-1}\alpha)\eta^X_F
   &=\eta^X_{F'}\alpha,\\
 \beta\varepsilon^X_G
   &=\varepsilon^X_{G'}(\epsilon_X^{-1}\epsilon_{X,*}\beta).
\end{align*}
\end{lemma}
\begin{proof}
\uses{def:weil-pbc-epsilon-comparison,def:weil-pbc-epsilon-unit-counit}
At an etale object $U$, both sides of the first equation send $s\in F(U)$
to the class $(1_{\iota_XU},\alpha_U(s))$.  At a pro-etale object $W$, apply
the second equation to a representative $(W\xrightarrow{a}\iota_XU,s)$.
Both sides give
$\beta_W(G(a)(s))=G'(a)(\beta_{\iota_XU}(s))$ by naturality of $\beta$.
These equations survive sheafification and prove the asserted naturality.
\end{proof}

\begin{lemma}[Triangle identities for the comparison adjunction]
\label{lemma:weil-pbc-epsilon-triangle-identities}
\uses{def:weil-pbc-epsilon-unit-counit}
The unit and counit satisfy
\begin{align*}
 \varepsilon^X_{\epsilon_X^{-1}F}\,
   \epsilon_X^{-1}(\eta^X_F)&=1_{\epsilon_X^{-1}F},\\
 \epsilon_{X,*}(\varepsilon^X_G)\,
   \eta^X_{\epsilon_{X,*}G}&=1_{\epsilon_{X,*}G}.
\end{align*}
\end{lemma}
\begin{proof}
\uses{lemma:weil-pbc-epsilon-unit-counit-naturality}

A section of $\epsilon_X^{-1}F$ is locally represented by
$(W\xrightarrow{a}\iota_XU,s)$.  The first composite replaces $s$ by its
identity representative at $\iota_XU$ and then restricts it along $a$, which
returns the original class.  For $s\in(\epsilon_{X,*}G)(U)=G(\iota_XU)$,
the second composite forms $(1_{\iota_XU},s)$ and evaluates it along the
identity, returning $s$.  Local representatives determine maps of sheaves,
so both displayed composites are identities.
\end{proof}

\begin{lemma}[Naturality of the comparison in the scheme]
\label{lemma:weil-pbc-epsilon-scheme-naturality}
\uses{def:weil-pbc-etale-proetale-inclusion,def:weil-pbc-epsilon-comparison,def:weil-pbc-epsilon-unit-counit}
For a scheme morphism $u:X\to Y$, base change on the two sites gives a
canonical natural isomorphism
\[
 u_{\mathrm{proet}}^*\iota_Y\xrightarrow{\sim}
 \iota_Xu_{\mathrm{et}}^*.
\]
It is compatible with identities and composition in $u$, and the induced
square of geometric morphisms carries $\eta^Y,\varepsilon^Y$ to
$\eta^X,\varepsilon^X$.
\end{lemma}
\begin{proof}
\uses{def:weil-pbc-epsilon-comparison,def:weil-pbc-epsilon-unit-counit}
On an etale object $V\to Y$, both composites have underlying $X$-scheme
$X\times_YV\to X$; on a morphism they both use its base change.  The identity
map of this fiber product gives the component of the natural isomorphism.
The associativity and unit isomorphisms for fiber products give its
composition and identity coherence.  Under this identification, the unit is
still represented by an identity arrow and the counit is still restriction
along $W\to\iota_XU$.  The formulas in Definition
\ref{def:weil-pbc-epsilon-unit-counit} therefore give the two compatibility
squares.
\end{proof}

\begin{definition}[Witness for coming from the etale site]
\label{def:weil-pbc-comes-from-etale}
\uses{def:weil-pbc-epsilon-comparison}
For a coefficient ring $A$ and $K\in D^+(X_{\mathrm{proet}},A)$, a witness
that $K$ comes from the etale site is a pair
\[
 (L,\theta),\qquad L\in D^+(X_{\mathrm{et}},A),\quad
 \theta:\epsilon_X^{-1}L\xrightarrow{\sim}K,
\]
including the two inverse equations for $\theta$.
\end{definition}

\begin{lemma}[Comparison inverse image commutes with pullback]
\label{lemma:weil-pbc-epsilon-pullback}
\uses{def:weil-pbc-epsilon-comparison,def:weil-pbc-comes-from-etale}
For $u:X\to Y$ and $L\in D^+(Y_{\mathrm{et}},A)$ there is a natural
isomorphism
\[
 Lu^*\epsilon_Y^{-1}L\xrightarrow{\sim}
 \epsilon_X^{-1}Lu^*L.
\]
It is compatible with identities, composition, and both horizontal maps in a
Cartesian square.
\end{lemma}
\begin{proof}
\uses{lemma:weil-pbc-epsilon-scheme-naturality,lemma:weil-pbc-epsilon-inverse-image-exact}

Lemma \ref{lemma:weil-pbc-epsilon-scheme-naturality} identifies the two
composite site functors.  Their inverse images are therefore sheafifications
of the same presheaf left Kan extension, and the identity on that presheaf
induces the displayed isomorphism.  The identity and composition coherence
proved there gives the corresponding equations and the two horizontal maps
in a Cartesian square.  Exactness of $\epsilon^{-1}$ and K-flat replacement
extend these equations to derived pullback.
\end{proof}

\begin{lemma}[The comparison unit is an isomorphism]
\label{lemma:weil-pbc-epsilon-unit-isomorphism}
\uses{def:weil-pbc-epsilon-unit-counit,def:weil-pbc-epsilon-comparison}
For every $A$-module sheaf $F$ on $X_{\mathrm{et}}$, the unit
\[
 \eta^X_F:F\longrightarrow\epsilon_{X,*}\epsilon_X^{-1}F
\]
is an isomorphism.
\end{lemma}
\begin{proof}
\uses{def:weil-pbc-epsilon-comparison,def:weil-pbc-epsilon-unit-counit}
Evaluate at an etale object $U\to X$.  Restriction of the left-Kan-extension
presheaf in Definition \ref{def:weil-pbc-epsilon-comparison} back to
$\iota_XU$ has the identity factorization
$\iota_XU\to\iota_XU$.  Refining any other factorization by its fiber product
with $U$ shows that the associated sheaf value is $F(U)$.  Under this
identification, the formula in Definition
\ref{def:weil-pbc-epsilon-unit-counit} sends $s$ to the same section $s$.
Thus every component of $\eta^X_F$ is the identity.  Etale objects detect
maps of etale sheaves, so the unit is an isomorphism.
\end{proof}

% Atomic construction leaf for EFI SITE-08.

\begin{definition}[Derived etale-to-pro-etale cohomology comparison]
\label{def:weil-efi-site-derived-cohomology-comparison}
\uses{def:weil-pbc-epsilon-comparison,
  def:weil-pbc-epsilon-unit-counit}
For a commutative ring $A$ and $L\in D^+(X_{\mathrm{et}},A)$, let
\[
 u^R_{X,L}:L\longrightarrow
 R\epsilon_{X,*}\mathbf L\epsilon_X^{-1}L
\]
be the derived adjunction unit.  Define
\[
 c_{X,L}:R\Gamma(X_{\mathrm{et}},L)
 \longrightarrow
 R\Gamma(X_{\mathrm{proet}},\mathbf L\epsilon_X^{-1}L)
\]
as $R\Gamma(X_{\mathrm{et}},u^R_{X,L})$ followed by the canonical identity
$R\Gamma(X_{\mathrm{et}},R\epsilon_{X,*}(-))
=R\Gamma(X_{\mathrm{proet}},-)$.  Write $c^n_{X,L}$ for its map on
$n$th cohomology.
\end{definition}

\begin{lemma}[Acyclicity of pulled-back injectives]
\label{lemma:weil-pbc-epsilon-injective-acyclicity}
\uses{}

Let $U$ be affine and etale over $X$, and let $I$ be an injective
$A$-module sheaf on $U_{\mathrm{et}}$.  Then
\[
 H^r(U_{\mathrm{proet}},\epsilon_U^{-1}I)=0\qquad(r>0).
\]
Consequently $R^r\epsilon_{X,*}(\epsilon_X^{-1}J)=0$ for every injective
$J$ on $X_{\mathrm{et}}$ and every $r>0$.
\end{lemma}
\begin{proof}
\uses{lemma:weil-pbc-epsilon-unit-isomorphism,lemma:weil-pbc-epsilon-scheme-naturality}

On $U_{\mathrm{proet}}$, take the basis $\mathcal B$ of affine objects
$V=\varprojlim_jV_j$ for which every $V_j\to U$ is affine etale and the
transition maps are affine.  Use finite pro-etale covers
$\{V_a\to V\}_{a=1}^m$ inside this basis.  Write each $V_a$ as an inverse
limit of affine etale $V_{a,j}\to V$.  Fiber products commute with these
limits, so every term of the Cech complex of the cover is the filtered colimit
of the corresponding terms for the finite-stage families
$\{V_{a,j}\to V_j\}$.

After increasing $j$, all finitely many arrows and their joint surjectivity
descend, so $\{V_{a,j}\to V_j\}$ is an etale covering.  The comparison unit
identifies the restriction of $\epsilon_U^{-1}I$ on every finite-stage etale
object with $I$.  Hence
\[
 \check H^r(\{V_a\to V\},\epsilon_U^{-1}I)
 =\varinjlim_j\check H^r(\{V_{a,j}\to V_j\},I)=0
 \qquad(r>0).
\]
The last equality uses exactness of filtered colimits and Cech acyclicity of
an injective sheaf.  The basis is closed under fiber products, and every
pro-etale object is covered by its members; the basis Cech criterion therefore
gives the displayed cohomology vanishing.

For an affine etale $U\to X$, scheme naturality identifies the restriction
of $\epsilon_X^{-1}J$ with $\epsilon_U^{-1}(J|_U)$.  The restriction of an
injective to the localization site is injective because extension by zero is
its exact left adjoint.  Thus the first part applies.  Affine etale objects
cover every etale object, so the sheaf associated to
$U\mapsto H^r(U_{\mathrm{proet}},\epsilon_U^{-1}(J|_U))$ is zero.  This sheaf
is $R^r\epsilon_{X,*}(\epsilon_X^{-1}J)$.
\end{proof}

\begin{lemma}[Derived comparison for the epsilon unit]
\label{lemma:weil-pbc-epsilon-derived-unit}
\uses{}

For $L\in D^+(X_{\mathrm{et}},A)$, the derived unit
\[
 L\longrightarrow R\epsilon_{X,*}\epsilon_X^{-1}L
\]
is an isomorphism, naturally in $L$.
\end{lemma}
\begin{proof}
\uses{lemma:weil-pbc-epsilon-injective-acyclicity,lemma:weil-pbc-epsilon-unit-isomorphism,lemma:weil-pbc-epsilon-inverse-image-exact}

Choose a bounded-below injective resolution $L\to I^\bullet$.  Exactness of
$\epsilon_X^{-1}$ represents the inverse image by
$\epsilon_X^{-1}I^\bullet$.  Lemma
\ref{lemma:weil-pbc-epsilon-injective-acyclicity} makes every term
$\epsilon_{X,*}$-acyclic.  The bounded-below acyclic-resolution lemma then
computes the derived direct image termwise:
\[
 R\epsilon_{X,*}\epsilon_X^{-1}L
 \simeq\epsilon_{X,*}\epsilon_X^{-1}I^\bullet.
\]
The unit isomorphism identifies the last complex with $I^\bullet$, including
its differentials.  A morphism of bounded-below complexes lifts to the chosen
injective resolutions, and two lifts are homotopic; termwise naturality of the
unit makes the resulting derived isomorphism independent of the lift.
\end{proof}

% Atomic proof leaf for EFI SITE-09.  This proves an isomorphism property;
% it does not treat a comparison Definition as evidence of invertibility.

\begin{proposition}[Finite-level etale/pro-etale cohomology comparison]
\label{prop:weil-efi-site-derived-comparison-isomorphism}
\uses{def:weil-efi-site-derived-cohomology-comparison}
Let $\ell$ be prime, let $r\geq1$, and let $\mathcal A$ be a constructible
$\mathbf Z/\ell^r$-module sheaf on $X_{\mathrm{et}}$.  For $n\geq0$, the map
\[
 c^n_{X,\mathcal A}:H^n(X_{\mathrm{et}},\mathcal A)
 \longrightarrow
 H^n(X_{\mathrm{proet}},\epsilon_X^{-1}\mathcal A)
\]
is an isomorphism.
\end{proposition}
\begin{proof}
\uses{lemma:weil-pbc-epsilon-derived-unit,
  lemma:weil-efi-site-inverse-image-exact}
Exactness identifies
$\mathbf L\epsilon_X^{-1}(\mathcal A[0])$ with
$\epsilon_X^{-1}\mathcal A[0]$.  The derived-unit lemma makes
\[
 \mathcal A[0]\longrightarrow
 R\epsilon_{X,*}\epsilon_X^{-1}\mathcal A[0]
\]
an isomorphism.  Derived global sections preserve this isomorphism, and the
canonical identity
$R\Gamma(X_{\mathrm{et}},R\epsilon_{X,*}(-))
=R\Gamma(X_{\mathrm{proet}},-)$ identifies the resulting arrow with
$c_{X,\mathcal A[0]}$.  Applying $H^n$ gives the displayed isomorphism.
\end{proof}

\begin{definition}[Canonical direct-image comparison]
\label{def:weil-pbc-epsilon-direct-image-comparison}
\uses{def:weil-pbc-epsilon-unit-counit}
Let $f:X\to Y$ be quasi-compact and quasi-separated.  For an $A$-module
sheaf $F$ on $X_{\mathrm{et}}$, define
\[
 \delta_{f,F}:\epsilon_Y^{-1}f_{\mathrm{et},*}F
 \longrightarrow f_{\mathrm{proet},*}\epsilon_X^{-1}F
\]
to be adjoint to
\[
 f_{\mathrm{et},*}F\xrightarrow{f_{\mathrm{et},*}\eta^X_F}
 f_{\mathrm{et},*}\epsilon_{X,*}\epsilon_X^{-1}F
 =\epsilon_{Y,*}f_{\mathrm{proet},*}\epsilon_X^{-1}F.
\]
The equality follows by evaluating both restrictions on an etale object of
$Y$.  Right deriving this natural transformation and using exactness of
$\epsilon_Y^{-1}$ defines
\[
 D_f(L):\epsilon_Y^{-1}Rf_{\mathrm{et},*}L
 \longrightarrow Rf_{\mathrm{proet},*}\epsilon_X^{-1}L
\]
for $L\in D^+(X_{\mathrm{et}},A)$.
\end{definition}


\begin{lemma}[Well-definedness certificate for def:weil-pbc-epsilon-direct-image-comparison]
\label{lemma:weil-sf-well-defined-def-weil-pbc-epsilon-direct-image-comparison}
\uses{def:weil-pbc-epsilon-direct-image-comparison}
The construction in Definition \ref{def:weil-pbc-epsilon-direct-image-comparison} is a typed raw datum.  Its theorem-valued naturality, comparison, descent, or coherence fields are isolated in this property lemma.
\end{lemma}
\begin{proof}
\uses{lemma:weil-pbc-epsilon-scheme-naturality,
  lemma:weil-pbc-epsilon-adjunction-equations-v1,
  lemma:weil-pbc-epsilon-inverse-image-exact}
The underived transformation is the adjoint of $f_*\eta^X$.  Naturality of
$\eta^X$ proves naturality in $F$; the identity formula and the triangle
identities prove the identity and composition equations.  Scheme naturality
identifies the two restrictions in the displayed equality
$f_{\mathrm{et},*}\epsilon_{X,*}=\epsilon_{Y,*}f_{\mathrm{proet},*}$ and
makes the construction compatible with Cartesian pullback.

For the derived map choose a bounded-below injective resolution
$L\to I^\bullet$ on $X_{\mathrm{et}}$, and then a bounded-below injective
resolution $\epsilon_X^{-1}I^\bullet\to J^\bullet$ on
$X_{\mathrm{proet}}$.  Exactness of $\epsilon_X^{-1}$ and
$\epsilon_Y^{-1}$ identifies the two derived source representatives, and the
complex map
\[
 \epsilon_Y^{-1}f_{\mathrm{et},*}I^\bullet
 \xrightarrow{\delta_{f,I^\bullet}}
 f_{\mathrm{proet},*}\epsilon_X^{-1}I^\bullet
 \longrightarrow f_{\mathrm{proet},*}J^\bullet
\]
represents $D_f(L)$.  Comparisons between two choices of either resolution
are unique up to homotopy; naturality of $\delta_f$ carries both comparison
homotopies to the target, so the localized map is independent of every
choice.  The underived
identity, composition, and base-change equations persist after localization.
Thus $D_f(L)$ and all its stated coherence fields are well-defined.
\end{proof}
\begin{lemma}[Site-local components of direct-image comparison]
\label{lemma:weil-pbc-epsilon-direct-image-site-comparison}
\uses{def:weil-pbc-epsilon-direct-image-comparison}
Suppose $Y$ is affine and let
$W=\varprojlim_iW_i$ be an affine pro-etale object, with $W_i\to Y$ affine
etale and affine transition maps.  Put
$X_i=X\times_YW_i$ and $X_W=X\times_YW$.  Then
\begin{align*}
 (\epsilon_Y^{-1}f_{\mathrm{et},*}F)(W)
   &=\varinjlim_i\Gamma(X_i,F),\\
 (f_{\mathrm{proet},*}\epsilon_X^{-1}F)(W)
   &=\Gamma(X_W,\epsilon_{X_W}^{-1}g_{\mathrm{et}}^{-1}F),
\end{align*}
where $g:X_W\to X$.  Under these identifications, $\delta_{f,F}(W)$ is the
map induced by pullback from the finite stages.  Such $W$ form a covering
basis of $Y_{\mathrm{proet}}$.
\end{lemma}
\begin{proof}
\uses{lemma:weil-pbc-epsilon-scheme-naturality}

The first equality is the left-Kan-extension formula for
$\epsilon_Y^{-1}$: the subsystem of the displayed affine etale stages is
cofinal among etale objects receiving a map from $W$.  The second equality is
the definition of pro-etale direct image, followed by the scheme-naturality
identification of the restricted inverse image.  Expanding the adjoint of
$f_{\mathrm{et},*}\eta^X_F$ shows that a section over $X_i$ is pulled back to
$X_W$; hence its filtered-colimit map is exactly $\delta_{f,F}(W)$.  Every
pro-etale object over affine $Y$ admits a cover by inverse limits of this
form, so these evaluations detect sheaf isomorphisms.
\end{proof}

\begin{lemma}[Higher direct-image comparison on finite sheaves]
\label{lemma:weil-pbc-epsilon-direct-image-higher-sheaves}
\uses{}

Let $A$ be finite, let $f:X\to Y$ be quasi-compact and quasi-separated, and
let $F$ be a constructible $A$-module sheaf on $X_{\mathrm{et}}$.  For every
$r\geq0$, the canonical map
\[
 \epsilon_Y^{-1}R^rf_{\mathrm{et},*}F
 \longrightarrow R^rf_{\mathrm{proet},*}\epsilon_X^{-1}F
\]
is an isomorphism.
\end{lemma}
\begin{proof}
\uses{lemma:weil-pbc-epsilon-direct-image-site-comparison,lemma:weil-pbc-epsilon-derived-unit}

The assertion is Zariski local on $Y$, so take $Y$ affine and use the basis
$W=\varprojlim_iW_i$ from Lemma
\ref{lemma:weil-pbc-epsilon-direct-image-site-comparison}.  The left side at
$W$ is
\[
 \varinjlim_iH^r_{\mathrm{et}}(X_i,F).
\]
For the right side, apply Lemma \ref{lemma:weil-pbc-epsilon-derived-unit} on
$X_W$ to $g_{\mathrm{et}}^{-1}F$.  It identifies its value at $W$ with
\[
 H^r_{\mathrm{proet}}(X_W,\epsilon_{X_W}^{-1}g_{\mathrm{et}}^{-1}F)
 =H^r_{\mathrm{et}}(X_W,g_{\mathrm{et}}^{-1}F).
\]
Because $f$ is quasi-compact and quasi-separated, $X_W$ is the inverse limit
of the $X_i$ with affine transition maps and etale cohomology is continuous
for this system.  Hence
\[
 \varinjlim_iH^r_{\mathrm{et}}(X_i,F)
 \xrightarrow{\sim}H^r_{\mathrm{et}}(X_W,g_{\mathrm{et}}^{-1}F).
\]
The component description in the preceding lemma identifies this
isomorphism with the canonical comparison map.  The basis evaluations are
conservative, proving the sheaf assertion.
\end{proof}

\begin{lemma}[Finite constructibility and lower bounds]
\label{lemma:weil-pbc-epsilon-direct-image-constructible-bounds}
\uses{}

Let $A$ be finite and let $L\in D_c^+(X_{\mathrm{et}},A)$ vanish in degrees
$<a$.  Then $\epsilon_X^{-1}L$ has constructible cohomology and also vanishes
in degrees $<a$.  Both $Rf_{\mathrm{et},*}L$ and
$Rf_{\mathrm{proet},*}\epsilon_X^{-1}L$ vanish in degrees $<a$.
\end{lemma}
\begin{proof}
\uses{lemma:weil-pbc-epsilon-inverse-image-exact,lemma:weil-pbc-epsilon-scheme-naturality}

Exactness gives
\[
 \mathcal H^q(\epsilon_X^{-1}L)
 =\epsilon_X^{-1}\mathcal H^q(L).
\]
For each $q$, choose a finite stratification on which $\mathcal H^q(L)$ is
locally constant with finite stalks.  Scheme naturality identifies its inverse
image on each pro-etale stratum with the same finite locally constant sheaf,
so that stratification witnesses constructibility of the corresponding
cohomology sheaf.  The equality also preserves the lower vanishing bound.
Right-derived direct image has no
negative derived degree; its hypercohomology terms have $p\geq0$ and
$q\geq a$, which gives the same lower bound for both direct images.
\end{proof}

\begin{lemma}[Naturality of direct-image comparison]
\label{lemma:weil-pbc-epsilon-direct-image-naturality}
\uses{def:weil-pbc-epsilon-direct-image-comparison}
The maps $D_f(L)$ are natural in $L$, equal the identity for $f=1_X$, and
respect composition of quasi-compact quasi-separated morphisms.  For a
Cartesian square
\[
\begin{CD}
X'@>{v}>>X\\
@V{f'}VV @VV{f}V\\
Y'@>{u}>>Y,
\end{CD}
\]
the square formed by $u^*D_f(L)$, $D_{f'}(v^*L)$, and the two canonical
base-change transformations commutes.
\end{lemma}
\begin{proof}
\uses{lemma:weil-pbc-epsilon-unit-counit-naturality,lemma:weil-pbc-epsilon-scheme-naturality}

The underived map is adjoint to $f_*\eta^X$.  Naturality of $\eta^X$ gives
naturality in $F$, its identity formula gives the identity case, and the
triangle identity shows that the adjoint for a composite is the ordered
composite of the two adjoints.  Deriving preserves these equations.

For the Cartesian square, expand both routes before taking adjoints.  Scheme
naturality identifies $v^*\eta^X_L$ with $\eta^{X'}_{v^*L}$, and both routes
then consist of this unit followed by restriction along the same projection
$X'\to X$ and the same base-change map.  They are equal before derivation,
hence the derived base-change square commutes.
\end{proof}

\begin{lemma}[Comparison direct-image transport]
\label{lemma:weil-pbc-epsilon-direct-image}
\uses{def:weil-pbc-epsilon-direct-image-comparison}
Let $A$ be finite, let $f:X\to Y$ be quasi-compact and quasi-separated, and
let $L\in D_c^+(X_{\mathrm{et}},A)$.  The canonical map
\[
 D_f(L):\epsilon_Y^{-1}Rf_{\mathrm{et},*}L
 \longrightarrow Rf_{\mathrm{proet},*}\epsilon_X^{-1}L
\]
is an isomorphism, natural in $L$, in $f$, and in Cartesian base change.
\end{lemma}
\begin{proof}
\uses{lemma:weil-pbc-epsilon-direct-image-higher-sheaves,lemma:weil-pbc-epsilon-direct-image-constructible-bounds,lemma:weil-pbc-epsilon-direct-image-naturality,lemma:weil-pbc-hypercohomology-convergence}

Choose $a$ below the cohomology of $L$.  The truncation towers on both sides
give morphisms of hypercohomology spectral sequences whose second pages are
\[
 \epsilon_Y^{-1}R^pf_{\mathrm{et},*}\mathcal H^q(L)
 \longrightarrow
 R^pf_{\mathrm{proet},*}\epsilon_X^{-1}\mathcal H^q(L).
\]
Every component is an isomorphism by Lemma
\ref{lemma:weil-pbc-epsilon-direct-image-higher-sheaves}.  The bounds lemma
puts both spectral sequences in $p\geq0$, $q\geq a$.  Thus each total degree
has a finite exhaustive separated filtration, by the numerical argument of
Lemma \ref{lemma:weil-pbc-hypercohomology-convergence}.  Page induction and
finite-filtration induction make the abutment map an isomorphism in every
degree.  Hence $D_f(L)$ is an isomorphism in the derived category.  Its
naturality statements are Lemma
\ref{lemma:weil-pbc-epsilon-direct-image-naturality}.
\end{proof}

\begin{proposition}[Finite-level pro-etale proper base change]
\label{prop:weil-pbc-proet-finite-level-transport}
\uses{def:weil-pbc-comes-from-etale}
Let $A$ be a finite torsion ring and let $(L,\theta)$ witness that
$K\in D^+(X_{\mathrm{proet}},A)$ comes from the etale site, with
$L\in D_c^+(X_{\mathrm{et}},A)$.  For a proper Cartesian square, the
pro-etale exchange map
\[
 Lg^*Rf_*K\longrightarrow Rf'_*L(g')^*K
\]
is an isomorphism.
\end{proposition}
\begin{proof}
\uses{lemma:weil-pbc-epsilon-pullback,lemma:weil-pbc-epsilon-direct-image,prop:weil-proet-finite-derived-proper-base-change}

Use $\theta$ to replace $K$ by $\epsilon_X^{-1}L$.  Commute
$\epsilon^{-1}$ past the two pullbacks and two direct images using the
preceding lemmas.  The resulting middle arrow is
$\epsilon_{Y'}^{-1}$ applied to the S7 etale exchange map and is therefore an
isomorphism.  Expanding all four comparison isomorphisms shows the transported
arrow is the pro-etale unit--counit exchange map; naturality follows from the
two comparison squares.
\end{proof}

\subsubsection{S9. Derived completion and inverse limits}

\begin{definition}[Adic coefficient datum and reduction tower]
\label{def:weil-pbc-complete-coefficient-datum}
\uses{}
Let $\Lambda$ be Noetherian, let $I=(f_1,\ldots,f_r)\subseteq\Lambda$ have
$\Lambda/I$ torsion, and suppose every $P_n:=\Lambda/I^n$ is perfect over
$\Lambda$.  For $K\in D(X_{\mathrm{proet}},\Lambda)$ put
\[
 K_n:=K\otimes^L_\Lambda P_n.
\]
For $b\geq a$, let $q_{b,a}:K_b\to K_a$ be induced by
$P_b\to P_a$.  The quotient identities give
$q_{a,a}=1$ and $q_{c,a}=q_{b,a}q_{c,b}$.  Associativity also gives the
normalization isomorphisms
\[
 K_b\otimes^L_{P_b}P_a\xrightarrow{\sim}K_a,
\]
and these obey the same cocycle.
\end{definition}

\begin{definition}[Explicit adic derived inverse limit]
\label{def:weil-pbc-adic-derived-inverse-limit}
\uses{def:weil-pbc-complete-coefficient-datum}
For a tower $M_\bullet=(M_n,t_{n+1,n})$, put
\[
 s((x_n)_n)=(t_{n+1,n}(x_{n+1}))_n,
 \qquad
 R\varprojlim_nM_n:=\operatorname{Cone}
 \left(1-s:\prod_nM_n\longrightarrow\prod_nM_n\right)[-1].
\]
Thus there is a functorial distinguished triangle
\[
 R\varprojlim_nM_n\longrightarrow\prod_nM_n
 \xrightarrow{1-s}\prod_nM_n\longrightarrow .
\]
Maps of towers act componentwise on the two products and hence on the cone.
\end{definition}

\begin{lemma}[Milnor sequence on pro-etale test objects]
\label{lemma:weil-pbc-adic-milnor-sequence}
\uses{def:weil-pbc-adic-derived-inverse-limit}
Let $U$ be a weakly contractible quasi-compact object of the pro-etale site.
For every tower $M_\bullet$ and every $i$ there is a functorial exact
sequence
\[
 0\longrightarrow\varprojlim
 {}^{1}_nH^{i-1}(U,M_n)\longrightarrow
 H^i\left(U,R\varprojlim_nM_n\right)\longrightarrow
 \varprojlim_nH^i(U,M_n)\longrightarrow0.
\]
\end{lemma}
\begin{proof}
\uses{} 
Weak contractibility makes $\Gamma(U,-)$ exact, and evaluation on $U$
commutes with products.  Apply it to the product-cone triangle.  The long
exact cohomology sequence identifies the kernel of $1-s$ on
$\prod_nH^i(U,M_n)$ with $\varprojlim H^i(U,M_n)$ and its cokernel in
degree $i-1$ with $\varprojlim{}^1H^{i-1}(U,M_n)$.  Kernel--cokernel
exactness gives the displayed sequence, and every construction is natural
in a map of towers.
\end{proof}

\begin{lemma}[Finite towers satisfy Mittag--Leffler]
\label{lemma:weil-pbc-adic-mittag-leffler}
\uses{}
For an inverse system of finite abelian groups $(A_n)$, the images
$\operatorname{Im}(A_m\to A_n)$ are constant for all sufficiently large
$m\geq n$; hence $(A_n)$ is Mittag--Leffler.
\end{lemma}
\begin{proof}
\uses{} 
For fixed $n$ these images form a descending chain of subgroups of the
finite group $A_n$.  Such a chain stabilizes.  Doing this for every $n$ is
exactly the Mittag--Leffler condition.
\end{proof}

\begin{lemma}[Mittag--Leffler kills the Milnor obstruction]
\label{lemma:weil-pbc-adic-lim-one-zero}
\uses{}

For a countable Mittag--Leffler system $(A_n)$ one has
$\varprojlim{}^1_nA_n=0$.  In particular this holds for a tower of finite
abelian groups.
\end{lemma}
\begin{proof}
\uses{lemma:weil-pbc-adic-mittag-leffler}

Replace $A_n$ by the stable image in $A_n$; this does not change either
$\varprojlim$ or the cokernel description of $\varprojlim{}^1$, and the new
transition maps are surjective.  Given $(y_n)\in\prod_nA_n$, choose $x_1$
arbitrarily and then choose $x_{n+1}$ recursively so that
$x_n-t_{n+1,n}(x_{n+1})=y_n$.  Thus $1-s$ is surjective on the product and
its cokernel, namely $\varprojlim{}^1A_n$, is zero.
\end{proof}

\begin{definition}[Intrinsic Koszul completion]
\label{def:weil-pbc-intrinsic-koszul-completion}
\uses{def:weil-pbc-adic-derived-inverse-limit}
Let
\[
 Q_m:=K(\Lambda;f_1^m,\ldots,f_r^m)
\]
be the Koszul complex in degrees $[-r,0]$.  Its map $Q_{m+1}\to Q_m$ is
the tensor product, over $j$, of the two-term chain maps whose degree-zero
component is the identity and whose degree-minus-one component is
multiplication by $f_j$.  Define
\[
 K^{\wedge,\mathrm{Kos}}:=
 R\varprojlim_m(K\otimes^L_\Lambda Q_m).
\]
The augmentations $\Lambda\to Q_m$ give
$\kappa_K:K\to K^{\wedge,\mathrm{Kos}}$.  Call $K$ intrinsically derived
$I$-complete when $\kappa_K$ is an isomorphism; below, ``derived
$I$-complete'' always means this intrinsic condition.
\end{definition}

\begin{definition}[Completion map]
\label{def:weil-pbc-completion-map}
\uses{def:weil-pbc-adic-derived-inverse-limit}
Define
\[
 \operatorname{comp}_K:K\longrightarrow R\varprojlim_nK_n
\]
as the map induced by the compatible quotient arrows $K\to K_n$.
\end{definition}

\begin{lemma}[Koszul and quotient towers have the same limit]
\label{lemma:weil-pbc-koszul-quotient-limit}
\uses{def:weil-pbc-intrinsic-koszul-completion,def:weil-pbc-completion-map}
There is a functorial isomorphism
\[
 c_K:K^{\wedge,\mathrm{Kos}}\xrightarrow{\sim}
 R\varprojlim_nK_n
\]
such that $c_K\kappa_K=\operatorname{comp}_K$.
\end{lemma}
\begin{proof}
\uses{} 
Put $J_m=(f_1^m,\ldots,f_r^m)$.  The augmentation
$Q_m\to\Lambda/J_m$ has cone $\tau_{<0}Q_m$.  For $p<0$ the modules
$H^p(Q_m)$ are finite and killed by $J_m$.  The transition in degree $p$
has entries in $(f_1^{b-a},\ldots,f_r^{b-a})$ from level $b$ to level $a$.
Artin--Rees, applied to the cycles inside the finite free term $Q_a^p$,
therefore makes $H^p(Q_b)\to H^p(Q_a)$ zero for $b\gg a$.  Hence
$(\tau_{<0}Q_m)_m$ is pro-zero and
$(Q_m)_m\to(\Lambda/J_m)_m$ is a pro-isomorphism.

The inclusions
\[
 I^{r(m-1)+1}\subseteq J_m\subseteq I^m
\]
show that $(\Lambda/J_m)_m$ and $(P_n)_n$ are cofinal pro-systems.  Tensor
the comparison maps with $K$ and apply the product-cone limit.  Pro-zero
cones have zero product-cone limit after a cofinal reindexing, so the result
is $c_K$.  Both units are induced by the same augmentations
$\Lambda\to Q_m\to\Lambda/J_m\to P_n$; this proves the triangle equation.
\end{proof}

\begin{lemma}[Derived-complete reconstruction]
\label{lemma:weil-pbc-derived-complete-reconstruction}
\uses{def:weil-pbc-intrinsic-koszul-completion}
If $K$ is intrinsically derived $I$-complete, then
\[
 \operatorname{comp}_K:K\xrightarrow{\sim}R\varprojlim_nK_n
\]
is an isomorphism, functorially in $K$.
\end{lemma}
\begin{proof}
\uses{lemma:weil-pbc-koszul-quotient-limit}

By intrinsic completeness, $\kappa_K$ is an isomorphism.  The preceding
lemma identifies $\operatorname{comp}_K$ with the composite
$c_K\kappa_K$, so it too is an isomorphism.  For $u:K\to L$, naturality of
the Koszul augmentations and of the product-cone construction gives
\[
 (R\varprojlim u_n)\operatorname{comp}_K
 =\operatorname{comp}_Lu.
\]
Conjugating by the two reconstruction isomorphisms proves functoriality,
including functoriality of their inverses.  This is the explicit
Koszul-to-adic argument behind Tag \texttt{099M}.
\end{proof}

\begin{lemma}[Finite reduction of adic completion]
\label{lemma:weil-pbc-adic-completion-reduction}
\uses{}

For every $a\geq1$, projection to level $a$ induces a functorial
isomorphism
\[
 \left(R\varprojlim_nK_n\right)\otimes^L_\Lambda P_a
 \xrightarrow{\sim}K_a.
\]
\end{lemma}
\begin{proof}
\uses{lemma:weil-pbc-koszul-quotient-limit}

Use the preceding comparison to replace the quotient tower by
$K\otimes^LQ_m$.  Since $P_a$ is perfect, tensoring with $P_a$ commutes
with both products in the product-cone model.  For $m\geq a$, the
differential of $Q_m\otimes^L P_a$ is zero.  Its degree-zero transition is
the identity, whereas every composite of $a$ successive transitions in a
negative degree contains a factor $f_j^a\in I^a$ and is zero over $P_a$.
Thus the negative-degree towers are pro-zero and the degree-zero tower is
the constant tower $P_a$.  After tensoring with $K$, its product-cone limit
is $K\otimes^LP_a=K_a$.  The surviving degree-zero projection is precisely
the displayed map.  This also exposes the finite Tor-amplitude input: the
bounded finite-projective complex representing $P_a$ is what permits it to
pass through the two countable products.
\end{proof}

\begin{definition}[Uniformly bounded constructible adic tower]
\label{def:weil-pbc-bounded-constructible-adic-tower}
\uses{def:weil-pbc-complete-coefficient-datum}
Say that the reduction tower $K_\bullet$ attached to $K$ is uniformly
bounded constructible if there are integers $d_0\leq d_1$ and one finite
stratification of $X$, each of whose strata has finitely many connected
components, such that, for every $n$, $H^i(K_n)=0$ outside
$[d_0,d_1]$ and each $H^i(K_n)$ is finite lisse on every stratum.  The
normalization isomorphisms and their cocycle are part of this datum.
\end{definition}

\begin{lemma}[Bounded constructible towers have no Milnor obstruction]
\label{lemma:weil-pbc-constructible-adic-lim-one-zero}
\uses{def:weil-pbc-bounded-constructible-adic-tower}
If the reduction tower $K_\bullet$ attached to $K$ is uniformly bounded
constructible, then on every weakly contractible quasi-compact $U$ lying
over one of the displayed strata,
\[
 H^i\left(U,R\varprojlim_nK_n\right)
 \xrightarrow{\sim}\varprojlim_nH^i(U,K_n),
\]
and both sides vanish outside $[d_0,d_1]$.
\end{lemma}
\begin{proof}
\uses{lemma:weil-pbc-adic-milnor-sequence,lemma:weil-pbc-adic-lim-one-zero}

Fix $n$ and $i$.  As $m$ increases, the image sheaves of
 $H^i(K_m)\to H^i(K_n)$ form a descending chain of finite lisse subsheaves on
each displayed stratum.  On each connected component their fibers are
descending chains of subgroups of a finite group, so they stabilize.  There
are only finitely many strata and, by hypothesis, finitely many connected
components on each stratum; hence one index works globally.  Exactness of
$\Gamma(U,-)$ identifies the
image of $H^i(U,K_m)\to H^i(U,K_n)$ with the sections of that image sheaf.
Consequently $(H^i(U,K_n))_n$ is Mittag--Leffler even when $U$ has
infinitely many connected components.  The vanishing lemma now kills its
$\varprojlim{}^1$.  The Milnor sequence is therefore the displayed
isomorphism, and uniform finite-level vanishing gives the last assertion.
\end{proof}

\begin{lemma}[Bounded constructibility survives reconstruction]
\label{lemma:weil-pbc-bounded-constructible-reconstruction}
\uses{}

If the reduction tower $K_\bullet$ attached to $K$ is uniformly bounded
constructible, then
$\widehat K:=R\varprojlim_nK_n$ is intrinsically derived $I$-complete, has
finite reduction $\widehat K\otimes^LP_n\simeq K_n$, and retains the common
cohomological interval $[d_0,d_1]$ on the conservative weakly contractible
basis subordinate to the displayed stratification.
\end{lemma}
\begin{proof}
\uses{lemma:weil-pbc-adic-completion-reduction,lemma:weil-pbc-constructible-adic-lim-one-zero,lemma:weil-pbc-koszul-quotient-limit}

Finite reduction is the preceding reduction lemma.  Hence the quotient
completion of $\widehat K$ is the limit of the same tower and its completion
map is the identity under the product-cone model.  The Koszul--quotient
comparison then makes its intrinsic Koszul unit an isomorphism.  The common
cohomological bounds and the absence of a hidden degree from
$\varprojlim{}^1$ are exactly the conclusion of the preceding lemma.
Weakly contractible objects subordinate to a finite stratification form a
conservative basis, so these bounds detect the bounds of $\widehat K$.
\end{proof}

\begin{definition}[Completed pullback]
\label{def:weil-pbc-completed-pullback}
\uses{def:weil-pbc-complete-coefficient-datum}
For $g:Y'\to Y$ define
\[
 Lg^*_{\mathrm{comp}}K:=
 R\varprojlim_n\bigl(Lg^*K\otimes^L_\Lambda\Lambda/I^n\bigr).
\]
The comparison from ordinary pullback is the completion map of $Lg^*K$.
\end{definition}

\begin{lemma}[Reduction commutes with completed pullback]
\label{lemma:weil-pbc-completed-pullback-reduction}
\uses{def:weil-pbc-completed-pullback}
For every $n$,
\[
 (Lg^*_{\mathrm{comp}}K)\otimes^L_\Lambda\Lambda/I^n
 \xrightarrow{\sim}Lg^*K_n.
\]
\end{lemma}
\begin{proof}
\uses{lemma:weil-pbc-adic-completion-reduction}

Derived pullback commutes with derived tensor, so its defining tower is the
reduction tower of $Lg^*K$.  Apply finite reduction of adic completion to
this object.  Its level $n$ is
$Lg^*K\otimes^L_\Lambda P_n\simeq Lg^*(K\otimes^L_\Lambda P_n)=Lg^*K_n$.
\end{proof}

\begin{definition}[Adic operator and base-change datum]
\label{def:weil-pbc-adic-operator-base-change-datum}
\uses{def:weil-pbc-completed-pullback}
Let $g:Y'\to Y$, let $K$ be on $Y_{\mathrm{proet}}$, and let $L$ be on
$Y'_{\mathrm{proet}}$.  An adic operator and base-change datum consists of
compatible tower endomorphisms
\[
 \phi_n:K_n\to K_n,\qquad \psi_n:L_n\to L_n,
\]
and compatible maps $b_n:Lg^*K_n\to L_n$ such that
\[
 b_n\,Lg^*(\phi_n)=\psi_n b_n
\]
for every $n$.  Compatibility includes the transition squares, the
normalization squares, and their cocycles.  This is a datum of ordinary
endomorphisms.  A genuine Frobenius or Weil structure has the different type
$\operatorname{Fr}^*K_n\to K_n$ and is not identified with this datum.
\end{definition}

\begin{lemma}[Reconstruction preserves operators and base change]
\label{lemma:weil-pbc-adic-reconstruction-compatibility}
\uses{def:weil-pbc-adic-operator-base-change-datum}
An adic operator and base-change datum induces maps
\[
 \widehat\phi=R\varprojlim_n\phi_n,\quad
 \widehat\psi=R\varprojlim_n\psi_n,\quad
 \widehat{Lg^*\phi}=R\varprojlim_nLg^*\phi_n,\quad
 \widehat b=R\varprojlim_nb_n:
 Lg^*_{\mathrm{comp}}K\longrightarrow\widehat L
\]
satisfying
$\widehat b\,\widehat{Lg^*\phi}
=\widehat\psi\widehat b$.  Under the canonical finite-reduction
isomorphisms, reduction modulo $I^n$ gives exactly
$\phi_n,\psi_n,b_n$.  If $K$ and $L$ are intrinsically derived complete,
the reconstruction isomorphisms identify these with operators and a
base-change map on $K$ and $L$ themselves.
\end{lemma}
\begin{proof}
\uses{lemma:weil-pbc-adic-completion-reduction}

Compatibility with transitions makes the three families maps of towers.
They therefore give maps between the two product-cone diagrams.  The
levelwise equations give an equality on both product terms, hence on their
fibres, proving the displayed completed equation.  The finite-reduction
lemma identifies projection to the $n$th factor with reduction modulo
$I^n$, so the reduced maps are precisely the named level maps.  Finally,
conjugate by the functorial reconstruction isomorphisms for complete $K$
and $L$; the same equality survives conjugation.
\end{proof}

\begin{definition}[Local strict-perfect coefficient package]
\label{def:weil-pbc-strict-perfect-coefficient-package}
\uses{}
On a localization of a ringed topos, let $Q$ be represented by a bounded
complex $F^\bullet$ of finite projective $\Lambda$-modules.  A strict-perfect
package for this representation consists, for every $r$, of an integer
$N_r$ and maps
\[
 F^r\xrightarrow{i_r}\Lambda^{N_r}
 \xrightarrow{p_r}F^r,\qquad p_ri_r=1.
\]
Write $e_r=i_rp_r$; then $e_r^2=e_r$, $p_re_r=p_r$, and $e_ri_r=i_r$.
Perfectness means that such packages exist after a covering.  When
$Q\in D(\Lambda)$ is a perfect coefficient complex, one global package and
its associated constant sheaf may be used.
\end{definition}

\begin{definition}[Projection-formula arrows]
\label{def:weil-pbc-projection-formula-arrows}
\uses{def:weil-pbc-strict-perfect-coefficient-package}
For a morphism of ringed topoi $h:\mathcal T\to\mathcal S$, an object
$E\in D(\mathcal T,\Lambda)$, and a strict-perfect package $F^\bullet$ on
$\mathcal S$, define
\[
 \pi_{h,E,F}:Rh_*E\otimes^L_\Lambda F^\bullet
 \longrightarrow Rh_*(E\otimes^L_\Lambda Lh^*F^\bullet)
\]
to be adjoint to
\[
 Lh^*(Rh_*E\otimes^LF^\bullet)
 \simeq Lh^*Rh_*E\otimes^LLh^*F^\bullet
 \xrightarrow{\epsilon_E\otimes1}
 E\otimes^LLh^*F^\bullet.
\]
For $F^r=\Lambda^{N_r}$, let $\rho_r$ be the inverse of the finite-sum
identification
\[
 (Rh_*E)^{N_r}\xrightarrow{\pi_r}Rh_*(E^{N_r}).
\]
For a projective term with the displayed splitting, define
\[
 \rho_{F^r}:=
 (Rh_*E\otimes p_r)\,\rho_r\,
 Rh_*(1_E\otimes h^*i_r).
\]
The bounded totalization of the $\rho_{F^r}$ is denoted
$\rho_{h,E,F}$.
\end{definition}

\begin{lemma}[Finite-free projection formula]
\label{lemma:weil-pbc-projection-formula-finite-free}
\uses{def:weil-pbc-projection-formula-arrows}
For $F=\Lambda^N[a]$, the maps $\pi_{h,E,F}$ and $\rho_{h,E,F}$ are inverse:
\[
 \rho_{h,E,F}\pi_{h,E,F}=1,\qquad
 \pi_{h,E,F}\rho_{h,E,F}=1.
\]
\end{lemma}
\begin{proof}
\uses{} 
Both source and target are the shift by $a$ of a finite direct sum of $N$
copies of $Rh_*E$.  A right adjoint preserves finite products, and in an
additive category finite products are finite coproducts.  Under these two
identifications, the counit defining $\pi$ is the identity in every summand.
The map $\rho$ is the reverse finite-sum identification, so the two displayed
composites are componentwise identities.
\end{proof}

\begin{lemma}[Finite-projective summand equations]
\label{lemma:weil-pbc-projection-formula-projective-summand}
\uses{def:weil-pbc-projection-formula-arrows}
For a finite projective $F^r$ and a displayed splitting $(i_r,p_r)$, the
arrows $\pi_{h,E,F^r}$ and $\rho_{F^r}$ are inverse and are independent of
the chosen finite-free summand presentation.
\end{lemma}
\begin{proof}
\uses{lemma:weil-pbc-projection-formula-finite-free}

Let $A(u)=Rh_*E\otimes u$ and
$B(u)=Rh_*(E\otimes h^*u)$.  Naturality of the counit gives
\[
 B(i_r)\pi_{F^r}=\pi_rA(i_r),\qquad
 B(p_r)\pi_r=\pi_{F^r}A(p_r).
\]
Using $p_ri_r=1$ and the finite-free inverse equations, compute
\[
 \rho_{F^r}\pi_{F^r}
 =A(p_r)\rho_rB(i_r)\pi_{F^r}
 =A(p_r)\rho_r\pi_rA(i_r)=1,
\]
and
\[
 \pi_{F^r}\rho_{F^r}
 =\pi_{F^r}A(p_r)\rho_rB(i_r)
 =B(p_r)\pi_r\rho_rB(i_r)=1.
\]
Thus $\rho_{F^r}$ is the unique inverse of the intrinsic map $\pi_{F^r}$;
uniqueness removes the dependence on $(i_r,p_r)$.
\end{proof}

\begin{lemma}[Signed totalization of the perfect formula]
\label{lemma:weil-pbc-projection-formula-totalization}
\uses{def:weil-pbc-projection-formula-arrows}
For a bounded finite-projective complex $F^\bullet$, the total arrows
$\pi_{h,E,F}$ and $\rho_{h,E,F}$ are inverse isomorphisms.
\end{lemma}
\begin{proof}
\uses{lemma:weil-pbc-projection-formula-projective-summand}

Choose a complex representing $Rh_*E$.  On a homogeneous tensor
$x\otimes y$ with $x$ of degree $a$, both total complexes use
\[
 d(x\otimes y)=d(x)\otimes y+(-1)^a x\otimes d_F(y).
\]
Naturality of $\pi_{F^r}$ in a map of finite projective modules and uniqueness
of its inverse give naturality of $\rho_{F^r}$.  Apply these two naturality
statements to every differential $d_F^r:F^r\to F^{r+1}$.  They say the term
maps commute with the horizontal differentials; their degree zero says they
commute with the signed vertical differentials.  Hence their finite
totalizations are chain maps.  On every bidegree the two composites are the
identities proved in the summand lemma.  Finite totalization preserves these
equalities, so
\[
 \rho_{h,E,F}\pi_{h,E,F}=1,
 \qquad \pi_{h,E,F}\rho_{h,E,F}=1.
\]
\end{proof}

\begin{definition}[Chosen compactification and compact-support image]
\label{def:weil-pbc-projection-compactification}
\uses{}
Let $f:X\to Y$ admit a chosen compactification
\[
 X\xrightarrow{j}\overline X\xrightarrow{p}Y,
 \qquad f=pj,
\]
with $j$ a quasi-compact open immersion and $p$ proper.  For this datum set
\[
 Rf_!^{(j,p)}E:=Rp_*j_!E.
\]
For proper $f$, its identity compactification is $(j,p)=(1_X,f)$.
\end{definition}

\begin{lemma}[Open-immersion tensor piece]
\label{lemma:weil-pbc-projection-formula-open-piece}
\uses{def:weil-pbc-projection-compactification}
For a bounded finite-projective complex $F^\bullet$ on $\overline X$, the
canonical map
\[
 (j_!E)\otimes^LF^\bullet
 \xrightarrow{\sim}j_!(E\otimes^LLj^*F^\bullet)
\]
is an isomorphism, natural in $E$ and $F^\bullet$.
\end{lemma}
\begin{proof}
\uses{} 
Extension by zero along an open immersion is exact.  Resolve only $E$;
tensoring with each finite projective term remains exact.  At a point of the
pro-etale topos lying over $X$, both sides are $E_x\otimes F_x^\bullet$ and
the map is the identity.  At a point lying over
$\overline X\setminus X$, both sides are zero.  These points form a
conservative family, so the term map is an isomorphism.  It commutes with
$d_F$ and with the signed total differential, hence gives the displayed
derived isomorphism.
\end{proof}

\begin{lemma}[Proper compactification piece]
\label{lemma:weil-pbc-projection-formula-proper-piece}
\uses{def:weil-pbc-projection-compactification}
For the proper piece $p$ and every strict-perfect $F^\bullet$ on $Y$,
\[
 (Rp_*j_!E)\otimes^LF^\bullet
 \xrightarrow[\pi_{p,j_!E,F}]{\sim}
 Rp_*(j_!E\otimes^LLp^*F^\bullet),
\]
with inverse $\rho_{p,j_!E,F}$.
\end{lemma}
\begin{proof}
\uses{lemma:weil-pbc-projection-formula-totalization}

Apply the signed-totalization lemma to the morphism of ringed pro-etale
topoi induced by $p$ and the object $j_!E$.
\end{proof}

\begin{lemma}[Compactified finite-level projection formula]
\label{lemma:weil-pbc-projection-formula-compactified}
\uses{}

For a chosen compactification and a perfect $Q$ on $Y$, the ordered composite
\[
\begin{aligned}
 (Rf_!^{(j,p)}E)\otimes^LQ
 &=(Rp_*j_!E)\otimes^LQ\\
 &\xrightarrow{\pi_{p,j_!E,Q}}
 Rp_*(j_!E\otimes^LLp^*Q)\\
 &\xrightarrow{Rp_*\theta_j}
 Rp_*j_!(E\otimes^LLf^*Q)
 =Rf_!^{(j,p)}(E\otimes^LLf^*Q)
\end{aligned}
\]
is an isomorphism.  Here $\theta_j$ is the open-piece isomorphism and the
inverse is $\rho_{p,j_!E,Q}\,Rp_*(\theta_j^{-1})$ in reverse order.
\end{lemma}
\begin{proof}
\uses{lemma:weil-pbc-projection-formula-open-piece,lemma:weil-pbc-projection-formula-proper-piece}

The assertion is local on $Y$.  On a cover choose a strict-perfect package
for $Q$.  The first arrow and its displayed inverse are the proper-piece
inverse pair.  The second arrow and its inverse come from the open-piece
inverse pair, using
$Lj^*Lp^*Q=Lf^*Q$.  Associativity gives
\[
 \rho_{p,j_!E,Q}(Rp_*\theta_j^{-1})
 (Rp_*\theta_j)\pi_{p,j_!E,Q}=1
\]
and
\[
 (Rp_*\theta_j)\pi_{p,j_!E,Q}
 \rho_{p,j_!E,Q}(Rp_*\theta_j^{-1})=1.
\]
These are precisely the two inverse equations for the ordered composite.
They prove that the globally defined forward map is an isomorphism; uniqueness
of an inverse glues the local reverse maps.
\end{proof}

\begin{lemma}[Identity compactification recovers direct image]
\label{lemma:weil-pbc-projection-formula-identity-compactification}
\uses{}

If $f$ is proper, the identity compactification identifies $Rf_!^{(1,f)}$
with $Rf_*$, makes $\theta_1$ the identity, and identifies the compactified
formula with $\pi_{f,E,Q}$ and its inverse with $\rho_{f,E,Q}$.
\end{lemma}
\begin{proof}
\uses{lemma:weil-pbc-projection-formula-compactified}

For $j=1_X$, extension by zero, inverse image, and both unit constraints are
identities.  Substitution into the ordered composite leaves
$\pi_{f,E,Q}$; substitution into its inverse leaves $\rho_{f,E,Q}$.
\end{proof}

\begin{lemma}[Projection formulas form a coefficient tower]
\label{lemma:weil-pbc-projection-formula-tower-compatibility}
\uses{def:weil-pbc-complete-coefficient-datum,def:weil-pbc-projection-formula-arrows}
For every $n$, let $q_n:P_{n+1}\to P_n$ be the quotient.  The square
\[
\begin{CD}
 (Rf_*K)\otimes^LP_{n+1} @>{\pi_{n+1}}>>
 Rf_*(K\otimes^LP_{n+1})\\
 @V{1\otimes q_n}VV @VV{Rf_*(1\otimes q_n)}V\\
 (Rf_*K)\otimes^LP_n @>{\pi_n}>> Rf_*(K\otimes^LP_n)
\end{CD}
\]
commutes, and so does the square with every $\pi$ replaced by $\rho$ and
horizontal arrows reversed; its vertical arrows remain the two quotient
maps.
\end{lemma}
\begin{proof}
\uses{} 
The two adjoints of the forward composites are both
\[
 Lf^*Rf_*K\otimes^LP_{n+1}
 \xrightarrow{\epsilon_K\otimes1}K\otimes^LP_{n+1}
 \xrightarrow{1\otimes q_n}K\otimes^LP_n.
\]
Thus the forward square commutes.  Invert the two horizontal isomorphisms in
the equality
$Rf_*(1\otimes q_n)\pi_{n+1}=\pi_n(1\otimes q_n)$ to obtain the asserted
compatibility for $\rho$.
\end{proof}

\begin{lemma}[Operator and base-change naturality of the formula]
\label{lemma:weil-pbc-projection-formula-base-change-naturality}
\uses{def:weil-pbc-cartesian-square,def:weil-pbc-projection-formula-arrows}
The projection-formula arrows are natural in $K$; hence for every
endomorphism $\phi:K\to K$ with reductions $\phi_n:K_n\to K_n$,
\[
 Rf_*\phi_n\,\pi_n=\pi_n((Rf_*\phi)\otimes1),
 \qquad
 ((Rf_*\phi)\otimes1)\rho_n=\rho_nRf_*\phi_n.
\]
For a Cartesian square, the following two routes are equal:
\[
\begin{aligned}
 Lg^*((Rf_*K)\otimes^LQ)
 &\xrightarrow{Lg^*\pi_f}Lg^*Rf_*(K\otimes^LLf^*Q)
 \xrightarrow{\operatorname{Ex}_{K\otimes Lf^*Q}}
 Rf'_*L(g')^*(K\otimes^LLf^*Q),\\
 Lg^*((Rf_*K)\otimes^LQ)
 &\xrightarrow{\sim}(Lg^*Rf_*K)\otimes^LLg^*Q
 \xrightarrow{\operatorname{Ex}_K\otimes1}
 (Rf'_*L(g')^*K)\otimes^LLg^*Q\\
 &\xrightarrow{\pi_{f'}}
 Rf'_*(L(g')^*K\otimes^LL(f')^*Lg^*Q).
\end{aligned}
\]
The targets are identified by the chosen Cartesian isomorphism.
\end{lemma}
\begin{proof}
\uses{} 
For a map $u:K\to L$, both naturality composites are adjoint to
\[
 Lf^*Rf_*K\otimes Lf^*Q
 \xrightarrow{\epsilon_K\otimes1}K\otimes Lf^*Q
 \xrightarrow{u\otimes1}L\otimes Lf^*Q.
\]
This proves the forward operator equation; conjugation by $\rho_n=\pi_n^{-1}$
gives the second.  For base change, expand both exchange maps as unit followed
by counit.  After adjunction to $X'$, both displayed routes become
\[
 L(g')^*Lf^*Rf_*K\otimes L(g')^*Lf^*Q
 \xrightarrow{L(g')^*(\epsilon_K)\otimes1}
 L(g')^*K\otimes L(g')^*Lf^*Q.
\]
The Cartesian identification
$L(g')^*Lf^*=L(f')^*Lg^*$ identifies the coefficient factors.  Hence the
adjoints, and therefore the two routes, are equal.
\end{proof}

\begin{lemma}[Perfect projection formula for reductions]
\label{lemma:weil-pbc-perfect-projection-formula}
\uses{def:weil-pbc-complete-coefficient-datum}
For every $n$ there is a natural isomorphism
\[
 (Rf_*K)\otimes^L_\Lambda\Lambda/I^n
 \xrightarrow[\pi_n]{\sim}Rf_*K_n
\]
whose inverse is $\rho_n$.  The families $(\pi_n)$ and $(\rho_n)$ commute
with the reduction tower, compatible coefficient endomorphisms, and base
change.
\end{lemma}
\begin{proof}
\uses{lemma:weil-pbc-projection-formula-totalization,lemma:weil-pbc-projection-formula-identity-compactification,lemma:weil-pbc-projection-formula-tower-compatibility,lemma:weil-pbc-projection-formula-base-change-naturality}

Choose a strict-perfect package $F_n^\bullet\to P_n$.  Signed totalization
gives inverse arrows
\[
 (Rf_*K)\otimes^LF_n^\bullet
 \mathrel{\substack{\xrightarrow{\ \pi_n\ }\\[-2pt]
 \xleftarrow[\ \rho_n\ ]{}}}
 Rf_*(K\otimes^LLf^*F_n^\bullet).
\]
Because $P_n$ is constant, $Lf^*P_n=P_n$, and the target is $Rf_*K_n$.
The identity-compactification lemma shows that this is also the open/proper
compactification composite for proper $f$.  The totalization lemma gives
$\rho_n\pi_n=1$ and $\pi_n\rho_n=1$.  Uniqueness of an inverse makes the map
independent of the strict-perfect package.  Tower, operator, and base-change
compatibility are exactly the preceding two lemmas.  This expands the
finite-projective argument located at Tag \texttt{0944}.
\end{proof}

\begin{lemma}[Direct image commutes with sequential homotopy limits]
\label{lemma:weil-pbc-direct-image-r-lim}
\uses{}
For an inverse tower $M_n$ on $X_{\mathrm{proet}}$,
\[
 Rf_*R\varprojlim_nM_n\xrightarrow{\sim}
 R\varprojlim_nRf_*M_n.
\]
\end{lemma}
\begin{proof}
\uses{} 
Use the product/equalizer model
\[
 R\varprojlim M_n\longrightarrow\prod_nM_n
 \xrightarrow{1-\mathrm{shift}}\prod_nM_n
\]
as a distinguished triangle.  Since $Rf_*$ is a right adjoint, it commutes
with products.  Applying it to the triangle gives the product/equalizer model
for the right-hand limit, and the identity maps on the two products give the
isomorphism of fibres.
\end{proof}

\begin{definition}[Completed tensor product]
\label{def:weil-pbc-completed-tensor-product}
\uses{def:weil-pbc-complete-coefficient-datum,def:weil-pbc-adic-derived-inverse-limit}
For $A,B$ on the same ringed pro-etale topos, define
\[
 A\widehat\otimes^L_\Lambda B:=
 R\varprojlim_n(A\otimes^L_\Lambda B\otimes^L_\Lambda P_n).
\]
Thus $A\widehat\otimes^L\Lambda$ is the derived $I$-completion of $A$.
\end{definition}

\begin{lemma}[Finite reduction of completed tensor]
\label{lemma:weil-pbc-completed-tensor-reduction}
\uses{def:weil-pbc-completed-tensor-product}
For every $a$ there is a functorial isomorphism
\[
 (A\widehat\otimes^LB)\otimes^LP_a
 \xrightarrow{\sim}A\otimes^LB\otimes^LP_a.
\]
\end{lemma}
\begin{proof}
\uses{lemma:weil-pbc-adic-completion-reduction}

The completed tensor is the adic completion of $A\otimes^LB$.  Apply finite
reduction of adic completion to that object and use associativity of derived
tensor product.
\end{proof}

\begin{lemma}[Completed projection-formula composite]
\label{lemma:weil-pbc-completed-projection-formula-composite}
\uses{def:weil-pbc-completed-tensor-product}
Let $B$ be an object on $Y_{\mathrm{proet}}$ such that
$B_n:=B\otimes^LP_n$ is perfect for every $n$.  Locally choose a
strict-perfect package for each $B_n$.  Then the ordered composite
\[
\begin{aligned}
 (Rf_*E)\widehat\otimes^LB
 &=R\varprojlim_n((Rf_*E)\otimes^LB_n)\\
 &\xrightarrow{R\varprojlim\pi_{E,B_n}}
 R\varprojlim_nRf_*(E\otimes^LLf^*B_n)\\
 &\xrightarrow{(Rf_*\text{-}R\varprojlim)^{-1}}
 Rf_*R\varprojlim_n(E\otimes^LLf^*B\otimes^LP_n)\\
 &=Rf_*(E\widehat\otimes^LLf^*B)
\end{aligned}
\]
is an isomorphism, natural in $E,B$, compatible endomorphisms, and base-change
squares.
Its inverse is the reverse composite using direct-image/limit commutation
and $R\varprojlim\rho_{E,B_n}$.
\end{lemma}
\begin{proof}
\uses{lemma:weil-pbc-projection-formula-totalization,lemma:weil-pbc-projection-formula-tower-compatibility,lemma:weil-pbc-direct-image-r-lim}

Naturality in the coefficient maps $B_{n+1}\to B_n$ makes the levelwise
projection formulas a map of towers.  Every $\pi_{E,B_n}$ is an isomorphism
locally by signed totalization, with inverse $\rho_{E,B_n}$; uniqueness of
an inverse glues the local inverses.  The second displayed arrow is an
isomorphism because direct image commutes with the product-cone limit.  The
forward composite followed by the stated reverse composite is
\[
 R\varprojlim(\rho_{E,B_n}\pi_{E,B_n})=R\varprojlim1=1;
\]
the reverse order uses
$R\varprojlim(\pi_{E,B_n}\rho_{E,B_n})=1$ and the two inverse equations for
direct-image/limit commutation.  Levelwise operator and base-change
naturality passes through products and the product-cone fibre, proving the
last assertion.
\end{proof}

\begin{lemma}[Direct image commutes with derived completion]
\label{lemma:weil-pbc-direct-image-derived-completion}
\uses{}

There is a functorial isomorphism
\[
 \widehat{Rf_*E}\xrightarrow{\sim}Rf_*\widehat E.
\]
If $E$ is intrinsically derived $I$-complete, then $Rf_*E$ is intrinsically
derived $I$-complete.  Both statements commute with compatible endomorphisms
and base change.
\end{lemma}
\begin{proof}
\uses{lemma:weil-pbc-completed-projection-formula-composite,lemma:weil-pbc-derived-complete-reconstruction}

Apply the completed projection-formula composite with $B=\Lambda$; its
finite reductions are $P_n$, which are perfect by the coefficient datum.
This gives the displayed map and its explicit inverse.  If $E$ is complete,
reconstruction identifies $\widehat E$ with $E$.  Naturality of the composite
identifies the resulting inverse $Rf_*E\to\widehat{Rf_*E}$ with the
completion unit, so that unit is an isomorphism.  The naturality assertion is
inherited from the completed composite.
\end{proof}

\begin{lemma}[Limits preserve compatible levelwise isomorphisms]
\label{lemma:weil-pbc-r-lim-levelwise-isomorphism}
\uses{}
If $u_n:M_n\to N_n$ is a compatible map of towers and every $u_n$ is an
isomorphism, then $R\varprojlim u_n$ is an isomorphism.
\end{lemma}
\begin{proof}
\uses{} 
The maps $\prod_nu_n$ are isomorphisms and commute with
$1-\mathrm{shift}$.  Taking fibres in the homotopy-limit triangles preserves
the isomorphism.
\end{proof}

\begin{lemma}[Complete cone detected by reductions]
\label{lemma:weil-pbc-complete-cone-zero}
\uses{}

If $C$ is derived complete and $C\otimes^L\Lambda/I^n=0$ for every $n$, then
$C=0$.
\end{lemma}
\begin{proof}
\uses{lemma:weil-pbc-derived-complete-reconstruction}

Reconstruction identifies $C$ with the derived limit of its zero reductions;
the product/equalizer model of the zero tower is zero.
\end{proof}

\begin{lemma}[Perfectness of $\mathbf Z/\ell^n$]
\label{lemma:weil-pbc-zl-reduction-perfect}
\uses{}
For $\Lambda=\mathbf Z_\ell$ and $I=(\ell)$, every
$\Lambda/I^n$ is perfect over $\Lambda$.
\end{lemma}
\begin{proof}
\uses{} 
The two-term complex
\[
 [\mathbf Z_\ell\xrightarrow{\ell^n}\mathbf Z_\ell]
\]
in degrees $-1,0$ consists of finite free modules.  Multiplication by
$\ell^n$ is injective and its cokernel is $\mathbf Z/\ell^n$, so the complex
is a finite projective resolution.
\end{proof}

% Closed intrinsic-completion leaves relevant to EFI COEFF-07.
% They do not close the comparison with the package-dependent A2 category.

\begin{definition}[Canonical finite reduction tower over $\mathbf Z_\ell$]
\label{def:weil-efi-coeff-zl-reduction-tower}
\uses{def:weil-a2-derived-category,
  def:weil-pbc-complete-coefficient-datum}
\constructedby{lemma:weil-pbc-zl-reduction-perfect}
Specialize the adic coefficient construction to
$\Lambda=\mathbf Z_\ell$, $I=(\ell)$, and
$P_a=\mathbf Z_\ell/\ell^a$.  For $K\in\mathcal D(X,\Lambda)$ set
\[
 K_a:=K\otimes^L_\Lambda P_a,
 \qquad q_{b,a}:K_b\longrightarrow K_a\quad(b\geq a)
\]
using the quotient map $P_b\to P_a$.
\end{definition}

\begin{lemma}[Identity equation for the canonical reduction tower]
\label{lemma:weil-efi-coeff-zl-reduction-identity}
\uses{def:weil-efi-coeff-zl-reduction-tower}
For every $a$, $q_{a,a}=\operatorname{id}_{K_a}$.
\end{lemma}
\begin{proof}
\uses{def:weil-efi-coeff-zl-reduction-tower}
The coefficient map $P_a\to P_a$ is the identity, and derived tensor sends
it to the identity of $K\otimes^L_\Lambda P_a$.
\end{proof}

\begin{lemma}[Composition equation for the canonical reduction tower]
\label{lemma:weil-efi-coeff-zl-reduction-composition}
\uses{def:weil-efi-coeff-zl-reduction-tower}
For $c\geq b\geq a$, $q_{c,a}=q_{b,a}q_{c,b}$.
\end{lemma}
\begin{proof}
\uses{def:weil-efi-coeff-zl-reduction-tower}
The two coefficient maps are induced by the equal quotient homomorphisms
$P_c\to P_a$.  Functoriality of derived tensor preserves that equality.
\end{proof}

\begin{definition}[Normalization comparison for the canonical reduction tower]
\label{def:weil-efi-coeff-zl-reduction-normalization}
\uses{def:weil-efi-coeff-zl-reduction-tower}
For $b\geq a$, let
\[
 \nu_{b,a,K}:
 K_b\otimes^L_{P_b}P_a\longrightarrow K_a
\]
be the derived-tensor associator followed by the multiplication isomorphism
$P_b\otimes^L_{P_b}P_a\simeq P_a$.
\end{definition}

\begin{lemma}[The normalization comparison is invertible]
\label{lemma:weil-efi-coeff-zl-reduction-normalization-iso}
\uses{def:weil-efi-coeff-zl-reduction-normalization}
The map $\nu_{b,a,K}$ is an isomorphism.
\end{lemma}
\begin{proof}
\uses{lemma:weil-pbc-zl-reduction-perfect}
The derived associator is an isomorphism.  The second factor is the standard
derived unitor
$P_b\otimes^L_{P_b}P_a\simeq P_a$; composing the two isomorphisms gives the
claim.
\end{proof}

\begin{lemma}[Cocycle equation for reduction normalization]
\label{lemma:weil-efi-coeff-zl-reduction-normalization-cocycle}
\uses{def:weil-efi-coeff-zl-reduction-normalization}
For $c\geq b\geq a$, the direct normalization from level $c$ to level $a$
equals the composite through level $b$.
\end{lemma}
\begin{proof}
\uses{lemma:weil-efi-coeff-zl-reduction-composition}
Expand the two maps into derived-tensor associators and quotient maps.  The
associator pentagon makes the reassociations equal, and the cited quotient
composition equation identifies the remaining coefficient map.
\end{proof}

\begin{definition}[Intrinsic constructible integral category candidate]
\label{def:weil-efi-coeff-intrinsic-constructible-zl-category}
\uses{def:weil-a2-derived-category,
  def:weil-pbc-intrinsic-koszul-completion,
  def:weil-pbc-bounded-constructible-adic-tower,
  def:weil-efi-coeff-zl-reduction-tower}
Define $\mathscr A^{\mathrm{intr}}_{\mathbf Z_\ell}(X)$ to be the full
subcategory of $\mathcal D(X,\mathbf Z_\ell)$ on those $K$ for which the
intrinsic Koszul completion unit is an isomorphism and the canonical finite
reduction tower is uniformly bounded constructible.  This definition has no
K-flat or K-injective package parameter.
\end{definition}

\begin{lemma}[Adic completion lands in the intrinsic candidate]
\label{lemma:weil-efi-coeff-completion-in-intrinsic-category}
\uses{def:weil-efi-coeff-intrinsic-constructible-zl-category}
If the canonical reduction tower of $K$ is uniformly bounded constructible,
then its explicit adic completion is an object of
$\mathscr A^{\mathrm{intr}}_{\mathbf Z_\ell}(X)$.
\end{lemma}
\begin{proof}
\uses{lemma:weil-pbc-bounded-constructible-reconstruction,
  lemma:weil-pbc-adic-completion-reduction}
The reconstruction lemma proves intrinsic completeness and preserves the
common constructible bounds.  The finite-reduction lemma identifies the
reductions of the completion with the original tower, so the same
constructibility witness applies.
\end{proof}

\subsubsection{S10. Complete pro-etale proper base change}

\begin{definition}[Four finite-level rectangle corners]
\label{def:weil-pbc-level-map-corners}
\uses{def:weil-pbc-cartesian-square,def:weil-pbc-complete-coefficient-datum}
For the proper Cartesian square $\square$, put $E=L(g')^*K$ and, for every
$n\geq1$, define
\[
\begin{aligned}
 A_n&:=Lg^*((Rf_*K)\otimes^LP_n),&
 B_n&:=Lg^*Rf_*K_n,\\
 C_n&:=(Rf'_*E)\otimes^LP_n,&
 D_n&:=Rf'_*E_n,
\end{aligned}
\]
where $E_n=L(g')^*K_n=E\otimes^LP_n$.  The chosen Cartesian
identification supplies
$L(g')^*Lf^*=L(f')^*Lg^*$ throughout.
\end{definition}

\begin{definition}[The four finite-level arrows]
\label{def:weil-pbc-level-map-arrows}
\uses{def:weil-pbc-level-map-corners}
Let $\xi_K:Lg^*Rf_*K\to Rf'_*E$ be the unit--counit exchange map
\[
 Lg^*Rf_*K\longrightarrow
 Rf'_*L(f')^*Lg^*Rf_*K
 \simeq Rf'_*L(g')^*Lf^*Rf_*K
 \xrightarrow{Rf'_*L(g')^*\epsilon_K}Rf'_*E.
\]
Define the four edges
\[
\begin{aligned}
 h_n&:=Lg^*(\pi_{f,K,n}):A_n\longrightarrow B_n,\\
 v_n&:=\xi_K\otimes^L1_{P_n}:A_n\longrightarrow C_n,\\
 h'_n&:=\pi_{f',E,n}:C_n\longrightarrow D_n,\\
 v'_n&:=\xi_{K_n}:B_n\longrightarrow D_n.
\end{aligned}
\]
For $v_n$ use
$Lg^*((Rf_*K)\otimes P_n)\simeq(Lg^*Rf_*K)\otimes P_n$.
The transported S8 map is $v'_n$, because its construction was identified
there with this unit--counit exchange map.
\end{definition}


\begin{lemma}[Well-definedness certificate for def:weil-pbc-level-map-arrows]
\label{lemma:weil-sf-well-defined-def-weil-pbc-level-map-arrows}
\uses{def:weil-pbc-level-map-arrows}
For $n\geq1$, the tuple $(h_n,v_n,h'_n,v'_n)$ belongs to
\[
 \operatorname{Hom}(A_n,B_n)\times
 \operatorname{Hom}(A_n,C_n)\times
 \operatorname{Hom}(C_n,D_n)\times
 \operatorname{Hom}(B_n,D_n),
\]
and $h_n,h'_n$ are isomorphisms.  Under finite-level pro-etale transport,
the fourth component is the unit--counit exchange map $\xi_{K_n}$.
\end{lemma}
\begin{proof}
\uses{lemma:weil-pbc-perfect-projection-formula,prop:weil-pbc-proet-finite-level-transport}
Perfectness of $P_n$ makes
$\pi_{f,K,n}:(Rf_*K)\otimes^LP_n\to Rf_*K_n$ and
$\pi_{f',E,n}:(Rf'_*E)\otimes^LP_n\to Rf'_*E_n$ isomorphisms.
Derived pullback preserves the first isomorphism, giving the stated type and
invertibility of $h_n$; the second gives the same facts for $h'_n$.
Symmetric monoidality of derived pullback gives
\[
 Lg^*((Rf_*K)\otimes^LP_n)\simeq
 (Lg^*Rf_*K)\otimes^LP_n,
\]
so tensoring $\xi_K$ with $1_{P_n}$ has type $A_n\to C_n$.
Proposition~\ref{prop:weil-pbc-proet-finite-level-transport} transports the
finite etale-origin exchange to a map $B_n\to D_n$ and identifies that map
with the unit--counit formula $\xi_{K_n}$.  These are exactly the four
components asserted above; no rectangle or tower equation is used here.
\end{proof}
\begin{lemma}[The finite-level exchange rectangle]
\label{lemma:weil-pbc-level-map-rectangle}
\uses{def:weil-pbc-level-map-arrows}
For every $n$, the diagram
\[
\begin{CD}
 A_n @>{h_n}>> B_n\\
 @V{v_n}VV @VV{v'_n}V\\
 C_n @>{h'_n}>> D_n
\end{CD}
\]
commutes; equivalently, its two composites satisfy
\[
 v'_nh_n=h'_nv_n:A_n\longrightarrow D_n.
\]
\end{lemma}
\begin{proof}
\uses{lemma:weil-pbc-projection-formula-base-change-naturality}

This is the base-change naturality rectangle for the projection-formula map
with coefficient $P_n$.  To expose the equality, adjoint both composites to
$X'$.  Their common adjoint is
\[
 L(g')^*Lf^*Rf_*K\otimes^LP_n
 \xrightarrow{L(g')^*(\epsilon_K)\otimes1}
 L(g')^*K\otimes^LP_n=E_n.
\]
Thus every unit, counit, tensor associator, and Cartesian comparison occurs
in the same order in the two adjoints.  Equality after adjunction proves the
displayed component equality.
\end{proof}

\begin{lemma}[Transition compatibility of all four edges]
\label{lemma:weil-pbc-level-map-tower-compatibility}
\uses{def:weil-pbc-level-map-arrows}
The four corner transitions from level $n+1$ to level $n$ are
\[
\begin{aligned}
 a_{n+1,n}&=Lg^*(1\otimes q_n):A_{n+1}\to A_n,\\
 b_{n+1,n}&=Lg^*Rf_*(K_{n+1}\to K_n):B_{n+1}\to B_n,\\
 c_{n+1,n}&=1\otimes q_n:C_{n+1}\to C_n,\\
 d_{n+1,n}&=Rf'_*L(g')^*(K_{n+1}\to K_n):D_{n+1}\to D_n.
\end{aligned}
\]
Every edge is a map of towers:
\[
\begin{aligned}
 b_{n+1,n}h_{n+1}&=h_na_{n+1,n},&
 c_{n+1,n}v_{n+1}&=v_na_{n+1,n},\\
 d_{n+1,n}h'_{n+1}&=h'_nc_{n+1,n},&
 d_{n+1,n}v'_{n+1}&=v'_nb_{n+1,n}.
\end{aligned}
\]
\end{lemma}
\begin{proof}
\uses{lemma:weil-pbc-projection-formula-tower-compatibility}

The first and third equations are the coefficient-tower equations for the
two projection formulas.  The second follows by tensor naturality of
$\xi_K$ with respect to $q_n$.  The last is naturality of the unit--counit
exchange map with respect to $K_{n+1}\to K_n$.  Expanded after adjunction,
both sides of the last equation are the counit of $K$ followed by the same
coefficient quotient.  Hence all four equations hold componentwise.
\end{proof}

\begin{lemma}[Operator naturality of the finite rectangle]
\label{lemma:weil-pbc-level-map-operator-naturality}
\uses{def:weil-pbc-level-map-arrows}
Let $\phi:K\to K$ be a tower-compatible endomorphism.  Give
$A_n,B_n,C_n,D_n$ the endomorphisms induced by
$Rf_*\phi$, $Rf_*\phi_n$, $Rf'_*L(g')^*\phi$, and
$Rf'_*L(g')^*\phi_n$, respectively.  Then $h_n,v_n,h'_n,v'_n$ intertwine
these endomorphisms, and the equality $v'_nh_n=h'_nv_n$ is an equality of
operator-equivariant maps.
\end{lemma}
\begin{proof}
\uses{lemma:weil-pbc-projection-formula-base-change-naturality}

The horizontal statements are operator naturality of the projection
formula.  The vertical statements follow by naturality of the unit and
counit defining exchange.  Composing either pair of equivariant edges gives
an equivariant map; the preceding rectangle identifies those two maps.
\end{proof}

\begin{definition}[Frobenius-equivariant square and Weil tower]
\label{def:weil-pbc-frobenius-tower-datum}
\uses{def:weil-pbc-level-map-corners,def:weil-pbc-level-map-arrows}
A Frobenius-equivariant datum for the proper Cartesian square $\square$ and
the tower $K_\bullet$ includes displayed witnesses
\[
 \epsilon_X^{-1}L_n\xrightarrow{\sim}K_n,
 \qquad L_n\in D_c^+(X_{\mathrm{et}},P_n),
\]
compatible with the tower transitions.  Pullback supplies the corresponding
etale-origin witnesses for $E_n=L(g')^*K_n$.  The datum then consists of
automorphisms
\[
 \operatorname{Fr}_X,\quad \operatorname{Fr}_Y,\quad
 \operatorname{Fr}_{X'},\quad \operatorname{Fr}_{Y'}
\]
such that
\[
\begin{aligned}
 f\operatorname{Fr}_X&=\operatorname{Fr}_Yf,&
 g\operatorname{Fr}_{Y'}&=\operatorname{Fr}_Yg,\\
 g'\operatorname{Fr}_{X'}&=\operatorname{Fr}_Xg',&
 f'\operatorname{Fr}_{X'}&=\operatorname{Fr}_{Y'}f'.
\end{aligned}
\]
The four resulting Frobenius faces are equipped with their Cartesian
identifications, and these identifications make the resulting cube commute
with the chosen Cartesian identification for $\square$.  Put
$\Phi_Z=L\operatorname{Fr}_Z^*$.  Since every $\operatorname{Fr}_Z$ is an
automorphism, every $\Phi_Z$ is an exact equivalence and preserves products
and fibres.

The pullback faces give canonical isomorphisms
\[
 \kappa_g:\Phi_{Y'}Lg^*\xrightarrow{\sim}Lg^*\Phi_Y,
 \qquad
 \kappa_{g'}:\Phi_{X'}L(g')^*\xrightarrow{\sim}L(g')^*\Phi_X.
\]
The proper Frobenius faces give finite-level exchange isomorphisms
\[
 \beta_{f,n}:\Phi_YRf_*K_n\xrightarrow{\sim}Rf_*\Phi_XK_n,
 \qquad
 \beta_{f',M}:\Phi_{Y'}Rf'_*M\xrightarrow{\sim}Rf'_*\Phi_{X'}M.
\]
Finally the datum contains isomorphisms
\[
 \vartheta_n:\Phi_XK_n\xrightarrow{\sim}K_n
\]
such that, for $r_{n+1,n}:K_{n+1}\to K_n$,
\[
 r_{n+1,n}\vartheta_{n+1}
 =\vartheta_n\Phi_X(r_{n+1,n}).
\]
The constant coefficient quotients carry their identity Frobenius
structures, and all displayed comparisons are required to respect their
normalizations and cocycles.
\end{definition}


\begin{lemma}[Well-definedness certificate for def:weil-pbc-frobenius-tower-datum]
\label{lemma:weil-sf-well-defined-def-weil-pbc-frobenius-tower-datum}
\uses{def:weil-pbc-frobenius-tower-datum}
The displayed Frobenius record has fields
\[
\begin{gathered}
 \kappa_g:\Phi_{Y'}Lg^*\xrightarrow{\sim}Lg^*\Phi_Y,
 \qquad
 \kappa_{g'}:\Phi_{X'}L(g')^*\xrightarrow{\sim}L(g')^*\Phi_X,\\
 \beta_{f,n}:\Phi_YRf_*K_n\xrightarrow{\sim}Rf_*\Phi_XK_n,
 \qquad
 \beta_{f',E_n}:\Phi_{Y'}Rf'_*E_n
       \xrightarrow{\sim}Rf'_*\Phi_{X'}E_n,\\
 \vartheta_n:\Phi_XK_n\xrightarrow{\sim}K_n.
\end{gathered}
\]
The two sides of the tower relation belong to
$\operatorname{Hom}(\Phi_XK_{n+1},K_n)$, so the relation is a typed field of
the record.
\end{lemma}
\begin{proof}
\uses{lemma:weil-pbc-derived-exchange-coherence,prop:weil-pbc-proet-finite-level-transport}
Pullback by the inverse of $\operatorname{Fr}_Z$ is a quasi-inverse to
$\Phi_Z=L\operatorname{Fr}_Z^*$; hence $\Phi_Z$ is an exact equivalence.
The two pullback faces in the Frobenius cube give $\kappa_g$ and
$\kappa_{g'}$ with the displayed endpoints.  Applied to the two proper
faces, finite-level pro-etale transport gives the displayed exchange
isomorphisms because $K_n$ and $E_n$ carry their etale-origin witnesses.
The identity and pasting equations required of these exchanges are those of
Lemma~\ref{lemma:weil-pbc-derived-exchange-coherence}.

Finally, $\Phi_X(r_{n+1,n})$ has source $\Phi_XK_{n+1}$ and target
$\Phi_XK_n$.  Thus
$r_{n+1,n}\vartheta_{n+1}$ and
$\vartheta_n\Phi_X(r_{n+1,n})$ have the common source and target stated in
the lemma.  This verifies the record typing without asserting a new
Frobenius-equivariance theorem for the finite rectangle.
\end{proof}
\begin{definition}[Weil structures on the four finite corners]
\label{def:weil-pbc-frobenius-corner-structures}
\uses{def:weil-pbc-frobenius-tower-datum}
For $E_n=L(g')^*K_n$, define
\[
 \vartheta_{E,n}:=
 L(g')^*(\vartheta_n)\,\kappa_{g',K_n}:
 \Phi_{X'}E_n\xrightarrow{\sim}E_n.
\]
The Weil structures on $B_n$ and $D_n$ are the typed composites
\[
\begin{aligned}
 \Theta_{B,n}:\Phi_{Y'}B_n
 &\xrightarrow{\kappa_g}Lg^*\Phi_YRf_*K_n
 \xrightarrow{Lg^*\beta_{f,n}}Lg^*Rf_*\Phi_XK_n
 \xrightarrow{Lg^*Rf_*\vartheta_n}B_n,\\
 \Theta_{D,n}:\Phi_{Y'}D_n
 &\xrightarrow{\beta_{f',E_n}}Rf'_*\Phi_{X'}E_n
 \xrightarrow{Rf'_*\vartheta_{E,n}}D_n.
\end{aligned}
\]
Transport these structures across the two projection-formula isomorphisms:
\[
 \Theta_{A,n}:=h_n^{-1}\Theta_{B,n}\Phi_{Y'}(h_n),
 \qquad
 \Theta_{C,n}:=(h'_n)^{-1}\Theta_{D,n}\Phi_{Y'}(h'_n).
\]
Thus every $\Theta_{Z,n}$ has the genuine Weil type
$\Phi_{Y'}Z_n\xrightarrow{\sim}Z_n$ for
$Z=A,B,C,D$.
\end{definition}


\begin{lemma}[Well-definedness certificate for def:weil-pbc-frobenius-corner-structures]
\label{lemma:weil-sf-well-defined-def-weil-pbc-frobenius-corner-structures}
\uses{def:weil-pbc-frobenius-corner-structures}
The tuple $(\Theta_{A,n},\Theta_{B,n},\Theta_{C,n},\Theta_{D,n})$ belongs to
\[
 \operatorname{Iso}(\Phi_{Y'}A_n,A_n)\times
 \operatorname{Iso}(\Phi_{Y'}B_n,B_n)\times
 \operatorname{Iso}(\Phi_{Y'}C_n,C_n)\times
 \operatorname{Iso}(\Phi_{Y'}D_n,D_n).
\]
This assertion records the four corner types and does not assert edge
equivariance.
\end{lemma}
\begin{proof}
\uses{lemma:weil-sf-well-defined-def-weil-pbc-level-map-arrows,lemma:weil-sf-well-defined-def-weil-pbc-frobenius-tower-datum,lemma:weil-pbc-perfect-projection-formula}
The certified tower datum makes $\kappa_g,\kappa_{g'},\beta_{f,n},
\beta_{f',E_n}$ and $\vartheta_n$ isomorphisms with the displayed
endpoints.  Their ordered composites therefore give
$\vartheta_{E,n}$, $\Theta_{B,n}$ and $\Theta_{D,n}$ with the asserted
types.  The level-arrow certificate and the perfect projection formula make
$h_n$ and $h'_n$ invertible.  Conjugating $\Theta_{B,n}$ by $h_n$ and
$\Theta_{D,n}$ by $h'_n$ gives the asserted types for $\Theta_{A,n}$ and
$\Theta_{C,n}$.  Edge equivariance remains the next lemma.
\end{proof}
\begin{lemma}[The finite rectangle is a Weil rectangle]
\label{lemma:weil-pbc-frobenius-level-rectangle}
\uses{def:weil-pbc-frobenius-corner-structures}
All four finite-level edges are Weil-equivariant:
\[
\begin{aligned}
 h_n\Theta_{A,n}&=\Theta_{B,n}\Phi_{Y'}(h_n),&
 v_n\Theta_{A,n}&=\Theta_{C,n}\Phi_{Y'}(v_n),\\
 h'_n\Theta_{C,n}&=\Theta_{D,n}\Phi_{Y'}(h'_n),&
 v'_n\Theta_{B,n}&=\Theta_{D,n}\Phi_{Y'}(v'_n).
\end{aligned}
\]
In particular the two explicitly typed composites
\[
\begin{aligned}
 \mathcal P_n&:=
 \Phi_{Y'}A_n\xrightarrow{\Theta_{A,n}}A_n
 \xrightarrow{v'_nh_n}D_n,\\
 \mathcal Q_n&:=
 \Phi_{Y'}A_n\xrightarrow{\Phi_{Y'}(v'_nh_n)}\Phi_{Y'}D_n
 \xrightarrow{\Theta_{D,n}}D_n
\end{aligned}
\]
are equal.  The same common map is obtained after replacing $v'_nh_n$ by
$h'_nv_n$.
\end{lemma}
\begin{proof}
\uses{lemma:weil-pbc-level-map-rectangle,lemma:weil-pbc-derived-exchange-coherence}

The equations for $h_n$ and $h'_n$ are the two transport definitions.  For
$v'_n$, expand the upper-right and lower-left routes in the Frobenius cube.
Pasting coherence for exchange and naturality in $\vartheta_n$ identify both
with
\[
\begin{aligned}
 \Phi_{Y'}Lg^*Rf_*K_n
 &\xrightarrow{\kappa_g}Lg^*\Phi_YRf_*K_n
 \xrightarrow{Lg^*\beta_{f,n}}Lg^*Rf_*\Phi_XK_n\\
 &\xrightarrow{\xi_{\Phi_XK_n}}Rf'_*L(g')^*\Phi_XK_n
 \xrightarrow{Rf'_*L(g')^*\vartheta_n}D_n.
\end{aligned}
\]
This proves the fourth equation with every pullback and exchange map typed.
For the remaining vertical edge, use the ordinary rectangle and the three
equations already proved:
\[
\begin{aligned}
 h'_nv_n\Theta_{A,n}
 &=v'_nh_n\Theta_{A,n}
 =v'_n\Theta_{B,n}\Phi_{Y'}(h_n)\\
 &=\Theta_{D,n}\Phi_{Y'}(v'_nh_n)
 =h'_n\Theta_{C,n}\Phi_{Y'}(v_n).
\end{aligned}
\]
Cancel the isomorphism $h'_n$.  Composing either pair of edge equations and
using $v'_nh_n=h'_nv_n$ gives
$\mathcal P_n=\mathcal Q_n$ and the final assertion.
\end{proof}

\begin{lemma}[Frobenius compatibility with tower transitions]
\label{lemma:weil-pbc-frobenius-tower-compatibility}
\uses{def:weil-pbc-frobenius-corner-structures}
For the four transitions $a_{n+1,n},b_{n+1,n},c_{n+1,n},d_{n+1,n}$,
\[
\begin{aligned}
 a_{n+1,n}\Theta_{A,n+1}
 &=\Theta_{A,n}\Phi_{Y'}(a_{n+1,n}),\\
 b_{n+1,n}\Theta_{B,n+1}
 &=\Theta_{B,n}\Phi_{Y'}(b_{n+1,n}),\\
 c_{n+1,n}\Theta_{C,n+1}
 &=\Theta_{C,n}\Phi_{Y'}(c_{n+1,n}),\\
 d_{n+1,n}\Theta_{D,n+1}
 &=\Theta_{D,n}\Phi_{Y'}(d_{n+1,n}).
\end{aligned}
\]
Hence the four corner Weil structures and all four edges are maps of towers.
\end{lemma}
\begin{proof}
\uses{lemma:weil-pbc-frobenius-level-rectangle,lemma:weil-pbc-level-map-tower-compatibility}

The equation for $B$ follows by naturality of $\kappa_g$ and $\beta_f$ and
the defining equation
$r_{n+1,n}\vartheta_{n+1}=\vartheta_n\Phi_X(r_{n+1,n})$.
The same argument for $\kappa_{g'}$, $\beta_{f'}$, and
$\vartheta_{E,n}$ proves the equation for $D$.  For $A$, combine
$b_{n+1,n}h_{n+1}=h_na_{n+1,n}$ with the transport formula for
$\Theta_A$ and the equation for $B$; cancelling $h_n$ gives the first
displayed equation.  The equation for $C$ follows identically from
$d_{n+1,n}h'_{n+1}=h'_nc_{n+1,n}$ and the equation for $D$.  The finite Weil
rectangle and the four transition equations now say precisely that every
edge is a morphism of Weil towers.
\end{proof}

\begin{lemma}[The rectangle passes to the derived inverse limit]
\label{lemma:weil-pbc-level-map-derived-limit}
\uses{def:weil-pbc-adic-derived-inverse-limit}
Write $A_\infty=R\varprojlim A_n$ and likewise for the other three towers,
and write $H,V,H',V'$ for the limits of $h,v,h',v'$.  Then
\[
 V'H=H'V:A_\infty\longrightarrow D_\infty.
\]
For every $n$, both maps have the same $n$th component
\[
 \operatorname{pr}_{D,n}V'H
 =v'_nh_n\operatorname{pr}_{A,n}
 =h'_nv_n\operatorname{pr}_{A,n}
 =\operatorname{pr}_{D,n}H'V.
\]
The equality is natural for compatible tower endomorphisms.
\end{lemma}
\begin{proof}
\uses{lemma:weil-pbc-level-map-rectangle,lemma:weil-pbc-level-map-tower-compatibility,lemma:weil-pbc-level-map-operator-naturality}

Transition compatibility makes each family a map between the product-cone
diagrams.  On both product terms, the $n$th coordinate of $V'H$ is
$v'_nh_n$, while that of $H'V$ is $h'_nv_n$; these are equal by the
finite-level rectangle.  Hence the product maps, their $1-\mathrm{shift}$
squares, and the induced maps on the fibres are equal.  This proves both the
limit equality and the three displayed component equations.  Operator
naturality holds on every product coordinate, so it also holds on the fibre.
\end{proof}

\begin{lemma}[Frobenius structures pass through the product cone]
\label{lemma:weil-pbc-frobenius-product-cone-limit}
\uses{def:weil-pbc-adic-derived-inverse-limit}
For $Z=A,B,C,D$, the product-cone limit carries the Weil structure
\[
 \Theta_{Z,\infty}:
 \Phi_{Y'}Z_\infty
 \xrightarrow{\sim}R\varprojlim_n\Phi_{Y'}Z_n
 \xrightarrow{R\varprojlim_n\Theta_{Z,n}}Z_\infty.
\]
The first arrow is the canonical comparison, which is an isomorphism because
$\Phi_{Y'}$ is an equivalence and therefore preserves the two products and
the fibre in the product-cone model.  Each of $H,V,H',V'$ is Weil-equivariant.
For the common limit map $W:=V'H=H'V:A_\infty\to D_\infty$, the two typed
composites
\[
\begin{aligned}
 \Phi_{Y'}A_\infty&\xrightarrow{\Theta_{A,\infty}}A_\infty
 \xrightarrow{W}D_\infty,\\
 \Phi_{Y'}A_\infty&\xrightarrow{\Phi_{Y'}(W)}\Phi_{Y'}D_\infty
 \xrightarrow{\Theta_{D,\infty}}D_\infty
\end{aligned}
\]
are equal.
\end{lemma}
\begin{proof}
\uses{lemma:weil-pbc-frobenius-level-rectangle,lemma:weil-pbc-frobenius-tower-compatibility,lemma:weil-pbc-level-map-derived-limit}

Apply $\Phi_{Y'}$ to the product-cone triangle for each tower.  Preservation
of products identifies the two product terms with the products of
$\Phi_{Y'}Z_n$, and transition compatibility identifies the two
$1-\mathrm{shift}$ arrows.  Thus the displayed comparison is an isomorphism
of fibres.  On its $n$th product coordinate, equivariance of $H,V,H',V'$ is
exactly one of the four finite edge equations.  Hence the product maps and
their maps on fibres are equivariant.  Composing the equations for $H,V'$ or
for $V,H'$ gives
\[
 W\Theta_{A,\infty}
 =\Theta_{D,\infty}\Phi_{Y'}(W).
\]
The equality $V'H=H'V$ is already an equality of maps of product-cone
diagrams, so both factorizations give the same typed Weil morphism.
\end{proof}

\begin{definition}[Completed base-change map]
\label{def:weil-pbc-completed-base-change-map}
\uses{def:weil-pbc-completed-pullback}
Let $K$ be derived complete and let every $K_n$ have a displayed witness
$\epsilon_X^{-1}L_n\simeq K_n$ with
$L_n\in D_c^+(X_{\mathrm{et}},\Lambda/I^n)$.  Define
\[
 \operatorname{Ex}^{\mathrm{comp}}_{\square,K}:
 Lg^*_{\mathrm{comp}}Rf_*K\longrightarrow
 Rf'_*L(g')^*_{\mathrm{comp}}K
\]
to be the ordered composite
\[
\begin{aligned}
 Lg^*_{\mathrm{comp}}Rf_*K
 &\xrightarrow{\sim}R\varprojlim_n Lg^*Rf_*K_n\\
 &\xrightarrow{R\varprojlim_n\operatorname{Ex}_{\square,K_n}}
 R\varprojlim_nRf'_*L(g')^*K_n\\
 &\xrightarrow{\sim}Rf'_*R\varprojlim_nL(g')^*K_n
 =Rf'_*L(g')^*_{\mathrm{comp}}K.
\end{aligned}
\]
In the notation above, the first two arrows are $H$ followed by $V'$.
The first isomorphism uses reconstruction and the perfect projection formula;
the second displayed isomorphism is the inverse of direct-image/limit
commutation.  This definition fills the map-construction omission in Tag
\texttt{09C9}.
\end{definition}


\begin{lemma}[Well-definedness certificate for def:weil-pbc-completed-base-change-map]
\label{lemma:weil-sf-well-defined-def-weil-pbc-completed-base-change-map}
\uses{def:weil-pbc-completed-base-change-map}
With $S=Lg^*_{\mathrm{comp}}Rf_*K$ and
$T=Rf'_*L(g')^*_{\mathrm{comp}}K$, the ordered expression
\[
 S\xrightarrow{s}A_\infty\xrightarrow{H}B_\infty
 \xrightarrow{V'}D_\infty\xrightarrow{t}T
\]
is a morphism $S\to T$ and is the map
$\operatorname{Ex}^{\mathrm{comp}}_{\square,K}$ of Definition
\ref{def:weil-pbc-completed-base-change-map}.  The endpoint comparisons
$s,t$ and the limit projection-formula map $H$ are isomorphisms.
\end{lemma}
\begin{proof}
\uses{lemma:weil-pbc-perfect-projection-formula,lemma:weil-pbc-level-map-derived-limit,lemma:weil-pbc-direct-image-r-lim}
Symmetric monoidality of $Lg^*$ identifies the defining completed-pullback
tower with the $A_n$ tower.  The perfect projection formula then gives the
map $H:A_\infty\to B_\infty$, while finite exchange gives the compatible
tower map $V':B_\infty\to D_\infty$ recorded by
Lemma~\ref{lemma:weil-pbc-level-map-derived-limit}.  Direct image commutation
with the product-cone limit gives
$t:D_\infty\xrightarrow{\sim}T$.  The target of each factor is the source
of the next one.

At level $n$, $h_n$ is invertible by the perfect projection formula.  From
$b_{n+1,n}h_{n+1}=h_na_{n+1,n}$, multiplication by the two inverses gives
$a_{n+1,n}h_{n+1}^{-1}=h_n^{-1}b_{n+1,n}$.  Hence the inverse family is a
map of towers, and its derived limit is inverse to $H$.  Thus $s,H,t$ are
isomorphisms.  Expanding the ordered definition gives the displayed
composite.  No invertibility or Frobenius assertion for the whole composite
is made in this certificate.
\end{proof}
\begin{definition}[Weil structures on the completed source and target]
\label{def:weil-pbc-completed-frobenius-structures}
\uses{def:weil-pbc-completed-base-change-map}
Write
\[
 S=Lg^*_{\mathrm{comp}}Rf_*K,
 \qquad T=Rf'_*L(g')^*_{\mathrm{comp}}K.
\]
Let $s:S\xrightarrow{\sim}A_\infty$ be the pullback/tensor comparison in the
first part of the completed construction, and let
$t:D_\infty\xrightarrow{\sim}T$ be the inverse of direct-image/limit
commutation.  Define genuine completed Weil structures by
\[
\begin{aligned}
 \Theta_S&:=s^{-1}\Theta_{A,\infty}\Phi_{Y'}(s):
 \Phi_{Y'}S\xrightarrow{\sim}S,\\
 \Theta_T&:=t\Theta_{D,\infty}\Phi_{Y'}(t^{-1}):
 \Phi_{Y'}T\xrightarrow{\sim}T.
\end{aligned}
\]
Under these identifications the completed base-change map is the ordered
composite
\[
 \operatorname{Ex}^{\mathrm{comp}}_{\square,K}=tWs,
 \qquad W=V'H=H'V.
\]
\end{definition}


\begin{lemma}[Well-definedness certificate for def:weil-pbc-completed-frobenius-structures]
\label{lemma:weil-sf-well-defined-def-weil-pbc-completed-frobenius-structures}
\uses{def:weil-pbc-completed-frobenius-structures}
The pair $(\Theta_S,\Theta_T)$ belongs to
\[
 \operatorname{Iso}(\Phi_{Y'}S,S)\times
 \operatorname{Iso}(\Phi_{Y'}T,T),
\]
and $tWs$ belongs to $\operatorname{Hom}(S,T)$.  This certificate records
the completed types but not the equivariance equation for $tWs$.
\end{lemma}
\begin{proof}
\uses{lemma:weil-sf-well-defined-def-weil-pbc-completed-base-change-map,lemma:weil-pbc-frobenius-product-cone-limit}
The product-cone lemma supplies isomorphisms
$\Theta_{A,\infty}:\Phi_{Y'}A_\infty\to A_\infty$ and
$\Theta_{D,\infty}:\Phi_{Y'}D_\infty\to D_\infty$, together with
$W:A_\infty\to D_\infty$.  The preceding certificate gives the
isomorphisms $s:S\to A_\infty$ and $t:D_\infty\to T$.  Conjugating the two
limit structures by $s$ and $t$ yields the displayed pair of isomorphism
types, and the endpoints of $s,W,t$ give $tWs:S\to T$.  The next lemma
proves the equivariance equation.
\end{proof}
\begin{lemma}[Completed proper base change is Frobenius-equivariant]
\label{lemma:weil-pbc-completed-frobenius-equivariance}
\uses{def:weil-pbc-completed-frobenius-structures}
For every Frobenius-equivariant square and Weil tower datum,
\[
 \operatorname{Ex}^{\mathrm{comp}}_{\square,K}\Theta_S
 =\Theta_T\Phi_{Y'}
   (\operatorname{Ex}^{\mathrm{comp}}_{\square,K}):
 \Phi_{Y'}S\longrightarrow T.
\]
If all four Frobenius automorphisms and all pullback/exchange comparisons are
identities, then $\vartheta_n$ is an ordinary compatible endomorphism and
this equation specializes exactly to the operator naturality of the finite
rectangle and its product-cone limit.  For nonidentity Frobenius, no
identification of $\Phi_XK_n\to K_n$ with an endomorphism $K_n\to K_n$ is
made.
\end{lemma}
\begin{proof}
\uses{lemma:weil-pbc-frobenius-product-cone-limit,lemma:weil-pbc-level-map-operator-naturality}

Using the displayed definitions and cancelling $s^{-1}s$ and $t^{-1}t$,
the left side is
\[
 tW\Theta_{A,\infty}\Phi_{Y'}(s),
\]
whereas the right side is
\[
 t\Theta_{D,\infty}\Phi_{Y'}(W)\Phi_{Y'}(s).
\]
They are equal by the product-cone Frobenius equation.  When every
$\Phi_Z$ is the identity, the four formulas defining $\Theta_{Z,n}$ become
the endomorphisms induced from $\vartheta_n$ and the equation just proved is
the fibrewise limit of the equations in operator naturality.  This proves
the specialization statement without using it to type the genuine
Frobenius case.
\end{proof}

\begin{definition}[The four edges of the completed reduction square]
\label{def:weil-pbc-completed-reduction-square}
\uses{def:weil-pbc-completed-base-change-map}
Put
\[
 S=Lg^*_{\mathrm{comp}}Rf_*K,
 \qquad T=Rf'_*L(g')^*_{\mathrm{comp}}K.
\]
Let
\[
 \sigma_n:S\otimes^LP_n\xrightarrow{\sim}A_n
\]
be finite reduction of completed pullback.  For $E=L(g')^*K$, let
\[
 \tau_n:T\otimes^LP_n\xrightarrow{\sim}C_n
\]
be the inverse of
$\widehat{Rf'_*E}\xrightarrow{\sim}Rf'_*\widehat E=T$, tensored with
$P_n$, followed by finite reduction of $\widehat{Rf'_*E}$.  Define
\[
 \alpha_n=h_n\sigma_n:S\otimes^LP_n\to B_n,
 \qquad
 \delta_n=h'_n\tau_n:T\otimes^LP_n\to D_n,
\]
and let
\[
 \overline\xi_n=
 \operatorname{Ex}^{\mathrm{comp}}_{\square,K}\otimes^L1_{P_n}:
 S\otimes^LP_n\longrightarrow T\otimes^LP_n.
\]
Thus every arrow in the completed reduction square is named, typed, and
expressed as an ordered composite of reduction, direct-image, projection-
formula, or base-change maps.
\end{definition}


\begin{lemma}[Well-definedness certificate for def:weil-pbc-completed-reduction-square]
\label{lemma:weil-sf-well-defined-def-weil-pbc-completed-reduction-square}
\uses{def:weil-pbc-completed-reduction-square}
For $n\geq1$, $\sigma_n:S\otimes^LP_n\to A_n$ and
$\tau_n:T\otimes^LP_n\to C_n$ are isomorphisms, while
\[
 (\alpha_n,\delta_n,\overline\xi_n)\in
 \operatorname{Hom}(S\otimes^LP_n,B_n)\times
 \operatorname{Hom}(T\otimes^LP_n,D_n)\times
 \operatorname{Hom}(S\otimes^LP_n,T\otimes^LP_n).
\]
The first two components of this triple are isomorphisms.  No commutativity
claim for the reduction square is part of this certificate.
\end{lemma}
\begin{proof}
\uses{lemma:weil-sf-well-defined-def-weil-pbc-level-map-arrows,lemma:weil-sf-well-defined-def-weil-pbc-completed-base-change-map,lemma:weil-pbc-completed-pullback-reduction,lemma:weil-pbc-direct-image-derived-completion,lemma:weil-pbc-completed-tensor-reduction}
Finite reduction of completed pullback gives $\sigma_n$ with the stated
source and target.  Direct-image completion gives
$\widehat{Rf'_*E}\xrightarrow{\sim}Rf'_*(\widehat E)=T$; tensoring its
inverse with $P_n$ and applying finite reduction of completed tensor gives
$\tau_n:T\otimes^LP_n\xrightarrow{\sim}C_n$.

The level-arrow certificate makes $h_n:A_n\to B_n$ and
$h'_n:C_n\to D_n$ isomorphisms.  Hence
$\alpha_n=h_n\sigma_n$ and $\delta_n=h'_n\tau_n$ have the asserted types
and are invertible.  Finally the completed-map certificate gives
$\operatorname{Ex}^{\mathrm{comp}}_{\square,K}:S\to T$, so derived tensor
with $1_{P_n}$ gives the asserted type of $\overline\xi_n$.  The next lemma
contains the square's commutativity and its compatibility fields.
\end{proof}
\begin{lemma}[Finite reduction of the completed map]
\label{lemma:weil-pbc-completed-map-reduction}
\uses{def:weil-pbc-completed-reduction-square,def:weil-pbc-completed-frobenius-structures}
For every $n$, the diagram
\[
\begin{CD}
 S\otimes^LP_n @>{\overline\xi_n}>> T\otimes^LP_n\\
 @V{\alpha_n}V{\sim}V @VV{\delta_n}V{\sim}\\
 B_n @>{v'_n}>> D_n
\end{CD}
\]
commutes.  Hence reduction of
$\operatorname{Ex}^{\mathrm{comp}}_{\square,K}$ modulo $I^n$, conjugated
by the displayed canonical reduction isomorphisms, is exactly the S8
transported finite-level exchange map $v'_n=\operatorname{Ex}_{\square,K_n}$.
The squares for $n+1\to n$ commute, and the assertion is natural for
compatible tower endomorphisms.  Under a Frobenius-equivariant datum, use
$\Phi_{Y'}P_n=P_n$ to give the two reductions their induced Weil structures.
Then
\[
\begin{aligned}
 \alpha_n(\Theta_S\otimes1)
 &=\Theta_{B,n}\Phi_{Y'}(\alpha_n),\\
 \delta_n(\Theta_T\otimes1)
 &=\Theta_{D,n}\Phi_{Y'}(\delta_n),
\end{aligned}
\]
and the displayed reduction square is a square of Weil objects.
\end{lemma}
\begin{proof}
\uses{lemma:weil-pbc-level-map-derived-limit,lemma:weil-pbc-frobenius-tower-compatibility,lemma:weil-pbc-completed-frobenius-equivariance}

Under $S=A_\infty$ and the direct-image completion identification
$T\simeq C_\infty$, the completed exchange map is both limit composites
\[
 V'H=H'V:A_\infty\longrightarrow D_\infty,
\]
followed by the inverse of direct-image/limit commutation.  More explicitly,
if $d:T\to D_\infty$ is direct-image/limit commutation and
$c:C_\infty\to T$ is direct-image completion, then its construction gives
$dc=H'$.  Functorial finite
reduction and the component equations of the limit rectangle give
\[
 \tau_n\overline\xi_n=v_n\sigma_n.
\]
Consequently the two composites around the completed reduction square are
\[
 \delta_n\overline\xi_n
 =h'_n\tau_n\overline\xi_n=h'_nv_n\sigma_n
\]
and
\[
 v'_n\alpha_n=v'_nh_n\sigma_n.
\]
They are equal term by term because $h'_nv_n=v'_nh_n$ in the finite-level
rectangle.  Naturality of finite reduction makes $\sigma_n$ and $\tau_n$
commute with $n+1\to n$, and functoriality of tensor reduction does the same
for $\overline\xi_n$.  Together with the four transition equations, this
shows that the completed squares commute with $K_{n+1}\to K_n$ on every
corner.  The limit construction and all four reduction edges are
operator-equivariant, so the conjugated equality is operator-natural as well.

For Frobenius, the definition of $\Theta_S$ transports
$\Theta_{A,\infty}$ across $s$, and the $n$th product-cone component followed
by $h_n$ is $\alpha_n$.  This gives the first displayed Weil equation.  The
definition of $\Theta_T$ transports $\Theta_{D,\infty}$ across $t$, and its
$n$th component is $\delta_n$, giving the second.  The completed Frobenius
equation for $\overline\xi_n$, together with finite-level equivariance of
$v'_n$, now says that every arrow of the reduction square is a Weil morphism.
\end{proof}

\begin{lemma}[Completed exchange coherence]
\label{lemma:weil-pbc-completed-exchange-coherence}
\uses{def:weil-pbc-completed-base-change-map}
The completed exchange map is natural in $K$, is the identity for identity
squares, and commutes with horizontal and vertical pasting.
\end{lemma}
\begin{proof}
\uses{lemma:weil-pbc-completed-map-reduction,lemma:weil-pbc-r-lim-levelwise-isomorphism}

At every finite level these equations are the S0/S8 equations.  They commute
with the tower transitions by naturality of coefficient reduction.  Apply
$R\varprojlim$ to each equality and conjugate by the same reconstruction,
projection-formula, and direct-image/limit isomorphisms used in the definition.
The resulting equations are precisely the completed equations.
\end{proof}

\begin{proposition}[Complete pro-etale proper base change]
\label{prop:weil-proet-complete-proper-base-change}
\uses{def:weil-pbc-complete-coefficient-datum,def:weil-pbc-completed-base-change-map}
Let $f:X\to Y$ be proper and let $\square$ be any base change.  Suppose
$K\in D(X_{\mathrm{proet}},\Lambda)$ is derived $I$-complete, each $K_n$ is
equipped with a displayed witness
$\epsilon_X^{-1}L_n\simeq K_n$ for some
$L_n\in D_c^+(X_{\mathrm{et}},\Lambda/I^n)$, and every $\Lambda/I^n$ is
perfect.  Then
\[
 Lg^*_{\mathrm{comp}}Rf_*K\xrightarrow{\sim}
 Rf'_*L(g')^*_{\mathrm{comp}}K.
\]
The isomorphism is natural and satisfies identity and pasting coherence.  If
the square and the reductions carry a Frobenius-equivariant square and Weil
tower datum, this is an isomorphism of the completed Weil objects
$(S,\Theta_S)$ and $(T,\Theta_T)$.
\end{proposition}
\begin{proof}
\uses{def:weil-pbc-completed-base-change-map,prop:weil-pbc-proet-finite-level-transport,lemma:weil-pbc-r-lim-levelwise-isomorphism,lemma:weil-pbc-completed-exchange-coherence,lemma:weil-pbc-completed-frobenius-equivariance}

Every finite-level middle arrow in Definition
\ref{def:weil-pbc-completed-base-change-map} is an isomorphism by S8.  The
derived limit preserves this compatible levelwise isomorphism.  The outer
arrows in the defining composite are reconstruction, perfect projection
formula, and direct-image/limit isomorphisms.  Hence the completed map is an
isomorphism.  Coherence is the preceding lemma.  This proves the mathematical
content located at Tag \texttt{09C9} without using that theorem as a leaf.
The completed Frobenius-equivariance lemma gives the last assertion with the
genuine pullback type $\Phi_{Y'}S\to T$.
\end{proof}

\begin{definition}[A2 constructible rebinding certificate]
\label{def:weil-pbc-a2-constructible-rebind}
\uses{def:weil-a2-dbc-zl,def:weil-a2-dbc-torsion,def:weil-a2-proper-complete-base-change-datum,def:weil-a2-finite-level-etale-origin-realization,def:weil-a2-torsion-constructibility-witness,def:weil-a2-constructible-stratification-witness,def:weil-a1-torsion-derived-category,def:weil-a1-rfshriek-selected,def:weil-pbc-level-map-corners,def:weil-pbc-completed-reduction-square}
Specialize A2 to $\Lambda=\mathbf Z_\ell$, $I=(\ell)$, and
$P_n=\Lambda/\ell^n$.  For a scheme $Z$ write
\[
 \mathsf C(Z):=D^b_c(Z_{\mathrm{proet}},\mathbf Z_\ell),
 \qquad
 \mathsf C_n(Z):=D^b_c(Z_{\mathrm{proet}},\mathbf Z/\ell^n)
\]
for the full subcategories defined, respectively, by
\texttt{def:weil-a2-dbc-zl} and \texttt{def:weil-a2-dbc-torsion}.

Fix a proper Cartesian square $\square$ as in S10.  An A2 constructible
rebinding certificate for $K\in\mathsf C(X)$ consists of the following
displayed data.
\begin{enumerate}
\item An A2 proper complete-base-change datum, hence compatible etale-origin
      objects
      \[
       L_n\in D^b_c(X_{\mathrm{et}},\mathbf Z/\ell^n),\qquad
       \alpha_n:\epsilon_X^{-1}L_n\xrightarrow{\sim}K_n,
      \]
      including their transition equations.  The superscript $b$ is part of
      the certificate: one interval $[d_0,d_1]$ works for every $L_n$ and
      every $K_n$.
\item On each of $X,X',Y,Y'$, one finite stratification, each stratum having
      finitely many connected components, and A2 finite-torsion
      constructibility witnesses, uniform in $n$, for every finite object
      occurring in S7--S10.  On $Y'$ these are, in particular, the four
      objects $A_n,B_n,C_n,D_n$ of Definition
      \ref{def:weil-pbc-level-map-corners}; their witnesses use one common
      interval and one common refinement.
\item The finite exceptional-direct-image values for the identity
      compactifications of the proper maps,
      \[
       Rf_!^{(1,f)}K_n\in\mathsf C_n(Y),\qquad
       Rf'_!{}^{(1,f')}E_n\in\mathsf C_n(Y'),
      \]
      with the A2 torsion witnesses retained by the typed A1 input and with
      their transition squares.  The certificate records their displayed
      identifications with $Rf_*K_n$ and $Rf'_*E_n$.
\item A2 constructible-stratification witnesses for
      \[
       S(K):=Lg^*_{\mathrm{comp}}Rf_*K,
       \qquad
       T(K):=Rf'_*L(g')^*_{\mathrm{comp}}K
      \]
      in $\mathsf C(Y')$.  Their finite reductions are required to be the
      objects identified by the maps $\alpha_n$ and $\delta_n$ of Definition
      \ref{def:weil-pbc-completed-reduction-square}.
\end{enumerate}
Let $\mathsf C_\square(X)$ be the full subcategory of $\mathsf C(X)$ on
objects for which the type of such certificates is nonempty.  Certificates
are membership witnesses only; morphisms in $\mathsf C_\square(X)$ are the
ambient A2 morphisms, so no selected witness enters their equality.
\end{definition}

\begin{lemma}[Reduction and completed-tensor transport to A2]
\label{lemma:weil-pbc-a2-reduction-tensor-transport}
\uses{def:weil-pbc-a2-constructible-rebind,def:weil-a2-finite-reduction,def:weil-a2-reduction-transition,def:weil-a2-completed-tensor,def:weil-pbc-complete-coefficient-datum,def:weil-pbc-completed-tensor-product}
For every $K$ and $n$, there is a canonical natural isomorphism
\[
 \chi_{K,n}:\operatorname{Red}^{\mathrm{A2}}_n(K)
 \xrightarrow{\sim}K\otimes^L_{\mathbf Z_\ell}\mathbf Z/\ell^n=K_n.
\]
It intertwines the A2 transition with the quotient transition
$q_{n+1,n}$ of S9.  For complete $A,B$ there is likewise a canonical natural
isomorphism
\[
 \omega_{A,B}:
 A\widehat\otimes^{L,\mathrm{A2}}B
 \xrightarrow{\sim}
 A\widehat\otimes^{L,\mathrm{S9}}B,
\]
and its $n$-th reduction is the identity of
$A\otimes^LB\otimes^L\mathbf Z/\ell^n$ under $\chi$.  Both isomorphisms are
independent of the K-flat, quotient-complex, and completion aliases.
\end{lemma}
\begin{proof}
\uses{lemma:weil-a2-finite-reduction-package-independent,prop:weil-a2-finite-reduction-functor,lemma:weil-pbc-zl-reduction-perfect,lemma:weil-pbc-completed-tensor-reduction}

A2 represents $\mathbf Z/\ell^n$ by the two-term complex
\[
 L_n=[\mathbf Z_\ell\xrightarrow{\ell^n}\mathbf Z_\ell].
\]
The augmentation $L_n\to\mathbf Z/\ell^n$ is a quasi-isomorphism and S9 uses
the same complex as its strict perfect resolution.  Tensor it with the A2
K-flat model of $K$ and localize.  This gives $\chi_{K,n}$.  If the K-flat
model is changed, the two maps are obtained from a common comparison roof;
A2 package independence says that they are equal.  Thus $\chi$ is canonical,
and A2 functoriality makes it natural.

The A2 transition is represented by multiplication by $\ell$ in degree
$-1$ and the identity in degree $0$.  After the augmentation it is the
quotient map
$\mathbf Z/\ell^{n+1}\to\mathbf Z/\ell^n$.  Therefore
\[
 \chi_{K,n}\,r^{\mathrm{A2}}_{n+1,n}
 =q_{n+1,n}\,\chi_{K,n+1}.
\]

Finally, A2 completes the localized K-flat tensor of $A$ and $B$, whereas S9
takes the product-cone limit of its reductions.  Apply $\chi$ levelwise to
these two identical towers.  The transition equation just proved gives a
map of product-cone models and hence $\omega_{A,B}$.  Reduction of completion
and the S9 completed-tensor reduction identify its $n$-th component with the
identity of the displayed triple tensor.  Common K-flat roofs give identical
level maps, and equality into the two products is componentwise; consequently
the induced cone map is independent of every alias.
\end{proof}

\begin{lemma}[Site-comparison and exceptional-image transport to A2]
\label{lemma:weil-pbc-a2-comparison-rfshriek-transport}
\uses{def:weil-pbc-a2-constructible-rebind,def:weil-a2-etale-origin-datum,def:weil-a2-etale-proetale-pullback-comparison,def:weil-a2-etale-proetale-direct-image-comparison,def:weil-a2-six-operation-coefficient-target,def:weil-a1-torsion-derived-category,def:weil-a1-rfshriek-selected,def:weil-pbc-epsilon-comparison,def:weil-pbc-projection-compactification}
Under the identity-on-objects site functor of S8, the A2 notation
$\nu_Z^*$ is the functor $\epsilon_Z^{-1}$.  The A2 pullback comparison
$c_g$, direct-image comparison $d_f$, and etale-origin isomorphism therefore
transport to the S8 maps with the same components.  Moreover the A1 typed
input category $\mathcal D_n(Z)$ transports by the identity on objects and
morphisms to $\mathsf C_n(Z)$, and on a proper map
\[
 Rf_!^{(1,f)}=Rf_*:\mathsf C_n(X)\longrightarrow\mathsf C_n(Y).
\]
Thus the $Rf_!$ field of the A2 six-operation target has exactly the displayed
source and target on $\mathsf C_\square(X)$.  These transports are independent
of whether the A1, A2, S8, or S9 alias is used.
\end{lemma}
\begin{proof}
\uses{lemma:weil-pbc-epsilon-pullback,lemma:weil-pbc-epsilon-direct-image,lemma:weil-pbc-projection-formula-identity-compactification}

Both $\nu_Z$ and $\epsilon_Z$ are the geometric morphism induced by the
functor that sends an etale object to the same object of the pro-etale site.
Their inverse images are the same sheafified left Kan extension, so the
identity of the indexing categories identifies them.  Expanding $c_g$ gives
the S8 pullback square, and expanding $d_f$ gives the S8 unit followed by the
same derived direct image.  The component formulas and their naturality
squares are therefore literally equal, not merely abstractly isomorphic.

A1 denoted by $\mathcal D_n(Z)$ the typed finite torsion constructible
category whose construction was assigned to A2.  A2 realizes that input as
the full subcategory defined by its finite-torsion constructibility witness.
Consequently the transport functor and its inverse are both the identity on
underlying complexes and maps; their two composites are definitional
identities.  The selected A1 exceptional image is $Rp_*j_!$.  This is the S9
definition of $Rf_!^{(j,p)}$.  When $f$ is proper, the selected
compactification is $(1_X,f)$, so $j_!=1$ and both expressions are $Rf_*$.
The identity-compactification lemma shows that the projection-formula arrows
also agree.  Since all transports are identities on representatives, changing
only the displayed alias cannot change an object, a morphism, or a coherence
equation.
\end{proof}

\begin{lemma}[A2 compatibility with every S7--S10 endpoint]
\label{lemma:weil-pbc-a2-s7-s10-endpoint-compatibility}
\uses{def:weil-a2-finite-etale-proper-base-change-map,def:weil-a2-transported-finite-level-base-change-map,def:weil-a2-proetale-proper-base-change-map}
Let $K\in\mathsf C_\square(X)$.  After the transports $\chi$ and $\omega$:
\begin{enumerate}
\item the A2 finite etale map $b_n$ is the S7 unit--counit exchange map;
\item the A2 transported pro-etale map $\widetilde\beta_n$ is the S8 map
      $v'_n$;
\item the A2 product-cone map $\beta_K^{\mathrm{proet}}$ is the S10 completed
      map $\xi_K$;
\item the $n$-th reduction of either completed alias is exactly
      $\widetilde\beta_n$, and the equations are compatible with every tower
      transition, compatible tower endomorphism, identity square, and pasted
      square.
\end{enumerate}
All four equalities are independent of the rebinding certificate chosen for
$K$.
\end{lemma}
\begin{proof}
\uses{lemma:weil-pbc-a2-reduction-tensor-transport,lemma:weil-pbc-a2-comparison-rfshriek-transport,prop:weil-proet-finite-derived-proper-base-change,prop:weil-pbc-proet-finite-level-transport,lemma:weil-pbc-bounded-constructible-reconstruction,lemma:weil-pbc-completed-pullback-reduction,lemma:weil-pbc-completed-tensor-reduction,lemma:weil-pbc-projection-formula-identity-compactification,lemma:weil-pbc-completed-map-reduction,lemma:weil-pbc-completed-exchange-coherence}

The A2 definition of $b_n$ and the S7 definition both use the adjunction unit,
the chosen Cartesian equality $fg'=gf'$, and the adjunction counit, in that
order.  Hence their resolved representatives are the same chain map.  By the
preceding lemma, the four A2 etale/pro-etale comparison arrows are the four S8
arrows.  Pasting the same comparison squares around the common middle map
$b_n$ therefore makes $\widetilde\beta_n=v'_n$.

The A2 completed map and the S10 completed map are both induced from the
tower $(v'_n)_n$ on the product-cone model.  The reduction and tensor
transports commute with the transitions, so the induced maps on both products
are componentwise equal and their cone maps are equal.  The completed
reduction square then gives
\[
 \delta_n\,\xi_K=v'_n\,\alpha_n,
\]
which is precisely the fourth assertion under $\chi$.  S7 naturality, S8
comparison naturality, the S9 tower equations, and S10 completed coherence
give the transition, operator, identity, and pasting equations.  Finally,
the certificate supplies only membership and bound witnesses: none occurs in
the unit--counit formula or in the product-cone map.  Changing it therefore
leaves every displayed morphism unchanged.
\end{proof}

\begin{lemma}[Constructible-subcategory restriction]
\label{lemma:weil-pbc-complete-constructible-restriction}
\uses{def:weil-pbc-a2-constructible-rebind}
For every proper Cartesian square, the complete base-change isomorphism
restricts to a natural isomorphism on the selected A2 full subcategory:
\[
 \xi_K:S(K)=Lg^*_{\mathrm{comp}}Rf_*K
 \xrightarrow{\sim}
 T(K)=Rf'_*L(g')^*_{\mathrm{comp}}K,
 \qquad K\in\mathsf C_\square(X),
\]
where both endpoints lie in $\mathsf C(Y')$.  Its finite reductions, $Rf_!$
proper aliases, etale/pro-etale transports, compatible tower endomorphisms,
and identity and pasting equations are the A2 transports of the S7--S10
endpoints.  The
restricted isomorphism is independent of all certificate and notation
choices.
\end{lemma}
\begin{proof}
\uses{lemma:weil-pbc-a2-reduction-tensor-transport,lemma:weil-pbc-a2-comparison-rfshriek-transport,lemma:weil-pbc-a2-s7-s10-endpoint-compatibility,prop:weil-proet-complete-proper-base-change}

A rebinding certificate supplies the compatible bounded constructible etale
witnesses required in the complete proposition.  Its uniform intervals and
stratifications are exactly the input of S9 bounded reconstruction, while the
perfect coefficient is the two-term resolution used in the reduction
transport lemma.  Thus Proposition
\ref{prop:weil-proet-complete-proper-base-change} gives an ambient isomorphism
$\xi_K:S(K)\to T(K)$ and its explicit ambient inverse.

The last field of the certificate exhibits $S(K)$ and $T(K)$ as objects of
the A2 full subcategory $\mathsf C(Y')$.  Fullness puts $\xi_K$ and its inverse
in that subcategory, and faithfulness preserves their two inverse equations.
The three transport lemmas identify all A2 and A3 aliases and prove every
finite-reduction and coherence equation.  Since the ambient map is constructed
without consulting a membership witness, replacing the certificate changes
neither the map nor its inverse.  This proves the claimed restriction with no
additional external hypothesis.
\end{proof}

\section{Compact-support continuation}\label{chap:compact-support}
\statusnote{Ordinary proper base change is used here only as the proper part of the raw compact-support comparison.  Densification, open exchange, refinement independence, and all four displayed coherence diagrams are additional inputs.}
\subsection*{A43. Dense compactifications and composition}

Let $f:X\to Y$ be separated and finitely presented.  A dense
compactification is data
\[
 C=(X\xrightarrow{j_C}\overline X_C\xrightarrow{p_C}Y,
       \epsilon_C:p_C\circ j_C=f),
\tag{BA43-DENSE-DATA}
\]
where $j_C$ is a schematically dense open immersion and $p_C$ is proper and
finitely presented.  The equality $\epsilon_C$ is stored structure data.  A
refinement $q:D\to C$ is a $Y$-map with $qj_D=j_C$ and
$q^{-1}(X)=X$.  Its properness is a separate geometric atom.

Before any constructible restriction, set
\[
 F_C^{\mathrm{amb}}=Rp_{C*}Rj_{C!}:
 D^+(X_{\acute et},\Lambda_n)\longrightarrow
 D^+(Y_{\acute et},\Lambda_n).
 \tag{BA43-AMBIENT-F}
\]
All $F_C$ and all selected $Rf_!$ in this fragment denote these ambient
functors.  A restriction to $D_c^+$ is a separate theorem.  Consequently
$Rg_!^{\mathrm{amb}}Rf_!^{\mathrm{amb}}$ is typed before any assertion that
$Rf_!K$ is constructible.

The refinement middle owner constructs
\[
 \gamma_q:Rq_*Rj_{D!}\xrightarrow{\sim}Rj_{C!},
 \qquad
 \rho_q:R(p_Cq)_*Rj_{D!}\xrightarrow{\sim}Rp_{C*}Rj_{C!},
\tag{BA43-RHO}
\]
and proves identity, composition, and required naturality.  It also constructs
the same proper-support comparison $\rho_q^{\mathrm{raw}}$ for a proper map
between typed raw models that is the identity over the source open; density is
not used for that construction.  Only dense models enter the following
$\kappa$ system.  For a common
dense refinement $E$, set
\[
 \kappa_{C,D}=\rho_{E\to D}\,\rho_{E\to C}^{-1}.
\]
Schematic-density uniqueness, common refinement, and the equations for
$\rho$ must prove independence of $E$ and the $\kappa$ cocycle.

For $X\xrightarrow fY\xrightarrow gZ$, a successive witness includes a
Cartesian square
\[
\begin{array}{ccc}
\overline X_f&\xrightarrow{u\ \mathrm{schematically\ dense\ open}}&W\\
\mathllap{p_f}\downarrow&&\downarrow\mathrlap{\bar p\ \mathrm{proper}}\\
Y&\xrightarrow{j_g}&\overline Y_g
\end{array}
\]
and a stored equality proving that $u\circ j_f:X\to W$ is schematically
dense.  The open--proper exchange
\[
 Rj_{g!}Rp_{f*}\xrightarrow{\sim}R\bar p_*Ru_!
\tag{BA43-OPX}
\]
must satisfy refinement naturality and typed two-square and three-square
pasting.  The separately constructed isomorphism
\[
 Ru_!Rj_{f!}\xrightarrow{\sim}R(u\circ j_f)_!
\]
has unit, composition, and pentagon laws.  Two-witness independence and triple
pasting then reassemble the choice-independent comparison
$Rg_!Rf_!\to R(gf)_!$.  Thus
\texttt{OBL-BA43-RELATIVE-COHERENT-CORE} remains
\texttt{OPEN/INCOMPLETE}.

\subsection*{A43/33. Densification after base change}

Let $g:Y'\to Y$ be an arbitrary morphism with both $Y$ and $Y'$ noetherian
and finite-dimensional.  Put $X'=X\times_Y Y'$ and
\[
 \overline X_{C,g}^{\mathrm{raw}}=\overline X_C\times_Y Y'.
\]
The pulled-back open need not be dense.  Define the dense model by schematic
closure
\[
 \overline X_{C,g}^{\mathrm{den}}
 =\overline{X'}^{\,\mathrm{sch}}
 \subset\overline X_{C,g}^{\mathrm{raw}},
 \qquad
 d_{C,g}:\overline X_{C,g}^{\mathrm{den}}
       \longrightarrow\overline X_{C,g}^{\mathrm{raw}}.
\tag{BA43-DEN}
\]
The schematic-closure leaf proves that $d_{C,g}$ is the resulting closed
proper map and that it restricts to the identity on $X'$.

A raw compactification is data
\[
 R=(X'\xrightarrow{j_R}\overline X_R\xrightarrow{p_R}Y',
       \epsilon_R:p_R\circ j_R=f'),
\tag{BA43-RAW-DATA}
\]
with $j_R$ open and $p_R$ proper and finitely presented, but with no density
condition.  Its stored equation $\epsilon_R$ is part of the type, and
$F_R^{\mathrm{amb}}=Rp_{R*}Rj_{R!}$ is a functor
$D^+(X'_{\acute et},\Lambda_n)\to D^+(Y'_{\acute et},\Lambda_n)$.
Its $D_c^+$ restriction is asserted only after a constructibility theorem.
Raw models are not objects of the dense
refinement category and have no $\kappa$ arrows.

The terminal leaf \texttt{CAT-GEO-BA43-DENSE-TO-RAW-SUPPORT} has exactly one
output:
\[
 \gamma^{\mathrm{raw}}_{d_{C,g}}:
 Rd_{C,g*}Rj_{C_g^{\mathrm{den}}!}
 \xrightarrow{\sim}Rj_{C_g^{\mathrm{raw}}!}.
\tag{BA43-GAMMA-RAW}
\]
It outputs neither $\sigma$, nor refinement naturality, nor identity,
iterated-base-change, or open-exchange coherence.  Those are all conclusions
of the middle reassembly \texttt{OBL-BA43-BASECHANGE-DENSIFICATION}.  Using
proper-direct-image composition, that middle owner first defines
\[
 \sigma_{C,g}:F_{C_g^{\mathrm{den}}}
 \xrightarrow{\sim}F_{C_g^{\mathrm{raw}}}.
\tag{BA43-SIGMA}
\]
It must then prove the following four separately owned diagrams.  For a
refinement $q:D\to C$, let $\rho^{\mathrm{den}}_{q,g}$ and
$\rho^{\mathrm{raw}}_{q,g}$ be the proper-support comparisons already
constructed by \texttt{OBL-BA43-REFINEMENT-EXCHANGE} on the densified and raw
pullbacks.  Refinement naturality is the
commutative square
\[
\begin{CD}
 F_{D_g^{\mathrm{den}}} @>{\sigma_{D,g}}>> F_{D_g^{\mathrm{raw}}}\\
 @V{\rho^{\mathrm{den}}_{q,g}}VV @VV{\rho^{\mathrm{raw}}_{q,g}}V\\
 F_{C_g^{\mathrm{den}}} @>{\sigma_{C,g}}>> F_{C_g^{\mathrm{raw}}}.
\end{CD}
\tag{BA43-DEN-REF}
\]
After the canonical identifications
$C_{\mathrm{id}}^{\mathrm{den}}=C_{\mathrm{id}}^{\mathrm{raw}}=C$, identity
coherence is the equation
\[
 \sigma_{C,\mathrm{id}_Y}=1_{F_C}.
 \tag{BA43-DEN-ID}
\]

For $Y''\xrightarrow{h}Y'\xrightarrow{g}Y$, let
$\delta_{C;g,h}:(C_g^{\mathrm{den}})_h^{\mathrm{den}}
\xrightarrow{\sim}C_{gh}^{\mathrm{den}}$ be the canonical comparison of the
two schematic closures.  Let $\tau_{C;g,h}$ be the proper-support comparison
from $(C_g^{\mathrm{den}})_h^{\mathrm{raw}}$ to $C_{gh}^{\mathrm{raw}}$
induced by the base change of $d_{C,g}$.  Iterated coherence is
\[
\begin{CD}
 F_{(C_g^{\mathrm{den}})_h^{\mathrm{den}}}
   @>{F(\delta_{C;g,h})}>> F_{C_{gh}^{\mathrm{den}}}\\
 @V{\sigma_{C_g^{\mathrm{den}},h}}VV @VV{\sigma_{C,gh}}V\\
 F_{(C_g^{\mathrm{den}})_h^{\mathrm{raw}}}
   @>{\tau_{C;g,h}}>> F_{C_{gh}^{\mathrm{raw}}}.
\end{CD}
\tag{BA43-DEN-ITER}
\]

Finally let
$A=L\bar g_{\mathrm{raw}}^*Rj_{C!}K$ and let
$s_d:A\xrightarrow{\sim}Rd_{C,g*}Ld_{C,g}^*A$ be the closed-support
identification.  Construct the raw and dense open-exchange maps independently.
Open-exchange coherence is precisely
\[
\begin{CD}
 A @>{\beta^{\mathrm{raw}}_{j_C,g,K}}>>
     Rj_{C_g^{\mathrm{raw}}!}Lg'^*K\\
 @V{s_d}VV @VV{(\gamma^{\mathrm{raw}}_{d_{C,g}})^{-1}}V\\
 Rd_{C,g*}Ld_{C,g}^*A
 @>{Rd_{C,g*}(\beta^{\mathrm{den}}_{j_C,g,K})}>>
 Rd_{C,g*}Rj_{C_g^{\mathrm{den}}!}Lg'^*K.
\end{CD}
\tag{BA43-DEN-OPX}
\]
None of these equations is a field of \textup{(BA43-GAMMA-RAW)}.

\subsection*{A33. Finite compact-support base change}

Continue with the same fixed $\Lambda_n$, ambient categories, and
noetherian finite-dimensional $Y,Y'$.  Let $f:X\to Y$ be separated and
finitely presented, let the square in A31 be Cartesian, and let
$K\in D_c^+(X,\Lambda_n)$.  For the dense compactification $C$ of
\textup{(BA43-DENSE-DATA)}, construct the raw one-model map as the following
two-arrow composite:
\[
\begin{split}
 \operatorname{BC}^{C,\mathrm{raw}}_{f,g,K}:Lg^*F_C(K)
 &\xrightarrow{\operatorname{BC}_{p_C,g,Rj_{C!}K}}
 Rp^{\mathrm{raw}}_{C,g,*}L\bar g^*Rj_{C!}K\\
 &\xrightarrow{Rp^{\mathrm{raw}}_{C,g,*}
      (\beta^{\mathrm{open}}_{j_C,\bar g,K})}
 Rp^{\mathrm{raw}}_{C,g,*}Rj^{\mathrm{raw}}_{C,g,!}Lg'^*K.
\end{split}
\tag{BA33-RAW-BC}
\]
Here the first arrow is the already constructed proper base-change mate and
the second is the derived open-extension exchange.  By definition the name
$\operatorname{BC}^{C,\mathrm{raw}}$ denotes their composite, never the first
arrow alone.  Before finite PBC is applied, the separate leaf
\texttt{CAT-GEO-BA33-OPEN-EXTENSION-CONSTRUCTIBILITY} must prove
$Rj_{C!}K\in D_c^+(\overline X_C,\Lambda_n)$.

The dense representative is
\[
 \operatorname{BC}^{C,\mathrm{den}}_{f,g,K}
 =\sigma_{C,g}^{-1}\circ
   \operatorname{BC}^{C,\mathrm{raw}}_{f,g,K}:
 Lg^*F_C(K)\longrightarrow
 F_{C_g^{\mathrm{den}}}(Lg'^*K).
\tag{BA33-DEN-BC}
\]
For a refinement $q:D\to C$,
\texttt{OBL-BA33-REFINEMENT-BASECHANGE-NATURALITY} proves equality of the
two explicit composites built from $Lg^*(\rho_q)$,
$\rho_{q_g^{\mathrm{den}}}$, the two maps
\textup{(BA33-DEN-BC)}, and the densification comparison.  The proof must
expand proper-base-change naturality, open-exchange naturality, and all four
densification diagrams; the equality is not supplied in a comparison record.

Only then may the $\kappa$ cocycle descend the representatives to
\[
 Lg^*Rf_!K\longrightarrow Rf'_!Lg'^*K.
\]
Identity and horizontal/vertical pasting consume the corresponding proper,
open, densification, and $\kappa$ equations.  Therefore
\texttt{OBL-BA33-FINITE-COMPACT-SUPPORT-BASE-CHANGE} remains
\texttt{OPEN/INCOMPLETE}.  Geometric-fibre, Frobenius-equivariant, tower,
integral, and rational exports are outside this finite-level owner.

\section{Exact unresolved interfaces}\label{app:gap-ledger}
\statusnote{Every entry in this appendix remains OPEN.  The middle entries are reductions; the leaf entries are the exact terminal proof obligations.  No entry is discharged by the extracted prose alone.  The eight obligations \nolinkurl{OBL-A4-RC-FORMAL-SERRE}, \nolinkurl{OBL-A4-RC-FORMAL-FINITE-DIRECT-IMAGE}, \nolinkurl{OBL-A4-RC-FORMAL-SEED-LEAVES}, \nolinkurl{OBL-A4-RC-FORMAL-AFFINE-LINE-DEVISSAGE}, \nolinkurl{OBL-A4-RC-FORMAL-AFFINE-SPACE}, \nolinkurl{OBL-A4-RC-FORMAL-CLOSED-IMMERSION}, \nolinkurl{OBL-A4-RC-FORMAL-AFFINE-MORPHISM}, \nolinkurl{OBL-A4-RC-FORMAL-PROPER-CECH} and the obligation \nolinkurl{OBL-SIX-FUNCTOR-ADIC} are referenced in the extracted prose but own no entry in this chapter: they lie outside the extracted S0--S10 block, and this chapter states no interface for them.}
\section{Middle interfaces}
\Needspace{8\baselineskip}
\begin{openproblem}[Exact unresolved interface]
\label{OBL-BA30-PROPER-FINITE-CONSTRUCTIBLE}
\textbf{Owner. }\nolinkurl{OBL-BA30-PROPER-FINITE-CONSTRUCTIBLE}\
\textbf{Status. }\texttt{OPEN}; \textbf{reduction status. }\texttt{INCOMPLETE}; \textbf{class. }\texttt{MIDDLE}.
\uses{OBL-BA30-A4-SEED-TO-PROPER-FINITE,EXT-BA30-RELATIVE-COHOMOLOGICAL-DIMENSION,FA-BA30-BOUNDED-POSTNIKOV-SERRE}
\textbf{Exact output required. }For fixed \nolinkurl{Lambda_n}, a proper finitely presented \nolinkurl{p:Z->Y} of fibre dimension at most \nolinkurl{d} sends \nolinkurl{D_c^[a,b](Z,Lambda_n)} to \nolinkurl{D_c^[a,b+2d](Y,Lambda_n)}.
\openobligation{OBL-BA30-PROPER-FINITE-CONSTRUCTIBLE}
\end{openproblem}

\Needspace{8\baselineskip}
\begin{openproblem}[Exact unresolved interface]
\label{OBL-BA30-A4-SEED-TO-PROPER-FINITE}
\textbf{Owner. }\nolinkurl{OBL-BA30-A4-SEED-TO-PROPER-FINITE}\
\textbf{Status. }\texttt{OPEN}; \textbf{reduction status. }\texttt{INCOMPLETE}; \textbf{class. }\texttt{MIDDLE}.
\uses{EXT-BA30-A4-SEED-CONSTRUCTOR,CAT-BA30-RFSTAR-REBIND,EXT-BA30-RELATIVE-COHOMOLOGICAL-DIMENSION,FA-BA30-BOUNDED-POSTNIKOV-SERRE}
\textbf{Exact output required. }An actual literal A4 seed, the ambient \nolinkurl{Rf_*} rebind, the formal A4 chain, the relative \nolinkurl{2d} bound, and finite Postnikov/Serre closure imply the finite proper theorem.
\openobligation{OBL-BA30-A4-SEED-TO-PROPER-FINITE}
\end{openproblem}

\Needspace{8\baselineskip}
\begin{openproblem}[Exact unresolved interface]
\label{OBL-BA31-FINITE-PBC}
\textbf{Owner. }\nolinkurl{OBL-BA31-FINITE-PBC}\
\textbf{Status. }\texttt{OPEN}; \textbf{reduction status. }\texttt{INCOMPLETE}; \textbf{class. }\texttt{MIDDLE}.
\uses{OBL-BA31-S0-GLOBAL-MATE,OBL-BA31-S1-FINITE,OBL-BA31-S2-STABILITY,OBL-BA31-S3-PROJECTIVE-REDUCTION,OBL-BA31-S4-HENSELIAN-P1,OBL-BA31-S5-STALK-COMPARISON,OBL-BA31-S6-SHEAF-ASSEMBLY,OBL-BA31-S7-BOUNDED-BELOW}
\textbf{Exact output required. }For fixed \nolinkurl{Lambda_n}, the globally constructed unit-counit mate is invertible on \nolinkurl{D_c^+} for a proper finitely presented Cartesian square and obeys mate-calculus identity/horizontal/vertical coherence.
\openobligation{OBL-BA31-FINITE-PBC}
\end{openproblem}

\Needspace{8\baselineskip}
\begin{openproblem}[Exact unresolved interface]
\label{OBL-BA31-S0-GLOBAL-MATE}
\textbf{Owner. }\nolinkurl{OBL-BA31-S0-GLOBAL-MATE}\
\textbf{Status. }\texttt{OPEN}; \textbf{reduction status. }\texttt{INCOMPLETE}; \textbf{class. }\texttt{MIDDLE}.
\uses{CAT-BA31-SMALL-ETALE-MATE,CAT-BA31-MATE-REPLACEMENT-INDEPENDENCE,CAT-BA31-MATE-COHERENCE}
\textbf{Exact output required. }Construct the single ambient mate and prove replacement independence, naturality, identity, and the two pasting equations.
\openobligation{OBL-BA31-S0-GLOBAL-MATE}
\end{openproblem}

\Needspace{8\baselineskip}
\begin{openproblem}[Exact unresolved interface]
\label{OBL-BA31-S1-FINITE}
\textbf{Owner. }\nolinkurl{OBL-BA31-S1-FINITE}\
\textbf{Status. }\texttt{OPEN}; \textbf{reduction status. }\texttt{INCOMPLETE}; \textbf{class. }\texttt{MIDDLE}.
\uses{OBL-BA31-S0-GLOBAL-MATE,CAT-GEO-BA31-FINITE-FIBRE-PRODUCT}
\textbf{Exact output required. }Prove exact finite pushforward and the finite fibre product formula, and identify that map with the stalk of the fixed mate.
\openobligation{OBL-BA31-S1-FINITE}
\end{openproblem}

\Needspace{8\baselineskip}
\begin{openproblem}[Exact unresolved interface]
\label{OBL-BA31-S2-STABILITY}
\textbf{Owner. }\nolinkurl{OBL-BA31-S2-STABILITY}\
\textbf{Status. }\texttt{OPEN}; \textbf{reduction status. }\texttt{INCOMPLETE}; \textbf{class. }\texttt{MIDDLE}.
\uses{OBL-BA31-S0-GLOBAL-MATE,CAT-BA30-RFSTAR-REBIND,FA-CAT-BA31-LERAY-CANCELLATION}
\textbf{Exact output required. }Prove typed composition and proper-surjective cancellation for the fixed mate and actual \nolinkurl{Rf_*}; no inverse is input data.
\openobligation{OBL-BA31-S2-STABILITY}
\end{openproblem}

\Needspace{8\baselineskip}
\begin{openproblem}[Exact unresolved interface]
\label{OBL-BA31-S3-PROJECTIVE-REDUCTION}
\textbf{Owner. }\nolinkurl{OBL-BA31-S3-PROJECTIVE-REDUCTION}\
\textbf{Status. }\texttt{OPEN}; \textbf{reduction status. }\texttt{INCOMPLETE}; \textbf{class. }\texttt{MIDDLE}.
\uses{OBL-BA31-S1-FINITE,OBL-BA31-S2-STABILITY,EXT-GEO-BA31-CHOW-PROPER-FP,GEO-BA31-SYMMETRIC-P1-FINITE-SURJECTIVE}
\textbf{Exact output required. }Give the proper-finitely-presented Chow datum and the finite-surjective retract/cancellation reduction through projective space to products of projective lines.
\openobligation{OBL-BA31-S3-PROJECTIVE-REDUCTION}
\end{openproblem}

\Needspace{8\baselineskip}
\begin{openproblem}[Exact unresolved interface]
\label{OBL-BA31-S4-HENSELIAN-P1}
\textbf{Owner. }\nolinkurl{OBL-BA31-S4-HENSELIAN-P1}\
\textbf{Status. }\texttt{OPEN}; \textbf{reduction status. }\texttt{INCOMPLETE}; \textbf{class. }\texttt{MIDDLE}.
\uses{EXT-GEO-BA31-HENSELIAN-H0,EXT-GEO-BA31-GABBER-AFFINE,EXT-GEO-BA31-TORSOR-H1}
\textbf{Exact output required. }Assemble independently owned H0, affine, and torsor/H1 inputs into strict-henselian projective-line comparison.
\openobligation{OBL-BA31-S4-HENSELIAN-P1}
\end{openproblem}

\Needspace{8\baselineskip}
\begin{openproblem}[Exact unresolved interface]
\label{OBL-BA31-S5-STALK-COMPARISON}
\textbf{Owner. }\nolinkurl{OBL-BA31-S5-STALK-COMPARISON}\
\textbf{Status. }\texttt{OPEN}; \textbf{reduction status. }\texttt{INCOMPLETE}; \textbf{class. }\texttt{MIDDLE}.
\uses{OBL-BA31-S0-GLOBAL-MATE,OBL-BA31-S4-HENSELIAN-P1,CAT-GEO-BA31-STRICT-HENSEL-STALK,EXT-GEO-BA31-CURVE-FIELD-INVARIANCE,CAT-BA31-SHEAF-STALK-FAITHFULNESS}
\textbf{Exact output required. }Identify the strict-henselian comparison with each \nolinkurl{H^i(BC)}; use sheaf-stalk faithfulness only for equality of already constructed maps of cohomology sheaves.
\openobligation{OBL-BA31-S5-STALK-COMPARISON}
\end{openproblem}

\Needspace{8\baselineskip}
\begin{openproblem}[Exact unresolved interface]
\label{OBL-BA31-S6-SHEAF-ASSEMBLY}
\textbf{Owner. }\nolinkurl{OBL-BA31-S6-SHEAF-ASSEMBLY}\
\textbf{Status. }\texttt{OPEN}; \textbf{reduction status. }\texttt{INCOMPLETE}; \textbf{class. }\texttt{MIDDLE}.
\uses{OBL-BA31-S0-GLOBAL-MATE,OBL-BA31-S1-FINITE,OBL-BA31-S2-STABILITY,OBL-BA31-S3-PROJECTIVE-REDUCTION,OBL-BA31-S5-STALK-COMPARISON,CAT-BA31-DERIVED-STALK-ISO-CONSERVATIVITY}
\textbf{Exact output required. }Apply derived-stalk isomorphism conservativity only to the fixed mate after all its cohomology-stalk maps are invertible.
\openobligation{OBL-BA31-S6-SHEAF-ASSEMBLY}
\end{openproblem}

\Needspace{8\baselineskip}
\begin{openproblem}[Exact unresolved interface]
\label{OBL-BA31-S7-BOUNDED-BELOW}
\textbf{Owner. }\nolinkurl{OBL-BA31-S7-BOUNDED-BELOW}\
\textbf{Status. }\texttt{OPEN}; \textbf{reduction status. }\texttt{INCOMPLETE}; \textbf{class. }\texttt{MIDDLE}.
\uses{OBL-BA31-S0-GLOBAL-MATE,OBL-BA31-S6-SHEAF-ASSEMBLY,FA-CAT-BA31-BOUNDED-BELOW-POSTNIKOV}
\textbf{Exact output required. }Use a selected bounded-below Postnikov/Cartan-Eilenberg construction, convergence, and comparison identification to upgrade the bounded result to \nolinkurl{D_c^+}.
\openobligation{OBL-BA31-S7-BOUNDED-BELOW}
\end{openproblem}

\Needspace{8\baselineskip}
\begin{openproblem}[Exact unresolved interface]
\label{OBL-BA43-RELATIVE-COHERENT-CORE}
\textbf{Owner. }\nolinkurl{OBL-BA43-RELATIVE-COHERENT-CORE}\
\textbf{Status. }\texttt{OPEN}; \textbf{reduction status. }\texttt{INCOMPLETE}; \textbf{class. }\texttt{MIDDLE}.
\uses{GEO-BA43-REFINEMENT-PROPER,GEO-BA43-DENSE-REFINEMENT-UNIQUENESS,GEO-BA43-PAIRWISE-COMMON-REFINEMENT,OBL-BA43-REFINEMENT-EXCHANGE,CAT-BA30-RFSTAR-REBIND,EXT-GEO-BA43-DENSE-SUCCESSIVE-WITNESS,CAT-GEO-BA43-OPEN-PROPER-EXCHANGE,CAT-BA43-OPEN-EXTENSION-COMPOSITION,CAT-BA43-SUCCESSIVE-WITNESS-INDEPENDENCE,CAT-BA43-TRIPLE-PASTING}
\textbf{Exact output required. }For typed dense compactifications, build refinement comparisons, the choice-independent selected \nolinkurl{Rf_!}, and unit/composition/associativity from successive witnesses.
\openobligation{OBL-BA43-RELATIVE-COHERENT-CORE}
\end{openproblem}

\Needspace{8\baselineskip}
\begin{openproblem}[Exact unresolved interface]
\label{OBL-BA43-REFINEMENT-EXCHANGE}
\textbf{Owner. }\nolinkurl{OBL-BA43-REFINEMENT-EXCHANGE}\
\textbf{Status. }\texttt{OPEN}; \textbf{reduction status. }\texttt{INCOMPLETE}; \textbf{class. }\texttt{MIDDLE}.
\uses{OBL-BA31-FINITE-PBC,CAT-BA30-RFSTAR-REBIND,CAT-BA33-DERIVED-OPEN-EXCHANGE}
\textbf{Exact output required. }Construct \nolinkurl{gamma_q} and \nolinkurl{rho_q}, their inverses, identity/composition, and downstream naturality for dense refinements; also construct the proper-support \nolinkurl{rho_raw} for a proper typed raw-model map that is the identity over the source open. This is a middle node, not a leaf.
\openobligation{OBL-BA43-REFINEMENT-EXCHANGE}
\end{openproblem}

\Needspace{8\baselineskip}
\begin{openproblem}[Exact unresolved interface]
\label{OBL-BA43-BASECHANGE-DENSIFICATION}
\textbf{Owner. }\nolinkurl{OBL-BA43-BASECHANGE-DENSIFICATION}\
\textbf{Status. }\texttt{OPEN}; \textbf{reduction status. }\texttt{INCOMPLETE}; \textbf{class. }\texttt{MIDDLE}.
\uses{GEO-BA43-SCHEMATIC-CLOSURE-DENSIFICATION,CAT-GEO-BA43-DENSE-TO-RAW-SUPPORT,CAT-BA30-RFSTAR-REBIND,OBL-BA43-REFINEMENT-EXCHANGE,CAT-BA43-DENSIFICATION-REFINEMENT-NATURALITY,CAT-BA43-DENSIFICATION-IDENTITY,CAT-BA43-DENSIFICATION-ITERATED-COHERENCE,CAT-BA43-DENSIFICATION-OPEN-EXCHANGE}
\textbf{Exact output required. }For noetherian finite-dimensional \nolinkurl{Y,Y'}, build \nolinkurl{sigma} from schematic closure and the gamma-raw atom, then prove refinement, identity, iterated-base-change, and open-exchange coherence.
\openobligation{OBL-BA43-BASECHANGE-DENSIFICATION}
\end{openproblem}

\Needspace{8\baselineskip}
\begin{openproblem}[Exact unresolved interface]
\label{OBL-BA33-REFINEMENT-BASECHANGE-NATURALITY}
\textbf{Owner. }\nolinkurl{OBL-BA33-REFINEMENT-BASECHANGE-NATURALITY}\
\textbf{Status. }\texttt{OPEN}; \textbf{reduction status. }\texttt{INCOMPLETE}; \textbf{class. }\texttt{MIDDLE}.
\uses{OBL-BA31-S0-GLOBAL-MATE,CAT-BA33-DERIVED-OPEN-EXCHANGE,CAT-GEO-BA33-OPEN-EXTENSION-CONSTRUCTIBILITY,OBL-BA43-REFINEMENT-EXCHANGE,OBL-BA43-BASECHANGE-DENSIFICATION}
\textbf{Exact output required. }Prove equality of the two explicit densified one-model BC composites under refinement after expanding proper BC, open exchange, \nolinkurl{rho}, and \nolinkurl{sigma}.
\openobligation{OBL-BA33-REFINEMENT-BASECHANGE-NATURALITY}
\end{openproblem}

\Needspace{8\baselineskip}
\begin{openproblem}[Exact unresolved interface]
\label{OBL-BA33-FINITE-COMPACT-SUPPORT-BASE-CHANGE}
\textbf{Owner. }\nolinkurl{OBL-BA33-FINITE-COMPACT-SUPPORT-BASE-CHANGE}\
\textbf{Status. }\texttt{OPEN}; \textbf{reduction status. }\texttt{INCOMPLETE}; \textbf{class. }\texttt{MIDDLE}.
\uses{OBL-BA31-FINITE-PBC,CAT-BA33-DERIVED-OPEN-EXCHANGE,CAT-GEO-BA33-OPEN-EXTENSION-CONSTRUCTIBILITY,OBL-BA43-RELATIVE-COHERENT-CORE,OBL-BA43-BASECHANGE-DENSIFICATION,OBL-BA33-REFINEMENT-BASECHANGE-NATURALITY,CAT-BA33-FINITE-BC-PASTING}
\textbf{Exact output required. }The raw representative is proper BC followed by open exchange; densify, descend through \nolinkurl{kappa}, and prove identity/horizontal/vertical pasting.
\openobligation{OBL-BA33-FINITE-COMPACT-SUPPORT-BASE-CHANGE}
\end{openproblem}

\section{Terminal leaves}
\Needspace{8\baselineskip}
\begin{openproblem}[Exact unresolved interface]
\label{EXT-BA30-A4-SEED-CONSTRUCTOR}
\textbf{Owner. }\nolinkurl{EXT-BA30-A4-SEED-CONSTRUCTOR}\
\textbf{Status. }\texttt{OPEN}; \textbf{reduction status. }\texttt{LEAF}; \textbf{class. }\texttt{MIXED}.
\uses{}
\textbf{Exact output required. }Construct the literal \nolinkurl{RelativeCurveConstructibilitySeed(T,Lambda_n)} inhabitant: specialization, finite monodromy, compactification, descent/refinement independence, finite affine-line devissage tree, and all coherence fields.
\openobligation{EXT-BA30-A4-SEED-CONSTRUCTOR}
\end{openproblem}

\Needspace{8\baselineskip}
\begin{openproblem}[Exact unresolved interface]
\label{CAT-BA30-RFSTAR-REBIND}
\textbf{Owner. }\nolinkurl{CAT-BA30-RFSTAR-REBIND}\
\textbf{Status. }\texttt{OPEN}; \textbf{reduction status. }\texttt{LEAF}; \textbf{class. }\texttt{CAT/FA}.
\uses{}
\textbf{Exact output required. }Identify every consumed historical direct image with the ambient K-injective \nolinkurl{Rf_*}; prove naturality, unit, composition, and pentagon.
\openobligation{CAT-BA30-RFSTAR-REBIND}
\end{openproblem}

\Needspace{8\baselineskip}
\begin{openproblem}[Exact unresolved interface]
\label{EXT-BA30-RELATIVE-COHOMOLOGICAL-DIMENSION}
\textbf{Owner. }\nolinkurl{EXT-BA30-RELATIVE-COHOMOLOGICAL-DIMENSION}\
\textbf{Status. }\texttt{OPEN}; \textbf{reduction status. }\texttt{LEAF}; \textbf{class. }\texttt{EXT-GEO}.
\uses{}
\textbf{Exact output required. }Under the displayed noetherian finite-dimensional and relative-dimension hypotheses, prove the exact vanishing above degree \nolinkurl{2d}.
\openobligation{EXT-BA30-RELATIVE-COHOMOLOGICAL-DIMENSION}
\end{openproblem}

\Needspace{8\baselineskip}
\begin{openproblem}[Exact unresolved interface]
\label{FA-BA30-BOUNDED-POSTNIKOV-SERRE}
\textbf{Owner. }\nolinkurl{FA-BA30-BOUNDED-POSTNIKOV-SERRE}\
\textbf{Status. }\texttt{OPEN}; \textbf{reduction status. }\texttt{LEAF}; \textbf{class. }\texttt{FA}.
\uses{}
\textbf{Exact output required. }From constructible cohomology sheaves and the independent degree bound, prove membership in \nolinkurl{D_c^[a,b+2d]}.
\openobligation{FA-BA30-BOUNDED-POSTNIKOV-SERRE}
\end{openproblem}

\Needspace{8\baselineskip}
\begin{openproblem}[Exact unresolved interface]
\label{CAT-BA31-SMALL-ETALE-MATE}
\textbf{Owner. }\nolinkurl{CAT-BA31-SMALL-ETALE-MATE}\
\textbf{Status. }\texttt{OPEN}; \textbf{reduction status. }\texttt{LEAF}; \textbf{class. }\texttt{CAT}.
\uses{}
\textbf{Exact output required. }Construct the unit-cartesian-identification-counit mate in the fixed ambient etale derived categories and prove naturality in \nolinkurl{K}.
\openobligation{CAT-BA31-SMALL-ETALE-MATE}
\end{openproblem}

\Needspace{8\baselineskip}
\begin{openproblem}[Exact unresolved interface]
\label{CAT-BA31-MATE-REPLACEMENT-INDEPENDENCE}
\textbf{Owner. }\nolinkurl{CAT-BA31-MATE-REPLACEMENT-INDEPENDENCE}\
\textbf{Status. }\texttt{OPEN}; \textbf{reduction status. }\texttt{LEAF}; \textbf{class. }\texttt{CAT/FA}.
\uses{}
\textbf{Exact output required. }Prove the mate is independent of the selected K-injective/K-flat representatives.
\openobligation{CAT-BA31-MATE-REPLACEMENT-INDEPENDENCE}
\end{openproblem}

\Needspace{8\baselineskip}
\begin{openproblem}[Exact unresolved interface]
\label{CAT-BA31-MATE-COHERENCE}
\textbf{Owner. }\nolinkurl{CAT-BA31-MATE-COHERENCE}\
\textbf{Status. }\texttt{OPEN}; \textbf{reduction status. }\texttt{LEAF}; \textbf{class. }\texttt{CAT}.
\uses{}
\textbf{Exact output required. }Prove identity and horizontal/vertical pasting for the fixed global mate by adjunction calculus.
\openobligation{CAT-BA31-MATE-COHERENCE}
\end{openproblem}

\Needspace{8\baselineskip}
\begin{openproblem}[Exact unresolved interface]
\label{CAT-GEO-BA31-FINITE-FIBRE-PRODUCT}
\textbf{Owner. }\nolinkurl{CAT-GEO-BA31-FINITE-FIBRE-PRODUCT}\
\textbf{Status. }\texttt{OPEN}; \textbf{reduction status. }\texttt{LEAF}; \textbf{class. }\texttt{CAT/GEO}.
\uses{}
\textbf{Exact output required. }For finite \nolinkurl{f}, prove exact pushforward, the finite geometric-fibre product formula, and equality with the stalk map of the global mate.
\openobligation{CAT-GEO-BA31-FINITE-FIBRE-PRODUCT}
\end{openproblem}

\Needspace{8\baselineskip}
\begin{openproblem}[Exact unresolved interface]
\label{FA-CAT-BA31-LERAY-CANCELLATION}
\textbf{Owner. }\nolinkurl{FA-CAT-BA31-LERAY-CANCELLATION}\
\textbf{Status. }\texttt{OPEN}; \textbf{reduction status. }\texttt{LEAF}; \textbf{class. }\texttt{FA/CAT}.
\uses{}
\textbf{Exact output required. }State and prove the exact Leray/composition cancellation diagram used in proper-surjective descent.
\openobligation{FA-CAT-BA31-LERAY-CANCELLATION}
\end{openproblem}

\Needspace{8\baselineskip}
\begin{openproblem}[Exact unresolved interface]
\label{EXT-GEO-BA31-CHOW-PROPER-FP}
\textbf{Owner. }\nolinkurl{EXT-GEO-BA31-CHOW-PROPER-FP}\
\textbf{Status. }\texttt{OPEN}; \textbf{reduction status. }\texttt{LEAF}; \textbf{class. }\texttt{EXT-GEO}.
\uses{}
\textbf{Exact output required. }Supply the exact Chow datum for proper finitely presented maps in the stated noetherian finite-dimensional class.
\openobligation{EXT-GEO-BA31-CHOW-PROPER-FP}
\end{openproblem}

\Needspace{8\baselineskip}
\begin{openproblem}[Exact unresolved interface]
\label{GEO-BA31-SYMMETRIC-P1-FINITE-SURJECTIVE}
\textbf{Owner. }\nolinkurl{GEO-BA31-SYMMETRIC-P1-FINITE-SURJECTIVE}\
\textbf{Status. }\texttt{OPEN}; \textbf{reduction status. }\texttt{LEAF}; \textbf{class. }\texttt{GEO}.
\uses{}
\textbf{Exact output required. }Construct the finite-surjective symmetric cover reducing the projective-space case to a product of projective lines, with the required retraction datum.
\openobligation{GEO-BA31-SYMMETRIC-P1-FINITE-SURJECTIVE}
\end{openproblem}

\Needspace{8\baselineskip}
\begin{openproblem}[Exact unresolved interface]
\label{EXT-GEO-BA31-HENSELIAN-H0}
\textbf{Owner. }\nolinkurl{EXT-GEO-BA31-HENSELIAN-H0}\
\textbf{Status. }\texttt{OPEN}; \textbf{reduction status. }\texttt{LEAF}; \textbf{class. }\texttt{EXT-GEO}.
\uses{}
\textbf{Exact output required. }Prove the degree-zero strict-henselian projective-line comparison.
\openobligation{EXT-GEO-BA31-HENSELIAN-H0}
\end{openproblem}

\Needspace{8\baselineskip}
\begin{openproblem}[Exact unresolved interface]
\label{EXT-GEO-BA31-GABBER-AFFINE}
\textbf{Owner. }\nolinkurl{EXT-GEO-BA31-GABBER-AFFINE}\
\textbf{Status. }\texttt{OPEN}; \textbf{reduction status. }\texttt{LEAF}; \textbf{class. }\texttt{EXT-GEO}.
\uses{}
\textbf{Exact output required. }Prove the exact affine effacement/comparison statement consumed by S4.
\openobligation{EXT-GEO-BA31-GABBER-AFFINE}
\end{openproblem}

\Needspace{8\baselineskip}
\begin{openproblem}[Exact unresolved interface]
\label{EXT-GEO-BA31-TORSOR-H1}
\textbf{Owner. }\nolinkurl{EXT-GEO-BA31-TORSOR-H1}\
\textbf{Status. }\texttt{OPEN}; \textbf{reduction status. }\texttt{LEAF}; \textbf{class. }\texttt{EXT-GEO}.
\uses{}
\textbf{Exact output required. }Prove the finite-etale torsor lifting and degree-one comparison consumed by S4.
\openobligation{EXT-GEO-BA31-TORSOR-H1}
\end{openproblem}

\Needspace{8\baselineskip}
\begin{openproblem}[Exact unresolved interface]
\label{CAT-GEO-BA31-STRICT-HENSEL-STALK}
\textbf{Owner. }\nolinkurl{CAT-GEO-BA31-STRICT-HENSEL-STALK}\
\textbf{Status. }\texttt{OPEN}; \textbf{reduction status. }\texttt{LEAF}; \textbf{class. }\texttt{CAT/GEO}.
\uses{}
\textbf{Exact output required. }Identify the geometric stalk of \nolinkurl{Rf_*K} with cohomology over the strict-henselization, naturally in the fixed mate.
\openobligation{CAT-GEO-BA31-STRICT-HENSEL-STALK}
\end{openproblem}

\Needspace{8\baselineskip}
\begin{openproblem}[Exact unresolved interface]
\label{EXT-GEO-BA31-CURVE-FIELD-INVARIANCE}
\textbf{Owner. }\nolinkurl{EXT-GEO-BA31-CURVE-FIELD-INVARIANCE}\
\textbf{Status. }\texttt{OPEN}; \textbf{reduction status. }\texttt{LEAF}; \textbf{class. }\texttt{EXT-GEO}.
\uses{}
\textbf{Exact output required. }Prove the proper-curve separably-closed field-extension invariance needed by S5.
\openobligation{EXT-GEO-BA31-CURVE-FIELD-INVARIANCE}
\end{openproblem}

\Needspace{8\baselineskip}
\begin{openproblem}[Exact unresolved interface]
\label{CAT-BA31-SHEAF-STALK-FAITHFULNESS}
\textbf{Owner. }\nolinkurl{CAT-BA31-SHEAF-STALK-FAITHFULNESS}\
\textbf{Status. }\texttt{OPEN}; \textbf{reduction status. }\texttt{LEAF}; \textbf{class. }\texttt{CAT}.
\uses{}
\textbf{Exact output required. }Equality of morphisms in the abelian category of \nolinkurl{Lambda_n}-sheaves is detected on all geometric stalks; no derived equality claim.
\openobligation{CAT-BA31-SHEAF-STALK-FAITHFULNESS}
\end{openproblem}

\Needspace{8\baselineskip}
\begin{openproblem}[Exact unresolved interface]
\label{CAT-BA31-DERIVED-STALK-ISO-CONSERVATIVITY}
\textbf{Owner. }\nolinkurl{CAT-BA31-DERIVED-STALK-ISO-CONSERVATIVITY}\
\textbf{Status. }\texttt{OPEN}; \textbf{reduction status. }\texttt{LEAF}; \textbf{class. }\texttt{CAT}.
\uses{}
\textbf{Exact output required. }An already constructed map in \nolinkurl{D^+} is invertible if every cohomology-sheaf stalk map is invertible; no fullness or inverse construction.
\openobligation{CAT-BA31-DERIVED-STALK-ISO-CONSERVATIVITY}
\end{openproblem}

\Needspace{8\baselineskip}
\begin{openproblem}[Exact unresolved interface]
\label{FA-CAT-BA31-BOUNDED-BELOW-POSTNIKOV}
\textbf{Owner. }\nolinkurl{FA-CAT-BA31-BOUNDED-BELOW-POSTNIKOV}\
\textbf{Status. }\texttt{OPEN}; \textbf{reduction status. }\texttt{LEAF}; \textbf{class. }\texttt{FA/CAT}.
\uses{}
\textbf{Exact output required. }Select bounded truncation/Postnikov or Cartan-Eilenberg models, prove convergence, independence, and identify the induced comparison with the fixed global mate.
\openobligation{FA-CAT-BA31-BOUNDED-BELOW-POSTNIKOV}
\end{openproblem}

\Needspace{8\baselineskip}
\begin{openproblem}[Exact unresolved interface]
\label{GEO-BA43-REFINEMENT-PROPER}
\textbf{Owner. }\nolinkurl{GEO-BA43-REFINEMENT-PROPER}\
\textbf{Status. }\texttt{OPEN}; \textbf{reduction status. }\texttt{LEAF}; \textbf{class. }\texttt{GEO}.
\uses{}
\textbf{Exact output required. }A refinement between the typed models is proper because its source is proper and its target separated over the base.
\openobligation{GEO-BA43-REFINEMENT-PROPER}
\end{openproblem}

\Needspace{8\baselineskip}
\begin{openproblem}[Exact unresolved interface]
\label{GEO-BA43-DENSE-REFINEMENT-UNIQUENESS}
\textbf{Owner. }\nolinkurl{GEO-BA43-DENSE-REFINEMENT-UNIQUENESS}\
\textbf{Status. }\texttt{OPEN}; \textbf{reduction status. }\texttt{LEAF}; \textbf{class. }\texttt{GEO}.
\uses{}
\textbf{Exact output required. }Two refinement maps agreeing on the schematically dense source open agree globally when the target is separated.
\openobligation{GEO-BA43-DENSE-REFINEMENT-UNIQUENESS}
\end{openproblem}

\Needspace{8\baselineskip}
\begin{openproblem}[Exact unresolved interface]
\label{GEO-BA43-PAIRWISE-COMMON-REFINEMENT}
\textbf{Owner. }\nolinkurl{GEO-BA43-PAIRWISE-COMMON-REFINEMENT}\
\textbf{Status. }\texttt{OPEN}; \textbf{reduction status. }\texttt{LEAF}; \textbf{class. }\texttt{GEO}.
\uses{}
\textbf{Exact output required. }Construct a typed common dense refinement with finitely presented proper structure map.
\openobligation{GEO-BA43-PAIRWISE-COMMON-REFINEMENT}
\end{openproblem}

\Needspace{8\baselineskip}
\begin{openproblem}[Exact unresolved interface]
\label{CAT-GEO-BA43-OPEN-PROPER-EXCHANGE}
\textbf{Owner. }\nolinkurl{CAT-GEO-BA43-OPEN-PROPER-EXCHANGE}\
\textbf{Status. }\texttt{OPEN}; \textbf{reduction status. }\texttt{LEAF}; \textbf{class. }\texttt{CAT/GEO}.
\uses{}
\textbf{Exact output required. }Construct \nolinkurl{Rj_!Rp_* ~= Rbarp_*Ru_!} for the stated Cartesian open-proper square and prove refinement naturality and two-/three-square pasting.
\openobligation{CAT-GEO-BA43-OPEN-PROPER-EXCHANGE}
\end{openproblem}

\Needspace{8\baselineskip}
\begin{openproblem}[Exact unresolved interface]
\label{EXT-GEO-BA43-DENSE-SUCCESSIVE-WITNESS}
\textbf{Owner. }\nolinkurl{EXT-GEO-BA43-DENSE-SUCCESSIVE-WITNESS}\
\textbf{Status. }\texttt{OPEN}; \textbf{reduction status. }\texttt{LEAF}; \textbf{class. }\texttt{EXT-GEO}.
\uses{}
\textbf{Exact output required. }Produce the typed Cartesian successive witness with schematically dense \nolinkurl{u} and schematically dense composite \nolinkurl{u j_f}.
\openobligation{EXT-GEO-BA43-DENSE-SUCCESSIVE-WITNESS}
\end{openproblem}

\Needspace{8\baselineskip}
\begin{openproblem}[Exact unresolved interface]
\label{CAT-BA43-OPEN-EXTENSION-COMPOSITION}
\textbf{Owner. }\nolinkurl{CAT-BA43-OPEN-EXTENSION-COMPOSITION}\
\textbf{Status. }\texttt{OPEN}; \textbf{reduction status. }\texttt{LEAF}; \textbf{class. }\texttt{CAT}.
\uses{}
\textbf{Exact output required. }Construct \nolinkurl{Ru_!Rj_! ~= R(uj)_!} and prove unit, composition, and pentagon.
\openobligation{CAT-BA43-OPEN-EXTENSION-COMPOSITION}
\end{openproblem}

\Needspace{8\baselineskip}
\begin{openproblem}[Exact unresolved interface]
\label{CAT-BA43-SUCCESSIVE-WITNESS-INDEPENDENCE}
\textbf{Owner. }\nolinkurl{CAT-BA43-SUCCESSIVE-WITNESS-INDEPENDENCE}\
\textbf{Status. }\texttt{OPEN}; \textbf{reduction status. }\texttt{LEAF}; \textbf{class. }\texttt{CAT}.
\uses{}
\textbf{Exact output required. }Compare two typed successive witnesses through the refinement system and prove the explicit comparison square.
\openobligation{CAT-BA43-SUCCESSIVE-WITNESS-INDEPENDENCE}
\end{openproblem}

\Needspace{8\baselineskip}
\begin{openproblem}[Exact unresolved interface]
\label{CAT-BA43-TRIPLE-PASTING}
\textbf{Owner. }\nolinkurl{CAT-BA43-TRIPLE-PASTING}\
\textbf{Status. }\texttt{OPEN}; \textbf{reduction status. }\texttt{LEAF}; \textbf{class. }\texttt{CAT}.
\uses{}
\textbf{Exact output required. }Prove the explicit triple-composition pasting equation yielding associativity.
\openobligation{CAT-BA43-TRIPLE-PASTING}
\end{openproblem}

\Needspace{8\baselineskip}
\begin{openproblem}[Exact unresolved interface]
\label{CAT-BA33-DERIVED-OPEN-EXCHANGE}
\textbf{Owner. }\nolinkurl{CAT-BA33-DERIVED-OPEN-EXCHANGE}\
\textbf{Status. }\texttt{OPEN}; \textbf{reduction status. }\texttt{LEAF}; \textbf{class. }\texttt{CAT}.
\uses{}
\textbf{Exact output required. }Construct the derived open-extension base-change map \nolinkurl{Lbar g^* Rj_! -> Rj'_! Lg'^*} and prove naturality and its identity/horizontal/vertical pasting equations.
\openobligation{CAT-BA33-DERIVED-OPEN-EXCHANGE}
\end{openproblem}

\Needspace{8\baselineskip}
\begin{openproblem}[Exact unresolved interface]
\label{GEO-BA43-SCHEMATIC-CLOSURE-DENSIFICATION}
\textbf{Owner. }\nolinkurl{GEO-BA43-SCHEMATIC-CLOSURE-DENSIFICATION}\
\textbf{Status. }\texttt{OPEN}; \textbf{reduction status. }\texttt{LEAF}; \textbf{class. }\texttt{GEO}.
\uses{}
\textbf{Exact output required. }For noetherian finite-dimensional \nolinkurl{Y,Y'}, the schematic closure of \nolinkurl{X'} in the raw pullback is proper finitely presented over \nolinkurl{Y'}, has dense source open, and maps closed-properly to the raw model.
\openobligation{GEO-BA43-SCHEMATIC-CLOSURE-DENSIFICATION}
\end{openproblem}

\Needspace{8\baselineskip}
\begin{openproblem}[Exact unresolved interface]
\label{CAT-GEO-BA43-DENSE-TO-RAW-SUPPORT}
\textbf{Owner. }\nolinkurl{CAT-GEO-BA43-DENSE-TO-RAW-SUPPORT}\
\textbf{Status. }\texttt{OPEN}; \textbf{reduction status. }\texttt{LEAF}; \textbf{class. }\texttt{CAT/GEO}.
\uses{}
\textbf{Exact output required. }Construct only \nolinkurl{gamma_raw: Rd_*Rj_den! ~= Rj_raw!} for the schematic-closure closed immersion \nolinkurl{d_(C,g)}, which is the identity over the source open. It supplies no \nolinkurl{sigma} and no naturality/coherence fields.
\openobligation{CAT-GEO-BA43-DENSE-TO-RAW-SUPPORT}
\end{openproblem}

\Needspace{8\baselineskip}
\begin{openproblem}[Exact unresolved interface]
\label{CAT-BA43-DENSIFICATION-REFINEMENT-NATURALITY}
\textbf{Owner. }\nolinkurl{CAT-BA43-DENSIFICATION-REFINEMENT-NATURALITY}\
\textbf{Status. }\texttt{OPEN}; \textbf{reduction status. }\texttt{LEAF}; \textbf{class. }\texttt{CAT}.
\uses{}
\textbf{Exact output required. }Prove square \nolinkurl{(BA43-DEN-REF)}: \nolinkurl{rho_raw sigma_D = sigma_C rho_den} for every typed refinement \nolinkurl{q:D->C}.
\openobligation{CAT-BA43-DENSIFICATION-REFINEMENT-NATURALITY}
\end{openproblem}

\Needspace{8\baselineskip}
\begin{openproblem}[Exact unresolved interface]
\label{CAT-BA43-DENSIFICATION-IDENTITY}
\textbf{Owner. }\nolinkurl{CAT-BA43-DENSIFICATION-IDENTITY}\
\textbf{Status. }\texttt{OPEN}; \textbf{reduction status. }\texttt{LEAF}; \textbf{class. }\texttt{CAT}.
\uses{}
\textbf{Exact output required. }Prove equation \nolinkurl{(BA43-DEN-ID)}: \nolinkurl{sigma_(C,id_Y)=1_(F_C)} under the canonical raw/dense identity identifications.
\openobligation{CAT-BA43-DENSIFICATION-IDENTITY}
\end{openproblem}

\Needspace{8\baselineskip}
\begin{openproblem}[Exact unresolved interface]
\label{CAT-BA43-DENSIFICATION-ITERATED-COHERENCE}
\textbf{Owner. }\nolinkurl{CAT-BA43-DENSIFICATION-ITERATED-COHERENCE}\
\textbf{Status. }\texttt{OPEN}; \textbf{reduction status. }\texttt{LEAF}; \textbf{class. }\texttt{CAT/GEO}.
\uses{}
\textbf{Exact output required. }Construct \nolinkurl{delta_(C;g,h)} and \nolinkurl{tau_(C;g,h)} with the displayed endpoints and prove square \nolinkurl{(BA43-DEN-ITER)}.
\openobligation{CAT-BA43-DENSIFICATION-ITERATED-COHERENCE}
\end{openproblem}

\Needspace{8\baselineskip}
\begin{openproblem}[Exact unresolved interface]
\label{CAT-BA43-DENSIFICATION-OPEN-EXCHANGE}
\textbf{Owner. }\nolinkurl{CAT-BA43-DENSIFICATION-OPEN-EXCHANGE}\
\textbf{Status. }\texttt{OPEN}; \textbf{reduction status. }\texttt{LEAF}; \textbf{class. }\texttt{CAT}.
\uses{}
\textbf{Exact output required. }For \nolinkurl{A=L bar(g)_raw^* Rj_C!K}, construct the closed-support map \nolinkurl{s_d} and independently constructed raw/dense beta maps, then prove square \nolinkurl{(BA43-DEN-OPX)}.
\openobligation{CAT-BA43-DENSIFICATION-OPEN-EXCHANGE}
\end{openproblem}

\Needspace{8\baselineskip}
\begin{openproblem}[Exact unresolved interface]
\label{CAT-GEO-BA33-OPEN-EXTENSION-CONSTRUCTIBILITY}
\textbf{Owner. }\nolinkurl{CAT-GEO-BA33-OPEN-EXTENSION-CONSTRUCTIBILITY}\
\textbf{Status. }\texttt{OPEN}; \textbf{reduction status. }\texttt{LEAF}; \textbf{class. }\texttt{CAT/GEO}.
\uses{}
\textbf{Exact output required. }For the typed open immersion and \nolinkurl{K in D_c^+}, prove \nolinkurl{Rj_!K in D_c^+} in the exact \nolinkurl{Lambda_n} site consumed by finite PBC.
\openobligation{CAT-GEO-BA33-OPEN-EXTENSION-CONSTRUCTIBILITY}
\end{openproblem}

\Needspace{8\baselineskip}
\begin{openproblem}[Exact unresolved interface]
\label{CAT-BA33-FINITE-BC-PASTING}
\textbf{Owner. }\nolinkurl{CAT-BA33-FINITE-BC-PASTING}\
\textbf{Status. }\texttt{OPEN}; \textbf{reduction status. }\texttt{LEAF}; \textbf{class. }\texttt{CAT}.
\uses{}
\textbf{Exact output required. }Prove identity and horizontal/vertical pasting for the final compact-support map by combining the already proved proper, open, densification, and \nolinkurl{kappa} diagrams.
\openobligation{CAT-BA33-FINITE-BC-PASTING}
\end{openproblem}

\section{Imported dependency registry}
\statusnote{These labels occur in an original \string\uses or cross-reference but are defined outside the extracted S0--S10 block.  A red status is an additional explicit proof gap; a green status means only that the frozen manifest records the imported declaration as defined or proved.}
\refstepcounter{importedinterface}\label{def:weil-a1-rfshriek-selected}
\noindent\textbf{Imported interface \theimportedinterface.} \nolinkurl{def:weil-a1-rfshriek-selected}\
\textbf{Kind. }\texttt{definition}; \textbf{manifest status. }\textcolor{red!65!black}{\texttt{CONDITIONAL-ON-OPEN-ANCESTOR}}.
\par\smallskip
\refstepcounter{importedinterface}\label{def:weil-a1-torsion-derived-category}
\noindent\textbf{Imported interface \theimportedinterface.} \nolinkurl{def:weil-a1-torsion-derived-category}\
\textbf{Kind. }\texttt{definition}; \textbf{manifest status. }\textcolor{green!40!black}{\texttt{DEFINED-IN-PREFIX}}.
\par\smallskip
\refstepcounter{importedinterface}\label{def:weil-a2-completed-tensor}
\noindent\textbf{Imported interface \theimportedinterface.} \nolinkurl{def:weil-a2-completed-tensor}\
\textbf{Kind. }\texttt{definition}; \textbf{manifest status. }\textcolor{green!40!black}{\texttt{DEFINED-IN-PREFIX}}.
\par\smallskip
\refstepcounter{importedinterface}\label{def:weil-a2-constructible-stratification-witness}
\noindent\textbf{Imported interface \theimportedinterface.} \nolinkurl{def:weil-a2-constructible-stratification-witness}\
\textbf{Kind. }\texttt{definition}; \textbf{manifest status. }\textcolor{green!40!black}{\texttt{DEFINED-IN-PREFIX}}.
\par\smallskip
\refstepcounter{importedinterface}\label{def:weil-a2-dbc-torsion}
\noindent\textbf{Imported interface \theimportedinterface.} \nolinkurl{def:weil-a2-dbc-torsion}\
\textbf{Kind. }\texttt{definition}; \textbf{manifest status. }\textcolor{green!40!black}{\texttt{DEFINED-IN-PREFIX}}.
\par\smallskip
\refstepcounter{importedinterface}\label{def:weil-a2-dbc-zl}
\noindent\textbf{Imported interface \theimportedinterface.} \nolinkurl{def:weil-a2-dbc-zl}\
\textbf{Kind. }\texttt{definition}; \textbf{manifest status. }\textcolor{green!40!black}{\texttt{DEFINED-IN-PREFIX}}.
\par\smallskip
\refstepcounter{importedinterface}\label{def:weil-a2-derived-category}
\noindent\textbf{Imported interface \theimportedinterface.} \nolinkurl{def:weil-a2-derived-category}\
\textbf{Kind. }\texttt{definition}; \textbf{manifest status. }\textcolor{green!40!black}{\texttt{DEFINED-IN-PREFIX}}.
\par\smallskip
\refstepcounter{importedinterface}\label{def:weil-a2-etale-origin-datum}
\noindent\textbf{Imported interface \theimportedinterface.} \nolinkurl{def:weil-a2-etale-origin-datum}\
\textbf{Kind. }\texttt{definition}; \textbf{manifest status. }\textcolor{green!40!black}{\texttt{DEFINED-IN-PREFIX}}.
\par\smallskip
\refstepcounter{importedinterface}\label{def:weil-a2-etale-proetale-direct-image-comparison}
\noindent\textbf{Imported interface \theimportedinterface.} \nolinkurl{def:weil-a2-etale-proetale-direct-image-comparison}\
\textbf{Kind. }\texttt{definition}; \textbf{manifest status. }\textcolor{green!40!black}{\texttt{DEFINED-IN-PREFIX}}.
\par\smallskip
\refstepcounter{importedinterface}\label{def:weil-a2-etale-proetale-pullback-comparison}
\noindent\textbf{Imported interface \theimportedinterface.} \nolinkurl{def:weil-a2-etale-proetale-pullback-comparison}\
\textbf{Kind. }\texttt{definition}; \textbf{manifest status. }\textcolor{green!40!black}{\texttt{DEFINED-IN-PREFIX}}.
\par\smallskip
\refstepcounter{importedinterface}\label{def:weil-a2-finite-etale-proper-base-change-map}
\noindent\textbf{Imported interface \theimportedinterface.} \nolinkurl{def:weil-a2-finite-etale-proper-base-change-map}\
\textbf{Kind. }\texttt{definition}; \textbf{manifest status. }\textcolor{green!40!black}{\texttt{DEFINED-IN-PREFIX}}.
\par\smallskip
\refstepcounter{importedinterface}\label{def:weil-a2-finite-level-etale-origin-realization}
\noindent\textbf{Imported interface \theimportedinterface.} \nolinkurl{def:weil-a2-finite-level-etale-origin-realization}\
\textbf{Kind. }\texttt{definition}; \textbf{manifest status. }\textcolor{green!40!black}{\texttt{DEFINED-IN-PREFIX}}.
\par\smallskip
\refstepcounter{importedinterface}\label{def:weil-a2-finite-reduction}
\noindent\textbf{Imported interface \theimportedinterface.} \nolinkurl{def:weil-a2-finite-reduction}\
\textbf{Kind. }\texttt{definition}; \textbf{manifest status. }\textcolor{green!40!black}{\texttt{DEFINED-IN-PREFIX}}.
\par\smallskip
\refstepcounter{importedinterface}\label{def:weil-a2-proetale-proper-base-change-map}
\noindent\textbf{Imported interface \theimportedinterface.} \nolinkurl{def:weil-a2-proetale-proper-base-change-map}\
\textbf{Kind. }\texttt{definition}; \textbf{manifest status. }\textcolor{green!40!black}{\texttt{DEFINED-IN-PREFIX}}.
\par\smallskip
\refstepcounter{importedinterface}\label{def:weil-a2-proper-complete-base-change-datum}
\noindent\textbf{Imported interface \theimportedinterface.} \nolinkurl{def:weil-a2-proper-complete-base-change-datum}\
\textbf{Kind. }\texttt{definition}; \textbf{manifest status. }\textcolor{green!40!black}{\texttt{DEFINED-IN-PREFIX}}.
\par\smallskip
\refstepcounter{importedinterface}\label{def:weil-a2-reduction-transition}
\noindent\textbf{Imported interface \theimportedinterface.} \nolinkurl{def:weil-a2-reduction-transition}\
\textbf{Kind. }\texttt{definition}; \textbf{manifest status. }\textcolor{green!40!black}{\texttt{DEFINED-IN-PREFIX}}.
\par\smallskip
\refstepcounter{importedinterface}\label{def:weil-a2-six-operation-coefficient-target}
\noindent\textbf{Imported interface \theimportedinterface.} \nolinkurl{def:weil-a2-six-operation-coefficient-target}\
\textbf{Kind. }\texttt{definition}; \textbf{manifest status. }\textcolor{green!40!black}{\texttt{DEFINED-IN-PREFIX}}.
\par\smallskip
\refstepcounter{importedinterface}\label{def:weil-a2-torsion-constructibility-witness}
\noindent\textbf{Imported interface \theimportedinterface.} \nolinkurl{def:weil-a2-torsion-constructibility-witness}\
\textbf{Kind. }\texttt{definition}; \textbf{manifest status. }\textcolor{green!40!black}{\texttt{DEFINED-IN-PREFIX}}.
\par\smallskip
\refstepcounter{importedinterface}\label{def:weil-a2-transported-finite-level-base-change-map}
\noindent\textbf{Imported interface \theimportedinterface.} \nolinkurl{def:weil-a2-transported-finite-level-base-change-map}\
\textbf{Kind. }\texttt{definition}; \textbf{manifest status. }\textcolor{green!40!black}{\texttt{DEFINED-IN-PREFIX}}.
\par\smallskip
\refstepcounter{importedinterface}\label{def:weil-small-etale-pretopology}
\noindent\textbf{Imported interface \theimportedinterface.} \nolinkurl{def:weil-small-etale-pretopology}\
\textbf{Kind. }\texttt{definition}; \textbf{manifest status. }\textcolor{green!40!black}{\texttt{DEFINED-IN-PREFIX}}.
\par\smallskip
\refstepcounter{importedinterface}\label{def:weil-small-etale-topology}
\noindent\textbf{Imported interface \theimportedinterface.} \nolinkurl{def:weil-small-etale-topology}\
\textbf{Kind. }\texttt{definition}; \textbf{manifest status. }\textcolor{green!40!black}{\texttt{DEFINED-IN-PREFIX}}.
\par\smallskip
\refstepcounter{importedinterface}\label{def:weil-small-proetale-category}
\noindent\textbf{Imported interface \theimportedinterface.} \nolinkurl{def:weil-small-proetale-category}\
\textbf{Kind. }\texttt{definition}; \textbf{manifest status. }\textcolor{green!40!black}{\texttt{DEFINED-IN-PREFIX}}.
\par\smallskip
\refstepcounter{importedinterface}\label{def:weil-small-proetale-precoverage}
\noindent\textbf{Imported interface \theimportedinterface.} \nolinkurl{def:weil-small-proetale-precoverage}\
\textbf{Kind. }\texttt{definition}; \textbf{manifest status. }\textcolor{green!40!black}{\texttt{DEFINED-IN-PREFIX}}.
\par\smallskip
\refstepcounter{importedinterface}\label{def:weil-small-proetale-topology}
\noindent\textbf{Imported interface \theimportedinterface.} \nolinkurl{def:weil-small-proetale-topology}\
\textbf{Kind. }\texttt{definition}; \textbf{manifest status. }\textcolor{green!40!black}{\texttt{DEFINED-IN-PREFIX}}.
\par\smallskip
\refstepcounter{importedinterface}\label{lemma:weil-a2-finite-reduction-package-independent}
\noindent\textbf{Imported interface \theimportedinterface.} \nolinkurl{lemma:weil-a2-finite-reduction-package-independent}\
\textbf{Kind. }\texttt{lemma}; \textbf{manifest status. }\textcolor{green!40!black}{\texttt{PROVED-IN-PREFIX}}.
\par\smallskip
\refstepcounter{importedinterface}\label{prop:weil-a2-finite-reduction-functor}
\noindent\textbf{Imported interface \theimportedinterface.} \nolinkurl{prop:weil-a2-finite-reduction-functor}\
\textbf{Kind. }\texttt{proposition}; \textbf{manifest status. }\textcolor{green!40!black}{\texttt{PROVED-IN-PREFIX}}.
\par\smallskip

\section{Retired or invalid proof routes}

\begin{warning}[Stalk-family descent is not a proof of PBC]
The historical certificate that starts with arbitrary specialization-compatible
stalk maps and asks them to descend to a derived morphism is not valid.
Geometric stalks detect whether an already constructed morphism is an
isomorphism; they are not a fullness theorem for the derived category.
\retiredobligation{OBL-PBC-FINITE-STALK-DESCENT}
\end{warning}

\begin{warning}[Certificate projection is conditional only]
An input record containing an inverse to the Beck--Chevalley map, both inverse
equations, and all coherence fields merely repackages the theorem.  Projection
from such a record proves only a conditional implication and cannot discharge
\nolinkurl{OBL-A1-PROPER-BASE-CHANGE-FINITE-TORSION}.
\end{warning}

\begin{warning}[Completed coefficient shortcut]
Coefficient reduction does not commute with proper direct image merely by
functoriality.  The completed proof must explicitly use a perfect quotient,
the projection-formula comparison, compatible finite-level maps, the
product-cone model for $R\varprojlim$, and reconstruction.  The older A5
shortcut is not an admissible proof source.
\end{warning}

\statusnote{The detailed body was retained because it exposes the intended
argument.  The warnings above and the gap ledger govern its logical status.}
